% Generated mechanically from frozen input p07/ja/algebraization.tex.
% Do not edit this generated file; edit the bound locale source instead.

% OK, start here.
%

\setcounter{chapter}{51}
\chapter{代数幾何と形式幾何}



\phantomsection
\label{algebraization-section-phantom}


\section{序論}
\label{algebraization-section-introduction}

\noindent
本章では形式代数幾何の研究を続け、とりわけ形式的対象が代数的対象の
完備化であるかという問題を扱う。基本文献は \cite{SGA2} である。
Stacks Project ですでに論じた結果を以下に列挙する。
\begin{enumerate}
\item 形式関数定理。Cohomology of Schemes, 第
\ref{coherent-section-theorem-formal-functions} 節を参照。
\item 連接形式加群。Cohomology of Schemes, 第
\ref{coherent-section-coherent-formal} 節を参照。
\item Grothendieck の存在定理。Cohomology of Schemes, 第
\ref{coherent-section-existence}、
\ref{coherent-section-existence-proper}、および
\ref{coherent-section-existence-proper-support} 節を参照。
\item Grothendieck の代数化定理。Cohomology of Schemes, 第
\ref{coherent-section-algebraization} 節を参照。
\item より一般的な Grothendieck の存在定理。More on Flatness, 第
\ref{flat-section-existence} および
\ref{flat-section-existence-derived} 節を参照。
\end{enumerate}
以下、本章の内容を概観する。

\medskip\noindent
スキーム $X$ と有限型の準連接イデアル層
$\mathcal{I} \subset \mathcal{O}_X$ を取る。本章で扱う問題の多くは、逆系
$(\mathcal{F}_n)$ であって、その各項が準連接 $\mathcal{O}_X$-加群であり
$\mathcal{F}_n = \mathcal{F}_{n + 1}/\mathcal{I}^n\mathcal{F}_{n + 1}$
を満たすものに関係する。重要な特別な場合は、$X$ が Noether 環 $A$ 上の
スキームであり、$\mathcal{I} = I \mathcal{O}_X$ となる場合である。ここで
$I \subset A$ はあるイデアルである。Cohomology, 第
\ref{cohomology-section-inverse-systems}、
\ref{cohomology-section-inverse-systems-bis}、および
\ref{cohomology-section-inverse-systems-tri} 節にはいくつかの一般的結果がある。
本章の第 \ref{algebraization-section-ML-degree-zero} 節と第
\ref{algebraization-section-formal-functions-principal} 節には、スキームと準連接加群に固有の
結果を収める。第 \ref{algebraization-section-formal-sections-cd-one} 節では、逆極限
$\lim H^p(X, \mathcal{F}_n)$ 上の位相が $I$-進位相であることを、
$\text{cd}(A, I) = 1$ の場合に証明する。本章の一つの主題は、主イデアルの場合
$I = (f)$ に証明された結果が、$\text{cd}(A, I) = 1$ のみを仮定する場合にも
成り立つことを示すことである。

\medskip\noindent
第 \ref{algebraization-section-derived-completion} 節では、環付きサイト
$(\mathcal{C}, \mathcal{O})$ 上の加群について、有限型イデアル層
$\mathcal{I}$ に関する導来完備化を論じる。この節は、More on Algebra, 第
\ref{more-algebra-section-derived-completion} 節で述べた可換代数における
導来完備化の理論の自然な続編である。第一の主要結果は、導来完備化が存在する
ことである。第二の主要結果は、環付きサイトの射 $f$ に対して導来完備化が
導来順像と可換するというものである：
$$
(Rf_*K)^\wedge = Rf_*(K^\wedge)
$$
ただし、上側のイデアル層が下側のイデアルから来る切断によって局所的に
生成されるものとする。補題
\ref{algebraization-lemma-pushforward-commutes-with-derived-completion} を参照されたい。
対象となるイデアルが大域的に有限生成である場合、これら二つの主要結果は
きわめて初等的であることを強調しておく。本章におけるこの理論の応用では
常にこの場合に当たる。上の等式は形式関数定理の「正しい」形である。
第 \ref{algebraization-section-formal-functions} 節の議論を参照されたい。

\medskip\noindent
Noether 環 $A$ を取り、$I, J$ を $A$ の二つのイデアルとする。
$M$ を有限 $A$-加群とする。本章で次に扱うのは写像
$$
R\Gamma_J(M) \longrightarrow R\Gamma_J(M)^\wedge
$$
すなわち $M$ の局所コホモロジーからその導来 $I$-進完備化への写像である。
適切な深さの条件を課すと、この写像はある次数範囲でコホモロジー上の同型に
なる。第 \ref{algebraization-section-algebraization-sections-general} 節では、本質的に今述べた
一般性のもとで議論する。第 \ref{algebraization-section-algebraization-punctured} 節では、$A$ が
局所環であり $J = \mathfrak m$ が極大イデアルであると仮定する。本章のこの部分に
属する他の二節より先に、この節を読むことを勧める。最後に第
\ref{algebraization-section-bootstrap} 節では、局所的な場合を足掛かりとして、一般の場合に
戻ってより強い結果を得る。

\medskip\noindent
本章の次の部分では、局所コホモロジーの完備化に関する結果を用いて、連接加群の
完備化のコホモロジーに関する、網羅的ではない一連の結果を得る。より正確には、
$A$ を Noether 環、$I \subset A$ をイデアル、$U \subset \Spec(A)$ を開部分スキーム
とする。$\mathcal{F}$ が連接 $\mathcal{O}_U$-加群ならば、写像
$$
H^i(U, \mathcal{F}) \longrightarrow \lim H^i(U, \mathcal{F}/I^n\mathcal{F})
$$
を考え、ある次数範囲で同型が得られるかを問うことができる。第
\ref{algebraization-section-algebraization-sections} 節では、$U$ が局所環の穴あきスペクトルである
いくつかの例を詳しく調べる。第
\ref{algebraization-section-algebraization-sections-coherent} 節では一般の場合を論じる。
第 \ref{algebraization-section-connected} 節では、得られた結果のいくつかを代数幾何における
連結性の問題へ応用する。

\medskip\noindent
本章の残りの節では、連接形式加群の代数化を論じる。言い換えると、上のような
連接加群の逆系 $(\mathcal{F}_n)$、すなわち $U$ 上の逆系であって
$\mathcal{F}_n = \mathcal{F}_{n + 1}/I^n\mathcal{F}_{n + 1}$
を満たすものが与えられたとき、連接 $\mathcal{O}_U$-加群 $\mathcal{F}$ であって
$\mathcal{F}_n = \mathcal{F}/I^n\mathcal{F}$
がすべての $n$ に対して成り立つものが存在するかを問う。第
\ref{algebraization-section-algebraization-modules} 節を読むことを勧める。そこには、この問題の
正確な定式化、有用な一般結果（補題 \ref{algebraization-lemma-when-done}）、および非自明な応用
（補題 \ref{algebraization-lemma-algebraization-principal-variant}）がある。この場合を本質的に
越える結果を証明するには、さらにかなりの理論を展開しなければならない。
本章で得られるこの種の最も強い結果については、第
\ref{algebraization-section-algebraization-modules-conclusion} 節を参照されたい。



\section{形式切断 I}
\label{algebraization-section-ML-degree-zero}

\noindent
まず Cohomology, 第 \ref{cohomology-section-inverse-systems} 節を
参照することを勧める。

\begin{lemma}
\label{algebraization-lemma-properties-system}
$X$ をスキーム、$\mathcal{I} \subset \mathcal{O}_X$ を準連接イデアル層とする。逆系
$$
\ldots \to \mathcal{F}_3 \to \mathcal{F}_2 \to \mathcal{F}_1
$$
の各項が準連接 $\mathcal{O}_X$-加群であり、かつ
$\mathcal{F}_n = \mathcal{F}_{n + 1}/\mathcal{I}^n\mathcal{F}_{n + 1}$
を満たすとする。$\mathcal{F} = \lim \mathcal{F}_n$ と置く。このとき、
\begin{enumerate}
\item $\mathcal{F} = R\lim \mathcal{F}_n$ であり、
\item 任意のアフィン開集合 $U \subset X$ に対して
$H^p(U, \mathcal{F}) = 0$ が $p > 0$ のとき成り立ち、
\item 各 $p$ に対して短完全系列
$0 \to R^1\lim H^{p - 1}(X, \mathcal{F}_n) \to
H^p(X, \mathcal{F}) \to \lim H^p(X, \mathcal{F}_n) \to 0$
がある。
\end{enumerate}
さらに $\mathcal{I}$ が有限型ならば、
\begin{enumerate}
\item[(4)]
$\mathcal{F}_n = \mathcal{F}/\mathcal{I}^n\mathcal{F}$ であり、
\item[(5)]
$\mathcal{I}^n \mathcal{F} = \lim_{m \geq n} \mathcal{I}^n\mathcal{F}_m$
である。
\end{enumerate}
\end{lemma}

\begin{proof}
(1)、(2)、(3) は、全射な遷移写像をもつ準連接加群の逆系に関する一般的事実で
ある。Derived Categories of Schemes, 補題
\ref{perfect-lemma-Rlim-quasi-coherent} および Cohomology, 補題
\ref{cohomology-lemma-RGamma-commutes-with-Rlim} を参照されたい。
次に $\mathcal{I}$ が有限型であると仮定する。$U \subset X$ をアフィン開集合とし、
$U = \Spec(A)$ で $\mathcal{I}|_U$ が $I \subset A$ に対応するとする。このとき
$I$ は有限生成イデアルである。準連接 $\mathcal{O}_U$-加群の圏と $A$-加群の圏の
同値（Schemes, 補題 \ref{schemes-lemma-equivalence-quasi-coherent}）により、
$M_n = \mathcal{F}_n(U)$ は $A$-加群の逆系であり、
$M_n = M_{n + 1}/I^nM_{n + 1}$ を満たす。したがって
$$
M = \mathcal{F}(U) = \lim \mathcal{F}_n(U) = \lim M_n
$$
これは Algebra, 補題 \ref{algebra-lemma-limit-complete} により $I$-進完備加群であり、
$M/I^nM = M_n$ を満たす。これで (4) が証明された。
(5) は
$\lim_{m \geq n} I^nM/I^mM = I^nM$
という主張に移される。$I^mM = I^{m - n} \cdot I^nM$ なので、これはちょうど
$I^mM$ が $I$-進完備であるという主張である。これは Algebra, 補題
\ref{algebra-lemma-hathat-finitely-generated} と $M$ が完備であることから従う。
\end{proof}





\section{形式切断 II}
\label{algebraization-section-formal-functions-principal}

\noindent
まず Cohomology, 第
\ref{cohomology-section-inverse-systems-bis} および
\ref{cohomology-section-inverse-systems-tri} 節を参照することを勧める。

\begin{lemma}
\label{algebraization-lemma-equivalent-f-good}
$X$ をスキームとし、$f \in \Gamma(X, \mathcal{O}_X)$ とする。逆系
$$
\ldots \to \mathcal{F}_3 \to \mathcal{F}_2 \to \mathcal{F}_1
$$
の各項が準連接 $\mathcal{O}_X$-加群であるとする。次の条件は同値である。
\begin{enumerate}
\item すべての $n \geq 1$ に対し、写像
$f : \mathcal{F}_{n + 1} \to \mathcal{F}_{n + 1}$ は
$\mathcal{F}_{n + 1} \to \mathcal{F}_n$ を経由して分解し、短完全系列
$0 \to \mathcal{F}_n \to \mathcal{F}_{n + 1} \to \mathcal{F}_1 \to 0$
を与える。
\item すべての $n \geq 1$ に対し、写像
$f^n : \mathcal{F}_{n + 1} \to \mathcal{F}_{n + 1}$
は $\mathcal{F}_{n + 1} \to \mathcal{F}_1$ を経由して分解し、短完全系列
$0 \to \mathcal{F}_1 \to \mathcal{F}_{n + 1} \to \mathcal{F}_n \to 0$
を与える。
\item ある $\mathcal{O}_X$-加群 $\mathcal{G}$ であって、$f$-可除かつ
$\mathcal{F}_n = \mathcal{G}[f^n]$ を満たすものが存在する。
\item ある $\mathcal{O}_X$-加群 $\mathcal{F}$ であって、$f$-捩れがなく、かつ
$\mathcal{F}_n = \mathcal{F}/f^n\mathcal{F}$ を満たすものが存在する。
\end{enumerate}
\end{lemma}

\begin{proof}
(1)、(2)、(3) の同値性と含意 (4) $\Rightarrow$ (1) は Cohomology, 補題
\ref{cohomology-lemma-equivalent-f-good} で証明されている。(1) が成り立つと仮定し、
$\mathcal{F} = \lim \mathcal{F}_n$ と置く。補題
\ref{algebraization-lemma-properties-system} の (4) により
$\mathcal{F}_n = \mathcal{F}/f^n\mathcal{F}$ である。$U \subset X$ を開集合とし、
$s = (s_n) \in \mathcal{F}(U) = \lim \mathcal{F}_n(U)$ とする。
$n \geq 1$ を選ぶ。$fs = 0$ ならば、条件 (1) により $s_{n + 1}$ は
$\mathcal{F}_{n + 1} \to \mathcal{F}_n$ の核に属する。したがって $s_n = 0$ である。
$n$ は任意だったので $s = 0$ である。ゆえに $\mathcal{F}$ は $f$-捩れをもたない。
\end{proof}

\begin{lemma}
\label{algebraization-lemma-formal-functions-principal}
\begin{reference}
\cite[補題 1.6]{Bhatt-local} をわずかに改善した形
\end{reference}
$A$ を環、$f \in A$ とし、$X$ を $A$ 上のスキームとする。
$\mathcal{F}$ を準連接 $\mathcal{O}_X$-加群とする。
$\mathcal{F}[f^n] = \Ker(f^n : \mathcal{F} \to \mathcal{F})$ が安定すると仮定する。
このとき、
$$
R\Gamma(X, \lim \mathcal{F}/f^n\mathcal{F}) =
R\Gamma(X, \mathcal{F})^\wedge
$$
である。ここで右辺はイデアル $(f) \subset A$ に関する導来完備化を表す。
したがって $p \in \mathbf{Z}$ に対して、行と列が完全な可換図式
$$
\xymatrix{
& 0 & 0 \\
0 \ar[r] &
\widehat{H^p(X, \mathcal{F})} \ar[r] \ar[u] &
\lim H^p(X, \mathcal{F}/f^n\mathcal{F}) \ar[r] \ar[u] &
T_f(H^{p + 1}(X, \mathcal{F})) \ar[r] &
0 \\
0 \ar[r] &
H^0(H^p(X, \mathcal{F})^\wedge) \ar[r] \ar[u] &
H^p(X, \lim \mathcal{F}/f^n\mathcal{F}) \ar[r] \ar[u] &
T_f(H^{p + 1}(X, \mathcal{F})) \ar[r] \ar@{=}[u] &
0 \\
&
R^1\lim H^p(X, \mathcal{F})[f^n] \ar[u] \ar[r]^\cong &
R^1\lim H^{p - 1}(X, \mathcal{F}/f^n\mathcal{F}) \ar[u] \\
& 0 \ar[u] & 0 \ar[u]
}
$$
を得る。ここで
$\widehat{H^p(X, \mathcal{F})} =
\lim H^p(X, \mathcal{F})/f^n H^p(X, \mathcal{F})$
は通常の $f$-進完備化であり、$T_f(-)$ は More on Algebra, 例
\ref{more-algebra-example-spectral-sequence-principal} における $f$-進 Tate 加群を
表す。
\end{lemma}

\begin{proof}
補題 \ref{algebraization-lemma-properties-system} により
$\lim \mathcal{F}/f^n\mathcal{F} = R\lim \mathcal{F}/f^n \mathcal{F}$
である。残りはすべて Cohomology, 例
\ref{cohomology-example-formal-functions-principal} から従う。
\end{proof}









\section{形式切断 III}
\label{algebraization-section-formal-sections-cd-one}

\noindent
本節では補題 \ref{algebraization-lemma-topology-I-adic} を証明する。この補題は（Noether
スキームと連接加群の設定において）Cohomology, 補題
\ref{cohomology-lemma-topology-I-adic-f} の類似である。ただし、イデアル $I$ は
主イデアルとは仮定せず、$\text{cd}(A, I) = 1$ を満たすものとする。

\begin{lemma}
\label{algebraization-lemma-cd-one}
$I = (f_1, \ldots, f_r)$ を Noether 環 $A$ のイデアルとする。
$\text{cd}(A, I) = 1$ ならば、$c \geq 1$ と写像
$\varphi_j : I^c \to A$ であって、$\sum f_j \varphi_j : I^c \to I$ が包含写像と
なるものが存在する。
\end{lemma}

\begin{proof}
$\text{cd}(A, I) = 1$ なので、補集合 $U = \Spec(A) \setminus V(I)$ はアフィンである
（Local Cohomology, 補題 \ref{local-cohomology-lemma-cd-is-one}）。
$U = \Spec(B)$ と書く。このとき $IB = B$ であるから、
$1 = \sum_{j = 1, \ldots, r} f_j b_j$ と書ける。ここである $b_j \in B$ を取った。
Cohomology of Schemes, 補題 \ref{coherent-lemma-homs-over-open} により、$b_j$ を写像
$\varphi_j : I^c \to A$ で表せる。これはある $c \geq 0$ に対して成り立つ。同じ
補題により、$\sum f_j \varphi_j : I^c \to I \subset A$ は標準的な埋め込みになる。
ただし、必要なら $c$ をより大きな整数に置き換える。
\end{proof}

\begin{lemma}
\label{algebraization-lemma-cd-one-extend}
$I = (f_1, \ldots, f_r)$ を Noether 環 $A$ のイデアルとし、
$\text{cd}(A, I) = 1$ とする。$c \geq 1$、$\varphi_j : I^c \to A$、
$j = 1, \ldots, r$ を補題 \ref{algebraization-lemma-cd-one} のものとする。このとき次数付き
$A$-代数写像
$$
\Phi : \bigoplus\nolimits_{n \geq 0} I^{nc} \to A[T_1, \ldots, T_r]
$$
であって、$\Phi(g) = \sum \varphi_j(g) T_j$ を満たすものが一意に存在する。ここで
$g \in I^c$ である。さらに、$\Phi$ と写像
$A[T_1, \ldots, T_r] \to \bigoplus_{n \geq 0} I^n$,
$T_j \mapsto f_j$ との合成は包含写像
$\bigoplus_{n \geq 0} I^{nc} \to \bigoplus_{n \geq 0} I^n$ である。
\end{lemma}

\begin{proof}
各 $j$ と $m \geq c$ に対し、$\varphi_j$ の $I^m$ への制限は写像
$\varphi_j : I^m \to I^{m - c}$ である。
$j_1, \ldots, j_n \in \{1, \ldots, r\}$ が与えられたとき、合成
$$
\varphi_{j_1} \ldots \varphi_{j_n} :
I^{nc} \to I^{(n - 1)c} \to \ldots \to I^c \to A
$$
は添字 $j_1, \ldots, j_n$ の順序に依存しないと主張する。実際、
$g = g_1 \ldots g_n$ かつ $g_i \in I^c$ ならば、
$$
(\varphi_{j_1} \ldots \varphi_{j_n})(g) =
\varphi_{j_1}(g_1) \ldots \varphi_{j_n}(g_n)
$$
は $A$ における乗法が可換なので順序に依存しない。したがって $\Phi$ を、
$g \in I^{nc}$ を
$$
\Phi(g) = \sum\nolimits_{e_1 + \ldots + e_r = n}
(\varphi_1^{e_1} \circ \ldots \circ \varphi_r^{e_r})(g)
T_1^{e_1} \ldots T_r^{e_r}
$$
へ送ることによって定義できる。これが所望の性質をもつ次数付き $A$-代数準同型で
あることは直ちに証明できる。一意性も最後の性質も明らかである。これで補題が
証明された。
\end{proof}

\begin{lemma}
\label{algebraization-lemma-cd-one-extend-to-module}
$I = (f_1, \ldots, f_r)$ を Noether 環 $A$ のイデアルとし、
$\text{cd}(A, I) = 1$ とする。$c \geq 1$、$\varphi_j : I^c \to A$、
$j = 1, \ldots, r$ を補題 \ref{algebraization-lemma-cd-one} のものとする。
$A \to B$ を環準同型とし、$B$ は Noether 環、$N$ は有限 $B$-加群とする。
このとき、必要なら $c$ を大きくし、それに応じて $\varphi_j$ を調整すれば、
一意な一意な次数付き $B$-加群写像
$$
\Phi_N : \bigoplus\nolimits_{n \geq 0} I^{nc}N \to N[T_1, \ldots, T_r]
$$
であって、$\Phi_N(g x) = \Phi(g) x$ を満たすものが存在する。ここで
$g \in I^{nc}$ かつ $x \in N$ であり、$\Phi$ は補題
\ref{algebraization-lemma-cd-one-extend} のものである。$\Phi_N$ と写像
$N[T_1, \ldots, T_r] \to \bigoplus_{n \geq 0} I^nN$,
$T_j \mapsto f_j$ との合成は包含写像
$\bigoplus_{n \geq 0} I^{nc}N \to \bigoplus_{n \geq 0} I^nN$ である。
\end{lemma}

\begin{proof}
公式と補題 \ref{algebraization-lemma-cd-one-extend} における $\Phi$ の一意性から、一意性は
明らかである。Noether $A$-代数 $B' = B \oplus N$ を考える。ここで $N$ は平方零の
イデアルである。$\Phi_N$ の存在を示すには、補題 \ref{algebraization-lemma-cd-one} により、
$\varphi_j$ が写像 $\varphi'_j : I^cB' \to B'$ に拡張されることを示せば十分で
ある。その際、必要なら $c$ をある $c'$ まで大きくする（そして
$\varphi_j$ を、包含 $I^{c'} \to I^c$ と $\varphi_j$ の合成で置き換える）。
$\varphi_j$ が切断
$$
h_j \in \Gamma(\Spec(A) \setminus V(I), \mathcal{O}_{\Spec(A)})
$$
に対応することを思い出す。Cohomology of Schemes, 補題
\ref{coherent-lemma-homs-over-open} を参照されたい。（実際、補題
\ref{algebraization-lemma-cd-one} の証明ではこのように $\varphi_j$ を選んだ。）同じ補題を用いて、
引戻し
$$
h'_j \in \Gamma(\Spec(B') \setminus V(IB'), \mathcal{O}_{\Spec(B')})
$$
（$h_j$ の引戻し）を $B'$-線形写像 $\varphi'_j : I^{c'}B' \to B'$ で表す。
これはある $c' \geq c$ に対して成り立つ。さらに同じ補題を適用すると、
$\varphi_j$ との一致が十分大きな $c'$ に対して成り立つ。実際、一致は $I^{c'}$ の有限個の
生成元上で検査できる。細部を一つ省略する。
\end{proof}

\begin{lemma}
\label{algebraization-lemma-cd-is-one-for-system}
$I = (f_1, \ldots, f_r)$ を Noether 環 $A$ のイデアルとし、
$\text{cd}(A, I) = 1$ とする。$c \geq 1$、$\varphi_j : I^c \to A$、
$j = 1, \ldots, r$ を補題 \ref{algebraization-lemma-cd-one} のものとする。
$X$ を $\Spec(A)$ 上の Noether スキームとする。逆系
$$
\ldots \to \mathcal{F}_3 \to \mathcal{F}_2 \to \mathcal{F}_1
$$
の各項が連接 $\mathcal{O}_X$-加群であり、
$\mathcal{F}_n = \mathcal{F}_{n + 1}/I^n\mathcal{F}_{n + 1}$ を満たすとする。
$\mathcal{F} = \lim \mathcal{F}_n$ と置く。このとき、必要なら $c$ を大きくし、
それに応じて $\varphi_j$ を調整すれば、次数付き $\mathcal{O}_X$-加群写像
$$
\Phi_\mathcal{F} :
\bigoplus\nolimits_{n \geq 0} I^{nc}\mathcal{F}
\longrightarrow
\mathcal{F}[T_1, \ldots, T_r]
$$
であって、$\Phi_\mathcal{F}(g s) = \Phi(g) s$ を満たすものが一意に存在する。ここで
$g \in I^{nc}$、$s$ は $\mathcal{F}$ の局所切断であり、
$\Phi$ は補題 \ref{algebraization-lemma-cd-one-extend} のものである。$\Phi_\mathcal{F}$ と写像
$\mathcal{F}[T_1, \ldots, T_r] \to \bigoplus_{n \geq 0} I^n\mathcal{F}$,
$T_j \mapsto f_j$
との合成は標準的な包含写像
$\bigoplus_{n \geq 0} I^{nc}\mathcal{F} \to
\bigoplus_{n \geq 0} I^n\mathcal{F}$ である。
\end{lemma}

\begin{proof}
$\mathcal{O}_X$-線形性と、$\Phi_\mathcal{F}(g s) = \Phi(g) s$ となるという要請から、
一意性は明らかである。ここで $g \in I^{nc}$ であり、$s$ は $\mathcal{F}$ の局所切断で
ある。したがって $X = \Spec(B)$ はアフィンであると仮定してよい。
$(\mathcal{F}_n)$ は、Cohomology of Schemes, 第
\ref{coherent-section-coherent-formal} 節で導入された圏
$\textit{Coh}(X, I\mathcal{O}_X)$ の対象であることに注意する。
$B' = B^\wedge$ を $I$-進完備化とする。これは $B$ の完備化である。
Cohomology of Schemes, 補題
\ref{coherent-lemma-inverse-systems-affine} により、対象 $(\mathcal{F}_n)$ は有限
$B'$-加群 $N$ に対応する。その意味は、$\mathcal{F}_n$ が有限 $B$-加群
$N/I^n N$ に付随する連接加群だということである。補題
\ref{algebraization-lemma-cd-one-extend-to-module} を $I \subset A \to B'$ と $N$ に適用すると、
必要なら $c$ を大きくして $\varphi_j$ をそれに応じて調整することで、一意な写像
$$
\Phi_N : \bigoplus\nolimits_{n \geq 0} I^{nc}N \to N[T_1, \ldots, T_r]
$$
であって対応する性質をもつものが得られる。次数 $n$ において、写像
$N/I^mN \to \bigoplus_{e_1 + \ldots + e_r = n}
N/I^{m - nc}N \cdot T_1^{e_1} \ldots T_r^{e_r}$ の逆系が得られることに注意する。
ここで $m \geq nc$ である。
これを連接層に戻して翻訳すると、$\Phi_N$ は写像の系
$$
\Phi^n_m :
I^{nc}\mathcal{F}_m
\longrightarrow
\bigoplus\nolimits_{e_1 + \ldots + e_r = n}
\mathcal{F}_{m - nc} \cdot T_1^{e_1} \ldots T_r^{e_r}
$$
に対応することが分かる。ここで $m \geq nc$ と $n \geq 1$ は変化する。
$m$ についてこれらの写像の逆極限を取れば
$\Phi_\mathcal{F} = \bigoplus_n \lim_m \Phi^n_m$ を得る。例えばアフィン上で評価すれば
$\lim_m I^t\mathcal{F}_m = I^t \mathcal{F}$ と分かることに注意する。しかし実際には、
写像 $I^t\mathcal{F} \to \lim_m I^t\mathcal{F}_m$ があることは明らかなので、この
事実は必要ない。
\end{proof}

\begin{lemma}
\label{algebraization-lemma-topology-I-adic}
$I$ を Noether 環 $A$ のイデアルとし、$X$ を $\Spec(A)$ 上の Noether スキームと
する。逆系
$$
\ldots \to \mathcal{F}_3 \to \mathcal{F}_2 \to \mathcal{F}_1
$$
の各項が連接 $\mathcal{O}_X$-加群であり、
$\mathcal{F}_n = \mathcal{F}_{n + 1}/I^n\mathcal{F}_{n + 1}$ を満たすとする。
$\text{cd}(A, I) = 1$ ならば、すべての $p \in \mathbf{Z}$ に対し
$\lim H^p(X, \mathcal{F}_n)$ 上の逆極限位相は $I$-進位相である。
\end{lemma}

\begin{proof}
まず、$I^t \lim H^p(X, \mathcal{F}_n)$ が $H^p(X, \mathcal{F}_t)$ において零へ写る
ことは明らかである。したがって $I$-進位相は逆極限位相より細かい。逆向きを示す
ため、$\mathcal{F} = \lim \mathcal{F}_n$ と置き、生成元
$f_1, \ldots, f_r$ を $I$ について選び、$c \geq 1$ を選び、補題
\ref{algebraization-lemma-cd-is-one-for-system} のように $\Phi_\mathcal{F}$ を選ぶ。以後、補題
\ref{algebraization-lemma-properties-system} の結果を断りなく用いる。特に短完全系列
$$
0 \to R^1\lim H^{p - 1}(X, \mathcal{F}_n) \to H^p(X, \mathcal{F})
\to \lim H^p(X, \mathcal{F}_n) \to 0
$$
がある。したがって任意の元 $\xi$ で $\lim H^p(X, \mathcal{F}_n)$ に属するものを
元 $\xi' \in H^p(X, \mathcal{F})$ へ持ち上げられる。$\xi$ が
$H^p(X, \mathcal{F}_{nc})$ において零へ写ると、ある $n$ に対して仮定する。
言い換えると、$\xi$ が
逆極限位相について「小さい」と仮定する。短完全系列
$$
0 \to I^{nc}\mathcal{F} \to \mathcal{F} \to \mathcal{F}_{nc} \to 0
$$
があるので、仮定は $\xi'$ を元 $\xi'' \in H^p(X, I^{nc}\mathcal{F})$ へ持ち上げ
られることを意味する。$\Phi_\mathcal{F}$ を適用すると
$$
\Phi_\mathcal{F}(\xi'') = \sum\nolimits_{e_1 + \ldots + e_r = n}
\xi'_{e_1, \ldots, e_r} \cdot T_1^{e_1} \ldots T_r^{e_r}
$$
を得る。ここで、ある $\xi'_{e_1, \ldots, e_r} \in H^p(X, \mathcal{F})$ が存在する。
$\xi_{e_1, \ldots, e_r} \in \lim H^p(X, \mathcal{F}_n)$ をそれらの像とし、補題
\ref{algebraization-lemma-cd-is-one-for-system} の最後の主張を用いると、
$$
\xi = \sum f_1^{e_1} \ldots f_r^{e_r} \xi_{e_1, \ldots, e_r}
$$
は所望のとおり $I^n \lim H^p(X, \mathcal{F}_n)$ に属すると結論できる。
\end{proof}

\begin{example}
\label{algebraization-example-not-I-adic}
$k$ を体とし、$A = k[x, y][[s, t]]/(xs - yt)$ とする。
$I = (s, t)$ および $\mathfrak a = (x, y, s, t)$ と置く。
$X = \Spec(A) - V(\mathfrak a)$ および
$\mathcal{F}_n = \mathcal{O}_X/I^n\mathcal{O}_X$ と置く。有理関数
$$
g = \frac{t}{x} = \frac{s}{y}
$$
はある開近傍 $V \subset X$ 上で正則である。この開近傍は $V(I\mathcal{O}_X)$ の
近傍である。したがって各冪
$g^e$ は切断 $g^e \in M = \lim H^0(X, \mathcal{F}_n)$ を定める。逆極限位相を入れた
$g^e \to 0$ が $e \to \infty$ のとき、$M$ の逆極限位相において成り立つことに
注意する。実際、
$g^e$ は $\mathcal{F}_e$ において零へ写る。一方、$g^e \not \in IM$ は任意の $e$ に
対して成り立つ。これは $H^0(U, \mathcal{F}_n)$ を計算すれば分かるが、
計算は省略する。$\text{cd}(A, I) = 2$ であることに注意する。したがって補題
\ref{algebraization-lemma-topology-I-adic} の結果は鋭い。
\end{example}








\section{Mittag--Leffler 条件}
\label{algebraization-section-ML}

\noindent
局所 Noether 環の極大イデアルに関する局所コホモロジーを取ると、しばしば
Mittag--Leffler 条件が自動的に得られる。このことから、穴あきスペクトル上で
全射な遷移写像をもつ連接層の逆系の高次コホモロジー群についても同じことが
成り立つ。

\begin{lemma}
\label{algebraization-lemma-descending-chain}
$(A, \mathfrak m)$ を Noether 局所環とする。
\begin{enumerate}
\item $M$ を有限 $A$-加群とする。このとき $A$-加群
$H^i_\mathfrak m(M)$ は任意の $i$ に対して降鎖条件を満たす。
\item $U = \Spec(A) \setminus \{\mathfrak m\}$ を $A$ の穴あきスペクトルとする。
$\mathcal{F}$ を連接 $\mathcal{O}_U$-加群とする。このとき $A$-加群
$H^i(U, \mathcal{F})$ は $i > 0$ に対して降鎖条件を満たす。
\end{enumerate}
\end{lemma}

\begin{proof}
(1) を $M$ の台の次元に関する帰納法で証明する。$M = 0$ ならば主張は成り立つ。
したがって $M$ は零でないと仮定してよいし、そう仮定する。

\medskip\noindent
帰納法の基底段階を考える。$\dim(\text{Supp}(M)) = 0$ ならば、$M$ の台は
$\{\mathfrak m\}$ であり、$H^0_\mathfrak m(M) = M$ かつ
$H^i_\mathfrak m(M) = 0$ が $i > 0$ に対して成り立つ。これは局所コホモロジーの
構成から明らかである。Dualizing Complexes, 第
\ref{dualizing-section-local-cohomology} 節を参照されたい。$M$ は有限長なので
（Algebra, 補題 \ref{algebra-lemma-length-finite}）、降鎖条件を満たす。

\medskip\noindent
帰納段階を考える。$\dim(\text{Supp}(M)) > 0$ と仮定する。基底段階により、有限加群
$H^0_\mathfrak m(M) \subset M$ は降鎖条件を満たす。Dualizing Complexes, 補題
\ref{dualizing-lemma-divide-by-torsion} により、$M$ を $M/H^0_\mathfrak m(M)$ で
置き換えてよい。このとき $H^0_\mathfrak m(M) = 0$ である。すなわち $M$ の深さは
$\geq 1$ である。Dualizing Complexes, 補題 \ref{dualizing-lemma-depth} を参照されたい。
$x \in \mathfrak m$ を、$x : M \to M$ が単射となるように選ぶ。Algebra, 補題
\ref{algebra-lemma-one-equation-module} により
$\dim(\text{Supp}(M/xM)) = \dim(\text{Supp}(M)) - 1$ であり、帰納法の仮定を適用
できる。添字 $i$ を一つ選び、完全系列
$$
H^{i - 1}_\mathfrak m(M/xM) \to H^i_\mathfrak m(M) \xrightarrow{x}
H^i_\mathfrak m(M)
$$
を考える。これは短完全系列 $0 \to M \xrightarrow{x} M \to M/xM \to 0$ から来る。
したがって $x$-捩れ $H^i_\mathfrak m(M)[x]$ は降鎖条件を満たす加群の商であり、
それ自身も降鎖条件を満たす。ゆえに $\mathfrak m$-捩れ部分加群
$H^i_\mathfrak m(M)[\mathfrak m]$ は降鎖条件を満たす（したがって
$A/\mathfrak m$ 上有限次元である）。そこで Dualizing Complexes, 補題
\ref{dualizing-lemma-describe-categories} により、$\mathfrak m$-冪捩れ加群
$H^i_\mathfrak m(M)$ は降鎖条件を満たすと結論する。

\medskip\noindent
(2) は (1) と Local Cohomology, 補題
\ref{local-cohomology-lemma-finiteness-pushforwards-and-H1-local} から従う。
\end{proof}

\begin{lemma}
\label{algebraization-lemma-ML-local}
$(A, \mathfrak m)$ を Noether 局所環とする。
\begin{enumerate}
\item $(M_n)$ を有限 $A$-加群の逆系とする。このとき逆系
$H^i_\mathfrak m(M_n)$ は任意の $i$ に対して Mittag--Leffler 条件を満たす。
\item $U = \Spec(A) \setminus \{\mathfrak m\}$ を $A$ の穴あきスペクトルとする。
$\mathcal{F}_n$ を連接 $\mathcal{O}_U$-加群の逆系とする。このとき逆系
$H^i(U, \mathcal{F}_n)$ は $i > 0$ に対して Mittag--Leffler 条件を満たす。
\end{enumerate}
\end{lemma}

\begin{proof}
補題 \ref{algebraization-lemma-descending-chain} から直ちに従う。
\end{proof}

\begin{lemma}
\label{algebraization-lemma-terrific}
$(A, \mathfrak m)$ を Noether 局所環とする。$(M_n)$ を有限 $A$-加群の逆系とする。
$M \to \lim M_n$ を写像とし、ここで $M$ は有限 $A$-加群であるとする。またある
$i$ に対して写像
$H^i_\mathfrak m(M) \to \lim H^i_\mathfrak m(M_n)$
が同型であるとする。このとき逆系 $H^i_\mathfrak m(M_n)$ は、本質的に定数値
$H^i_\mathfrak m(M)$ をとる。
\end{lemma}

\begin{proof}
補題 \ref{algebraization-lemma-ML-local} により逆系 $H^i_\mathfrak m(M_n)$ は Mittag--Leffler 条件を
満たす。$E_n \subset H^i_\mathfrak m(M_n)$ を $H^i_\mathfrak m(M_{n'})$ の像とする。
ここで $n' \gg n$ である。このとき $(E_n)$ は全射な遷移写像をもつ逆系で
あり、$H^i_\mathfrak m(M) = \lim E_n$ である。$H^i_\mathfrak m(M)$ は補題
\ref{algebraization-lemma-descending-chain} により降鎖条件を満たすので、全射
$H^i_\mathfrak m(M) \to E_n$ の非自明な核は有限個しかない。したがって
$E_n \to E_{n - 1}$ はすべての $n \gg 0$ に対して同型である。これは所望の結論である。
\end{proof}

\begin{lemma}
\label{algebraization-lemma-local-cohomology-derived-completion}
$(A, \mathfrak m)$ を Noether 局所環とする。$I \subset A$ をイデアルとし、$M$ を
有限 $A$-加群とする。このとき
$$
H^i(R\Gamma_\mathfrak m(M)^\wedge) = \lim H^i_\mathfrak m(M/I^nM)
$$
がすべての $i$ に対して成り立つ。ここで $R\Gamma_\mathfrak m(M)^\wedge$ は導来
$I$-進完備化を表す。
\end{lemma}

\begin{proof}
Dualizing Complexes, 補題 \ref{dualizing-lemma-completion-local} と補題
\ref{algebraization-lemma-ML-local} を適用すると、$R^1\lim$ の項が消えることが分かる。
\end{proof}








\section{環付きサイト上の導来完備化}
\label{algebraization-section-derived-completion}

\noindent
初読の際には本節を読み飛ばすことを勧める．

\medskip\noindent
この内容の代数版は More on Algebra, 第
\ref{more-algebra-section-derived-completion} 節にある．
$\mathcal{O}$ をサイト $\mathcal{C}$ 上の環の層とする．
$f$ を $\mathcal{O}$ の大域切断とする．局所化の前層に付随する層を
$\mathcal{O}_f$ と記す．この前層は
$U \mapsto \mathcal{O}(U)_f$ で与えられる．

\begin{lemma}
\label{algebraization-lemma-map-twice-localize}
$(\mathcal{C}, \mathcal{O})$ を環付きサイトとし，$f$ を
$\mathcal{O}$ の大域切断とする．
\begin{enumerate}
\item $L, N \in D(\mathcal{O}_f)$ に対して
$R\SheafHom_\mathcal{O}(L, N) = R\SheafHom_{\mathcal{O}_f}(L, N)$.
特に二つの $\mathcal{O}_f$-加群構造は
$R\SheafHom_\mathcal{O}(L, N)$ 上で一致する．
\item $K \in D(\mathcal{O})$ および
$L \in D(\mathcal{O}_f)$ に対して
$$
R\SheafHom_\mathcal{O}(L, K) =
R\SheafHom_{\mathcal{O}_f}(L, R\SheafHom_\mathcal{O}(\mathcal{O}_f, K))
$$
が成り立つ．特に
$R\SheafHom_\mathcal{O}(\mathcal{O}_f,
R\SheafHom_\mathcal{O}(\mathcal{O}_f, K)) =
R\SheafHom_\mathcal{O}(\mathcal{O}_f, K)$.
\item $g$ が $\mathcal{O}$ のもう一つの大域切断ならば
$$
R\SheafHom_\mathcal{O}(\mathcal{O}_f, R\SheafHom_\mathcal{O}(\mathcal{O}_g, K))
= R\SheafHom_\mathcal{O}(\mathcal{O}_{gf}, K).
$$
\end{enumerate}
\end{lemma}

\begin{proof}
(1) の証明．$\mathcal{J}^\bullet$ を $\mathcal{O}_f$-加群の
K-入射的複体であって $N$ を表すものとする．Cohomology on Sites, 補題
\ref{sites-cohomology-lemma-K-injective-flat} により，
$\mathcal{J}^\bullet$ は $\mathcal{O}$-加群の K-入射的複体でもある．
$\mathcal{F}^\bullet$ を $\mathcal{O}_f$-加群の複体であって
$L$ を表すものとする．このとき
$$
R\SheafHom_\mathcal{O}(L, N) =
R\SheafHom_\mathcal{O}(\mathcal{F}^\bullet, \mathcal{J}^\bullet) =
R\SheafHom_{\mathcal{O}_f}(\mathcal{F}^\bullet, \mathcal{J}^\bullet)
$$
である．実際，$\mathcal{J}^\bullet$ は $\mathcal{O}$-加群としても
$\mathcal{O}_f$-加群としても K-入射的複体なので，これは
Modules on Sites, 補題 \ref{sites-modules-lemma-epimorphism-modules}
から従う．

\medskip\noindent
(2) の証明．$\mathcal{I}^\bullet$ を $\mathcal{O}$-加群の
K-入射的複体であって $K$ を表すものとする．このとき
$R\SheafHom_\mathcal{O}(\mathcal{O}_f, K)$ は
$\SheafHom_\mathcal{O}(\mathcal{O}_f, \mathcal{I}^\bullet)$ で表される．
Cohomology on Sites, 補題 \ref{sites-cohomology-lemma-hom-K-injective} および
\ref{sites-cohomology-lemma-K-injective-flat} により，これは
$\mathcal{O}_f$-加群としても $\mathcal{O}$-加群としても
K-入射的な複体である．$\mathcal{F}^\bullet$ を $\mathcal{O}_f$-加群の
複体であって $L$ を表すものとする．すると
$$
R\SheafHom_\mathcal{O}(L, K) =
R\SheafHom_\mathcal{O}(\mathcal{F}^\bullet, \mathcal{I}^\bullet) =
R\SheafHom_{\mathcal{O}_f}(\mathcal{F}^\bullet,
\SheafHom_\mathcal{O}(\mathcal{O}_f, \mathcal{I}^\bullet))
$$
となる．これは Modules on Sites, 補題
\ref{sites-modules-lemma-adjoint-hom-restrict} と，
$\SheafHom_\mathcal{O}(\mathcal{O}_f, \mathcal{I}^\bullet)$ が
$\mathcal{O}_f$-加群の K-入射的複体であることによる．

\medskip\noindent
(3) の証明．これは
$R\SheafHom_\mathcal{O}(\mathcal{O}_g, \mathcal{I}^\bullet)$
が $\mathcal{O}$-加群の複体として K-入射的であること，および
$\SheafHom_\mathcal{O}(\mathcal{O}_f,
\SheafHom_\mathcal{O}(\mathcal{O}_g, \mathcal{H})) = 
\SheafHom_\mathcal{O}(\mathcal{O}_{gf}, \mathcal{H})$
がすべての $\mathcal{O}$-加群の層 $\mathcal{H}$ に対して成り立つことから従う．
\end{proof}

\noindent
$K \in D(\mathcal{O})$ とする．$T(K, f)$ を（Derived Categories, 定義
\ref{derived-definition-derived-limit} の意味で）逆系
$$
\ldots \to K \xrightarrow{f} K \xrightarrow{f} K
$$
の $D(\mathcal{O})$ における導来極限と記す．

\begin{lemma}
\label{algebraization-lemma-hom-from-Af}
$(\mathcal{C}, \mathcal{O})$ を環付きサイトとし，$f$ を
$\mathcal{O}$ の大域切断とする．$K \in D(\mathcal{O})$ とする．
次は同値である．
\begin{enumerate}
\item $R\SheafHom_\mathcal{O}(\mathcal{O}_f, K) = 0$,
\item $R\SheafHom_\mathcal{O}(L, K) = 0$ が，すべての $L$，すなわち
$D(\mathcal{O}_f)$ の対象について成り立つ，
\item $T(K, f) = 0$.
\end{enumerate}
\end{lemma}

\begin{proof}
(2) が (1) を含意することは明らかである．(1) $\Rightarrow$ (2) は
補題 \ref{algebraization-lemma-map-twice-localize} から従う．
$\mathcal{O}$-加群 $\mathcal{O}_f$ の自由分解は
$$
0 \to \bigoplus\nolimits_{n \in \mathbf{N}} \mathcal{O} \to
\bigoplus\nolimits_{n \in \mathbf{N}} \mathcal{O}
\to \mathcal{O}_f \to 0
$$
で与えられる．ここで第一の写像は局所切断 $(x_0, x_1, \ldots)$ を
$(x_0, x_1 - fx_0, x_2 - fx_1, \ldots)$ に送り，第二の写像は
$(x_0, x_1, \ldots)$ を $x_0 + x_1/f + x_2/f^2 + \ldots$ に送る．
$\SheafHom_\mathcal{O}(-, \mathcal{I}^\bullet)$ を適用する．ここで
$\mathcal{I}^\bullet$ は $\mathcal{O}$-加群の K-入射的複体であって
$K$ を表すものとする．すると複体の短完全列
$$
0 \to \SheafHom_\mathcal{O}(\mathcal{O}_f, \mathcal{I}^\bullet) \to
\prod \mathcal{I}^\bullet \to \prod \mathcal{I}^\bullet \to 0
$$
を得る．これは $\mathcal{I}^n$ が入射的 $\mathcal{O}$-加群だからである．
これらの積は $D(\mathcal{O})$ における積である（Injectives, 補題
\ref{injectives-lemma-derived-products} を参照）．したがって対象
$T(K, f)$ は $R\SheafHom_\mathcal{O}(\mathcal{O}_f, K)$ の
$D(\mathcal{O})$ における代表である．ゆえに (1) と (3) は同値である．
\end{proof}

\begin{lemma}
\label{algebraization-lemma-ideal-of-elements-complete-wrt}
$(\mathcal{C}, \mathcal{O})$ を環付きサイトとし，$K \in D(\mathcal{O})$
とする．$U$ に集合 $\mathcal{I}(U)$ を対応させる規則を考える．ここでこれは，
$f \in \mathcal{O}(U)$ であって $T(K|_U, f) = 0$ を満たす切断全体である．この規則は
$\mathcal{O}$ のイデアル層をなす．
\end{lemma}

\begin{proof}
以下，補題 \ref{algebraization-lemma-hom-from-Af} の結果を断りなく用いる．
$f \in \mathcal{I}(U)$ かつ $g \in \mathcal{O}(U)$ ならば，
$\mathcal{O}_{U, gf}$ は $\mathcal{O}_{U, f}$-加群であるから
$R\SheafHom_\mathcal{O}(\mathcal{O}_{U, gf}, K|_U) = 0$，したがって
$gf \in \mathcal{I}(U)$ である．$f, g \in \mathcal{O}(U)$ とする．
このとき短完全列
$$
0 \to \mathcal{O}_{U, f + g} \to
\mathcal{O}_{U, f(f + g)} \oplus \mathcal{O}_{U, g(f + g)} \to
\mathcal{O}_{U, gf(f + g)} \to 0
$$
がある．実際，$f, g$ は $\mathcal{O}(U)_{f + g}$ において単位イデアルを生成する．
これは Algebra, 補題 \ref{algebra-lemma-standard-covering} と，
最後の矢が全射であるという容易な事実から従う．
$R\SheafHom_\mathcal{O}( - , K|_U)$ は三角圏の間の完全関手なので，
$R\SheafHom_{\mathcal{O}_U}(\mathcal{O}_{U, f(f + g)}, K|_U)$,
$R\SheafHom_{\mathcal{O}_U}(\mathcal{O}_{U, g(f + g)}, K|_U)$，および
$R\SheafHom_{\mathcal{O}_U}(\mathcal{O}_{U, gf(f + g)}, K|_U)$,
の消滅は
$R\SheafHom_{\mathcal{O}_U}(\mathcal{O}_{U, f + g}, K|_U)$.
の消滅を含意する．層条件の確認は省略する．
\end{proof}

\noindent
任意の環付きサイトについて次の定義を設けることができる．

\begin{definition}
\label{algebraization-definition-derived-complete}
$(\mathcal{C}, \mathcal{O})$ を環付きサイトとする．
$\mathcal{I} \subset \mathcal{O}$ をイデアル層とし，
$K \in D(\mathcal{O})$ とする．$K$ が
{\it $\mathcal{I}$ に関して導来完備}であるとは，任意の対象
$U$（$\mathcal{C}$ の対象）と $f \in \mathcal{I}(U)$ に対して，対象
$T(K|_U, f)$ が $D(\mathcal{O}_U)$ において零であることをいう．
\end{definition}

\noindent
導来完備な対象からなる充満部分圏
$D_{comp}(\mathcal{O}) = D_{comp}(\mathcal{O}, \mathcal{I}) \subset
D(\mathcal{O})$ は飽和三角部分圏であることは明らかである（
Derived Categories, 定義
\ref{derived-definition-triangulated-subcategory} および
\ref{derived-definition-saturated} を参照）．この部分圏は
$D(\mathcal{O})$ における積とホモトピー極限のもとで保たれる．
しかし一般には可算直和のもとでは保たれない．

\begin{lemma}
\label{algebraization-lemma-derived-complete-internal-hom}
$(\mathcal{C}, \mathcal{O})$ を環付きサイトとする．
$\mathcal{I} \subset \mathcal{O}$ をイデアル層とする．
$K \in D(\mathcal{O})$ かつ $L \in D_{comp}(\mathcal{O})$ ならば
$R\SheafHom_\mathcal{O}(K, L) \in D_{comp}(\mathcal{O})$.
\end{lemma}

\begin{proof}
$U$ を $\mathcal{C}$ の対象，$f \in \mathcal{I}(U)$ とする．次を思い出そう．
$$
\Hom_{D(\mathcal{O}_U)}(\mathcal{O}_{U, f}, R\SheafHom_\mathcal{O}(K, L)|_U)
=
\Hom_{D(\mathcal{O}_U)}(
K|_U \otimes_{\mathcal{O}_U}^\mathbf{L} \mathcal{O}_{U, f}, L|_U)
$$
これは Cohomology on Sites, 補題 \ref{sites-cohomology-lemma-internal-hom}
による．右辺は補題 \ref{algebraization-lemma-hom-from-Af} と内部 Hom と実際の Hom の関係
（Cohomology on Sites, 補題 \ref{sites-cohomology-lemma-section-RHom-over-U}）
により零である．同じ消滅はすべての $U'/U$ について成り立つ．したがって
$R\SheafHom_{\mathcal{O}_U}(\mathcal{O}_{U, f},
R\SheafHom_\mathcal{O}(K, L)|_U)$ という $D(\mathcal{O}_U)$ の対象の
第 $0$ コホモロジー層は消滅する（同所引用）．他のコホモロジー層についても同様である．
すなわち $R\SheafHom_{\mathcal{O}_U}(\mathcal{O}_{U, f},
R\SheafHom_\mathcal{O}(K, L)|_U)$ は $D(\mathcal{O}_U)$ において零である．
補題 \ref{algebraization-lemma-hom-from-Af} により結論を得る．
\end{proof}

\begin{lemma}
\label{algebraization-lemma-restriction-derived-complete}
$\mathcal{C}$ をサイトとする．$\mathcal{O} \to \mathcal{O}'$ を環の層の準同型とし，
$\mathcal{I} \subset \mathcal{O}$ をイデアル層とする．
$D_{comp}(\mathcal{O}, \mathcal{I})$ の制限関手
$D(\mathcal{O}') \to D(\mathcal{O})$ による逆像は
$D_{comp}(\mathcal{O}', \mathcal{I}\mathcal{O}')$.
\end{lemma}

\begin{proof}
補題 \ref{algebraization-lemma-ideal-of-elements-complete-wrt} により，
$K' \in D(\mathcal{O}')$ が
$D_{comp}(\mathcal{O}', \mathcal{I}\mathcal{O}')$
に属することと，$T(K'|_U, f)$ がすべての局所切断
$f \in \mathcal{I}(U)$ に対して零であることとは同値である．
$T(K'|_U, f)$ のコホモロジー層はアーベル層の圏で計算されるので，
$f$ を $\mathcal{O}$ の切断とみなしても，その像を
$f$ の $\mathcal{O}'$ における切断として用いても違いはない．
これと導来完備対象の定義から補題は直ちに従う．
\end{proof}

\begin{lemma}
\label{algebraization-lemma-pushforward-derived-complete}
$f : (\Sh(\mathcal{D}), \mathcal{O}') \to (\Sh(\mathcal{C}), \mathcal{O})$
を環付きトポスの射とする．$\mathcal{I} \subset \mathcal{O}$ および
$\mathcal{I}' \subset \mathcal{O}'$ を，$f^\sharp$ が
$f^{-1}\mathcal{I}$ を $\mathcal{I}'$ に送るようなイデアル層とする．
このとき $Rf_*$ は $D_{comp}(\mathcal{O}', \mathcal{I}')$ を
$D_{comp}(\mathcal{O}, \mathcal{I})$ に送る．
\end{lemma}

\begin{proof}
$f$ は連続関手 $\mathcal{C} \to \mathcal{D}$ に対応する環付きサイトの射によって
与えられると仮定してよい（Modules on Sites, 補題
\ref{sites-modules-lemma-morphism-ringed-topoi-comes-from-morphism-ringed-sites}
）．$U$ を $\mathcal{C}$ の対象とし，$g$ を
$\mathcal{I}$ の $U$ 上の切断とする．
$\Hom_{D(\mathcal{O}_U)}(\mathcal{O}_{U, g}, Rf_*K|_U) = 0$
を，$K$ が $\mathcal{I}'$ に関して導来完備であるときに示せばよい．
実際，Cohomology on Sites, 補題
\ref{sites-cohomology-lemma-section-RHom-over-U}
により，これを $U$ 上のすべての対象と $K$ のすべてのシフトに適用すれば，
$R\SheafHom_{\mathcal{O}_U}(\mathcal{O}_{U, g}, Rf_*K|_U)$ が零となり，さらに
$T(Rf_*K|_U, g)$ が零となる（補題 \ref{algebraization-lemma-hom-from-Af}）．
これは示すべきことにほかならない（定義 \ref{algebraization-definition-derived-complete}）．
$V$ を $\mathcal{D}$ における $U$ の像とする．このとき
$$
\Hom_{D(\mathcal{O}_U)}(\mathcal{O}_{U, g}, Rf_*K|_U) =
\Hom_{D(\mathcal{O}'_V)}(\mathcal{O}'_{V, g'}, K|_V) = 0
$$
である．ここで $g' = f^\sharp(g) \in \mathcal{I}'(V)$ である．
第二の等号は $K$ が導来完備であることから，第一の等号は
$\mathcal{O}_{U, g}$ の導来逆像が $\mathcal{O}'_{V, g'}$ であることと
Cohomology on Sites, 補題 \ref{sites-cohomology-lemma-adjoint} から従う．
\end{proof}

\noindent
次の補題は導来完備化が存在する最も簡単な場合である．

\begin{lemma}
\label{algebraization-lemma-derived-completion}
$(\mathcal{C}, \mathcal{O})$ を環付きサイトとする．$f_1, \ldots, f_r$ を
$\mathcal{O}$ の大域切断とする．$\mathcal{I} \subset \mathcal{O}$ を
$f_1, \ldots, f_r$ で生成されるイデアル層とする．
このとき包含関手 $D_{comp}(\mathcal{O}) \to D(\mathcal{O})$ は左随伴をもつ．
すなわち，任意の対象 $K$（$D(\mathcal{O})$ の対象）に対して写像
$K \to K^\wedge$ が存在し，$K^\wedge$ は $D_{comp}(\mathcal{O})$ に属し，写像
$$
\Hom_{D(\mathcal{O})}(K^\wedge, E) \longrightarrow \Hom_{D(\mathcal{O})}(K, E)
$$
は $E$ が $D_{comp}(\mathcal{O})$ に属するとき常に全単射である．実際，
$$
K^\wedge =
R\SheafHom_\mathcal{O}
(\mathcal{O} \to \prod\nolimits_{i_0} \mathcal{O}_{f_{i_0}} \to
\prod\nolimits_{i_0 < i_1} \mathcal{O}_{f_{i_0}f_{i_1}} \to
\ldots \to \mathcal{O}_{f_1\ldots f_r}, K)
$$
が $K$ に関して関手的に成り立つ．
\end{lemma}

\begin{proof}
補題の最後の表示式によって $K^\wedge$ を定義する．複体の写像
$$
(\mathcal{O} \to \prod\nolimits_{i_0} \mathcal{O}_{f_{i_0}} \to
\prod\nolimits_{i_0 < i_1} \mathcal{O}_{f_{i_0}f_{i_1}} \to
\ldots \to \mathcal{O}_{f_1\ldots f_r}) \longrightarrow \mathcal{O}
$$
があり，これは写像 $K \to K^\wedge$ を誘導する．$K^\wedge$ が導来完備であり，
$K \to K^\wedge$ が，$K$ が導来完備ならば同型であることを示せば十分である．

\medskip\noindent
$f$ を $\mathcal{O}$ の大域切断とする．補題
\ref{algebraization-lemma-map-twice-localize} により，対象
$R\SheafHom_\mathcal{O}(\mathcal{O}_f, K^\wedge)$
は
$$
R\SheafHom_\mathcal{O}(
(\mathcal{O}_f \to \prod\nolimits_{i_0} \mathcal{O}_{ff_{i_0}} \to
\prod\nolimits_{i_0 < i_1} \mathcal{O}_{ff_{i_0}f_{i_1}} \to
\ldots \to \mathcal{O}_{ff_1\ldots f_r}), K)
$$
に等しい．$f = f_i$ となる $i$ が存在するならば，$f_1, \ldots, f_r$ は
$\mathcal{O}_f$ において単位イデアルを生成する．したがって拡張交代
{\v C}ech 複体
$$
\mathcal{O}_f \to \prod\nolimits_{i_0} \mathcal{O}_{ff_{i_0}} \to
\prod\nolimits_{i_0 < i_1} \mathcal{O}_{ff_{i_0}f_{i_1}} \to
\ldots \to \mathcal{O}_{ff_1\ldots f_r}
$$
は零である（実際，零にホモトピックである）．これにより $K^\wedge$ が
導来完備であることが分かる．

\medskip\noindent
$K$ が導来完備ならば，$R\SheafHom_\mathcal{O}(\mathcal{O}_f, K)$ は
すべての $f = f_{i_0} \ldots f_{i_p}$，$p \geq 0$ に対して零である．したがって
$K \to K^\wedge$ は $D(\mathcal{O})$ における同型である．
\end{proof}

\noindent
次に，導来完備化が実際に完備化である理由を説明する．

\begin{lemma}
\label{algebraization-lemma-derived-completion-koszul}
$(\mathcal{C}, \mathcal{O})$ を環付きサイトとする．$f_1, \ldots, f_r$ を
$\mathcal{O}$ の大域切断とする．$\mathcal{I} \subset \mathcal{O}$ を
$f_1, \ldots, f_r$ で生成されるイデアル層とし，$K \in D(\mathcal{O})$ とする．
補題 \ref{algebraization-lemma-derived-completion} の導来完備化 $K^\wedge$ は式
$$
K^\wedge = R\lim K \otimes^\mathbf{L}_\mathcal{O} K_n
$$
で与えられる．ここで $K_n = K(\mathcal{O}, f_1^n, \ldots, f_r^n)$ は
$f_1^n, \ldots, f_r^n$ に関する $\mathcal{O}$ 上の Koszul 複体である．
\end{lemma}

\begin{proof}
More on Algebra, 補題
\ref{more-algebra-lemma-extended-alternating-Cech-is-colimit-koszul}
で見たように，拡張交代 {\v C}ech 複体
$$
\mathcal{O} \to \prod\nolimits_{i_0} \mathcal{O}_{f_{i_0}} \to
\prod\nolimits_{i_0 < i_1} \mathcal{O}_{f_{i_0}f_{i_1}} \to
\ldots \to \mathcal{O}_{f_1\ldots f_r}
$$
は，Koszul 複体 $K^n = K(\mathcal{O}, f_1^n, \ldots, f_r^n)$ の余極限であり，
これらは次数 $0, \ldots, r$ に置かれる．$K^n$ は
有限自由 $\mathcal{O}$-加群からなる有限鎖複体であり，その双対は
$\SheafHom_\mathcal{O}(K^n, \mathcal{O}) = K_n$ であることに注意する．ここで
$K_n$ は通常どおり次数 $-r, \ldots, 0$ に置かれた Koszul 余鎖複体である．
補題 \ref{algebraization-lemma-derived-completion} により，関手 $E \mapsto E^\wedge$ は
拡張交代 {\v C}ech 複体からの $R\SheafHom$ を取り，その第二変数を
$E$ とすることで得られる：
$$
E^\wedge = R\SheafHom(\colim K^n, E)
$$
これは $R\lim (E \otimes_\mathcal{O}^\mathbf{L} K_n)$ に等しい．
Cohomology on Sites, 補題
\ref{sites-cohomology-lemma-colim-and-lim-of-duals} を参照されたい．
\end{proof}

\begin{lemma}
\label{algebraization-lemma-all-rings}
次のものを構成する方法が存在する：
\begin{enumerate}
\item 各対 $(A, I)$，すなわち環 $A$ と有限生成イデアル $I \subset A$ からなる対に対する，
複体 $K(A, I)$；これは $A$-加群の複体である；
\item 写像 $K(A, I) \to A$；これは $A$-加群の複体の写像である；
\item 各環準同型 $A \to B$ と有限生成イデアル $I \subset A$ に対する，
複体の写像 $K(A, I) \to K(B, IB)$．
\end{enumerate}
これらは次を満たす：
\begin{enumerate}
\item[(a)] $A \to B$ と有限生成な $I \subset A$ に対して図式
$$
\xymatrix{
K(A, I) \ar[r] \ar[d] & A \ar[d] \\
K(B, IB) \ar[r] & B
}
$$
は可換である；
\item[(b)] $A \to B \to C$ と有限生成な $I \subset A$ に対して，写像の合成
$K(A, I) \to K(B, IB) \to K(C, IC)$ は写像 $K(A, I) \to K(C, IC)$ である；
\item[(c)] $A \to B$ と有限生成イデアル $I \subset A$ に対して，誘導写像
$K(A, I) \otimes_A^\mathbf{L} B \to K(B, IB)$ は $D(B)$ における同型である；
\item[(d)] $I = (f_1, \ldots, f_r) \subset A$ ならば，可換図式
$$
\xymatrix{
(A \to \prod\nolimits_{i_0} A_{f_{i_0}} \to
\prod\nolimits_{i_0 < i_1} A_{f_{i_0}f_{i_1}} \to
\ldots \to A_{f_1\ldots f_r}) \ar[r] \ar[d] &  K(A, I) \ar[d] \\
A \ar[r]^1 & A
}
$$
が $D(A)$ に存在し，その水平矢印は同型である．
\end{enumerate}
\end{lemma}

\begin{proof}
 $S$ を，次の形の環 $A_0$ 全体の集合とする：
$A_0 = \mathbf{Z}[x_1, \ldots, x_n]/J$.
任意の有限型 $\mathbf{Z}$-代数は $S$ のある元と同型である．
$\mathcal{A}_0$ を次の圏とする．その対象は対 $(A_0, I_0)$ であり，ここで
$A_0 \in S$ かつ $I_0 \subset A_0$ はイデアルである．射
$(A_0, I_0) \to (B_0, J_0)$ は環準同型 $\varphi : A_0 \to B_0$ であり，
$J_0 = \varphi(I_0)B_0$ を満たす．

\medskip\noindent
$K(A_0, I_0) \to A_0$ を $\mathcal{A}_0$ の対象に対して関手的に構成でき，
性質 (a)，(b)，(c)，(d) を満たすと仮定する．このとき
$$
K(A, I) = \colim_{\varphi : (A_0, I_0) \to (A, I)} K(A_0, I_0)
$$
と置く．ここで余極限は，環準同型 $\varphi : A_0 \to A$ であって
$\varphi(I_0)A = I$ を満たすもの全体にわたり，$(A_0, I_0)$ は
$\mathcal{A}_0$ に属する．
$(A_0, I_0) \to (A, I)$ と $(A_0', I_0') \to (A, I)$ の間の射は，
射 $(A_0, I_0) \to (A_0', I_0')$，すなわち $\mathcal{A}_0$ の射であって
$A$ への写像と可換するもので与えられる．これらの
$(A_0, I_0) \to (A, I)$ の圏はフィルター付きである（詳細は省略する）．さらに
$\colim_{\varphi : (A_0, I_0) \to (A, I)} A_0 = A$ なので，
$K(A, I)$ は $A$-加群の複体である．最後に，$\varphi : A \to B$ と
$I \subset A$ が与えられたとする．余極限に現れるすべての
$(A_0, I_0) \to (A, I)$ に対して，合成
$(A_0, I_0) \to (B, IB)$ は $(B, IB)$ に対する余極限に現れる．
このようにして余極限上の写像を得る．性質 (a)，(b)，(c)，(d) は，これと複体
$K(A_0, I_0)$ の対応する性質から直ちに従う．

\medskip\noindent
$\mathcal{C}_0 = \mathcal{A}_0^{opp}$ にカオス位相を入れる．
$\mathcal{C}_0$ に環の層
$\mathcal{O} : (A, I) \mapsto A$ を入れる．イデアル $I$ は貼り合わさって
イデアル層 $\mathcal{I} \subset \mathcal{O}$ を与える．入射分解
$\mathcal{O} \to \mathcal{J}^\bullet$ を選び，対象
$$
\mathcal{F}^\bullet = \bigcup\nolimits_n \mathcal{J}^\bullet[\mathcal{I}^n]
$$
を考える．$U = (A, I) \in \Ob(\mathcal{C}_0)$ とする．
$\mathcal{C}_0$ の位相はカオス位相なので，値
$\mathcal{J}^\bullet(U)$ は $A$ の入射的 $A$-加群による分解である．したがって値
$\mathcal{F}^\bullet(U)$ は $D(A)$ の対象であり，$R\Gamma_I(A)$ の
$D(A)$ における像を表す．Dualizing Complexes, 節
\ref{dualizing-section-local-cohomology} を参照されたい．
$\mathcal{O}$-加群の複体 $\mathcal{K}^\bullet$ と可換図式
$$
\xymatrix{
\mathcal{O} \ar[r] & \mathcal{J}^\bullet \\
\mathcal{K}^\bullet \ar[r] \ar[u] & \mathcal{F}^\bullet \ar[u]
}
$$
を，水平矢印が擬同型となるように選ぶ．これは導来圏
$D(\mathcal{O})$ の構成によって可能である．$K(A, I) = \mathcal{K}^\bullet(U)$ と置く．
ここで $U = (A, I)$ である．性質 (a)，(b) は明らかであり，性質 (c)，(d) は
Dualizing Complexes, 補題
\ref{dualizing-lemma-compute-local-cohomology-noetherian} および
\ref{dualizing-lemma-local-cohomology-change-rings} から従う．
\end{proof}

\begin{lemma}
\label{algebraization-lemma-global-extended-cech-complex}
$(\mathcal{C}, \mathcal{O})$ を環付きサイトとし，
$\mathcal{I} \subset \mathcal{O}$ を有限型のイデアル層とする．
写像 $K \to \mathcal{O}$ が $D(\mathcal{O})$ に存在し，次を満たす．
$U \in \Ob(\mathcal{C})$ が，$\mathcal{I}|_U$ が
$f_1, \ldots, f_r \in \mathcal{I}(U)$ で生成される任意の対象ならば，同型
$$
(\mathcal{O}_U \to \prod\nolimits_{i_0} \mathcal{O}_{U, f_{i_0}} \to
\prod\nolimits_{i_0 < i_1} \mathcal{O}_{U, f_{i_0}f_{i_1}} \to
\ldots \to \mathcal{O}_{U, f_1\ldots f_r}) \longrightarrow K|_U
$$
があり，これは $\mathcal{O}_U$ への写像と両立する．
\end{lemma}

\begin{proof}
$\mathcal{C}' \subset \mathcal{C}$ を，対象 $U$ であって
$\mathcal{I}|_U$ が有限個の切断で生成されるものの充満部分圏とする．このとき
$\mathcal{C}' \to \mathcal{C}$ は特別な余連続関手である
（Sites, 定義 \ref{sites-definition-special-cocontinuous-functor}）．したがって
$\mathcal{C}'$ で考えれば十分である．Sites, 補題
\ref{sites-lemma-equivalence} を参照されたい．言い換えると，任意の対象
$U$（$\mathcal{C}$ の対象）に対して有限生成イデアル
$I \subset \mathcal{I}(U)$ が存在し，
$\mathcal{I}|_U = \Im(I \otimes \mathcal{O}_U \to \mathcal{O}_U)$
となると仮定してよい．このとき $I$ が $\mathcal{I}|_U$ を生成するという．
注意：$\mathcal{I}(U)$ が $\mathcal{O}(U)$ の有限生成イデアルであるとは限らない．

\medskip\noindent
$U$ を対象とし，$I \subset \mathcal{O}(U)$ を有限生成イデアルであって
$\mathcal{I}|_U$ を生成するものとする．圏 $\mathcal{C}/U$ 上で前層の複体
$$
K_{U, I}^\bullet : U'/U \longmapsto K(\mathcal{O}(U'), I\mathcal{O}(U'))
$$
を考える．ここで $K(-, -)$ は補題 \ref{algebraization-lemma-all-rings} のものである．
この層化は $I$ の選択によらないと主張する．実際，
$I' \subset \mathcal{O}(U)$ が有限生成イデアルであり，同じく
$\mathcal{I}|_U$ を生成するならば，被覆 $\{U_j \to U\}$ が存在して
$I\mathcal{O}(U_j) = I'\mathcal{O}(U_j)$ となる．（ヒント：これは $I$ と
$I'$ がともに有限生成であり，$\mathcal{I}|_U$ を生成することによる．）
したがって $K_{U, I}^\bullet$ と $K_{U, I'}^\bullet$ は，いずれかの
$U_j$ 上にある任意の対象について{\it 同一}である．ゆえに層化についての主張が従う．
共通の値を $K_U^\bullet$ と記す．

\medskip\noindent
選択が $I$ によらないことから，さらに
$K_U^\bullet|_{\mathcal{C}/U'} = K_{U'}^\bullet$
も分かる．ここで射 $U' \to U$，したがって局所化射
$\mathcal{C}/U' \to \mathcal{C}/U$ が与えられている．ゆえに複体
$K_U^\bullet$ は貼り合わさって，一意に定まる一つの複体
$K^\bullet$（$\mathcal{O}$-加群の複体）を与える．写像
$K^\bullet \to \mathcal{O}$ の存在と補題の擬同型は，複体
$K(-, -)$ の対応する性質から直ちに従う（補題 \ref{algebraization-lemma-all-rings}）．
\end{proof}

\begin{proposition}
\label{algebraization-proposition-derived-completion}
$(\mathcal{C}, \mathcal{O})$ を環付きサイトとする．
$\mathcal{I} \subset \mathcal{O}$ を有限型のイデアル層とする．包含関手
$D_{comp}(\mathcal{O}) \to D(\mathcal{O})$ には左随伴が存在する．
\end{proposition}

\begin{proof}
$K \to \mathcal{O}$ を $D(\mathcal{O})$ における，補題
\ref{algebraization-lemma-global-extended-cech-complex} で構成した写像とする．
$E \in D(\mathcal{O})$ とする．このとき $E^\wedge = R\SheafHom(K, E)$ と
写像 $E \to E^\wedge$ が求めるものになる．実際，サイト
$\mathcal{C}$ 上で局所的に，補題 \ref{algebraization-lemma-derived-completion} の随伴を回収する．
これにより $E^\wedge$ は常に導来完備であり，$E \to E^\wedge$ は
$E$ が導来完備ならば同型であることが分かる．
\end{proof}

\begin{remark}[完備化との比較]
\label{algebraization-remark-compare-with-completion}
$(\mathcal{C}, \mathcal{O})$ を環付きサイトとする．
$\mathcal{I} \subset \mathcal{O}$ を有限型のイデアル層とする．
$K \mapsto K^\wedge$ を命題 \ref{algebraization-proposition-derived-completion} の
導来完備化関手とする．任意の $n \geq 1$ に対して対象
$K \otimes_\mathcal{O}^\mathbf{L} \mathcal{O}/\mathcal{I}^n$
は導来完備である．実際，$\mathcal{I}$ の局所切断の冪で消される．したがって，
標準写像の標準的な分解
$$
K \to K^\wedge \to K \otimes_\mathcal{O}^\mathbf{L} \mathcal{O}/\mathcal{I}^n
$$
が存在する．ここで元の標準写像は
$K \to K \otimes_\mathcal{O}^\mathbf{L} \mathcal{O}/\mathcal{I}^n$.
である．これらの写像は $n$ の変化と両立し，比較写像
$$
K^\wedge
\longrightarrow
R\lim \left(K \otimes_\mathcal{O}^\mathbf{L} \mathcal{O}/\mathcal{I}^n\right)
$$
を得る．右辺は一種の完備化としてより見慣れた形である．
一般には，この比較写像は同型ではない．
\end{remark}

\begin{remark}[局所化と導来完備化]
\label{algebraization-remark-localization-and-completion}
$(\mathcal{C}, \mathcal{O})$ を環付きサイトとする．
$\mathcal{I} \subset \mathcal{O}$ を有限型のイデアル層とする．
$K \mapsto K^\wedge$ を命題 \ref{algebraization-proposition-derived-completion} の
導来完備化関手とする．命題の証明中の構成から，$K^\wedge|_U$ は
$K|_U$ の導来完備化である．これは任意の $U \in \Ob(\mathcal{C})$ に対して成り立つ．
これは次のように証明することもできる．導来完備対象の定義から
$K^\wedge|_U$ は導来完備である．したがって標準写像
$a : (K|_U)^\wedge \to K^\wedge|_U$ を得る．一方，$E$ が
$D(\mathcal{O}_U)$ の導来完備対象ならば，$Rj_*E$ は
$D(\mathcal{O})$ の導来完備対象である（補題
\ref{algebraization-lemma-pushforward-derived-complete}）．ここで $j$ は局所化射である
（Modules on Sites, 節 \ref{sites-modules-section-localize}）．したがって標準写像
$b : K^\wedge \to Rj_*((K|_U)^\wedge)$ も得る．$b$ の随伴が $a$ の逆写像で
あることの（形式的な）確認は省略する．
\end{remark}

\begin{remark}[完備化テンソル積]
\label{algebraization-remark-completed-tensor-product}
$(\mathcal{C}, \mathcal{O})$ を環付きサイトとし，
$\mathcal{I} \subset \mathcal{O}$ を有限型のイデアル層とする．
$K \mapsto K^\wedge$ を命題 \ref{algebraization-proposition-derived-completion} の
随伴と記す．このとき
$$
K \otimes^\wedge_\mathcal{O} L = (K \otimes_\mathcal{O}^\mathbf{L} L)^\wedge
$$
と置く．この{\it 完備化テンソル積}は関手
$D_{comp}(\mathcal{O}) \times D_{comp}(\mathcal{O}) \to D_{comp}(\mathcal{O})$
を定め，次が成り立つ：
$$
\Hom_{D_{comp}(\mathcal{O})}(K, R\SheafHom_\mathcal{O}(L, M))
=
\Hom_{D_{comp}(\mathcal{O})}(K \otimes_\mathcal{O}^\wedge L, M)
$$
これは $K, L, M \in D_{comp}(\mathcal{O})$ に対して成り立つ．なお，
$R\SheafHom_\mathcal{O}(L, M) \in D_{comp}(\mathcal{O})$ であることは
補題 \ref{algebraization-lemma-derived-complete-internal-hom} による．
\end{remark}

\begin{lemma}
\label{algebraization-lemma-map-identifies-koszul-and-cech-complexes}
$\mathcal{C}$ をサイトとする．$\varphi : \mathcal{O} \to \mathcal{O}'$ を
環の層の平坦な準同型と仮定する．$f_1, \ldots, f_r$ を
$\mathcal{O}$ の大域切断であって
$\mathcal{O}/(f_1, \ldots, f_r) \cong
\mathcal{O}'/(f_1, \ldots, f_r)\mathcal{O}'$ を満たすものとする．
このとき拡張交代 {\v C}ech 複体の写像
$$
\xymatrix{
\mathcal{O} \to
\prod_{i_0} \mathcal{O}_{f_{i_0}} \to
\prod_{i_0 < i_1} \mathcal{O}_{f_{i_0}f_{i_1}} \to \ldots \to
\mathcal{O}_{f_1\ldots f_r} \ar[d] \\
\mathcal{O}' \to
\prod_{i_0} \mathcal{O}'_{f_{i_0}} \to
\prod_{i_0 < i_1} \mathcal{O}'_{f_{i_0}f_{i_1}} \to \ldots \to
\mathcal{O}'_{f_1\ldots f_r}
}
$$
は擬同型である．
\end{lemma}

\begin{proof}
第二の複体は第一の複体と $\mathcal{O}'$ とのテンソル積である．第一の
拡張交代 {\v C}ech 複体は Koszul 複体
$K_n = K(\mathcal{O}, f_1^n, \ldots, f_r^n)$ の余極限として書ける．
More on Algebra, 補題
\ref{more-algebra-lemma-extended-alternating-Cech-is-colimit-koszul}.
したがって $K_n \to K_n \otimes_\mathcal{O} \mathcal{O}'$ が擬同型であることを
示せば十分である．$\mathcal{O} \to \mathcal{O}'$ は平坦なので，
$H^i \to H^i \otimes_\mathcal{O} \mathcal{O}'$ が同型であることを示せば十分である．
ここで $H^i$ は第 $i$ コホモロジー層 $H^i = H^i(K_n)$ である．これらの層は
$f_1^n, \ldots, f_r^n$ で消される．More on Algebra, 補題
\ref{more-algebra-lemma-homotopy-koszul} を参照されたい．したがって，これらの層は
$(f_1, \ldots, f_r)^m$ で消される．ここで $m \gg 0$ を適当に取る．ゆえに
$H^i \to H^i \otimes_\mathcal{O} \mathcal{O}'$ は同型である．
Modules on Sites, 補題 \ref{sites-modules-lemma-neighbourhood-isomorphism} を参照されたい．
\end{proof}

\begin{lemma}
\label{algebraization-lemma-restriction-derived-complete-equivalence}
$\mathcal{C}$ をサイトとする．$\mathcal{O} \to \mathcal{O}'$ を
環の層の準同型とする．$\mathcal{I} \subset \mathcal{O}$ を
有限型のイデアル層とする．$\mathcal{O} \to \mathcal{O}'$ が平坦であり，
$\mathcal{O}/\mathcal{I} \cong \mathcal{O}'/\mathcal{I}\mathcal{O}'$
ならば，制限関手 $D(\mathcal{O}') \to D(\mathcal{O})$ は圏同値
$D_{comp}(\mathcal{O}', \mathcal{I}\mathcal{O}') \to
D_{comp}(\mathcal{O}, \mathcal{I})$.
\end{lemma}

\begin{proof}
補題 \ref{algebraization-lemma-pushforward-derived-complete} により，制限
$r : D(\mathcal{O}') \to D(\mathcal{O})$ は
$D_{comp}(\mathcal{O}', \mathcal{I}\mathcal{O}')$ を
$D_{comp}(\mathcal{O}, \mathcal{I})$ に送る．擬逆 $E \mapsto E'$ を構成する．

\medskip\noindent
$K \to \mathcal{O}$ を $D(\mathcal{O})$ の射であって，補題
\ref{algebraization-lemma-global-extended-cech-complex} で構成したものとする．
$K' = K \otimes_\mathcal{O}^\mathbf{L} \mathcal{O}'$ と
$D(\mathcal{O}')$ において置く．このとき $K' \to \mathcal{O}'$ は
$D(\mathcal{O}')$ の写像であり，補題
\ref{algebraization-lemma-global-extended-cech-complex} の結論を
$\mathcal{I}' = \mathcal{I}\mathcal{O}'$ に関して満たす．写像
$K \to r(K')$ は補題
\ref{algebraization-lemma-map-identifies-koszul-and-cech-complexes} により擬同型である．
そこで $E \in D_{comp}(\mathcal{O}, \mathcal{I})$ に対して
$$
E' = R\SheafHom_\mathcal{O}(r(K'), E)
$$
と置く．これは $D(\mathcal{O}')$ の対象とみなす．そのために
$\mathcal{O}'$-加群構造（$K'$ 上の構造）を用いる．$E$ は導来完備なので
$E = R\SheafHom_\mathcal{O}(K, E)$ である．命題
\ref{algebraization-proposition-derived-completion} の証明を参照されたい．一方，
$K \to r(K')$ は同型であるから，同型
$E \to r(E')$ が $D(\mathcal{O})$ に存在する．証明を完了するには，
$E = r(M')$ であり，$M'$ が
$D_{comp}(\mathcal{O}', \mathcal{I}')$ の対象ならば
$E' \cong M'$ となることを示さなければならない．写像を得るために
$$
M' = R\SheafHom_{\mathcal{O}'}(\mathcal{O}', M') \to
R\SheafHom_\mathcal{O}(r(\mathcal{O}'), r(M')) \to
R\SheafHom_\mathcal{O}(r(K'), r(M')) = E'
$$
を用いる．第二の矢印では写像 $K' \to \mathcal{O}'$ を用いた．これが同型で
あることを見るには，この矢印に $r$ を適用したものが，上の同型
$E \to r(E')$ と同じであることを示せばよい．詳細は省略する．
\end{proof}

\begin{lemma}
\label{algebraization-lemma-pushforward-derived-complete-adjoint}
$f : (\Sh(\mathcal{D}), \mathcal{O}') \to (\Sh(\mathcal{C}), \mathcal{O})$
を環付きトポスの射とする．$\mathcal{I} \subset \mathcal{O}$ および
$\mathcal{I}' \subset \mathcal{O}'$ を有限型のイデアル層であって，
$f^\sharp$ が $f^{-1}\mathcal{I}$ を $\mathcal{I}'$ に送るものとする．
このとき $Rf_*$ は $D_{comp}(\mathcal{O}', \mathcal{I}')$ を
$D_{comp}(\mathcal{O}, \mathcal{I})$ に送り，左随伴 $Lf_{comp}^*$ をもつ．
これは $Lf^*$ の後に導来完備化を施すものである．
\end{lemma}

\begin{proof}
第一の主張は補題 \ref{algebraization-lemma-pushforward-derived-complete} で見た．
第二の主張が意味をもつことにも注意する．実際，導来完備化関手
$D(\mathcal{O}') \to D_{comp}(\mathcal{O}', \mathcal{I}')$ が存在する
（命題 \ref{algebraization-proposition-derived-completion}）．
そこで $K \in D_{comp}(\mathcal{O}, \mathcal{I})$ および
$M \in D_{comp}(\mathcal{O}', \mathcal{I}')$ とする．このとき
$$
\Hom(K, Rf_*M) = \Hom(Lf^*K, M) = \Hom(Lf_{comp}^*K, M)
$$
が導来完備化の普遍性により成り立つ．
\end{proof}

\begin{lemma}
\label{algebraization-lemma-pushforward-commutes-with-derived-completion}
\begin{reference}
\cite[補題 6.5.9 (2)]{BS} の一般化である．準連接加群および（導来）
代数スタックの射という設定では \cite[定理 6.5]{HL-P} と比較されたい．
\end{reference}
$f : (\Sh(\mathcal{D}), \mathcal{O}') \to (\Sh(\mathcal{C}), \mathcal{O})$
を環付きトポスの射とする．$\mathcal{I} \subset \mathcal{O}$ を
有限型のイデアル層とする．$\mathcal{I}' \subset \mathcal{O}'$ を
$f^\sharp(f^{-1}\mathcal{I})$ で生成されるイデアルとする．このとき
$Rf_*$ は導来完備化と可換する．すなわち，
$Rf_*(K^\wedge) = (Rf_*K)^\wedge$.
\end{lemma}

\begin{proof}
命題 \ref{algebraization-proposition-derived-completion} により導来完備化関手は存在する．
補題 \ref{algebraization-lemma-pushforward-derived-complete} により対象
$Rf_*(K^\wedge)$ は導来完備である．したがって導来完備化の普遍性から標準写像
$(Rf_*K)^\wedge \to Rf_*(K^\wedge)$ を得る．この写像が同型であることは
$\mathcal{C}$ 上で局所的に確認してよい．導来完備化は局所化と可換するので
（注 \ref{algebraization-remark-localization-and-completion}），$\mathcal{I}$ は大域切断
$f_1, \ldots, f_r$ で生成されると仮定してよい．このとき $\mathcal{I}'$ は
$g_i = f^\sharp(f_i)$ で生成される．補題
\ref{algebraization-lemma-derived-completion-koszul} により，示すべきことは
$$
R\lim \left(
Rf_*K \otimes^\mathbf{L}_\mathcal{O} K(\mathcal{O}, f_1^n, \ldots, f_r^n)
\right)
=
Rf_*\left(
R\lim
K \otimes^\mathbf{L}_{\mathcal{O}'} K(\mathcal{O}', g_1^n, \ldots, g_r^n)
\right)
$$
$Rf_*$ は $R\lim$ と可換するので（Cohomology on Sites, 補題
\ref{sites-cohomology-lemma-Rf-commutes-with-Rlim}），次を示せば十分である：
$$
Rf_*K \otimes^\mathbf{L}_\mathcal{O} K(\mathcal{O}, f_1^n, \ldots, f_r^n) =
Rf_*\left(
K \otimes^\mathbf{L}_{\mathcal{O}'} K(\mathcal{O}', g_1^n, \ldots, g_r^n)
\right)
$$
これは射影公式（Cohomology on Sites, 補題
\ref{sites-cohomology-lemma-projection-formula}）と，次の事実から従う：
$Lf^*K(\mathcal{O}, f_1^n, \ldots, f_r^n) =
K(\mathcal{O}', g_1^n, \ldots, g_r^n)$.
\end{proof}

\begin{lemma}
\label{algebraization-lemma-formal-functions-general}
$A$ を環とし，$I \subset A$ を有限生成イデアルとする．
$\mathcal{C}$ をサイトとし，$\mathcal{O}$ を $A$-代数の層とする．
$\mathcal{F}$ を $\mathcal{O}$-加群の層とする．このとき
$$
R\Gamma(\mathcal{C}, \mathcal{F})^\wedge =
R\Gamma(\mathcal{C}, \mathcal{F}^\wedge)
$$
が $D(A)$ において成り立つ．ここで $\mathcal{F}^\wedge$ は
$\mathcal{F}$ の $I\mathcal{O}$ に関する導来完備化であり，左辺では
$I$ に関する導来完備化を取る．これにより二つのスペクトル系列
$$
E_2^{i, j} = H^i(H^j(\mathcal{C}, \mathcal{F})^\wedge)
\quad\text{および}\quad
E_2^{p, q} = H^p(\mathcal{C}, H^q(\mathcal{F}^\wedge))
$$
を得る．いずれも
$H^*(R\Gamma(\mathcal{C}, \mathcal{F})^\wedge) =
H^*(\mathcal{C}, \mathcal{F}^\wedge)$
に収束する．
\end{lemma}

\begin{proof}
補題 \ref{algebraization-lemma-pushforward-commutes-with-derived-completion} を環付きトポスの射
$(\mathcal{C}, \mathcal{O}) \to (pt, A)$ に適用し，コホモロジーを取れば
第一の主張を得る．第二のスペクトル系列は Derived Categories, 補題
\ref{derived-lemma-two-ss-complex-functor} の第二のスペクトル系列である．
第一のスペクトル系列は More on Algebra, 例
\ref{more-algebra-example-derived-completion-spectral-sequence} のものを
$R\Gamma(\mathcal{C}, \mathcal{F})^\wedge$ に適用して得られる．
\end{proof}

\begin{remark}
\label{algebraization-remark-local-calculation-derived-completion}
$(\mathcal{C}, \mathcal{O})$ を環付きサイトとする．
$\mathcal{I} \subset \mathcal{O}$ を有限型のイデアル層とする．
$K \mapsto K^\wedge$ を命題 \ref{algebraization-proposition-derived-completion} の
導来完備化とする．$U \in \Ob(\mathcal{C})$ を，$\mathcal{I}$ がイデアル層として
$f_1, \ldots, f_r \in \mathcal{I}(U)$ によって生成されるような対象とする．
$A = \mathcal{O}(U)$ および $I = (f_1, \ldots, f_r) \subset A$ と置く．
注意：$I = \mathcal{I}(U)$ とは限らない．このとき
$$
R\Gamma(U, K^\wedge) = R\Gamma(U, K)^\wedge
$$
が成り立つ．ここで右辺は，対象 $R\Gamma(U, K)$（$D(A)$ の対象）の
$I$ に関する導来完備化である．これは導来完備化が局所化と可換すること
（注 \ref{algebraization-remark-localization-and-completion}）と，補題
\ref{algebraization-lemma-formal-functions-general} による．
\end{remark}








\section{形式関数定理}
\label{algebraization-section-formal-functions}

\noindent
ここで叙述の流れをいったん中断し，スキーム上の準連接加群という設定での
導来完備化について少し述べる．そしてこれを用いて，形式関数定理のやや異なる
証明を与える．文献への手引きは注 \ref{algebraization-remark-references} に掲げる．

\medskip\noindent
補題 \ref{algebraization-lemma-pushforward-commutes-with-derived-completion} は形式関数定理の
（きわめて形式的な）導来版である（Cohomology of Schemes, 定理
\ref{coherent-theorem-formal-functions}）．これをより明示的に述べよう．
$f : X \to S$ をスキームの射，$\mathcal{I} \subset \mathcal{O}_S$ を有限型の
準連接イデアル層，$\mathcal{F}$ を $X$ 上の準連接層とする．このとき補題は
\begin{equation}
\label{algebraization-equation-formal-functions}
Rf_*(\mathcal{F}^\wedge) = (Rf_*\mathcal{F})^\wedge
\end{equation}
を主張する．ここで $\mathcal{F}^\wedge$ は $\mathcal{F}$ の
$f^{-1}\mathcal{I} \cdot \mathcal{O}_X$ に関する導来完備化であり，右辺は
$Rf_*\mathcal{F}$ の $\mathcal{I}$ に関する導来完備化である．これから
形式関数定理を回収するには，多少の準備が必要である．

\begin{lemma}
\label{algebraization-lemma-sections-derived-completion-pseudo-coherent}
$X$ を局所 Noether スキームとする．$\mathcal{I} \subset \mathcal{O}_X$ を
準連接イデアル層とする．$K$ を $D(\mathcal{O}_X)$ の擬連接対象とし，
その導来完備化を $K^\wedge$ とする．このとき
$$
H^p(U, K^\wedge) = \lim H^p(U, K)/I^nH^p(U, K) =
H^p(U, K)^\wedge
$$
が任意のアフィン開集合 $U \subset X$ に対して成り立つ．ここで
$I = \mathcal{I}(U)$ であり，右端では $I$ に関する導来完備化を取る．
\end{lemma}

\begin{proof}
$U = \Spec(A)$ と書く．環 $A$ は Noether 環であり，したがって
$I \subset A$ は有限生成である．このとき
$$
R\Gamma(U, K^\wedge) = R\Gamma(U, K)^\wedge
$$
が注 \ref{algebraization-remark-local-calculation-derived-completion} により成り立つ．
ここで $R\Gamma(U, K)$ は $A$-加群の擬連接複体である
（Derived Categories of Schemes, 補題
\ref{perfect-lemma-pseudo-coherent-affine}）．More on Algebra, 補題
\ref{more-algebra-lemma-derived-completion-pseudo-coherent} により，第
$p$ コホモロジー加群，すなわち $R\Gamma(U, K^\wedge)$ のその加群は，
$I$-進完備化された $H^p(U, K)$ に等しい．これで第一の等号が示された．
第二の（それほど重要でない）等号は，いま用いた補題をもう一度適用すれば直ちに従う．
\end{proof}

\begin{lemma}
\label{algebraization-lemma-derived-completion-pseudo-coherent}
$X$ を局所 Noether スキームとし，$\mathcal{I} \subset \mathcal{O}_X$ を
準連接イデアル層とする．$K$ を $D(\mathcal{O}_X)$ の対象とする．このとき
\begin{enumerate}
\item 導来完備化 $K^\wedge$ は
$R\lim (K \otimes_{\mathcal{O}_X}^\mathbf{L} \mathcal{O}_X/\mathcal{I}^n)$.
に等しい．
\end{enumerate}
$K$ を $D(\mathcal{O}_X)$ の擬連接対象とする．このとき
\begin{enumerate}
\item[(2)] コホモロジー層 $H^q(K^\wedge)$ は
$\lim H^q(K)/\mathcal{I}^nH^q(K)$.
に等しい．
\end{enumerate}
$\mathcal{F}$ を連接 $\mathcal{O}_X$-加群とする\footnote{例えば
$H^q(K)$．ここで $K$ は局所 Noether スキーム $X$ 上で擬連接である．}．このとき
\begin{enumerate}
\item[(3)] 導来完備化 $\mathcal{F}^\wedge$ は
$\lim \mathcal{F}/\mathcal{I}^n\mathcal{F}$,
に等しい；
\item[(4)]
$\lim \mathcal{F}/\mathcal{I}^n \mathcal{F} =
R\lim \mathcal{F}/\mathcal{I}^n \mathcal{F}$,
\item[(5)] $H^p(U, \mathcal{F}^\wedge) = 0$ が，$p \not = 0$ および任意の
アフィン開集合 $U \subset X$ に対して成り立つ．
\end{enumerate}
\end{lemma}

\begin{proof}
(1) の証明．標準写像
$$
K \longrightarrow
R\lim (K \otimes_{\mathcal{O}_X}^\mathbf{L} \mathcal{O}_X/\mathcal{I}^n),
$$
が存在する．注 \ref{algebraization-remark-compare-with-completion} を参照されたい．
導来完備化は開部分スキームへの移行と可換する（注
\ref{algebraization-remark-localization-and-completion}）．$R\lim$ の形成も開部分スキームへの
移行と可換する．したがって，この写像が同型であることは局所的に確認してよい．
ゆえに $X = U = \Spec(A)$ と仮定してよい．$I = (f_1, \ldots, f_r)$ と書き，
$K_n = K(A, f_1^n, \ldots, f_r^n)$ を Koszul 複体とする．
More on Algebra, 補題 \ref{more-algebra-lemma-sequence-Koszul-complexes} により，
プロ系 $\{K_n\}$ と $\{A/I^n\}$ は $D(A)$ において同型である．
圏同値 $D(A) = D_{\QCoh}(\mathcal{O}_X)$（Derived Categories of Schemes, 補題
\ref{perfect-lemma-affine-compare-bounded}）を用いると，プロ系
$\{K(\mathcal{O}_X, f_1^n, \ldots, f_r^n)\}$ と
$\{\mathcal{O}_X/\mathcal{I}^n\}$ は $D(\mathcal{O}_X)$ において同型である．
これで次の式の第二の等号が示される：
$$
K^\wedge = R\lim \left(
K \otimes_{\mathcal{O}_X}^\mathbf{L} K(\mathcal{O}_X, f_1^n, \ldots, f_r^n)
\right) =
R\lim (K \otimes_{\mathcal{O}_X}^\mathbf{L} \mathcal{O}_X/\mathcal{I}^n)
$$
第一の等号は補題 \ref{algebraization-lemma-derived-completion-koszul} である．

\medskip\noindent
$K$ が擬連接であると仮定する．アフィン開集合 $U \subset X$ に対して
$H^q(U, K^\wedge) = \lim H^q(U, K)/\mathcal{I}^n(U)H^q(U, K)$
が補題 \ref{algebraization-lemma-sections-derived-completion-pseudo-coherent} により成り立つ．
これはすべての $U$ に対して成り立つので，層として
$H^q(K^\wedge) = \lim H^q(K)/\mathcal{I}^nH^q(K)$ となる．これで (2) が示された．

\medskip\noindent
(3) は (2) の特別な場合である．(4) と (5) は
Derived Categories of Schemes, 補題
\ref{perfect-lemma-Rlim-quasi-coherent} から従う．
\end{proof}

\begin{lemma}
\label{algebraization-lemma-formal-functions}
$A$ を Noether 環とし，$I \subset A$ をイデアルとする．$X$ を
$A$ 上の Noether スキームとする．$\mathcal{F}$ を連接
$\mathcal{O}_X$-加群とする．$H^p(X, \mathcal{F})$ が有限
$A$-加群であると，すべての $p$ に対して仮定する．このとき短完全列
$$
0 \to R^1\lim H^{p - 1}(X, \mathcal{F}/I^n\mathcal{F}) \to
H^p(X, \mathcal{F})^\wedge \to \lim H^p(X, \mathcal{F}/I^n\mathcal{F}) \to 0
$$
が $A$-加群について存在する．ここで $H^p(X, \mathcal{F})^\wedge$ は通常の
$I$-進完備化である．$f$ が固有ならば，$R^1\lim$ の項は零である．
\end{lemma}

\begin{proof}
補題 \ref{algebraization-lemma-formal-functions-general} の二つのスペクトル系列を考える．
第一は More on Algebra, 補題
\ref{more-algebra-lemma-derived-completion-pseudo-coherent} により退化する．
$H^p(X, \mathcal{F})^\wedge$ を次数 $p$ に得る．ここで
$H^p(X, \mathcal{F})$ が有限 $A$-加群であるという仮定を用いる．
第二は，次が層であるため退化する：
$$
\mathcal{F}^\wedge = \lim \mathcal{F}/I^n\mathcal{F} =
R\lim \mathcal{F}/I^n\mathcal{F}
$$
これは補題 \ref{algebraization-lemma-derived-completion-pseudo-coherent} による．
$H^p(X, \lim \mathcal{F}/I^n\mathcal{F})$ を次数 $p$ に得る．
$R\Gamma(X, -)$ は導来逆極限と可換するので（Injectives, 補題
\ref{injectives-lemma-RF-commutes-with-Rlim}），さらに
$$
R\Gamma(X, \lim \mathcal{F}/I^n\mathcal{F}) =
R\Gamma(X, R\lim \mathcal{F}/I^n\mathcal{F}) =
R\lim R\Gamma(X, \mathcal{F}/I^n\mathcal{F})
$$
More on Algebra, 注 \ref{more-algebra-remark-how-unique} により完全列
$$
0 \to
R^1\lim H^{p - 1}(X, \mathcal{F}/I^n\mathcal{F}) \to
H^p(X, \lim \mathcal{F}/I^n\mathcal{F}) \to
\lim H^p(X, \mathcal{F}/I^n\mathcal{F}) \to 0
$$
を $A$-加群について得る．以上を合わせれば補題の第一の主張を得る．
$R^1\lim$ の項の消滅は Cohomology of Schemes, 補題
\ref{coherent-lemma-ML-cohomology-powers-ideal} から従う．
\end{proof}

\begin{remark}
\label{algebraization-remark-references}
関連事項を論じる文献をいくつか挙げる．固有射に対する「導来形式関数定理」は
\cite[定理 6.3.1]{lurie-thesis} に遡るようである．\cite{dag12} にも議論があり，
特に第 4 章ではアフィンの場合が扱われる．More on Algebra, 節
\ref{more-algebra-section-derived-completion} も参照されたい．
\cite[節 2.9]{G-R} には，（ind）連接加群を用いた固有底変換と導来完備化の
議論がある．準連接加群の複体に対する
(\ref{algebraization-equation-formal-functions}) の類似は \cite[定理 6.5]{HL-P} に見いだせる．
\end{remark}







\section{局所コホモロジーの代数化 I}
\label{algebraization-section-algebraization-sections-general}

\noindent
$A$ を Noether 環とし，$I$ と $J$ を $A$ の二つのイデアルとする．
$M$ を有限 $A$-加群とする．本節では，対象
$$
R\Gamma_J(M)^\wedge
\quad\text{（次の圏の対象）}\quad
D(A)
$$
のコホモロジー群を調べる．ここで ${}^\wedge$ は導来 $I$-進完備化を表す．
Dualizing Complexes, 補題 \ref{dualizing-lemma-completion-local-H0} では，
$A$ が $I$ に関して完備ならば，同型
$$
\colim H^0_Z(M) \longrightarrow H^0(R\Gamma_J(M)^\wedge)
$$
が存在することを示した．ここで（有向）余極限は閉部分集合
$Z = V(J')$ であって，$J' \subset J$ かつ
$V(J') \cap V(I) = V(J) \cap V(I)$ を満たすもの全体にわたる．
これらの閉部分集合の合併は
\begin{equation}
\label{algebraization-equation-associated-subset}
T = \{\mathfrak p \in \Spec(A) :
V(\mathfrak p) \cap V(I) \subset V(J) \cap V(I)\}
\end{equation}
である．これは特殊化のもとで安定な $\Spec(A)$ の部分集合である．
上の結果は，写像
$$
H^0_T(M) \longrightarrow H^0(R\Gamma_J(M)^\wedge)
$$
が同型であるという主張になる．ただし $A$ は $I$ に関して完備とする．
Local Cohomology, 補題 \ref{local-cohomology-lemma-adjoint-ext} および注
\ref{local-cohomology-remark-upshot} を参照されたい．この同型を高次
コホモロジー群へ拡張する方法は，次の補題に基づく．

\begin{lemma}
\label{algebraization-lemma-kill-completion-general}
$I, J$ を Noether 環 $A$ のイデアルとする．$M$ を有限
$A$-加群，$\mathfrak p \subset A$ を素イデアルとする．$s$ と $d$ を整数とする．
次を仮定する：
\begin{enumerate}
\item $A$ は双対化複体をもつ；
\item $\mathfrak p \not \in V(J) \cap V(I)$；
\item $\text{cd}(A, I) \leq d$；
\item すべての素イデアル $\mathfrak p' \subset \mathfrak p$ に対して
$\text{depth}_{A_{\mathfrak p'}}(M_{\mathfrak p'}) +
\dim((A/\mathfrak p')_\mathfrak q) > d + s$
が，すべての $\mathfrak q \in V(\mathfrak p') \cap V(J) \cap V(I)$ に対して成り立つ．
\end{enumerate}
このとき $f \in A$，$f \not \in \mathfrak p$ を満たす元であって，
$H^i(R\Gamma_J(M)^\wedge)$ を $i \leq s$ に対して消すものが存在する．
ここで ${}^\wedge$ は $I$-進完備化を表す．
\end{lemma}

\begin{proof}
$R\Gamma_J = R\Gamma_{V(J)}$ を用いる．$I + J$ についても同様である．
Dualizing Complexes, 補題 \ref{dualizing-lemma-local-cohomology-noetherian}
を参照されたい．また
$R\Gamma_J(M)^\wedge = R\Gamma_I(R\Gamma_J(M))^\wedge =
R\Gamma_{I + J}(M)^\wedge$ である．Dualizing Complexes, 補題
\ref{dualizing-lemma-complete-and-local} および
\ref{dualizing-lemma-local-cohomology-ss} を参照されたい．
したがって $J$ を $I + J$ で置き換え，$I \subset J$ かつ
$\mathfrak p \not \in V(J)$ と仮定してよい．次を思い出そう：
$$
R\Gamma_J(M)^\wedge = R\Hom_A(R\Gamma_I(A), R\Gamma_J(M))
$$
である．これは More on Algebra, 補題
\ref{more-algebra-lemma-derived-completion} の導来完備化の記述と，
Dualizing Complexes, 補題
\ref{dualizing-lemma-compute-local-cohomology-noetherian} の局所コホモロジーの
記述を合わせたものである．仮定 (3) は，$R\Gamma_I(A)$ の非零コホモロジーが
次数 $\leq d$ にしか存在しないことを意味する．$R\Gamma_I(A)$ の標準切断を
用いると，次を示せば十分である：
$$
\text{Ext}^i(N, R\Gamma_J(M))
$$
これは，ある $f \in A$，$f \not \in \mathfrak p$ によって，
$i \leq s + d$ と任意の $A$-加群 $N$ に対して消される．次に
$R\Gamma_J(M)$ の標準切断を用いると，$H^i_J(M)$ がある
$f \in A$，$f \not \in \mathfrak p$ によって $i \leq s + d$ に対して
消されることを示せば十分である．
これは Local Cohomology, 補題
\ref{local-cohomology-lemma-kill-local-cohomology-at-prime} から従う．
\end{proof}

\begin{lemma}
\label{algebraization-lemma-kill-colimit-weak-general}
$I, J$ をある Noether 環のイデアルとし，$M$ を有限 $A$-加群とする．
$s$ と $d$ を整数とする．$T$ を
(\ref{algebraization-equation-associated-subset}) のものとして，次を仮定する：
\begin{enumerate}
\item $A$ は双対化複体をもつ；
\item $\mathfrak p \in V(I)$ ならば条件を課さない；
\item $\mathfrak p \not \in V(I)$ かつ $\mathfrak p \in T$ ならば，
$\dim((A/\mathfrak p)_\mathfrak q) \leq d$ となる
ある $\mathfrak q \in V(\mathfrak p) \cap V(J) \cap V(I)$ が存在する；
\item $\mathfrak p \not \in V(I)$ かつ $\mathfrak p \not \in T$ ならば，
$$
\text{depth}_{A_\mathfrak p}(M_\mathfrak p) \geq s
\quad\text{または}\quad
\text{depth}_{A_\mathfrak p}(M_\mathfrak p) +
\dim((A/\mathfrak p)_\mathfrak q) > d + s
$$
が，すべての $\mathfrak q \in V(\mathfrak p) \cap V(J) \cap V(I)$ に対して成り立つ．
\end{enumerate}
このとき，イデアル $J_0 \subset J$ であって
$V(J_0) \cap V(I) = V(J) \cap V(I)$ を満たすものが存在し，次の性質をもつ．
任意の $J' \subset J_0$ であって
$V(J') \cap V(I) = V(J) \cap V(I)$ を満たすものに対して写像
$$
R\Gamma_{J'}(M) \longrightarrow R\Gamma_{J_0}(M)
$$
は次数 $\leq s$ のコホモロジー上で同型を誘導し，さらにこれらの加群は
$J_0I$ のある冪で消される．
\end{lemma}

\begin{proof}
集合
$$
B = \{\mathfrak p \not \in V(I),\ \mathfrak p \in T,\text{ かつ }
\text{depth}(M_\mathfrak p) \leq s\}
$$
を考える．$J_0 \subset J$ を，$V(J_0)$ が $B \cup V(J)$ の閉包となるように選ぶ．

\medskip\noindent
主張 I：$V(J_0) \cap V(I) = V(J) \cap V(I)$．

\medskip\noindent
主張 I の証明．包含 $\supset$ は構成から成り立つ．
$\mathfrak p$ を $V(J_0)$ の極小素イデアルとする．
$\mathfrak p \in B \cup V(J)$ ならば，$\mathfrak p \in T$ または
$\mathfrak p \in V(J)$ であり，いずれの場合にも望みどおり
$V(\mathfrak p) \cap V(I) \subset V(J) \cap V(I)$ となる．
$\mathfrak p \not \in B \cup V(J)$ ならば，$V(\mathfrak p) \cap B$ は
稠密であり，したがって無限である．ゆえに
$\text{depth}(M_\mathfrak p) < s$ である（Local Cohomology, 補題
\ref{local-cohomology-lemma-depth-function}）．実際，
$V(\mathfrak p) \cap B = \{\mathfrak p_\lambda\}_{\lambda \in \Lambda}$.
と置く．(3) にあるような
$\mathfrak q_\lambda \in V(\mathfrak p_\lambda) \cap V(J) \cap V(I)$
を選ぶ．$\delta : \Spec(A) \to \mathbf{Z}$ を，双対化複体
$\omega_A^\bullet$（$A$ に対するもの）に付随する次元関数とする．$\Lambda$ は無限で
$\delta$ は有界なので，無限部分集合 $\Lambda' \subset \Lambda$ であって
$\delta(\mathfrak q_\lambda)$ がその上で一定となるものが存在する．
$\lambda \in \Lambda'$ に対して
$$
\text{depth}(M_{\mathfrak p_\lambda}) +
\delta(\mathfrak p_\lambda) - \delta(\mathfrak q_\lambda) =
\text{depth}(M_{\mathfrak p_\lambda}) +
\dim((A/\mathfrak p_\lambda)_{\mathfrak q_\lambda})
\leq d + s
$$
が (3) と $B$ の定義により成り立つ．関数
$\text{depth} + \delta$ の半連続性（Duality for Schemes, 補題
\ref{duality-lemma-sitting-in-degrees}）から
$$
\text{depth}(M_\mathfrak p) +
\dim((A/\mathfrak p)_{\mathfrak q_\lambda}) =
\text{depth}(M_\mathfrak p) + \delta(\mathfrak p) - \delta(\mathfrak q_\lambda)
\leq d + s
$$
を得る．さらに $\mathfrak p \not \in V(I)$ なので，(4) から
$\mathfrak p \in T$，すなわち
$V(\mathfrak p) \cap V(I) \subset V(J) \cap V(I)$ を読み取れる．
これで主張 I の証明が終わる．

\medskip\noindent
主張 II：$H^i_{J_0}(M) \to H^i_J(M)$ は，$i \leq s$ および
$J' \subset J_0$ かつ $V(J') \cap V(I) = V(J) \cap V(I)$ に対して同型である．

\medskip\noindent
主張 II の証明．$\mathfrak p \in V(J')$ であって $V(J_0)$ に属さないものを選ぶ．
$H^i_{\mathfrak pA_\mathfrak p}(M_\mathfrak p) = 0$ が $i \leq s$ に対して
成り立つことを示せば十分である．Local Cohomology, 補題
\ref{local-cohomology-lemma-isomorphism} を参照されたい．
$\mathfrak p \in T$ であることに注意する．したがって $\mathfrak p$ は
$B$ に属さないので，$\text{depth}(M_\mathfrak p) > s$ であり，群は
Dualizing Complexes, 補題 \ref{dualizing-lemma-depth} により消滅する．

\medskip\noindent
主張 III．補題の最後の主張が成り立つ．

\medskip\noindent
主張 II により，$i \leq s$ に対して
$$
H^i_T(M) = H^i_{J_0}(M) = H^i_{J'}(M)
$$
が，$J' \subset J_0$ かつ
$V(J')  \cap V(I) = V(J) \cap V(I)$ を満たすすべてのイデアルについて成り立つ．
Local Cohomology, 補題 \ref{local-cohomology-lemma-adjoint-ext} を参照されたい．
Local Cohomology, 命題 \ref{local-cohomology-proposition-annihilator} の仮定を，
部分集合 $T \subset T \cup V(I)$，加群 $M$，整数 $s$ について確認する．
$\mathfrak p \subset \mathfrak q$，$\mathfrak p \not \in T \cup V(I)$，
$\mathfrak q \in T$ が与えられたとき，次を示さなければならない：
$$
\text{depth}_{A_\mathfrak p}(M_\mathfrak p) +
\dim((A/\mathfrak p)_\mathfrak q) > s
$$
$\text{depth}(M_\mathfrak p) \geq s$ ならば，これは
$(A/\mathfrak p)_\mathfrak q$ の次元が少なくとも $1$ なので成り立つ．
したがって $\text{depth}(M_\mathfrak p) < s$ と仮定してよい．
$\mathfrak q \in V(I)$ ならば $\mathfrak q \in V(J) \cap V(I)$ であり，
不等式は (4) により成り立つ．$\mathfrak q \not \in V(I)$ ならば，(3) を用いて
$\mathfrak q' \in V(\mathfrak q) \cap V(J) \cap V(I)$ であって
$\dim((A/\mathfrak q)_{\mathfrak q'}) \leq d$ を満たすものを選べる．
このとき仮定 (4) から
$$
\text{depth}_{A_\mathfrak p}(M_\mathfrak p) +
\dim((A/\mathfrak p)_{\mathfrak q'}) > s + d
$$
を得る．$A$ はカテナリーなので，これは求める不等式を含意する．
Local Cohomology, 命題 \ref{local-cohomology-proposition-annihilator} を適用すると，
$J'' \subset A$ であって $V(J'') \subset T \cup V(I)$ を満たし，
$J''$ が $H^i_T(M)$ を $i \leq s$ に対して消すものが得られる．このとき
$V(J'') \cup V(J_0) \cup V(I) = V(J'I)$ と書ける．ここで
$J' \subset J_0$ かつ $V(J') \cap V(I) = V(J) \cap V(I)$ である．
$J_0$ を $J'$ で置き換えれば証明は完了する．
\end{proof}

\begin{lemma}
\label{algebraization-lemma-kill-colimit-general}
補題 \ref{algebraization-lemma-kill-colimit-weak-general} において，空の条件 (2) の代わりに
\begin{enumerate}
\item[(2')] $\mathfrak p \in V(I)$ かつ
$\mathfrak p \not \in V(J) \cap V(I)$ ならば，
$\text{depth}_{A_\mathfrak p}(M_\mathfrak p) +
\dim((A/\mathfrak p)_\mathfrak q) > s$
が，すべての $\mathfrak q \in V(\mathfrak p) \cap V(J) \cap V(I)$ に対して成り立つ，
\end{enumerate}
と仮定するならば，これらの条件はさらに $H^i_{J_0}(M)$ が有限
$A$-加群であることを $i \leq s$ に対して含意する．
\end{lemma}

\begin{proof}
$H^i_{J_0}(M) = H^i_T(M)$ であることを思い出そう．補題
\ref{algebraization-lemma-kill-colimit-weak-general} の証明を参照されたい．したがって，
$\mathfrak p \not \in T$，$\mathfrak q \in T$，かつ
$\mathfrak p \subset \mathfrak q$ ならば
$\text{depth}_{A_\mathfrak p}(M_\mathfrak p) +
\dim((A/\mathfrak p)_\mathfrak q) > s$ となることを確認すれば十分である．
Local Cohomology, 命題 \ref{local-cohomology-proposition-finiteness} を参照されたい．
条件 (2') により，これは $\mathfrak p \in V(I)$ のとき成り立つ．
$H^i_T(M)$ は $IJ_0$ のある冪で消されるから，
$\mathfrak p \not \in V(IJ_0)$ のときにも条件は成り立つ
（Local Cohomology, 命題 \ref{local-cohomology-proposition-annihilator}）．
これですべての場合が尽くされ，証明が完了する．
\end{proof}

\begin{lemma}
\label{algebraization-lemma-kill-colimit-support-general}
補題 \ref{algebraization-lemma-kill-colimit-weak-general} において，さらに
\begin{enumerate}
\item[(6)] $\mathfrak p \not \in V(I)$ かつ $\mathfrak p \in T$ ならば，
$\text{depth}_{A_\mathfrak p}(M_\mathfrak p) > s$,
\end{enumerate}
と仮定するならば，$H^i_{J_0}(M) = H^i_J(M) = H^i_{J + I}(M)$ が
$i \leq s$ に対して成り立ち，これらの加群は $I$ のある冪で消される．
\end{lemma}

\begin{proof}
$\mathfrak p \in V(J)$ または $\mathfrak p \in V(J_0)$ であるが，
$\mathfrak p \not \in V(J + I) = V(J_0 + I)$ となるものを選ぶ．
$H^i_{\mathfrak pA_\mathfrak p}(M_\mathfrak p) = 0$ が $i \leq s$ に対して
成り立つことを示せば十分である．Local Cohomology, 補題
\ref{local-cohomology-lemma-isomorphism} を参照されたい．これらの群は条件 (6) と
Dualizing Complexes, 補題 \ref{dualizing-lemma-depth} により消滅する．最後の主張は
Local Cohomology, 命題 \ref{local-cohomology-proposition-annihilator} から従う．
\end{proof}

\begin{lemma}
\label{algebraization-lemma-algebraize-local-cohomology-general}
$I, J$ を Noether 環 $A$ のイデアルとし，$M$ を有限 $A$-加群とする．
$s$ と $d$ を整数とする．$T$ を
(\ref{algebraization-equation-associated-subset}) のものとして，次を仮定する：
\begin{enumerate}
\item $A$ は $I$-進完備であり，双対化複体をもつ；
\item $\mathfrak p \in V(I)$ ならば条件を課さない；
\item $\text{cd}(A, I) \leq d$；
\item $\mathfrak p \not \in V(I)$ かつ $\mathfrak p \not \in T$ ならば，
$$
\text{depth}_{A_\mathfrak p}(M_\mathfrak p) \geq s
\quad\text{または}\quad
\text{depth}_{A_\mathfrak p}(M_\mathfrak p) +
\dim((A/\mathfrak p)_\mathfrak q) > d + s
$$
が，すべての $\mathfrak q \in V(\mathfrak p) \cap V(J) \cap V(I)$ に対して成り立つ；
\item $\mathfrak p \not \in V(I)$，$\mathfrak p \not \in T$，
$V(\mathfrak p) \cap V(J) \cap V(I) \not = \emptyset$，かつ
$\text{depth}(M_\mathfrak p) < s$ ならば，次のいずれかが成り立つ
\footnote{我々の方法ではこの追加条件が必要になる．これについては後で
（将来の参照を挿入）再び論じる．}：
\begin{enumerate}
\item $\dim(\text{Supp}(M_\mathfrak p)) < s + 2$\footnote{例えば，
$M$ が Serre の条件 $(S_s)$ を $V(I) \cup T$ の補集合上で満たす場合．}；または
\item  $\delta(\mathfrak p) > d + \delta_{max} - 1$
である．ここで $\delta$ は次元関数であり，$\delta_{max}$ は
$\delta$ の $V(J) \cap V(I)$ 上での最大値である；または
\item $\text{depth}_{A_\mathfrak p}(M_\mathfrak p) +
\dim((A/\mathfrak p)_\mathfrak q) > d + s + \delta_{max} - \delta_{min} - 2$
が，すべての $\mathfrak q \in V(\mathfrak p) \cap V(J) \cap V(I)$ に対して成り立つ．
\end{enumerate}
\end{enumerate}
このとき，イデアル $J_0 \subset J$ であって
$V(J_0) \cap V(I) = V(J) \cap V(I)$ を満たすものが存在し，次の性質をもつ．
任意の $J' \subset J_0$ であって
$V(J') \cap V(I) = V(J) \cap V(I)$ を満たすものに対して，写像
$$
R\Gamma_{J'}(M) \longrightarrow R\Gamma_J(M)^\wedge
$$
は次数 $\leq s$ のコホモロジー上で同型を誘導する．ここで
${}^\wedge$ は導来 $I$-進完備化を表す．
\end{lemma}

\noindent
読者にはまず局所の場合の証明（補題
\ref{algebraization-lemma-algebraize-local-cohomology}）を読むことを勧める．そこでは，
すべての不等式を扱わずに証明の構造が説明される．

\begin{proof}
イデアル $\mathfrak a \subset A$ に対して
$R\Gamma_\mathfrak a = R\Gamma_{V(\mathfrak a)}$ である．Dualizing Complexes, 補題
\ref{dualizing-lemma-local-cohomology-noetherian} を参照されたい．次に，
$$
R\Gamma_J(M)^\wedge =
R\Gamma_I(R\Gamma_J(M))^\wedge =
R\Gamma_{I + J}(M)^\wedge =
R\Gamma_{I + J'}(M)^\wedge =
R\Gamma_I(R\Gamma_{J'}(M))^\wedge =
R\Gamma_{J'}(M)^\wedge
$$
となることに注意する．これは Dualizing Complexes, 補題
\ref{dualizing-lemma-local-cohomology-ss} および
\ref{dualizing-lemma-complete-and-local} による．これで補題の主張にある
矢印をどのように定義するかが分かる．

\medskip\noindent
補題 \ref{algebraization-lemma-kill-colimit-weak-general} の仮定は，現在の $M$ に関する
仮定から従うと主張する．確認すべきは仮定 (3) だけである．そこで
$\mathfrak p \not \in V(I)$，$\mathfrak p \in T$ とする．
$V(\mathfrak p) \cap V(I)$ は空でない．実際，$I$ は $A$ の
Jacobson 根基に含まれる（Algebra, 補題
\ref{algebra-lemma-radical-completion}）．$\mathfrak p \in T$ なので
$V(\mathfrak p) \cap V(I) = V(\mathfrak p) \cap V(J) \cap V(I)$ である．
$\mathfrak q \in V(\mathfrak p) \cap V(I)$ を既約成分の一般点とする．
Local Cohomology, 補題 \ref{local-cohomology-lemma-cd-local} により
$\text{cd}(A_\mathfrak q, I_\mathfrak q) \leq d$ である．
$V(\mathfrak pA_\mathfrak q) \cap V(I_\mathfrak q) =
\{\mathfrak qA_\mathfrak q\}$ である．これは $\mathfrak q$ の選び方による．
さらに Local Cohomology, 補題
\ref{local-cohomology-lemma-cd-bound-dim-local} により
$\dim((A/\mathfrak p)_\mathfrak q) \leq d$ を得る．

\medskip\noindent
補題は $s < 0$ に対して成り立つことに注意する．これは自明ではない．実際，
$H^i(R\Gamma_J(M)^\wedge)$ が $i < 0$ に対して零であることは先験的には
明らかでない．しかしこの消滅は Dualizing Complexes, 補題
\ref{dualizing-lemma-completion-local} で示されている．
$s \geq 0$ に関して帰納法で補題を証明する．

\medskip\noindent
$s = 0$ の場合は，補題 \ref{algebraization-lemma-kill-colimit-weak-general} の結論と
Dualizing Complexes, 補題 \ref{dualizing-lemma-completion-local-H0} から直ちに従う．

\medskip\noindent
$s > 0$ とし，$s$ より小さい値について補題はすでに示されているとする．
$M' \subset M$ を，その台が $V(I) \cup T$ に含まれる最大の部分加群とする．
このとき $M'$ は有限 $A$ 加群であり，その台は $V(J') \cup V(I)$ に含まれる．
ここで，あるイデアル $J' \subset J$ は
$V(J') \cap V(I) = V(J) \cap V(I)$ を満たす．次を主張する：
$$
R\Gamma_{J'}(M') \to R\Gamma_J(M')^\wedge
$$
は $J'$ のどの選択に対しても同型である．実際，短完全列
$0 \to M_1 \oplus M_2 \to M' \to N \to 0$ であって，$M_1$ は
$J'$ のある冪で零化され，$M_2$ は $I$ のある冪で零化され，$N$ は
$I + J'$ のある冪で零化されるものを選べる．したがって，この主張を
$M_1$，$M_2$，$N$ について示せば十分である．$M_1$ の場合には
$R\Gamma_{J'}(M_1) = M_1$ である．また $M_1$ は有限 $A$ 加群かつ
$I$ 進完備なので $M_1^\wedge = M_1$ である．証明冒頭の考察により，
これは $M_1$ に対する主張を示す．$M_2$ の場合には
$H^i_J(M_2) = H^i_{I + J}(M) = H^i_{I + J'}(M) = H^i_{J'}(M_2)$
は $I$ のある冪で零化され，したがって導来完備である．よってこの場合にも
主張が従う．$N$ については，いま述べた二つの議論のいずれかを使える．
短完全列
$0 \to M' \to M \to M/M' \to 0$
を考えると，補題を $M/M'$ について証明すれば十分であることが分かる．
したがって $\text{Ass}(M) \cap (V(I) \cup T) = \emptyset$ と仮定してよい．

\medskip\noindent
$\mathfrak p \in \text{Ass}(M)$ で
$V(\mathfrak p) \cap V(J) \cap V(I) = \emptyset$ を満たすものを取る．
$I$ は $A$ の Jacobson 根基に含まれるので，
$V(\mathfrak p) \cap V(J') = \emptyset$ が，$J' \subset J$ で
$V(J') \cap V(I) = V(J) \cap V(I)$ を満たす任意のイデアルについて成り立つ．
そこで $N = H^0_\mathfrak p(M)$ とおけば，
$R\Gamma_J(N) = R\Gamma_{J'}(N) = 0$ が，$J' \subset J$ で
$V(J') \cap V(I) = V(J) \cap V(I)$ を満たすすべてのイデアルについて成り立つ．
特に $R\Gamma_J(N)^\wedge = 0$ である．したがって $M$ を $M/N$ で
置き換えてよい．というのも，これは $M$ の構造を条件 (4) または (5) に
関与しない素イデアルにおいてのみ変えるからである．これを繰り返せば，
$V(\mathfrak p) \cap V(J) \cap V(I) \not = \emptyset$ が，すべての
$\mathfrak p \in \text{Ass}(M)$ について成り立つと仮定してよい．

\medskip\noindent
$\text{Ass}(M) \cap (V(I) \cup T) = \emptyset$ と仮定し，さらに
$V(\mathfrak p) \cap V(J) \cap V(I) \not = \emptyset$ がすべての
$\mathfrak p \in \text{Ass}(M)$ について成り立つと仮定する．
$\mathfrak p \in \text{Ass}(M)$ を取る．補題
\ref{algebraization-lemma-kill-completion-general} を適用できることを示したい．
この検証において補足条件 (5) を用いる．$\mathfrak p' \subset \mathfrak p$
および $\mathfrak q' \subset V(\mathfrak p) \cap V(J) \cap V(I)$ を選ぶ．
\begin{enumerate}
\item $M_{\mathfrak p'} = 0$ ならば，
$\text{depth}(M_{\mathfrak p'}) = \infty$ であり，
$\text{depth}(M_{\mathfrak p'}) +
\dim((A/\mathfrak p')_{\mathfrak q'}) > d + s$.
\item $\text{depth}(M_{\mathfrak p'}) < s$ ならば，(4) により
$\text{depth}(M_{\mathfrak p'}) +
\dim((A/\mathfrak p')_{\mathfrak q'}) > d + s$ である．
\end{enumerate}
残る場合には $M_{\mathfrak p'} \not = 0$ かつ
$\text{depth}(M_{\mathfrak p'}) \geq s$ である．特に $\mathfrak p'$ は
$M$ の台に属し，$\mathfrak p'' \subset \mathfrak p'$ かつ
$\mathfrak p'' \in \text{Ass}(M)$ となるものを選べる．
\begin{enumerate}
\item[(a)] Algebra, 補題 \ref{algebra-lemma-depth-dim-associated-primes} により
$\dim((A/\mathfrak p'')_{\mathfrak p'}) \geq \text{depth}(M_{\mathfrak p'})$
であることに注意する．等号が成り立つならば，
$$
\text{depth}(M_{\mathfrak p'}) + \dim((A/\mathfrak p')_{\mathfrak q'}) =
\text{depth}(M_{\mathfrak p''}) + \dim((A/\mathfrak p'')_{\mathfrak q'})
> s + d
$$
を得る．これは $\mathfrak p''$ に (4) を適用した結果であり，これでよい．
したがって問題が起こり得るのは
$\dim((A/\mathfrak p'')_{\mathfrak p'}) > \text{depth}(M_{\mathfrak p'})$.
の場合だけである．これは $\dim(M_\mathfrak p) \geq s + 2$ を導く．
ゆえに (5)(a) が成り立つならば，この場合は起こらない．
\item[(b)] (5)(b) が成り立つならば，
$$
\text{depth}(M_{\mathfrak p'}) + \dim((A/\mathfrak p')_{\mathfrak q'})
\geq s + \delta(\mathfrak p') - \delta(\mathfrak q')
\geq s + 1 + \delta(\mathfrak p) - \delta_{max}
> s + d
$$
となり，所望の結論を得る．
\item[(c)] (5)(c) が成り立つならば，
\begin{align*}
\text{depth}(M_{\mathfrak p'}) + \dim((A/\mathfrak p')_{\mathfrak q'})
& \geq
s + \delta(\mathfrak p') - \delta(\mathfrak q') \\
& \geq
s + 1 + \delta(\mathfrak p) - \delta(\mathfrak q') \\
& =
s + 1 + \delta(\mathfrak p) - \delta(\mathfrak q) +
\delta(\mathfrak q) - \delta(\mathfrak q') \\
& >
s + 1 + (s + d + \delta_{max} - \delta_{min} - 2) +
\delta(\mathfrak q) - \delta(\mathfrak q') \\
& \geq 
2s + d - 1 \geq s + d
\end{align*}
となり，やはり所望の結論を得る．この議論が成り立つのは，素イデアル
$\mathfrak q \in V(\mathfrak p) \cap V(J) \cap V(I)$ が存在することを
知っているからである．
\end{enumerate}
これで帰納段階に進む準備が整った．

\medskip\noindent
補題 \ref{algebraization-lemma-kill-colimit-weak-general} にあるようなイデアル $J_0$ と，
整数 $t > 0$ であって，$(J_0I)^t$ が $H^s_J(M)$ を零化するものを選ぶ．
補題 \ref{algebraization-lemma-kill-completion-general} の仮定は，すべての
$\mathfrak p \in \text{Ass}(M)$ について満たされている
（直前の段落を参照）．したがって，零化イデアル $\mathfrak a \subset A$，すなわち
$H^s(R\Gamma_J(M)^\wedge)$ の零化イデアルは，$\mathfrak p$ に含まれない．
ここで $\mathfrak p \in \text{Ass}(M)$ である．ゆえに
$f \in \mathfrak a(J_0I)^t$ であって，$M$ のどの付随素イデアルにも属さず，
$H^s(R\Gamma_J(M)^\wedge)$ と $H^s_J(M)$ の両方を零化するものを見いだせる．
この $f$ は $M$ 上の非零因子なので，
短完全列
$$
0 \to M \xrightarrow{f} M \to M/fM \to 0
$$
を考えられる．$f$ の選択により，
$$
\xymatrix{
H^{s - 1}_{J'}(M) \ar[d] \ar[r] &
H^{s - 1}_{J'}(M/fM) \ar[d] \ar[r] &
H^s_{J'}(M) \ar[d] \ar[r] & 0 \\
H^{s - 1}(R\Gamma_J(M)^\wedge) \ar[r] &
H^{s - 1}(R\Gamma_J(M/fM)^\wedge) \ar[r] &
H^s(R\Gamma_J(M)^\wedge) \ar[r] & 0
}
$$
$J' \subset J_0$ で $V(J') \cap V(I) = V(J) \cap V(I)$ を満たす任意の
イデアルに対してこの図を得る．したがって，$J'$ を
$M$，$M/fM$，$s - 1$ に対して機能するように選べば（帰納法の仮定により
可能である；次の段落を参照），補題が成り立つと結論できる．

\medskip\noindent
証明を完了するには，加群 $M/fM$ が $s - 1$ に対する仮定 (4) と (5) を
満たすことを示さなければならない．そこで素イデアル $\mathfrak p$ であって，
$M/fM$ の台に属し，$\text{depth}((M/fM)_\mathfrak p) < s - 1$ かつ
$V(\mathfrak p) \cap V(J) \cap V(I)$ が空でないものを取る．このとき
$\dim(M_\mathfrak p) = \dim((M/fM)_\mathfrak p) + 1$ かつ
$\text{depth}(M_\mathfrak p) = \text{depth}((M/fM)_\mathfrak p) + 1$ である．
特に，(4) と (5) が $\mathfrak p$ および $M$ について元の値 $s$ で
成り立つことが分かっている．所望の不等式はこれらを見れば直ちに従う．
\end{proof}

\begin{example}
\label{algebraization-example-no-ML}
補題 \ref{algebraization-lemma-algebraize-local-cohomology-general} において，逆系
$H^i_J(M/I^nM)$ が Mittag--Leffler 条件を満たすかどうかは分からない．
例えば $A = \mathbf{Z}_p[[x, y]]$，$I = (p)$，$J = (p, x)$，
$M = A/(xy - p)$ とする．このとき
$H^0_J(M/p^nM) \to H^0_J(M/pM)$
の像は，$y^n$ が $M/pM = A/(p, xy)$ において生成するイデアルである．
\end{example}





\section{局所コホモロジーの代数化 II}
\label{algebraization-section-algebraization-punctured}

\noindent
この節では，節 \ref{algebraization-section-algebraization-sections-general} の議論を，
$(A, \mathfrak m)$ が局所環で，局所コホモロジー $R\Gamma_\mathfrak m$ を
$\mathfrak m$ に関して取る場合についてやり直す．前と同様，
主な道具は次の補題である．

\begin{lemma}
\label{algebraization-lemma-kill-completion}
$(A, \mathfrak m)$ を Noether 局所環とする．
$I \subset A$ をイデアルとする．$M$ を有限 $A$ 加群とし，
$\mathfrak p \subset A$ を素イデアルとする．$s$ および $d$ を整数とする．
次を仮定する．
\begin{enumerate}
\item $A$ は双対化複体をもつ．
\item $\text{cd}(A, I) \leq d$ である．
\item
$\text{depth}_{A_\mathfrak p}(M_\mathfrak p) + \dim(A/\mathfrak p) > d + s$
である．
\end{enumerate}
このとき，$f \in A \setminus \mathfrak p$ であって，
$H^i(R\Gamma_\mathfrak m(M)^\wedge)$ を $i \leq s$ に対して零化するものが
存在する．ここで ${}^\wedge$ は
$I$ 進完備化を表す．
\end{lemma}

\begin{proof}
Local Cohomology, 補題
\ref{local-cohomology-lemma-sitting-in-degrees}
によれば，関数
$$
\mathfrak p' \longmapsto
\text{depth}_{A_{\mathfrak p'}}(M_{\mathfrak p'}) + \dim(A/\mathfrak p')
$$
は $\Spec(A)$ 上で下半連続である．したがって，$\mathfrak p' \subset \mathfrak p$
におけるこの関数の値は $> s + d$ である．ゆえに
$\mathfrak p \not = \mathfrak m$ ならば，本補題は補題
\ref{algebraization-lemma-kill-completion-general} の特別な場合である．
$\mathfrak p = \mathfrak m$ ならば，$H^i_\mathfrak m(M) = 0$ が
$i \leq s + d$ に対して成り立つ．これは深さと局所コホモロジーの関係
（Dualizing Complexes, 補題 \ref{dualizing-lemma-depth}）による．したがって補題
\ref{algebraization-lemma-kill-completion-general} の証明で与えた議論から，この（退化した）
場合には $H^i(R\Gamma_\mathfrak m(M)^\wedge) = 0$ が $i \leq s$ に対して従う．
\end{proof}

\begin{lemma}
\label{algebraization-lemma-kill-colimit-weak}
$(A, \mathfrak m)$ を Noether 局所環とする．
$I \subset A$ をイデアルとし，$M$ を有限 $A$ 加群とする．
$s$ および $d$ を整数とする．次を仮定する．
\begin{enumerate}
\item $A$ は双対化複体をもつ．
\item $\mathfrak p \in V(I)$ ならば，条件を課さない．
\item $\mathfrak p \not \in V(I)$ かつ
$V(\mathfrak p) \cap V(I) = \{\mathfrak m\}$ ならば，
$\dim(A/\mathfrak p) \leq d$ である．
\item $\mathfrak p \not \in V(I)$ かつ
$V(\mathfrak p) \cap V(I) \not = \{\mathfrak m\}$ ならば，
$$
\text{depth}_{A_\mathfrak p}(M_\mathfrak p) \geq s
\quad\text{または}\quad
\text{depth}_{A_\mathfrak p}(M_\mathfrak p) + \dim(A/\mathfrak p) > d + s
$$
\end{enumerate}
このとき，イデアル $J_0 \subset A$ であって
$V(J_0) \cap V(I) = \{\mathfrak m\}$ を満たすものが存在し，任意の
$J \subset J_0$ で $V(J) \cap V(I) = \{\mathfrak m\}$ を満たすものに対して写像
$$
R\Gamma_J(M) \longrightarrow R\Gamma_{J_0}(M)
$$
は次数 $\leq s$ のコホモロジー上で同型を誘導し，さらにこれらの加群は
$J_0I$ のある冪で零化される．
\end{lemma}

\begin{proof}
これは補題 \ref{algebraization-lemma-kill-colimit-weak-general} の特別な場合である．
\end{proof}

\begin{lemma}
\label{algebraization-lemma-kill-colimit}
補題 \ref{algebraization-lemma-kill-colimit-weak} において，空の条件 (2) の代わりに
\begin{enumerate}
\item[(2')] $\mathfrak p \in V(I)$ かつ $\mathfrak p \not = \mathfrak m$ ならば，
$\text{depth}_{A_\mathfrak p}(M_\mathfrak p) + \dim(A/\mathfrak p) > s$ である．
\end{enumerate}
を仮定すれば，これらの条件はさらに $H^i_{J_0}(M)$ が有限 $A$ 加群であることを
$i \leq s$ に対して導く．
\end{lemma}

\begin{proof}
これは補題 \ref{algebraization-lemma-kill-colimit-general} の特別な場合である．
\end{proof}

\begin{lemma}
\label{algebraization-lemma-kill-colimit-support}
補題 \ref{algebraization-lemma-kill-colimit-weak} において，さらに
\begin{enumerate}
\item[(6)] $\mathfrak p \not \in V(I)$ かつ
$V(\mathfrak p) \cap V(I) = \{\mathfrak m\}$ ならば，
$\text{depth}_{A_\mathfrak p}(M_\mathfrak p) > s$ である．
\end{enumerate}
を仮定すれば，$H^i_{J_0}(M) = H^i_J(M) = H^i_\mathfrak m(M)$ が
$i \leq s$ に対して成り立ち，これらの加群は
$I$ のある冪で零化される．
\end{lemma}

\begin{proof}
これは補題 \ref{algebraization-lemma-kill-colimit-support-general} の特別な場合である．
\end{proof}

\begin{lemma}
\label{algebraization-lemma-algebraize-local-cohomology}
$(A, \mathfrak m)$ を Noether 局所環とする．
$I \subset A$ をイデアルとし，$M$ を有限 $A$ 加群とする．
$s$ および $d$ を整数とする．次を仮定する．
\begin{enumerate}
\item $A$ は $I$ 進完備で，双対化複体をもつ．
\item $\mathfrak p \in V(I)$ ならば，条件を課さない．
\item $\text{cd}(A, I) \leq d$ である．
\item $\mathfrak p \not \in V(I)$ かつ
$V(\mathfrak p) \cap V(I) \not = \{\mathfrak m\}$ ならば，
$$
\text{depth}_{A_\mathfrak p}(M_\mathfrak p) \geq s
\quad\text{または}\quad
\text{depth}_{A_\mathfrak p}(M_\mathfrak p) + \dim(A/\mathfrak p) > d + s
$$
\end{enumerate}
このとき，イデアル $J_0 \subset A$ であって
$V(J_0) \cap V(I) = \{\mathfrak m\}$ を満たすものが存在し，任意の
$J \subset J_0$ で $V(J) \cap V(I) = \{\mathfrak m\}$ を満たすものに対して写像
$$
R\Gamma_J(M) \longrightarrow
R\Gamma_J(M)^\wedge = R\Gamma_\mathfrak m(M)^\wedge
$$
は次数 $\leq s$ のコホモロジー上で同型を誘導する．ここで ${}^\wedge$ は
導来 $I$ 進完備化を表す．
\end{lemma}

\begin{proof}
本補題は補題 \ref{algebraization-lemma-algebraize-local-cohomology-general} の特別な場合である．
実際，$\delta_{max} = \delta_{min} = \delta(\mathfrak m)$ なので，条件 (4) は
条件 (5)(c) を導く．この重要な特別の場合は幾分容易である（確認事項が少ない）
ため，ここで証明を与える．

\medskip\noindent
現在の状況では $R\Gamma_\mathfrak a$ と $R\Gamma_{V(\mathfrak a)}$ の間に
違いはない．Dualizing Complexes, 補題
\ref{dualizing-lemma-local-cohomology-noetherian} を参照されたい．次に，
$$
R\Gamma_\mathfrak m(M)^\wedge =
R\Gamma_I(R\Gamma_J(M))^\wedge =
R\Gamma_J(M)^\wedge
$$
に注意する．これは Dualizing Complexes, 補題
\ref{dualizing-lemma-local-cohomology-ss} および
\ref{dualizing-lemma-complete-and-local} によるもので，補題の主張にある等号を
説明する．

\medskip\noindent
補題は $s < 0$ に対して成り立つことに注意する．これは自明な場合ではない．
実際，$H^s(R\Gamma_\mathfrak m(M)^\wedge)$ が負の $s$ に対して零であることは
先験的には明らかでない．しかしこの消滅は補題
\ref{algebraization-lemma-local-cohomology-derived-completion} で示されている．
$s \geq 0$ に関して帰納法で補題を証明する．

\medskip\noindent
補題 \ref{algebraization-lemma-kill-colimit-weak} の仮定は，Local Cohomology, 補題
\ref{local-cohomology-lemma-cd-bound-dim-local} により満たされる．
$s = 0$ の場合は，補題 \ref{algebraization-lemma-kill-colimit-weak} および
Dualizing Complexes, 補題 \ref{dualizing-lemma-completion-local-H0} から従う．

\medskip\noindent
$s > 0$ とし，$s$ より小さい値について補題が成り立つとする．
$M' \subset M$ を，その元の台が $V(I) \cup V(J)$ に含まれるような部分加群と
する．ここで，あるイデアル $J$ は $V(J) \cap V(I) = \{\mathfrak m\}$ を
満たすものとする．このとき $M'$ は有限 $A$ 加群である．次を主張する：
$$
R\Gamma_J(M') \to R\Gamma_\mathfrak m(M')^\wedge
$$
これは $J$ のどの選択に対しても同型である．実際，このような加群に対しては
短完全列 $0 \to M_1 \oplus M_2 \to M' \to N \to 0$ であって，$M_1$ は
$J$ のある冪で零化され，$M_2$ は $I$ のある冪で零化され，$N$ は
$\mathfrak m$ のある冪で零化されるものが存在する．$M_1$ の場合には
$R\Gamma_J(M_1) = M_1$ である．また $M_1$ は有限 $A$ 加群かつ $I$ 進完備なので
$M_1^\wedge = M_1$ である．よって $M_1$ に対する主張が成り立つ．
$M_2$ の場合には，$H^i_J(M_2)$ は $I$ のある冪で零化され，したがって導来完備で
ある．よって $M_2$ に対する主張も成り立つ．同じ議論により $N$ に対する主張も
成り立ち，主張全体が従う．短完全列 $0 \to M' \to M \to M/M' \to 0$ を
考えると，補題を $M/M'$ について証明すれば十分であることが分かる．したがって，
$\mathfrak p \in \text{Ass}(M)$ ならば
$V(\mathfrak p) \cap V(I) \not = \{\mathfrak m\}$，すなわち
$\mathfrak p$ は (4) に現れる素イデアルであると仮定してよい．

\medskip\noindent
補題 \ref{algebraization-lemma-kill-colimit-weak} にあるようなイデアル $J_0$ と，
整数 $t > 0$ であって $(J_0I)^t$ が $H^s_J(M)$ を零化するものを選ぶ．
ここで $J$ は，$J \subset J_0$ かつ
$V(J) \cap V(I) = \{\mathfrak m\}$ を満たす任意のイデアルを表す．
補題 \ref{algebraization-lemma-kill-completion} の仮定は，すべての
$\mathfrak p \in \text{Ass}(M)$ について満たされる（直前の段落を参照）．
したがって，零化イデアル $\mathfrak a \subset A$，すなわち
$H^s(R\Gamma_\mathfrak m(M)^\wedge)$ の零化イデアルは，$\mathfrak p$ に含まれない．ここで
$\mathfrak p \in \text{Ass}(M)$ である．ゆえに
$f \in \mathfrak a(J_0I)^t$ であって，$M$ のどの付随素イデアルにも属さず，
$H^s(R\Gamma_\mathfrak m(M)^\wedge)$ と $H^s_J(M)$ の両方を零化するものを
見いだせる．この $f$ は $M$ 上の非零因子なので，短完全列
$$
0 \to M \xrightarrow{f} M \to M/fM \to 0
$$
を考えられる．$f$ の選択により，
$$
\xymatrix{
H^{s - 1}_J(M) \ar[d] \ar[r] &
H^{s - 1}_J(M/fM) \ar[d] \ar[r] &
H^s_J(M) \ar[d] \ar[r] & 0 \\
H^{s - 1}(R\Gamma_\mathfrak m(M)^\wedge) \ar[r] &
H^{s - 1}(R\Gamma_\mathfrak m(M/fM)^\wedge) \ar[r] &
H^s(R\Gamma_\mathfrak m(M)^\wedge) \ar[r] & 0
}
$$
$J \subset J_0$ で $V(J) \cap V(I) = \{\mathfrak m\}$ を満たす任意の
イデアルに対してこの図を得る．したがって，$J$ を $M$，$M/fM$，
$s - 1$ に対して機能するように選べば（帰納法の仮定により可能である），
補題が成り立つと結論できる．
\end{proof}








\section{局所コホモロジーの代数化 III}
\label{algebraization-section-bootstrap}

\noindent
この節では，節 \ref{algebraization-section-algebraization-sections-general} および
\ref{algebraization-section-algebraization-punctured} の内容をブートストラップし，次の状況で
より強い結果を得る．

\begin{situation}
\label{algebraization-situation-bootstrap}
ここで $A$ は Noether 環である．イデアル $I \subset A$，有限 $A$ 加群 $M$，
および特殊化で安定な部分集合 $T \subset V(I)$ が与えられている．整数 $s$ と
$d$ が与えられている．次を仮定する．
\begin{enumerate}
\item[(1)] $A$ は双対化複体をもつ．
\item[(3)] $\text{cd}(A, I) \leq d$ である．
\item[(4)] 素イデアル $\mathfrak p \subset \mathfrak r \subset \mathfrak q$ が
$\mathfrak p \not \in V(I)$，$\mathfrak r \in V(I) \setminus T$，
$\mathfrak q \in T$ を満たすならば，
$$
\text{depth}_{A_\mathfrak p}(M_\mathfrak p) \geq s
\quad\text{または}\quad
\text{depth}_{A_\mathfrak p}(M_\mathfrak p) +
\dim((A/\mathfrak p)_\mathfrak q) > d + s
$$
\item[(6)] $\mathfrak q \in T$ が与えられたとき，
$A', \mathfrak m', I', M'$ がそれぞれ
$I$ 進完備化，すなわち
$A_\mathfrak q, \mathfrak qA_\mathfrak q, I_\mathfrak q, M_\mathfrak q$ の通常の
完備化を表すならば，
$$
\text{depth}(M'_{\mathfrak p'}) > s
$$
が，すべての $\mathfrak p' \in \Spec(A') \setminus V(I')$ で
$V(\mathfrak p') \cap V(I') = \{\mathfrak m'\}$ を満たすものに対して成り立つ．
\end{enumerate}
\end{situation}

\noindent
次の補題は，状況 \ref{algebraization-situation-bootstrap} において，実際の仮定には
(4) の素イデアルの三つ組 $\mathfrak p \subset \mathfrak r \subset \mathfrak q$
だけを調べれば十分である理由を説明する．実際の仮定には $\mathfrak p$ と
$\mathfrak q$ しか現れない．

\begin{lemma}
\label{algebraization-lemma-helper-bootstrap}
状況 \ref{algebraization-situation-bootstrap} において，$\mathfrak p \subset \mathfrak q$ を
$A$ の素イデアルで，$\mathfrak p \not \in V(I)$ かつ $\mathfrak q \in T$ を
満たすものとする．$\mathfrak r \in V(I) \setminus T$ で
$\mathfrak p \subset \mathfrak r \subset \mathfrak q$ を満たすものが存在しないならば，
$\text{depth}(M_\mathfrak p) > s$ である．
\end{lemma}

\begin{proof}
$\mathfrak q' \in T$ で，$\mathfrak p \subset \mathfrak q' \subset \mathfrak q$ を
満たし，$T$ に属する素イデアルが $\mathfrak p$ と $\mathfrak q'$ の間に真に
存在しないものを選ぶ．補題を証明するため，$\mathfrak q$ を $\mathfrak q'$ で
置き換えてよく，そうする．次に，$\mathfrak p' \subset A_\mathfrak q$ を
$\mathfrak p$ に対応する素イデアルとする．このとき
$V(\mathfrak p') \cap V(IA_\mathfrak q) = \{\mathfrak q A_\mathfrak q\}$
である．これは補題にあるような素イデアル $\mathfrak r$ が存在しないからである．
$A', I', \mathfrak m', M'$ を，それぞれ
$I$ 進完備化，すなわち
$A_\mathfrak q, I_\mathfrak q, \mathfrak qA_\mathfrak q, M_\mathfrak q$ の完備化とする．
$A_\mathfrak q \to A'$ は忠実平坦なので（Algebra, 補題
\ref{algebra-lemma-completion-faithfully-flat}），$\mathfrak p'' \subset A'$ で
$\mathfrak p'$ の上にあり，
$\dim(A'_{\mathfrak p''}/\mathfrak p' A'_{\mathfrak p''}) = 0$ を満たすものを
選べる．このとき
$$
\text{depth}(M'_{\mathfrak p''}) =
\text{depth}((M_\mathfrak q \otimes_{A_\mathfrak q} A')_{\mathfrak p''}) =
\text{depth}(M_\mathfrak p \otimes_{A_\mathfrak p} A'_{\mathfrak p''}) =
\text{depth}(M_\mathfrak p)
$$
を得る．これは $A \to A'$ の平坦性と $\mathfrak p''$ の選び方による．
Algebra, 補題 \ref{algebra-lemma-apply-grothendieck-module} を参照されたい．
$\mathfrak p''$ は $\mathfrak p'$ の上にあるので，
$V(\mathfrak p'') \cap V(I') = \{\mathfrak m'\}$ である．したがって状況
\ref{algebraization-situation-bootstrap} の条件 (6) から
$\text{depth}(M'_{\mathfrak p''}) > s$ が従い，証明が完了する．
\end{proof}

\noindent
次の煩雑な補題は，課し得るさまざまな条件の組の間の関係を説明する．

\begin{lemma}
\label{algebraization-lemma-bootstrap-inherited}
状況 \ref{algebraization-situation-bootstrap} において，次が成り立つ．
\begin{enumerate}
\item[(E)] $T' \subset T$ がより小さい特殊化で安定な部分集合ならば，
$A, I, T', M$ は状況 \ref{algebraization-situation-bootstrap} の仮定を満たす．
\item[(F)] $S \subset A$ が乗法的部分集合ならば，
$S^{-1}A, S^{-1}I, T', S^{-1}M$
は状況 \ref{algebraization-situation-bootstrap} の仮定を満たす．ここで
$T' \subset V(S^{-1}I)$ は $T$ の逆像である．
\item[(G)] 四つ組 $A', I', T', M'$ は状況 \ref{algebraization-situation-bootstrap} の仮定を
満たす．ここで $A', I', M'$ は通常の $I$ 進完備化，すなわち
$A, I, M$ の完備化であり，
$T' \subset V(I')$ は $T$ の逆像である．
\end{enumerate}
$I \subset \mathfrak a \subset A$ を，$V(\mathfrak a) \subset T$ を満たす
イデアルとする．このとき次が成り立つ．
\begin{enumerate}
\item[(A)] $I$ が $A$ の Jacobson 根基に含まれるならば，補題
\ref{algebraization-lemma-kill-colimit-weak-general} および
\ref{algebraization-lemma-kill-colimit-support-general} のすべての仮定が
$A, I, \mathfrak a, M$ に対して満たされる．
\item[(B)] $A$ が $I$ に関して完備ならば，補題
\ref{algebraization-lemma-algebraize-local-cohomology-general} の仮定のうち，おそらく (5) を
除くすべてが $A, I, \mathfrak a, M$ に対して満たされる．
\item[(C)] $A$ が極大イデアル $\mathfrak m = \mathfrak a$ をもつ局所環ならば，
補題 \ref{algebraization-lemma-kill-colimit-weak} および \ref{algebraization-lemma-kill-colimit-support} の
すべての仮定が $A, \mathfrak m, I, M$ に対して成り立つ．
\item[(D)] $A$ が極大イデアル $\mathfrak m = \mathfrak a$ をもつ局所環で，
$I$ 進完備ならば，補題 \ref{algebraization-lemma-algebraize-local-cohomology} のすべての仮定が
$A, \mathfrak m, I, M$ に対して成り立つ．
\end{enumerate}
\end{lemma}

\begin{proof}
（E）の証明．状況 \ref{algebraization-situation-bootstrap} の仮定 (1)，(3)，(4)，(6) が
$A, I, T, M$ に対して成り立つことを示さなければならない．$T$ を $T'$ に
縮小すると，仮定 (6) は弱くなり，仮定 (4) は強くなる．しかし，
$\mathfrak p \subset \mathfrak r \subset \mathfrak q$ が
$\mathfrak p \not \in V(I)$，$\mathfrak r \in V(I) \setminus T'$，
$\mathfrak q \in T'$ を満たすものが，$A, I, T', M$ に対する仮定 (4) に現れると
する．このとき，$\mathfrak r \in V(I) \setminus T$ を選べる場合には
$A, I, T, M$ に対する条件 (4) が適用される．そのような $\mathfrak r$ を
見いだせない場合には，補題 \ref{algebraization-lemma-helper-bootstrap} により
$\text{depth}(M_\mathfrak p) > s$ を得る．これで所望のように (4) が
$A, I, T', M$ に対して成り立つことが示された．

\medskip\noindent
（F）の証明．これは直接的なので，詳細を省略する．

\medskip\noindent
（G）の証明．状況 \ref{algebraization-situation-bootstrap} の仮定 (1)，(3)，(4)，(6) が
$I$ 進完備化 $A', I', T', M'$ に対して成り立つことを示さなければならない．
$\Spec(A') \to \Spec(A)$ が同型 $V(I') \to V(I)$ を誘導することに注意する．

\medskip\noindent
仮定 (1)：環 $A'$ は双対化複体をもつ．Dualizing Complexes, 補題
\ref{dualizing-lemma-ubiquity-dualizing} を参照されたい．

\medskip\noindent
仮定 (3)：$I' = IA'$ なので，これは Local Cohomology, 補題
\ref{local-cohomology-lemma-cd-change-rings} から従う．

\medskip\noindent
仮定 (4)：素イデアル $\mathfrak p' \subset \mathfrak r' \subset \mathfrak q'$ が $A'$ で
が $\mathfrak p' \not \in V(I')$，$\mathfrak r' \in V(I') \setminus T'$，
$\mathfrak q' \in T'$ を満たすとする．その像
$\mathfrak p \subset \mathfrak r \subset \mathfrak q$ は $A$ のスペクトルにあり，
$\mathfrak p \not \in V(I)$, $\mathfrak r \in V(I) \setminus T$,
$\mathfrak q \in T$ を満たす．したがって
$$
\text{depth}_{A_\mathfrak p}(M_\mathfrak p) \geq s
\quad\text{または}\quad
\text{depth}_{A_\mathfrak p}(M_\mathfrak p) +
\dim((A/\mathfrak p)_\mathfrak q) > d + s
$$
が，$A, I, T, M$ に対する仮定 (4) により成り立つ．また
$\text{depth}(M'_{\mathfrak p'}) \geq \text{depth}(M_\mathfrak p)$ かつ
$\text{depth}(M'_{\mathfrak p'}) +
\dim((A'/\mathfrak p')_{\mathfrak q'}) =
\text{depth}(M_\mathfrak p) +
\dim((A/\mathfrak p)_\mathfrak q)$
である．これは Local Cohomology, 補題
\ref{local-cohomology-lemma-change-completion} による．ゆえに仮定 (4) は
$A', I', T', M'$ に対して成り立つ．

\medskip\noindent
仮定 (6)：$\mathfrak q' \in T'$ を，素イデアル $\mathfrak q \in T$ の上に
あるものとする．このとき $A'_{\mathfrak q'}$ と $A_\mathfrak q$ は
同型な $I$ 進完備化をもち，$M_\mathfrak q$ と $M'_{\mathfrak q'}$ についても
同様である．したがって，$A', I', T', M'$ に対する仮定 (6) は
$A, I, T, M$ に対する仮定 (6) と同値である．

\medskip\noindent
（A）の証明．補題 \ref{algebraization-lemma-kill-colimit-weak-general} および
\ref{algebraization-lemma-kill-colimit-support-general} の条件 (1)，(2)，(3)，(4)，(6) を
$(A, I, \mathfrak a, M)$ について確認しなければならない．注意：これらの補題の
主張に現れる集合 $T$ は，上の集合 $T$ と同じではない．

\medskip\noindent
条件 (1)：状況 \ref{algebraization-situation-bootstrap} において $A$ は双対化複体をもつと
仮定したので，これは成り立つ．

\medskip\noindent
条件 (2)：これは空である．

\medskip\noindent
条件 (3)：$\mathfrak p \subset A$ が
$V(\mathfrak p) \cap V(I) \subset V(\mathfrak a)$ を満たすとする．
$I$ は $A$ の Jacobson 根基に含まれるので，
$V(\mathfrak p) \cap V(I) \not = \emptyset$ である．
$\mathfrak q \in V(\mathfrak p) \cap V(I)$ を一般点とする．
$\text{cd}(A_\mathfrak q, I_\mathfrak q) \leq d$ であり（Local Cohomology, 補題
\ref{local-cohomology-lemma-cd-local}），さらに
$V(\mathfrak p A_\mathfrak q) \cap V(I_\mathfrak q) =
\{\mathfrak q A_\mathfrak q\}$ なので，
$\dim((A/\mathfrak p)_\mathfrak q) \leq d$ を得る．これは Local Cohomology, 補題
\ref{local-cohomology-lemma-cd-bound-dim-local} による．これで (3) が示された．

\medskip\noindent
条件 (4)：$\mathfrak p \not \in V(I)$ かつ
$\mathfrak q \in V(\mathfrak p) \cap V(\mathfrak a)$ とする．次を示せば十分である：
$$
\text{depth}_{A_\mathfrak p}(M_\mathfrak p) \geq s
\quad\text{または}\quad
\text{depth}_{A_\mathfrak p}(M_\mathfrak p) +
\dim((A/\mathfrak p)_\mathfrak q) > d + s
$$
素イデアル $\mathfrak p \subset \mathfrak r \subset \mathfrak q$ で
$\mathfrak r \in V(I) \setminus T$ を満たすものが存在するならば，これは状況
\ref{algebraization-situation-bootstrap} の仮定 (4) から直ちに従う．存在しないならば，
$\text{depth}(M_\mathfrak p) > s$ が補題 \ref{algebraization-lemma-helper-bootstrap} から従う．

\medskip\noindent
条件 (6)：$\mathfrak p \not \in V(I)$ が
$V(\mathfrak p) \cap V(I) \subset V(\mathfrak a)$ を満たすとする．$I$ は $A$ の
Jacobson 根基に含まれるので，$V(\mathfrak p) \cap V(I) \not = \emptyset$ である．
$\mathfrak q \in V(\mathfrak p) \cap V(I) \subset V(\mathfrak a)$ を選ぶ．
素イデアル $\mathfrak p \subset \mathfrak r \subset \mathfrak q$ で
$\mathfrak r \in V(I) \setminus T$ を満たすものが存在しないことは明らかである．
補題 \ref{algebraization-lemma-helper-bootstrap} により $\text{depth}(M_\mathfrak p) > s$ であり，
これで (6) が示された．

\medskip\noindent
（B）の証明．補題 \ref{algebraization-lemma-algebraize-local-cohomology-general} の条件
(1)，(2)，(3)，(4) を確認しなければならない．注意：この補題の主張に現れる
集合 $T$ は，上の集合 $T$ と同じではない．

\medskip\noindent
条件 (1)：$A$ は完備で双対化複体をもつので，これは成り立つ．

\medskip\noindent
条件 (2)：これは空である．

\medskip\noindent
条件 (3)：これは状況 \ref{algebraization-situation-bootstrap} の仮定 (3) と同じである．

\medskip\noindent
条件 (4)：これは（A）で示した補題 \ref{algebraization-lemma-kill-colimit-weak-general} の
仮定 (4) と同じである．

\medskip\noindent
（C）の証明．局所の場合には，補題 \ref{algebraization-lemma-kill-colimit-weak} および
\ref{algebraization-lemma-kill-colimit-support} の仮定は，補題
\ref{algebraization-lemma-kill-colimit-weak-general} および
\ref{algebraization-lemma-kill-colimit-support-general} の仮定と同じであり，これらが成り立つことを
（A）で示したからである．

\medskip\noindent
（D）の証明．補題 \ref{algebraization-lemma-algebraize-local-cohomology} の仮定は，補題
\ref{algebraization-lemma-algebraize-local-cohomology-general} の仮定 (1)，(2)，(3)，(4) と
同じであり，これらが成り立つことを（B）で示したからである．
\end{proof}

\begin{lemma}
\label{algebraization-lemma-algebraize-local-cohomology-bis}
状況 \ref{algebraization-situation-bootstrap} において，$A$ は極大イデアル $\mathfrak m$ をもつ
局所環であり，$T = \{\mathfrak m\}$ と仮定する．このとき
$H^i_\mathfrak m(M) \to \lim H^i_\mathfrak m(M/I^nM)$
は $i \leq s$ に対して同型であり，これらの加群は $I$ のある冪で零化される．
\end{lemma}

\begin{proof}
$A', I', \mathfrak m', M'$ を，通常の $I$ 進完備化，すなわち
$A, I, \mathfrak m, M$ の完備化とする．次の等式を思い出そう：
$H^i_\mathfrak m(M) \otimes_A A' = H^i_{\mathfrak m'}(M')$
これは $A \to A'$ の平坦性と Dualizing Complexes, 補題
\ref{dualizing-lemma-torsion-change-rings} による．
$H^i_\mathfrak m(M)$ は $\mathfrak m$ の冪で捩れなので，
$H^i_\mathfrak m(M) = H^i_\mathfrak m(M) \otimes_A A'$ である．More on Algebra,
補題 \ref{more-algebra-lemma-neighbourhood-equivalence} を参照されたい．したがって
$H^i_\mathfrak m(M) = H^i_{\mathfrak m'}(M')$ を得る．まったく同じ議論から
$H^i_\mathfrak m(M/I^nM) =  H^i_{\mathfrak m'}(M'/(I')^nM')$
がすべての $n$ および $i$ に対して成り立つ．

\medskip\noindent
補題 \ref{algebraization-lemma-algebraize-local-cohomology}，
\ref{algebraization-lemma-kill-colimit-weak}，および
\ref{algebraization-lemma-kill-colimit-support}
は，補題 \ref{algebraization-lemma-bootstrap-inherited} の (C) と (D) により
$A', \mathfrak m', I', M'$ に適用できる．したがって，同型
$$
H^i_{\mathfrak m'}(M') \longrightarrow H^i(R\Gamma_{\mathfrak m'}(M')^\wedge)
$$
を $i \leq s$ に対して得る．ここで ${}^\wedge$ は導来 $I'$ 進完備化であり，
これらの加群は $I'$ のある冪で零化される．補題
\ref{algebraization-lemma-local-cohomology-derived-completion} により，同型
$$
H^i_{\mathfrak m'}(M') \longrightarrow
\lim H^i_{\mathfrak m'}(M'/(I')^nM'))
$$
を $i \leq s$ に対して得る．すでに確立した $A$ 上の局所コホモロジーとの比較と
合わせれば，補題が成り立つと結論できる．
\end{proof}

\begin{lemma}
\label{algebraization-lemma-bootstrap-bis-bis}
$I \subset \mathfrak a$ を Noether 環 $A$ のイデアルとする．
$M$ を有限 $A$ 加群とする．$s$ および $d$ を整数とする．次を仮定する．
\begin{enumerate}
\item[(a)] $A$ は双対化複体をもつ．
\item[(b)] $\text{cd}(A, I) \leq d$ である．
\item[(c)] $\mathfrak p \not \in V(I)$ かつ
$\mathfrak q \in V(\mathfrak p) \cap V(\mathfrak a)$ ならば，
$\text{depth}_{A_\mathfrak p}(M_\mathfrak p) > s$ または
$\text{depth}_{A_\mathfrak p}(M_\mathfrak p) +
\dim((A/\mathfrak p)_\mathfrak q) > d + s$ である．
\end{enumerate}
このとき $A, I, V(\mathfrak a), M, s, d$ は状況
\ref{algebraization-situation-bootstrap} にあるとおりである．
\end{lemma}

\begin{proof}
状況 \ref{algebraization-situation-bootstrap} の仮定 (1)，(3)，(4)，(6) が成り立つことを
示さなければならない．(a) $\Rightarrow$ (1)，(b) $\Rightarrow$ (3)，および
(c) $\Rightarrow$ (4) は明らかである．証明を完了するため，次の段落で
(6) が成り立つことを示す．

\medskip\noindent
$\mathfrak q \in V(\mathfrak a)$ とする．$A', I', \mathfrak m', M'$ を
$I$ 進完備化，すなわち
$A_\mathfrak q, I_\mathfrak q, \mathfrak qA_\mathfrak q, M_\mathfrak q$ の完備化と
する．$\mathfrak p' \subset A'$ を，
$V(\mathfrak p') \cap V(I') = \{\mathfrak m'\}$ を満たす非極大素イデアルとする．
このとき $\dim(A'/\mathfrak p') \leq d$ であることに注意する．これは
Local Cohomology, 補題 \ref{local-cohomology-lemma-cd-bound-dim-local} による．
$\mathfrak p \subset A$ を $\mathfrak p'$ の像とする．次が成り立つ：
$\text{depth}(M'_{\mathfrak p'}) \geq \text{depth}(M_\mathfrak p)$ かつ
$\text{depth}(M'_{\mathfrak p'}) +
\dim(A'/\mathfrak p') =
\text{depth}(M_\mathfrak p) +
\dim((A/\mathfrak p)_\mathfrak q)$
である．これは Local Cohomology, 補題
\ref{local-cohomology-lemma-change-completion} による．仮定 (c) により，
$\text{depth}(M'_{\mathfrak p'}) \geq \text{depth}(M_\mathfrak p) > s$
であってこれで終わるか，または
$\text{depth}(M'_{\mathfrak p'}) +
\dim(A'/\mathfrak p') > s + d$ である．後者は
$\text{depth}(M'_{\mathfrak p'}) > s$ を導く．これは，すでに示した不等式
$\dim(A'/\mathfrak p') \leq d$ による．いずれの場合にも所望の結論を得る．
\end{proof}

\begin{lemma}
\label{algebraization-lemma-bootstrap}
状況 \ref{algebraization-situation-bootstrap} において，逆系
$\{H^i_T(I^nM)\}_{n \geq 0}$ は $i \leq s$ に対してプロ零である．
さらに，整数 $m_0$ が存在し，すべての $m \geq m_0$ に対して，整数
$m'(m) \geq m$ であって，$k \geq m'(m)$ のとき
$H^{s + 1}_T(I^kM) \to H^{s + 1}_T(I^mM)$
の像が $H^{s + 1}_T(I^{m_0}M)$ へ単射的に写るものが存在する．
\end{lemma}

\begin{proof}
$m$ を固定する．$\mathfrak q \in T$ とする．補題
\ref{algebraization-lemma-bootstrap-inherited} および
\ref{algebraization-lemma-algebraize-local-cohomology-bis} により，
$$
H^i_\mathfrak q(M_\mathfrak q)
\longrightarrow
\lim H^i_\mathfrak q(M_\mathfrak q/I^nM_\mathfrak q)
$$
は $i \leq s$ に対して同型である．逆系
$\{H^i_\mathfrak q(I^nM_\mathfrak q)\}_{n \geq 0}$ および
$\{H^i_\mathfrak q(M/I^nM)\}_{n \geq 0}$
は，すべての $i$ に対して Mittag--Leffler 条件を満たす．補題
\ref{algebraization-lemma-ML-local} を参照されたい．したがって，長完全列の逆系
$$
0 \to H^0_\mathfrak q(I^nM_\mathfrak q) \to
H^0_\mathfrak q(M_\mathfrak q) \to
H^0_\mathfrak q(M_\mathfrak q/I^nM_\mathfrak q) \to
H^1_\mathfrak q(I^nM_\mathfrak q) \to
H^1_\mathfrak q(M_\mathfrak q) \to \ldots
$$
を調べると（いくつかの詳細は省略する），整数 $m'(m, \mathfrak q) \geq m$ が
存在して，すべての $k \geq m'(m, \mathfrak q)$ に対して写像
$H^i_\mathfrak q(I^kM_\mathfrak q) \to H^i_\mathfrak q(I^mM_\mathfrak q)$
は $i \leq s$ に対して零であり，さらに
$H^{s + 1}_\mathfrak q(I^kM_\mathfrak q) \to
H^{s + 1}_\mathfrak q(I^mM_\mathfrak q)$
の像は $k \geq m'(m, \mathfrak q)$ に依存せず，
$H^{s + 1}_\mathfrak q(M_\mathfrak q)$ へ単射的に写ることが分かる．

\medskip\noindent
$m'(m, \mathfrak q)$ を $\mathfrak q \in T$ に依存せず選べることを示せたとする．
このとき補題は Local Cohomology, 補題 \ref{local-cohomology-lemma-zero} および
\ref{local-cohomology-lemma-essential-image} から直ちに従う．

\medskip\noindent
$\omega_A^\bullet$ を双対化複体とし，
$\delta : \Spec(A) \to \mathbf{Z}$ を対応する次元関数とする．
$\delta$ が有限個の値しかとらないことを想起する．Dualizing Complexes, 補題
\ref{dualizing-lemma-universally-catenary} を参照されたい．次を主張する：各
$d \in \mathbf{Z}$ に対して，整数 $m'(m, \mathfrak q)$ は
$\mathfrak q \in T$ のうち $\delta(\mathfrak q) = d$ を満たすものに依存せず選べる．
上で述べたことから，この主張が補題を含意することは明らかである．

\medskip\noindent
$\mathfrak q \in T$ を，$\delta(\mathfrak q) = d$ を満たすように取る．Ext 加群
$$
E(n, j) = \text{Ext}^j_A(I^nM, \omega_A^\bullet)
$$
を考える．ここで用いる重要な点は，これらが有限 $A$-加群であることである．
双対化複体に付随する次元関数の定義により，
$(\omega_A^\bullet)_\mathfrak q[-d]$ は $A_\mathfrak q$ の正規化された
双対化複体であることを想起する．Dualizing Complexes, 第
\ref{dualizing-section-dimension-function} 節を参照されたい．局所双対定理
（Dualizing Complexes, 補題
\ref{dualizing-lemma-special-case-local-duality}）によれば，
$\mathfrak qA_\mathfrak q$-進完備化を $E(n, -d - i)_\mathfrak q$ に施したものは
$H^i_\mathfrak q(I^nM_\mathfrak q)$ の Matlis 双対である．したがって，
第1段落における $m'(m, \mathfrak q)$ の選び方から，$i \leq s$ に対して，
$k \geq m'(m, \mathfrak q)$ かつ $j \geq -d - s$ のとき，写像
$$
E(m, j)_\mathfrak q \to E(k, j)_\mathfrak q
$$
は零である．これらの加群は有限であり，零でないものが生じうる $j$ は有限個しかない
（細部は省略する）ので，$W \subset \Spec(A)$ を $\mathfrak q$ の開近傍として，
$$
E(m, j)_{\mathfrak q'} \to E(m'(m, \mathfrak q), j)_{\mathfrak q'}
$$
が $j \geq -d - s$ に対して，すべての $\mathfrak q' \in W$ で零となるものを取れる．
するともちろん，写像 $E(m, j)_{\mathfrak q'} \to E(k, j)_{\mathfrak q'}$ も
$k \geq m'(m, \mathfrak q)$ に対して零である．

\medskip\noindent
$i = s + 1$（これは $j = - d - s - 1$ に対応する）について，局所双対と第1段落の結果から，
$$
K_{k, \mathfrak q} =
\Ker(E(m, -d - s - 1)_\mathfrak q \to E(k, -d - s - 1)_\mathfrak q)
$$
は $k \geq m'(m, \mathfrak q)$ に依存せず，かつ写像
$$
E(0, -d - s - 1)_\mathfrak q \to
E(m, -d - s - 1)_\mathfrak q/K_{m'(m, \mathfrak q), \mathfrak q}
$$
は全射であることが分かる．$k \geq m'(m, \mathfrak q)$ に対して
$$
K_k = \Ker(E(m, -d - s - 1) \to E(k, -d - s - 1))
$$
と置く．$K_k$ は有限加群 $E(m, -d - s - 1)$ の部分加群の増大列なので，
$m'(m, \mathfrak q)$ を少し大きく取り直すことにより，
$K_{m'(m, \mathfrak q)} = K_k$ が $k \geq m'(m, \mathfrak q)$ に対して成り立つと仮定してよい．
必要なら $W$ をさらに縮小して，
$$
E(0, -d - s - 1)_{\mathfrak q'} \to
E(m, -d - s - 1)_{\mathfrak q'}/K_{m'(m, \mathfrak q), \mathfrak q'}
$$
がすべての $\mathfrak q' \in W$ に対して全射であると仮定することもできる
（先ほどと同様，これらの加群が有限であり，この写像が $\mathfrak q$ における
局所化の後で全射となることを用いる）．

\medskip\noindent
特に $T_d = \{\mathfrak q \in T \text{ で }\delta(\mathfrak q) = d\}$ のような，
Noether 位相空間 $\Spec(A)$ の任意の部分集合は，誘導位相に関して Noether であり，
したがって準コンパクトである．上で，すべての $\mathfrak q \in T_d$ に対して，
開近傍 $W$ が存在し，そこで $m'(m, \mathfrak q)$ は
すべての $\mathfrak q' \in T_d \cap W$ に対して有効であることを見た．
したがって，整数 $m'(m, d)$ を，すべての $\mathfrak q \in T_d$ に対して
$$
E(m, j)_\mathfrak q \to E(m'(m, d), j)_\mathfrak q
$$
が $j \geq -d - s$ に対して零となり，かつ
$K_{m'(m, d)} = \Ker(E(m, -d - s - 1) \to E(m'(m, d), -d - s - 1))$
と置けば
$$
K_{m'(m, d), \mathfrak q} =
\Ker(E(m, -d - s - 1)_{\mathfrak q} \to E(k, -d - s - 1)_{\mathfrak q})
$$
がすべての $k \geq m'(m, d)$ に対して成り立ち，さらに写像
$$
E(0, -d - s - 1)_\mathfrak q \to
E(m, -d - s - 1)_\mathfrak q/K_{m'(m, d), \mathfrak q}
$$
が全射となるように選べる．局所双対定理を再び
（逆向きに）用いると，主張が正しいことが分かる．これで証明は完了する．
\end{proof}

\begin{lemma}
\label{algebraization-lemma-final-bootstrap}
状況 \ref{algebraization-situation-bootstrap} において，次を満たす整数 $m_0 \geq 0$ が存在する：
\begin{enumerate}
\item $H^i_T(M) \to H^i_T(M/I^nM)$ は，すべての $i \leq s$ および
$n \geq m_0$ に対して単射である；
\item $\{H^i_T(M/I^nM)\}_{n \geq 0}$
は $i < s$ に対して Mittag--Leffler 条件を満たす；
\item $\{H^i_T(I^{m_0}M/I^nM)\}_{n \geq m_0}$
は $i \leq s$ に対して Mittag--Leffler 条件を満たす；
\item $H^i_T(M) \to \lim H^i_T(M/I^nM)$
は $i < s$ に対して同型である；
\item $H^s_T(I^{m_0}M) \to \lim H^s_T(I^{m_0}M/I^nM)$
は $i \leq s$ に対して同型である；
\item $H^s_T(M) \to \lim H^s_T(M/I^nM)$ は
単射であり，その余核は $I^{m_0}$ で零化される；
\item $R^1\lim H^s_T(M/I^nM)$ は $I^{m_0}$ で零化される．
\end{enumerate}
\end{lemma}

\begin{proof}
(1) を示すため，$0 \leq i \leq s$ とし，整数 $m > n > 0$ を，写像
$H^i_T(I^mM) \to H^i_T(I^nM)$ が零となるように選ぶ．
これは補題 \ref{algebraization-lemma-bootstrap} により可能である．行が完全な可換図式
$$
\xymatrix{
H^i_T(I^nM) \ar[r] &
H^i_T(M) \ar[r] &
H^i_T(M/I^nM) \\
H^i_T(I^mM) \ar[r] \ar[u]^0 &
H^i_T(M) \ar[r] \ar[u]^1 &
H^i_T(M/I^mM) \ar[u] \\
}
$$
から，写像 $H^i_T(M) \to H^i_T(M/I^mM)$ が単射であることが分かる．
したがって，十分大きな任意の $m_0$ に対して (1) が成り立つ．

\medskip\noindent
長完全列
$$
0 \to H^0_T(I^nM) \to H^0_T(M) \to
H^0_T(M/I^nM) \to H^1_T(I^nM) \to
H^1_T(M) \to \ldots
$$
を考える．(2) と (4) は，これと補題 \ref{algebraization-lemma-bootstrap} から従う．

\medskip\noindent
$m_0$ および $m'(-)$ を補題 \ref{algebraization-lemma-bootstrap} のものとする．
$m \geq m_0$ に対して長完全列
$$
H^s_T(I^mM) \to H^s_T(I^{m_0}M) \to
H^s_T(I^{m_0}M/I^mM) \to H^{s + 1}_T(I^mM) \to
H^1_T(I^{m_0}M)
$$
を考える．すると $k \geq m'(m)$ に対して，
$H^{s + 1}_T(I^kM) \to H^{s + 1}_T(I^mM)$
の像は $H^{s + 1}_T(I^{m_0}M)$ へ単射的に写る．したがって，
$H^s_T(I^{m_0}M/I^kM) \to H^s_T(I^{m_0}M/I^mM)$
の像は $H^{s + 1}_T(I^mM)$ へ零として写り，これはすべての $k \geq m'(m)$ に対して成り立つ．
よって (3) と (5) が成り立つ．

\medskip\noindent
短完全列
$0 \to I^{m_0}M \to M \to M/I^{m_0} M \to 0$ および
$0 \to I^{m_0}M/I^nM \to M/I^nM \to M/I^{m_0} M \to 0$ を考える．
図式
$$
\xymatrix{
H^{s - 1}_T(M/I^{m_0}M) \ar[r] &
\lim H^s_T(I^{m_0}M/I^nM) \ar[r] &
\lim H^s_T(M/I^nM) \ar[r] &
H^s_T(M/I^{m_0}M) \\
H^{s - 1}_T(M/I^{m_0}M) \ar[r] \ar@{=}[u] &
H^s_T(I^{m_0}M) \ar[r] \ar[u]_{\cong} &
H^s_T(M) \ar[r] \ar[u] &
H^s_T(M/I^{m_0}M) \ar@{=}[u]
}
$$
を得る．下の行は完全である．上の行も Homology, 補題
\ref{homology-lemma-apply-Mittag-Leffler} により（中央の二箇所で）完全である．
これから (6) が従う．

\medskip\noindent
$B_n = H^s_T(M/I^nM)$ と書く．$A_n \subset B_n$ を
$H^s_T(I^{m_0}M/I^nM) \to H^s_T(M/I^nM)$ の像とする．
すると $(A_n)$ は (3) および Homology, 補題
\ref{homology-lemma-Mittag-Leffler} により Mittag--Leffler 条件を満たす．
また $C_n = B_n/A_n$ は $I^{m_0}$ で零化される．したがって
$R^1\lim B_n \cong R^1\lim C_n$ も $I^{m_0}$ で零化され，(7) を得る．
\end{proof}

\begin{theorem}
\label{algebraization-theorem-final-bootstrap}
状況 \ref{algebraization-situation-bootstrap} において，$i \leq s$ に対する逆系
$\{H^i_T(M/I^nM)\}_{n \geq 0}$ は値 $H^i_T(M)$ をもつ本質的定値系である．
特に，$\{H^i_T(M/I^nM)\}_{n \geq 0}$ は Mittag--Leffler であり，
$H^i_T(M)$ は $I$ のある冪で零化され，かつ
$H^i_T(M) = \lim H^i_T(M/I^nM)$ である．
\end{theorem}

\begin{proof}
$0 \leq i \leq s$ とする．$\{H^i_T(M/I^nM)\}_{n \geq 0}$ が
Mittag--Leffler であり，$H^i_T(M) = \lim H^i_T(M/I^nM)$ であることを示せば，
補題 \ref{algebraization-lemma-final-bootstrap} で示した $H^i_T(M) \to H^i_T(M/I^nM)$ の
$n \gg 0$ に対する単射性により，
$\{H^i_T(M/I^nM)\}_{n \geq 0}$ は値 $H^i_T(M)$ をもつ本質的定値系となる．
細部は省略する．

\medskip\noindent
補題 \ref{algebraization-lemma-final-bootstrap} により，$\{H^i_T(M/I^nM)\}_{n \geq 0}$ が
Mittag--Leffler であり，$H^i_T(M) = \lim H^i_T(M/I^nM)$ が $i < s$ に対して
成り立つことは既に分かっている．また，これらの主張が $i = s$ および加群
$I^{m_0}M$ について，ある $m_0 > 0$ に対し成り立つことも分かっている．
証明を終えるため，これらの $i = s$ に関する主張が実際に $M$ について成り立つことを示す．

\medskip\noindent
$M' = H^0_I(M)$ および $M'' = M/M'$ と置くと，短完全列
$$
0 \to M' \to M \to M'' \to 0
$$
を得る．また Dualizing Complexes, 補題
\ref{dualizing-lemma-divide-by-torsion} により，$M''$ は $H^0_I(M') = 0$ を満たす．
Artin--Rees（Algebra, 補題 \ref{algebra-lemma-Artin-Rees}）により，短完全列
$$
0 \to M' \to M/I^n M \to M''/I^n M'' \to 0
$$
を十分大きな $n$ に対して得る．長完全列
$$
H^s_T(M') \to
H^s_T(M/I^nM) \to
H^s_T(M''/I^nM'') \to
H^{s + 1}_T(M')
$$
を考える．ここで，逆系 $\{H^s_T(M''/I^nM'')\}_{n \geq 0}$ が
Mittag--Leffler ならば，逆系 $\{H^s_T(M/I^nM)\}_{n \geq 0}$ も
Mittag--Leffler であることは容易に分かる．
（群の逆系 $G_n$ に対する Mittag--Leffler 条件は，十分大きな $n$ における逆系の値だけに依存することに注意する．）
さらに列
$$
H^s_T(M') \to
\lim H^s_T(M/I^nM) \to
\lim H^s_T(M''/I^nM'') \to
H^{s + 1}_T(M')
$$
は，必要な箇所で Mittag--Leffler 条件が成り立つため完全である．Homology, 補題
\ref{homology-lemma-apply-Mittag-Leffler} を参照されたい．したがって，
$H^s_T(M'') \to \lim H^s_T(M''/I^nM'')$ が同型ならば，五項補題
（Homology, 補題 \ref{homology-lemma-five-lemma}）により
$H^s_T(M) \to \lim H^s_T(M/I^nM)$ も同型である．
これにより，次の段落で扱う場合へ帰着される．

\medskip\noindent
$H^0_I(M) = 0$ と仮定する．$f_1, \ldots, f_r$ を $I^{m_0}$ の生成元として選ぶ．
ここで $m_0$ は，補題
\ref{algebraization-lemma-final-bootstrap} で $M$ に対して得られた整数である．すると，完全列
$$
0 \to M \xrightarrow{f_1, \ldots, f_r}
(I^{m_0}M)^{\oplus r} \to Q \to 0
$$
によって $Q$ を定義する．いくつか注意する．第1の写像は $H^0_I(M) = 0$ であることから単射である．
この単射の余核 $Q$ は有限 $A$-加群であり，すべての $1 \leq j \leq r$ に対して
$Q_{f_j} \cong (M_{f_j})^{\oplus r - 1}$ が成り立つ．特に，素イデアル
$\mathfrak p \subset A$ が $\mathfrak p \not \in V(I)$ を満たせば，
$Q_\mathfrak p \cong (M_\mathfrak p)^{\oplus r - 1}$ である．同様に，
$\mathfrak q \in T$ および $\mathfrak p' \subset A' = (A_\mathfrak q)^\wedge$ が与えられ，
$V(IA')$ に含まれないならば，
$Q'_{\mathfrak p'} \cong (M'_{\mathfrak p'})^{\oplus r - 1}$ が成り立つ．ここで
$Q' = (Q_\mathfrak q)^\wedge$ および $M' = (M_\mathfrak q)^\wedge$ である．
したがって，状況 \ref{algebraization-situation-bootstrap} の条件は $A, I, T, Q$ に対して成り立つ．
（$Q$ は非零の $H^0_I(Q)$ をもちうるが，これは問題にならないことに注意する．）

\medskip\noindent
任意の $n \geq 0$ に対して $F^nM = M \cap I^n(I^{m_0}M)^{\oplus r}$ と置くと，短完全列
$$
0 \to F^nM \to I^n(I^{m_0}M)^{\oplus r} \to I^nQ \to 0
$$
を得る．Artin--Rees（Algebra, 補題 \ref{algebra-lemma-Artin-Rees}）により，
$c \geq 0$ が存在して，$I^n M \subset F^nM \subset I^{n - c}M$ が
すべての $n \geq c$ に対して成り立つ．$m_0$ を整数とし，
$m'(m)$ を，$m \geq m_0$ に対して定義され，$M$ に補題
\ref{algebraization-lemma-bootstrap} を適用して得られる関数とする．整数 $m_0$ は上で選んだ整数 $m_0$ と同じである
（これは確認しなくてよい．望むなら二つの整数の最大値を取ればよい）．最後に，
$Q$ に補題 \ref{algebraization-lemma-bootstrap} を適用すると，すべての整数 $m$ に対して，
$m''(m) \geq m$ であって，写像 $H^s_T(I^kQ) \to H^s_T(I^mQ)$ が
すべての $k \geq m''(m)$ に対して零となる整数が存在する．

\medskip\noindent
$m \geq m_0$ を固定する．$k \geq m'(m''(m + c))$ を選ぶ．
$\xi \in H^{s + 1}_T(I^kM)$ を，$H^{s + 1}_T(M)$ で零に写るように選ぶ．
$\xi$ が $H^{s + 1}_T(I^mM)$ で零に写ることを示したい．実際，これにより，補題
$\{H^s_T(M/I^nM)\}_{n \geq 0}$ が，補題 \ref{algebraization-lemma-final-bootstrap} の証明と
まったく同様に Mittag--Leffler であることが示される．
議論を視覚化するための図式は次のとおりである：
$$
\xymatrix{
&
H^{s + 1}_T(I^kM) \ar[r] \ar[d] &
H^{s + 1}_T(I^k(I^{m_0}M)^{\oplus r}) \ar[d] &
\\
H^s_T(I^{m''(m + c)}Q) \ar[r]_-\delta \ar[d] &
H^{s + 1}_T(F^{m''(m + c)}M) \ar[r] \ar[d] &
H^{s + 1}_T(I^{m''(m + c)}(I^{m_0}M)^{\oplus r}) \\
H^s_T(I^{m + c}Q) \ar[r] &
H^{s + 1}_T(F^{m + c}M) \ar[d] &
\\
&
H^{s + 1}_T(I^mM)
}
$$
$\xi$ の $H^{s + 1}_T(I^k(I^{m_0}M)^{\oplus r})$ における像は
$H^{s + 1}_T((I^{m_0}M)^{\oplus r})$ で零に写り，したがって
$H^{s + 1}_T(I^{m''(m + c)}(I^{m_0}M)^{\oplus r})$ でも，$m'(-)$ の選び方により零に写る．
したがって，像 $\xi' \in H^{s + 1}_T(F^{m''(m + c)}M)$ は
$H^{s + 1}_T(I^{m''(m + c)}(I^{m_0}M)^{\oplus r})$ で零に写り，ゆえに
$\xi' = \delta(\eta)$ となる $\eta \in H^s_T(I^{m''(m + c)}Q)$ が存在する．
$m''(-)$ の選び方により，$\eta$ は $H^s_T(I^{m + c}Q)$ で零に写る．
したがって $\xi'$ は $H^{s + 1}_T(F^{m + c}M)$ で零に写る．
$F^{m + c}M \subset I^mM$ なので結論を得る．

\medskip\noindent
最後に，極限に関する主張を証明する．短完全列
$$
0 \to M/F^nM \to (I^{m_0}M)^{\oplus r}/I^n (I^{m_0}M)^{\oplus r}
\to Q/I^nQ \to 0
$$
を考える．これらの逆系は pro-同型なので，
$\lim H^s_T(M/I^nM) = \lim H^s_T(M/F^nM)$ である．可換図式
$$
\xymatrix{
H^{s - 1}_T(Q) \ar[r] \ar[d] &
\lim H^{s - 1}_T(Q/I^nQ) \ar[d] \\
H^s_T(M) \ar[r] \ar[d] &
\lim H^s_T(M/I^nM) \ar[d] \\
H^s_T((I^{m_0}M)^{\oplus r}) \ar[r] \ar[d] &
\lim H^s_T((I^{m_0}M)^{\oplus r}/I^n(I^{m_0}M)^{\oplus r}) \ar[d] \\
H^s_T(Q) \ar[r] &
\lim H^s_T(Q/I^nQ)
}
$$
を得る．必要な箇所で Mittag--Leffler 条件が成り立つため，右列は完全である．Homology, 補題
\ref{homology-lemma-apply-Mittag-Leffler} を参照されたい．最下段の水平射は補題
\ref{algebraization-lemma-final-bootstrap} の (6) により単射である（！）．その一つ上の水平射は，補題
\ref{algebraization-lemma-final-bootstrap} の (5) により全単射である．コホモロジー次数 $\leq s - 1$ における射は同型である．
したがって，五項補題（Homology, 補題 \ref{homology-lemma-five-lemma}）により，
$H^s_T(M) \to \lim H^s_T(M/I^nM)$ が同型であることが分かる．
これで定理の証明は完了する．
\end{proof}

\begin{lemma}
\label{algebraization-lemma-combine-two}
$I \subset \mathfrak a \subset A$ を Noether 環 $A$ のイデアルとし，
$M$ を有限 $A$-加群とする．$s$ および $d$ を整数とする．次を仮定する：
\begin{enumerate}
\item $A, I, V(\mathfrak a), M$ は，$s$ および $d$ に対して状況
\ref{algebraization-situation-bootstrap} の仮定を満たす；
\item $A, I, \mathfrak a, M$ は，$s + 1$ および $d$ に対して，
$J = \mathfrak a$ とした補題 \ref{algebraization-lemma-algebraize-local-cohomology-general} の条件を満たす．
\end{enumerate}
このとき，イデアル $J_0 \subset \mathfrak a$ で
$V(J_0) \cap V(I) = V(\mathfrak a)$ を満たすものが存在し，任意の
$J \subset J_0$ で $V(J) \cap V(I) = V(\mathfrak a)$ を満たすものに対して，写像
$$
H^{s + 1}_J(M) \longrightarrow \lim H^{s + 1}_\mathfrak a(M/I^nM)
$$
は同型である．
\end{lemma}

\begin{proof}
実際，補題 \ref{algebraization-lemma-algebraize-local-cohomology-general} により $J_0$ が存在し，同型
$H^{s + 1}_J(M) = H^{s + 1}(R\Gamma_\mathfrak a(M)^\wedge)$ を得る．また，短完全列
$$
0 \to R^1\lim H^s_\mathfrak a(M/I^nM) \to
H^{s + 1}(R\Gamma_\mathfrak a(M)^\wedge) \to
\lim H^{s + 1}_\mathfrak a(M/I^nM) \to 0
$$
を Dualizing Complexes, 補題 \ref{dualizing-lemma-completion-local} により得る．さらに，加群
$R^1\lim H^s_\mathfrak a(M/I^nM)$ は零である．実際，定理
\ref{algebraization-theorem-final-bootstrap} により
$\{H^s_\mathfrak a(M/I^nM)\}_{n \geq 0}$ は Mittag--Leffler である．
\end{proof}








\section{形式切断の代数化 I}
\label{algebraization-section-algebraization-sections}

\noindent
本節では，局所的な場合に形式切断を代数化する問題を調べる．
$(A, \mathfrak m)$ を Noether 局所環とする．
$I \subset A$ をイデアルとする．
$$
X = \Spec(A) \supset U = \Spec(A) \setminus \{\mathfrak m\}
$$
$Y = V(I)$ を $I$ に対応する閉部分スキームと記す．
$\mathcal{F}$ を連接 $\mathcal{O}_U$-加群とする．本節では極限
$$
\lim_n H^i(U, \mathcal{F}/I^n\mathcal{F})
$$
を考える．これは，$\mathcal{F}$ を $U$ の $Y$ に沿う形式完備化へ引き戻したものの
コホモロジーと密接に関係する．しかし，まだ形式スキームを導入していないので，
ここではこの用語を使えない．

\begin{lemma}
\label{algebraization-lemma-compare-with-derived-completion}
$U$ を Noether 局所環 $A$ の穴あきスペクトルとする．
$\mathcal{F}$ を連接 $\mathcal{O}_U$-加群とし，$I \subset A$ をイデアルとする．このとき
$$
H^i(R\Gamma(U, \mathcal{F})^\wedge) =
\lim H^i(U, \mathcal{F}/I^n\mathcal{F})
$$
がすべての $i$ に対して成り立つ．ここで $R\Gamma(U, \mathcal{F})^\wedge$ は
導来 $I$-進完備化を表す．
\end{lemma}

\begin{proof}
補題 \ref{algebraization-lemma-formal-functions-general} および
\ref{algebraization-lemma-derived-completion-pseudo-coherent} により
$$
R\Gamma(U, \mathcal{F})^\wedge =
R\Gamma(U, \mathcal{F}^\wedge) =
R\Gamma(U, R\lim \mathcal{F}/I^n\mathcal{F})
$$
である．したがって，短完全列
$$
0 \to R^1\lim H^{i - 1}(U, \mathcal{F}/I^n\mathcal{F}) \to
H^i(R\Gamma(U, \mathcal{F})^\wedge) \to
\lim H^i(U, \mathcal{F}/I^n\mathcal{F}) \to 0
$$
を Cohomology, 補題 \ref{cohomology-lemma-RGamma-commutes-with-Rlim} により得る．
$R^1\lim$ の項は消える．実際，群の逆系
$H^i(U, \mathcal{F}/I^n\mathcal{F})$ は補題 \ref{algebraization-lemma-ML-local} により
Mittag--Leffler 条件を満たす．
\end{proof}

\begin{theorem}
\label{algebraization-theorem-algebraization-formal-sections}
\begin{reference}
証明法はおおむね \cite[定理 1]{Faltings-algebraisation} および
\cite[Satz 2]{Faltings-uber} の証明法に従う．この結果は
\cite[定理 1.1]{MRaynaud-paper}（補集合がアフィンの場合）および
\cite[定理 3.9]{MRaynaud-book}（補集合が少数のアフィンの合併である場合）とほぼ同じである．
\end{reference}
$(A, \mathfrak m)$ を，双対化複体をもち，イデアル $I$ に関して完備な Noether 局所環とする．
$X = \Spec(A)$，$Y = V(I)$，および $U = X \setminus \{\mathfrak m\}$ と置く．
$\mathcal{F}$ を $U$ 上の連接層とする．次を仮定する：
\begin{enumerate}
\item $\text{cd}(A, I) \leq d$，すなわち，
$H^i(X \setminus Y, \mathcal{G}) = 0$ が，$i \geq d$ および任意の準連接層
$\mathcal{G}$（$X$ 上）に対して成り立つ；
\item $x \in X \setminus Y$ で，閉包 $\overline{\{x\}}$ が $X$ において
$U \cap Y$ と交わる任意のものに対して，
$$
\text{depth}_{\mathcal{O}_{X, x}}(\mathcal{F}_x) \geq s
\quad\text{または}\quad
\text{depth}_{\mathcal{O}_{X, x}}(\mathcal{F}_x)
+ \dim(\overline{\{x\}}) > d + s
$$
\end{enumerate}
このとき，開集合 $V_0 \subset U$ で $U \cap Y$ を含むものが存在し，
任意の開集合 $V \subset V_0$ で $U \cap Y$ を含むものに対して写像
$$
H^i(V, \mathcal{F}) \to \lim H^i(U, \mathcal{F}/I^n\mathcal{F})
$$
は $i < s$ に対して同型である．さらに
$
\text{depth}_{\mathcal{O}_{X, x}}(\mathcal{F}_x) +
\dim(\overline{\{x\}}) > s
$
がすべての $x \in U \cap Y$ に対して成り立つならば，これらのコホモロジー群は有限
$A$-加群である．
\end{theorem}

\begin{proof}
有限 $A$-加群 $M$ を，$\mathcal{F}$ が $U$ 上で，$\mathcal{O}_X$-加群として
$M$ に付随する連接加群の制限となるように選ぶ．Local Cohomology, 補題
\ref{local-cohomology-lemma-finiteness-pushforwards-and-H1-local} を参照されたい．
すると補題 \ref{algebraization-lemma-algebraize-local-cohomology} の仮定が満たされる．
同補題のように $J_0$ を選び，$V_0 = X \setminus V(J_0)$ と置く．
開集合 $V \subset V_0$ で $U \cap Y$ を含むものは，$1$ 対 $1$ で，
イデアル $J \subset J_0$ で $V(J) \cap V(I) = \{\mathfrak m\}$ を満たすものと対応する．
さらに，このような選択に対して識別三角
$$
R\Gamma_J(M) \to M \to R\Gamma(V, \mathcal{F}) \to
R\Gamma_J(M)[1]
$$
を得る．同様に，識別三角
$$
R\Gamma_\mathfrak m(M)^\wedge \to
M \to
R\Gamma(U, \mathcal{F})^\wedge \to
R\Gamma_\mathfrak m(M)^\wedge[1]
$$
を得る．これは導来 $I$-進完備化を含む．補題
\ref{algebraization-lemma-compare-with-derived-completion} により，
$R\Gamma(U, \mathcal{F})^\wedge$ のコホモロジー群は定理の主張にある極限に等しい．
これらの三角の間の標準写像と容易な議論から，本定理は主要な補題
\ref{algebraization-lemma-algebraize-local-cohomology} より従う
（ここでは $i < s$ であるのに対し，補題では $i \leq s$ であることに注意する．
これはシフトによる）．追加の仮定の下でのコホモロジー群の有限性は補題
\ref{algebraization-lemma-kill-colimit} から従う．
\end{proof}

\begin{lemma}
\label{algebraization-lemma-application-theorem}
$(A, \mathfrak m)$ を，双対化複体をもち，イデアル $I$ に関して完備な Noether 局所環とする．
$X = \Spec(A)$，$Y = V(I)$，および $U = X \setminus \{\mathfrak m\}$ と置く．
$\mathcal{F}$ を $U$ 上の連接層とする．任意の $x \in U$ が $\mathcal{F}$ の随伴点であるとき，
$\dim(\overline{\{x\}}) > \text{cd}(A, I) + 1$ と仮定する．ここで
$\overline{\{x\}}$ は $X$ における閉包である．このとき写像
$$
\colim H^0(V, \mathcal{F})
\longrightarrow
\lim H^0(U, \mathcal{F}/I^n\mathcal{F})
$$
は有限 $A$-加群の同型である．ここで余極限は，開集合 $V \subset U$ で
$U \cap Y$ を含むもの全体にわたる．
\end{lemma}

\begin{proof}
$s = 1$ として定理 \ref{algebraization-theorem-algebraization-formal-sections} を適用する
（有限性も得られる）．
\end{proof}




\section{形式切断の代数化 II}
\label{algebraization-section-algebraization-sections-coherent}

\noindent
第 \ref{algebraization-section-algebraization-sections-general} 節および第
\ref{algebraization-section-bootstrap} 節の結果が，準アフィンスキーム上の連接層のコホモロジーと，
あるイデアルに関するその完備化に対してもつ帰結をすべて簡潔に述べるのはやや難しい．
本節では $H^0$ への応用を網羅的ではない形で列挙する．次節では高次コホモロジーへの応用を扱う．

\begin{lemma}
\label{algebraization-lemma-ses-H0}
$I \subset \mathfrak a$ を Noether 環 $A$ のイデアルとする．
$0 \to \mathcal{F}' \to \mathcal{F} \to \mathcal{F}'' \to 0$ を
$U = \Spec(A) \setminus V(\mathfrak a)$ 上の連接加群の短完全列とする．
$\mathcal{V}$ を，$V \subset U$ であり，かつ $U \cap V(I)$ を含む開部分スキーム全体の集合とし，
逆包含で順序づける．可換図式
$$
\xymatrix{
\colim_\mathcal{V} H^0(V, \mathcal{F}') \ar[d] \ar[r] &
\colim_\mathcal{V} H^0(V, \mathcal{F}) \ar[d] \ar[r] &
\colim_\mathcal{V} H^0(V, \mathcal{F}'') \ar[d] \\
\lim H^0(U, \mathcal{F}'/I^n\mathcal{F}') \ar[r] &
\lim H^0(U, \mathcal{F}'/I^n\mathcal{F}) \ar[r] &
\lim H^0(U, \mathcal{F}'/I^n\mathcal{F}'')
}
$$
を考える．左と右の下向き射が同型ならば，中央の下向き射も同型である．
中央と左の下向き射が同型ならば，左の下向き射も同型である．
\end{lemma}

\begin{proof}
図式の列は中央で完全であり，第1の射は単射である．したがって最後の主張は容易な図式追跡から従う．
証明の残りでは，左と右の下向き射が同型であると仮定する．図式追跡により中央の下向き射は単射である．
残るのは，これが全射であることを示すことだけである．

\medskip\noindent
有限 $A$-加群 $M$ および $M'$ を，$\mathcal{F}$ と $\mathcal{F}'$ がそれぞれ
$\widetilde{M}$ と $\widetilde{M}'$ の $U$ への制限となるように選べる．Local Cohomology, 補題
\ref{local-cohomology-lemma-finiteness-pushforwards-and-H1-local} を参照されたい．
$M'$ を $\mathfrak a^n M'$（ある $n \geq 0$）で置き換えると，
$\mathcal{F}' \to \mathcal{F}$ が加群写像 $M' \to M$ に対応すると仮定してよい．
Cohomology of Schemes, 補題 \ref{coherent-lemma-homs-over-open} を参照されたい．
$M'$ を $M' \to M$ の像で置き換え，$M'' = M/M'$ と置くと，もとの短完全列は，
短完全列 $0 \to M' \to M \to M'' \to 0$ という $A$-加群列に付随する連接加群の
短完全列を制限したものに対応することが分かる．

\medskip\noindent
$\hat s \in \lim H^0(U, \mathcal{F}/I^n\mathcal{F})$ とし，その像を
$\hat s'' \in \lim H^0(U, \mathcal{F}''/I^n\mathcal{F}'')$ とする．仮定により，
$V \in \mathcal{V}$ と，切断 $s'' \in \mathcal{F}''(V)$ で $\hat s''$ に写るものが見つかる．
$J \subset A$ を $V(J) = \Spec(A) \setminus V$ を満たすイデアルとする．
Cohomology of Schemes, 補題 \ref{coherent-lemma-homs-over-open} により，
$J$ をある冪で置き換えれば，$A$-線形写像 $\varphi : J \to M''$ で
$s''$ に対応するものが存在すると仮定してよい．この $J$ の選択を固定する．
証明の残りでは $V$ を，より小さい $V$（$\mathcal{V}$ の元）で置き換える．
すなわち，$V \cap V(J) = \emptyset$ とする．

\medskip\noindent
表示 $A^{\oplus m} \to A^{\oplus n} \to J \to 0$ を選ぶ．
$g_1, \ldots, g_n \in J$ を $A^{\oplus n}$ の基底ベクトルの像と記す．すると
$J = (g_1, \ldots, g_n)$ である．写像 $A^{\oplus m} \to A^{\oplus n}$ が行列
$(a_{ji})$ で与えられ，$\sum a_{ji} g_i = 0$，$j = 1, \ldots, m$ となるものとする．
$M \to M''$ は全射なので，各 $i$ に対して，$m_i \in M$ で
$\varphi(g_i) \in M''$ に写るものを選べる．すると，元 $g_i \hat s - m_i$ は
$\lim H^0(U, \mathcal{F}/I^n\mathcal{F})$ の部分加群
$\lim H^0(U, \mathcal{F}'/I^n\mathcal{F}')$ に属する．
仮定により，$V$ を縮小すれば，$s'_i \in \mathcal{F}'(V)$，$i = 1, \ldots, n$ であって，
その $s'_i$ が $g_i \hat s - m_i$ に写るものが存在すると仮定してよい．
$s_i = s'_i + m_i$ と $\mathcal{F}(V)$ の中で置く．
$\sum a_{ji} s_i$ は $\sum a_{ji}g_i\hat s = 0$ に，写像
$$
\colim_\mathcal{V} \mathcal{F}(V')
\longrightarrow
\lim H^0(U, \mathcal{F}/I^n\mathcal{F})
$$
によって写ることに注意する．
この写像は単射なので（上を参照），$V$ を縮小した後，$\sum a_{ji}s_i = 0$ が
$\mathcal{F}(V)$ で，すべての $j = 1, \ldots, m$ に対して成り立つと仮定してよい．したがって，
$A$-加群写像 $J \to \mathcal{F}(V)$ で $g_i$ を $s_i$ に送るものを得る．
$\widetilde{J}$ の普遍性により，この $A$-加群写像は $\mathcal{O}_V$-加群写像
$\widetilde{J}|_V \to \mathcal{F}$ に対応する．しかし $V(J) \cap V = \emptyset$ なので，
$\widetilde{J}|_V = \mathcal{O}_V$ である．したがって切断 $s \in \mathcal{F}(V)$ を構成した．
$s$ が上に表示した写像によって $\hat s$ に写ることを示す計算は省略する．
\end{proof}

\noindent
次の補題は，命題 \ref{algebraization-proposition-application-H0} によって後に置き換えられる．

\begin{lemma}
\label{algebraization-lemma-application-H0-pre}
$I \subset \mathfrak a$ を Noether 環 $A$ のイデアルとする．
$\mathcal{F}$ を
$U = \Spec(A) \setminus V(\mathfrak a)$ 上の連接加群とする．
次を仮定する．
\begin{enumerate}
\item $A$ は $I$-進完備であり，双対化複体をもつ．
\item $x \in \text{Ass}(\mathcal{F})$，$x \not \in V(I)$，
$\overline{\{x\}} \cap V(I) \not \subset V(\mathfrak a)$,
かつ $z \in \overline{\{x\}} \cap V(\mathfrak a)$ ならば，
$\dim(\mathcal{O}_{\overline{\{x\}}, z}) > \text{cd}(A, I) + 1$ である．
\item 次のいずれかが成り立つ．
\begin{enumerate}
\item $\mathcal{F}$ の $U \setminus V(I)$ への制限は $(S_1)$ を満たす．
\item $V(\mathfrak a)$ の次元は高々 $2$ である\footnote{
$V(\mathfrak a)$ 上で，$\Spec(A)$ 上の次元関数がとる最大値と最小値との差が
高々 $2$ である，という意味である．}．
\end{enumerate}
\end{enumerate}
このとき同型
$$
\colim H^0(V, \mathcal{F})
\longrightarrow
\lim H^0(U, \mathcal{F}/I^n\mathcal{F})
$$
を得る．ここで余極限は，開部分スキーム $V \subset U$ で $U \cap V(I)$ を含むもの全体にわたる．
\end{lemma}

\begin{proof}
有限 $A$-加群 $M$ を，$\mathcal{F}$ が $U$ 上で $M$ に付随する連接加群の
制限となるように選ぶ．Local Cohomology, 補題
\ref{local-cohomology-lemma-finiteness-pushforwards-and-H1-local} を参照されたい．
$d = \text{cd}(A, I)$ と置く．
$\mathfrak p$ を $A$ の素イデアルで $V(I)$ に含まれないものとし，
$\mathfrak q \in V(\mathfrak p) \cap V(\mathfrak a)$ とする．
このとき，$\mathfrak p$ が $M$ の随伴素イデアルでなく，したがって
$\text{depth}(M_\mathfrak p) \geq 1$ であるか，または (2) により
$\dim((A/\mathfrak p)_\mathfrak q) > d + 1$ である．
したがって補題 \ref{algebraization-lemma-algebraize-local-cohomology-general} の仮定は
$s = 1$ および $d$ に対して満たされる．ここで条件 (3) を用いた．
ゆえに，イデアル $J_0 \subset \mathfrak a$ で $V(J_0) \cap V(I) = V(\mathfrak a)$ を満たすものが存在し，
任意の $J \subset J_0$ で $V(J) \cap V(I) = V(\mathfrak a)$ を満たすものに対して写像
$$
H^i_J(M) \longrightarrow H^i(R\Gamma_\mathfrak a(M)^\wedge)
$$
は $i = 0, 1$ に対して同型である．完全三角の射
$$
\xymatrix{
R\Gamma_J(M)  \ar[d] \ar[r] &
M \ar[r] \ar[d] &
R\Gamma(V, \mathcal{F}) \ar[d] \ar[r] &
R\Gamma_J(M)[1]  \ar[d] \\
R\Gamma_J(M)^\wedge \ar[r] &
M \ar[r] &
R\Gamma(V, \mathcal{F})^\wedge \ar[r] &
R\Gamma_J(M)^\wedge[1] \\
R\Gamma_\mathfrak a(M)^\wedge \ar[r] \ar[u] &
M \ar[r] \ar[u] &
R\Gamma(U, \mathcal{F})^\wedge \ar[r] \ar[u] &
R\Gamma_\mathfrak a(M)^\wedge[1] \ar[u]
}
$$
を考える．ここで $V = \Spec(A) \setminus V(J)$ である．
$R\Gamma_\mathfrak a(M)^\wedge \to R\Gamma_J(M)^\wedge$ は同型であることを想起する
（$\mathfrak a$，$\mathfrak a + I$，および $J + I$ が同じ閉部分スキームを定めるからである．
たとえば補題 \ref{algebraization-lemma-algebraize-local-cohomology-general} の証明を参照されたい）．
したがって
$R\Gamma(U, \mathcal{F})^\wedge = R\Gamma(V, \mathcal{F})^\wedge$.
これより可換図式
$$
\xymatrix{
0 \ar[r] &
H^0_J(M) \ar[r] \ar[d] &
M \ar[r] \ar[d] \ar[r] &
\Gamma(V, \mathcal{F}) \ar[d] \ar[r] &
H^1_J(M) \ar[d] \ar[r] &
0 \\
0 \ar[r] &
H^0(R\Gamma_J(M)^\wedge) \ar[r] &
M \ar[r] &
H^0(R\Gamma(V, \mathcal{F})^\wedge) \ar[r] &
H^1(R\Gamma_J(M)^\wedge) \ar[r] &
0 \\
0 \ar[r] &
H^0(R\Gamma_\mathfrak a(M)^\wedge) \ar[r] \ar[u] &
M \ar[r] \ar[u] &
H^0(R\Gamma(U, \mathcal{F})^\wedge) \ar[r] \ar[u] &
H^1(R\Gamma_\mathfrak a(M)^\wedge) \ar[r] \ar[u] &
0
}
$$
を得る．各行は完全であり，下側の垂直射は同型である．したがって同型
$\Gamma(V, \mathcal{F}) \to H^0(R\Gamma(U, \mathcal{F})^\wedge)$.
を得る．補題 \ref{algebraization-lemma-formal-functions-general} および
\ref{algebraization-lemma-derived-completion-pseudo-coherent} により
$$
R\Gamma(U, \mathcal{F})^\wedge =
R\Gamma(U, \mathcal{F}^\wedge) =
R\Gamma(U, R\lim \mathcal{F}/I^n\mathcal{F})
$$
であり，Cohomology, 補題 \ref{cohomology-lemma-RGamma-commutes-with-Rlim} によって
$H^0(R\Gamma(U, \mathcal{F})^\wedge) =
\lim H^0(U, \mathcal{F}/I^n\mathcal{F})$ を得る．
\end{proof}

\noindent
ここで前の補題をブートストラップして条件 (3) を取り除く．

\begin{proposition}
\label{algebraization-proposition-application-H0}
$I \subset \mathfrak a$ を Noether 環 $A$ のイデアルとする．
$\mathcal{F}$ を
$U = \Spec(A) \setminus V(\mathfrak a)$ 上の連接加群とする．
次を仮定する．
\begin{enumerate}
\item $A$ は $I$-進完備であり，双対化複体をもつ．
\item $x \in \text{Ass}(\mathcal{F})$，$x \not \in V(I)$，
$\overline{\{x\}} \cap V(I) \not \subset V(\mathfrak a)$,
かつ $z \in \overline{\{x\}} \cap V(\mathfrak a)$ ならば，
$\dim(\mathcal{O}_{\overline{\{x\}}, z}) > \text{cd}(A, I) + 1$ である．
\end{enumerate}
このとき同型
$$
\colim H^0(V, \mathcal{F})
\longrightarrow
\lim H^0(U, \mathcal{F}/I^n\mathcal{F})
$$
を得る．ここで余極限は，開部分スキーム $V \subset U$ で $U \cap V(I)$ を含むもの全体にわたる．
\end{proposition}

\begin{proof}
$T \subset U$ を，点 $x$ で $\overline{\{x\}} \cap V(I) \subset V(\mathfrak a)$ を満たすもの全体の集合とする．
$\mathcal{F} \to \mathcal{F}'$ を，Local Cohomology, 補題
\ref{local-cohomology-lemma-get-depth-1-along-Z} で構成された $U$ 上の連接加群の全射とする．
$\mathcal{F} \to \mathcal{F}'$ は，ある開集合 $V \subset U$ で $U \cap V(I)$ を含むものの上で同型なので，
$\mathcal{F}$ を $\mathcal{F}'$ で置き換えた場合に補題を証明すれば十分である．したがって，
$x \in U$ が $\overline{\{x\}} \cap V(I) \subset V(\mathfrak a)$ を満たすならば
$\text{depth}(\mathcal{F}_x) \geq 1$ であると仮定してよいし，そう仮定する．

\medskip\noindent
$\mathcal{V}$ を，$V \subset U$ であり $U \cap V(I)$ を含む開部分スキーム全体の集合とし，
包含関係の逆で順序づける．これは有向集合である．まず，写像
$$
\mathcal{F}(V)
\longrightarrow
\lim H^0(U, \mathcal{F}/I^n\mathcal{F})
$$
は任意の $V \in \mathcal{F}$ に対して単射であると主張する（特に補題の写像は単射となる）．
実際，前段落により，随伴点 $x$ は $\mathcal{F}$ に対して
$\overline{\{x\}} \cap U \cap Y \not = \emptyset$ を満たさなければならない．
$y \in \overline{\{x\}} \cap U \cap Y$ ならば，$\mathcal{F}_x$ は $\mathcal{F}_y$ の局所化であり，
Krull の共通部分定理（Algebra, 補題
\ref{algebra-lemma-intersect-powers-ideal-module-zero}）により
$\mathcal{F}_y \subset \lim \mathcal{F}_y/I^n \mathcal{F}_y$ である．
核に属する切断 $s \in \mathcal{F}(V)$ は台が空でなければならず，したがって零であるから，主張が従う．

\medskip\noindent
有限 $A$-加群 $M$ を，$\mathcal{F}$ が $\widetilde{M}$ の $U$ への制限となるように選ぶ．
Local Cohomology, 補題 \ref{local-cohomology-lemma-finiteness-pushforwards-and-H1-local} を参照されたい．
$H^0_\mathfrak a(M) = 0$ であると仮定してよいし，そう仮定する．
$\text{Ass}(M) \setminus V(I) = \{\mathfrak p_1, \ldots, \mathfrak p_n\}$ とする．
$n$ に関する帰納法で補題を証明する．順序を入れ替えれば，$\mathfrak p_n$ は
$\{\mathfrak p_1, \ldots, \mathfrak p_n\}$ の包含関係に関する極小元，すなわち
$\mathfrak p_n$ は $M$ の台の一般点であると仮定してよい．
次のように置く．
$$
M' = H^0_{\mathfrak p_1 \ldots \mathfrak p_{n - 1} I}(M)
$$
$M'' = M/M'$ と置く．$\mathcal{F}'$ および $\mathcal{F}''$ を，連接 $\mathcal{O}_U$-加群で，
それぞれ $M'$ および $M''$ に対応するものとする．Dualizing Complexes, 補題
\ref{dualizing-lemma-divide-by-torsion} により，$M''$ の随伴素イデアルはただ一つ，すなわち
$\mathfrak p_n$ だけである．したがって $\mathcal{F}''$ の随伴点もただ一つであり，
補題 \ref{algebraization-lemma-application-H0-pre} の条件 (3)(a) が成り立つ．ゆえに写像
$\colim H^0(V, \mathcal{F}'')
\to \lim H^0(U, \mathcal{F}''/I^n\mathcal{F}'')$ は同型である．一方，
$\mathfrak p_n \not \in V(\mathfrak p_1 \ldots \mathfrak p_{n - 1} I)$
なので，$\mathfrak p_n$ は $M'$ の随伴素イデアルではない．
したがって帰納法の仮定を $M'$ に適用できる．$\mathcal{F}' \subset \mathcal{F}$ であるから，
$\text{depth}(\mathcal{F}'_x) \geq 1$ という条件は，点 $x$ が
$\overline{\{x\}} \cap V(I) \subset V(\mathfrak a)$ を満たすときに成り立つことに注意する．
Algebra, 補題 \ref{algebra-lemma-depth-in-ses} を参照されたい．したがって写像
$\colim H^0(V, \mathcal{F}')
\to \lim H^0(U, \mathcal{F}'/I^n\mathcal{F}')$ も同型である．
補題 \ref{algebraization-lemma-ses-H0} により結論を得る．
\end{proof}

\begin{lemma}
\label{algebraization-lemma-application-H0}
$I \subset \mathfrak a$ を Noether 環 $A$ のイデアルとする．
$\mathcal{F}$ を
$U = \Spec(A) \setminus V(\mathfrak a)$ 上の連接加群とする．
次を仮定する．
\begin{enumerate}
\item $A$ は $I$-進完備であり，双対化複体をもつ．
\item $x \in \text{Ass}(\mathcal{F})$，$x \not \in V(I)$，
$\overline{\{x\}} \cap V(I) \not \subset V(\mathfrak a)$，かつ
$z \in V(\mathfrak a) \cap \overline{\{x\}}$ ならば，
$\dim(\mathcal{O}_{\overline{\{x\}}, z}) > \text{cd}(A, I) + 1$ である．
\item $x \in U$ が $\overline{\{x\}} \cap V(I) \subset V(\mathfrak a)$ を満たすならば，
$\text{depth}(\mathcal{F}_x) \geq 2$ である．
\end{enumerate}
このとき同型
$$
H^0(U, \mathcal{F})
\longrightarrow
\lim H^0(U, \mathcal{F}/I^n\mathcal{F})
$$
を得る．
\end{lemma}

\begin{proof}
$\hat s \in \lim H^0(U, \mathcal{F}/I^n\mathcal{F})$ とする．
命題 \ref{algebraization-proposition-application-H0} により，$\hat s$ はある元 $s \in \mathcal{F}(V)$ の像であり，
ここで $V \subset U$ は $U \cap V(I)$ を含む開集合である．
しかし条件 (3) により，$\text{depth}(\mathcal{F}_x) \geq 2$ が
すべての $x \in U \setminus V$ に対して成り立つ．したがって Divisors, 補題
\ref{divisors-lemma-depth-2-hartog} により
$\mathcal{F}(V) = \mathcal{F}(U)$ であり，証明が完了する．
\end{proof}

\begin{lemma}
\label{algebraization-lemma-alternative-colim-H0}
$A$ を Noether 環とする．$f \in \mathfrak a \subset A$ を $A$ のイデアルの元とする．
$M$ を有限 $A$-加群とする．次を仮定する．
\begin{enumerate}
\item $A$ は $f$-進完備である．
\item $f$ は $M$ 上の非零因子である．
\item $H^1_\mathfrak a(M/fM)$ は有限 $A$-加群である．
\end{enumerate}
このとき，$U = \Spec(A) \setminus V(\mathfrak a)$ と置けば，写像
$$
\colim_V \Gamma(V, \widetilde{M})
\longrightarrow
\lim \Gamma(U, \widetilde{M/f^nM})
$$
は同型である．ここで余極限は，$V \subset U$ で $U \cap V(f)$ を含む開集合全体にわたる．
\end{lemma}

\begin{proof}
$\mathcal{F} = \widetilde{M}|_U$ と置く．$H^1_\mathfrak a(M/fM)$ の有限性から
$H^0(U, \mathcal{F}/f\mathcal{F})$ は有限である．Local Cohomology, 補題
\ref{local-cohomology-lemma-finiteness-pushforwards-and-H1-local} を参照されたい．
Cohomology, 補題 \ref{cohomology-lemma-limit-finite}（$f$ は $\mathcal{F}$ 上の
非零因子なので適用できる）により，$N = \lim H^0(U, \mathcal{F}/f^n\mathcal{F})$ は
有限 $A$-加群であり，$f$-捩れをもたず，
$N/fN \subset H^0(U, \mathcal{F}/f\mathcal{F})$ である．一方，写像 $M \to N$ と，
それと両立する写像
$$
M/fM \longrightarrow H^0(U, \mathcal{F}/f\mathcal{F})
$$
がある．$g \in \mathfrak a$ に対し，$(M/fM)_g$ は
$H^0(U \cap D(f), \mathcal{F}/f\mathcal{F})$ へ同型に写る．なぜなら
$\mathcal{F}/f\mathcal{F}$ は $\widetilde{M/fM}$ の $U$ への制限だからである．
したがって $M_g \to N_g$ は同型
$$
M_g/fM_g = (M/fM)_g \to (N/fN)_g = N_g/fN_g
$$
を誘導する．$f$ は $N$ と $M$ の双方で非零因子なので，$M_g \to N_g$ は
$f$-進完備化の間に同型を誘導する．したがって $M_g \to N_g$ は
$V(f) \cap D(g)$ のある開近傍上で同型である．$g \in \mathfrak a$ は任意だったから，
$M$ と $N$ は補題の主張に現れるある開集合 $V$ 上で同型な連接加群を定める．
これで証明が完了する．
\end{proof}

\begin{proposition}
\label{algebraization-proposition-alternative-colim-H0}
$A$ を Noether 環とする．$f \in \mathfrak a \subset A$ を $A$ のイデアルの元とする．
$\mathcal{F}$ を $U = \Spec(A) \setminus V(\mathfrak a)$ 上の連接加群とする．次を仮定する．
\begin{enumerate}
\item $A$ は $f$-進完備であり，双対化複体をもつ．
\item $x \in \text{Ass}(\mathcal{F})$，$x \not \in V(f)$，
$\overline{\{x\}} \cap V(f) \not \subset V(\mathfrak a)$,
かつ $z \in \overline{\{x\}} \cap V(\mathfrak a)$ ならば，
$\dim(\mathcal{O}_{\overline{\{x\}}, z}) > 2$ である．
\end{enumerate}
このとき写像
$$
\colim_V \Gamma(V, \mathcal{F})
\longrightarrow
\lim \Gamma(U, \mathcal{F}/f^n\mathcal{F})
$$
は同型である．ここで余極限は，$V \subset U$ で $U \cap V(f)$ を含む開集合全体にわたる．
\end{proposition}

\begin{proof}[第1の証明]
$A$ は普遍鎖状であり，Gorenstein 形式的ファイバーをもつことを想起する．
Dualizing Complexes, 補題
\ref{dualizing-lemma-dualizing-gorenstein-formal-fibres} および
\ref{dualizing-lemma-universally-catenary} を参照されたい．したがって，
Local Cohomology, 補題 \ref{local-cohomology-lemma-make-S2-along-Z} で構成された写像
$\mathcal{F} \to \mathcal{F}'$ を考えられる．これは閉部分集合 $V(f) \cap U$ を
$U$ の中で考えた構成である．次を観察する．
\begin{enumerate}
\item $\mathcal{F} \to \mathcal{F}'$ の核と余核は $V(f) \cap U$ 上に台をもつ．
\item $\mathcal{F}'$ は $f$-捩れをもたない．実際，その茎の深さは $\geq 1$ であり，
これは $V(f) \cap U$ のすべての点で成り立つ．すなわち $\mathcal{F}'$ は $V(f) \cap U$ に随伴点をもたない．
\item $y \in V(f) \cap U$ が $\mathcal{F}'/f\mathcal{F}'$ の随伴点ならば，
$\text{depth}(\mathcal{F}'_y) = 1$ であり，したがって（$\mathcal{F}'$ の構成により）
直後の特殊化 $x \leadsto y$ で，$x \not \in V(f)$ が $\mathcal{F}$ の随伴点となるものが存在する．
仮定 (2) により，$y$ は $\Spec(A)$ において点 $z \in V(\mathfrak a)$ への直後の特殊化をもてない．
\item (3) より $H^0(U, \mathcal{F}'/f\mathcal{F}')$ は有限 $A$-加群である．
Local Cohomology, 補題 \ref{local-cohomology-lemma-finiteness-Rjstar} を参照されたい．
\end{enumerate}
これらの観察を用いて証明を終える．

\medskip\noindent
まず，命題の主張は $\mathcal{F}'$ に対して成り立つと主張する．
実際，有限 $A$-加群 $M'$ を，$\mathcal{F}'$ が $U$ 上で $M'$ に付随する連接加群の
制限となるように選ぶ．Local Cohomology, 補題
\ref{local-cohomology-lemma-finiteness-pushforwards-and-H1-local} を参照されたい．
$\mathcal{F}'$ は $f$-捩れをもたないので，$M'$ も $f$-捩れをもたないと仮定してよい．
上の観察 (4) より $H^1_\mathfrak a(M')$ は有限 $A$-加群である．Local Cohomology, 補題
\ref{local-cohomology-lemma-finiteness-pushforwards-and-H1-local} を参照されたい．
したがって補題 \ref{algebraization-lemma-alternative-colim-H0} により主張が従う．

\medskip\noindent
次に，命題の主張は任意の連接 $\mathcal{O}_U$-加群で $V(f) \cap U$ 上に台をもつものに対して
自明に成り立つことを観察する．$\mathcal{K}$，$\mathcal{G}$，$\mathcal{Q}$ をそれぞれ写像
$\mathcal{F} \to \mathcal{F}'$ の核，像，余核とする．短完全列
$0 \to \mathcal{G} \to \mathcal{F}' \to \mathcal{Q} \to 0$
と補題 \ref{algebraization-lemma-ses-H0} により，結論は $\mathcal{G}$ に対して成り立つ．次に短完全列
$0 \to \mathcal{K} \to \mathcal{F} \to \mathcal{G} \to 0$
について同じ議論を行えば，証明が完了する．
\end{proof}

\begin{proof}[第2の証明]
この命題は命題 \ref{algebraization-proposition-application-H0} の特別な場合である．
\end{proof}

\begin{lemma}
\label{algebraization-lemma-alternative-H0}
$A$ を Noether 環とする．$f \in \mathfrak a \subset A$ を $A$ のイデアルの元とする．
$M$ を有限 $A$-加群とする．次を仮定する．
\begin{enumerate}
\item $A$ は $f$-進完備である．
\item $H^1_\mathfrak a(M)$ および $H^2_\mathfrak a(M)$ は $f$ のある冪で零化される．
\end{enumerate}
このとき，$U = \Spec(A) \setminus V(\mathfrak a)$ と置けば，写像
$$
\Gamma(U, \widetilde{M})
\longrightarrow
\lim \Gamma(U, \widetilde{M/f^nM})
$$
は同型である．
\end{lemma}

\begin{proof}
補題 \ref{algebraization-lemma-formal-functions-principal} を $U$ および
$\mathcal{F} = \widetilde{M}|_U$ に適用できる．実際，$\mathcal{F}$ は連接
$\mathcal{O}_U$-加群の圏における Noether 対象である．
$H^1(U, \mathcal{F}) = H^2_\mathfrak a(M)$ は（Local Cohomology, 補題
\ref{local-cohomology-lemma-finiteness-pushforwards-and-H1-local})
により $f$ のある冪で零化されるので，その $f$-進 Tate 加群は零である．
したがって補題により，$\lim H^0(U, \mathcal{F}/f^n \mathcal{F})$ は，
$f$-進位相に関する $H^0(U, \mathcal{F})$ の通常の完備化に等しい．短完全列
$$
0 \to M/H^0_\mathfrak a(M) \to H^0(U, \mathcal{F}) \to H^1_\mathfrak a(M) \to 0
$$
がある．これは Local Cohomology, 補題
\ref{local-cohomology-lemma-finiteness-pushforwards-and-H1-local} による．
$M/H^0_\mathfrak a(M)$ は有限 $A$-加群なので完備である．Algebra, 補題
\ref{algebra-lemma-completion-tensor} を参照されたい．
$H^1_\mathfrak a(M)$ は $f$ のある冪で零化されるので，Algebra, 補題
\ref{algebra-lemma-completion-differ-by-torsion} により
$H^0(U, \mathcal{F})$ も完備である．これで証明が完了する．
\end{proof}







\section{形式切断の代数化 III}
\label{algebraization-section-algebraization-sections-coherent-III}

\noindent
次節では，局所コホモロジーの完備化に関する結果を，準アフィンスキーム上の
連接加群の高次コホモロジーおよびイデアルに関するその完備化へ応用する例を，
網羅的ではない形で列挙する．

\begin{proposition}
\label{algebraization-proposition-application-higher}
$I \subset \mathfrak a$ を Noether 環 $A$ のイデアルとする．
$\mathcal{F}$ を
$U = \Spec(A) \setminus V(\mathfrak a)$ 上の連接加群とする．
$s \geq 0$ とする．次を仮定する．
\begin{enumerate}
\item $A$ は $I$-進完備であり，双対化複体をもつ．
\item $x \in U \setminus V(I)$ ならば，$\text{depth}(\mathcal{F}_x) > s$ であるか，または
$$
\text{depth}(\mathcal{F}_x) +
\dim(\mathcal{O}_{\overline{\{x\}}, z}) > \text{cd}(A, I) + s + 1
$$
がすべての $z \in V(\mathfrak a) \cap \overline{\{x\}}$ に対して成り立つ．
\item 次の条件のいずれかが成り立つ．
\begin{enumerate}
\item $\mathcal{F}$ の $U \setminus V(I)$ への制限は $(S_{s + 1})$ を満たす．
\item $V(\mathfrak a)$ の次元は高々 $2$ である\footnote{
$V(\mathfrak a)$ 上で，$\Spec(A)$ 上の次元関数がとる最大値と最小値との差が
高々 $2$ である，という意味である．}．
\end{enumerate}
\end{enumerate}
このとき写像
$$
H^i(U, \mathcal{F})
\longrightarrow
\lim H^i(U, \mathcal{F}/I^n\mathcal{F})
$$
は $i < s$ に対して同型である．さらに同型
$$
\colim H^s(V, \mathcal{F})
\longrightarrow
\lim H^s(U, \mathcal{F}/I^n\mathcal{F})
$$
を得る．ここで余極限は，$V \subset U$ で $U \cap V(I)$ を含む開集合全体にわたる．
\end{proposition}

\begin{proof}
$s > 0$ と仮定してよい．$s = 0$ の場合は命題 \ref{algebraization-proposition-application-H0} で証明した．

\medskip\noindent
有限 $A$-加群 $M$ を，$\mathcal{F}$ が $U$ 上で $M$ に付随する連接加群の制限となるように選ぶ．
Local Cohomology, 補題 \ref{local-cohomology-lemma-finiteness-pushforwards-and-H1-local} を参照されたい．
$d = \text{cd}(A, I)$ と置く．$\mathfrak p$ を $A$ の素イデアルで $V(I)$ に含まれないものとし，
$\mathfrak q \in V(\mathfrak p) \cap V(\mathfrak a)$ とする．このとき
$\text{depth}(M_\mathfrak p) \geq s + 1 > s$ であるか，または (2) により
$\dim((A/\mathfrak p)_\mathfrak q) > d + s + 1$ である．補題 \ref{algebraization-lemma-bootstrap-bis-bis} により，
状況 \ref{algebraization-situation-bootstrap} の仮定は $A, I, V(\mathfrak a), M, s, d$ に対して満たされる．
一方，補題 \ref{algebraization-lemma-algebraize-local-cohomology-general} の仮定は $s + 1$ および $d$ に対して
満たされる．ここで条件 (3) を用いる．

\medskip\noindent
補題 \ref{algebraization-lemma-algebraize-local-cohomology-general} を適用すると，イデアル
$J_0 \subset \mathfrak a$ で $V(J_0) \cap V(I) = V(\mathfrak a)$ を満たすものが存在し，
任意の $J \subset J_0$ で $V(J) \cap V(I) = V(\mathfrak a)$ を満たすものに対して写像
$$
H^i_J(M) \longrightarrow H^i(R\Gamma_\mathfrak a(M)^\wedge)
$$
は $i \leq s + 1$ に対して同型である．

\medskip\noindent
$i \leq s$ に対して，写像 $H^i_\mathfrak a(M) \to H^i_J(M)$ は補題
\ref{algebraization-lemma-bootstrap-inherited} および \ref{algebraization-lemma-kill-colimit-support-general} により同型である．
コホモロジーと局所コホモロジーの比較（Local Cohomology, 補題
\ref{local-cohomology-lemma-local-cohomology}）を用いると，写像
$H^i(U, \mathcal{F}) \to H^i(V,\mathcal{F})$
は $V = \Spec(A) \setminus V(J)$ および $i < s$ に対して同型である．

\medskip\noindent
定理 \ref{algebraization-theorem-final-bootstrap} により
$H^i_\mathfrak a(M) = \lim H^i_\mathfrak a(M/I^nM)$
が $i \leq s$ に対して成り立つ．補題 \ref{algebraization-lemma-combine-two} により
$H^{s + 1}_\mathfrak a(M) = \lim H^{s + 1}_\mathfrak a(M/I^nM)$.

\medskip\noindent
同型 $H^0(U, \mathcal{F}) = H^0(V, \mathcal{F}) =
\lim H^0(U, \mathcal{F}/I^n\mathcal{F})$ は上の議論と命題
\ref{algebraization-proposition-application-H0} から従う．$0 < i < s$ に対しても，望む同型
$H^i(U, \mathcal{F}) = H^i(V, \mathcal{F}) =
\lim H^i(U, \mathcal{F}/I^n\mathcal{F})$ を，局所コホモロジーとコホモロジーの関係から
同様に得る．$i = 0$ の場合より容易である．なぜなら $i > 0$ に対して
$$
H^i(U, \mathcal{F}) = H^{i + 1}_\mathfrak a(M),
\quad
H^i(V, \mathcal{F}) = H^{i + 1}_J(M),
\quad
H^i(R\Gamma(U, \mathcal{F})^\wedge) = 
H^{i + 1}(R\Gamma_\mathfrak a(M)^\wedge)
$$
だからである．最後の主張についても同様である．
\end{proof}

\begin{lemma}
\label{algebraization-lemma-alternative-higher}
$A$ を Noether 環とする．$f \in \mathfrak a \subset A$ を $A$ のイデアルの元とする．
$M$ を有限 $A$-加群とする．$s \geq 0$ とする．次を仮定する．
\begin{enumerate}
\item $A$ は $f$-進完備である．
\item $H^i_\mathfrak a(M)$ は $f$ のある冪で零化される．これは $i \leq s + 1$ に対して成り立つ．
\end{enumerate}
このとき，$U = \Spec(A) \setminus V(\mathfrak a)$ と置けば，写像
$$
H^i(U, \widetilde{M})
\longrightarrow
\lim H^i(U, \widetilde{M/f^nM})
$$
は $i < s$ に対して同型である．
\end{lemma}

\begin{proof}
$s$ に関する帰納法による．$s = 0$ ならば主張は空である．$s = 1$ ならば，
結論は補題 \ref{algebraization-lemma-alternative-H0} である．$s > 1$ と仮定する．帰納法により，
$i = s - 1 \geq 1$ に対して結論を証明すれば十分である．
補題 \ref{algebraization-lemma-formal-functions-principal} を $U$ および
$\mathcal{F} = \widetilde{M}|_U$ に適用できる．実際，$\mathcal{F}$ は連接
$\mathcal{O}_U$-加群の圏における Noether 対象である．Local Cohomology, 補題
\ref{local-cohomology-lemma-finiteness-pushforwards-and-H1-local} により，
$H^j(U, \mathcal{F}) = H^{j + 1}_\mathfrak a(M)$ である．これはすべての $j$ に対して成り立つ．したがって，
$j = s = (s - 1) + 1$ の場合，これは仮定により $f$ のある冪で零化される．
ゆえに補題 \ref{algebraization-lemma-formal-functions-principal} から，
$\lim H^{s - 1}(U, \mathcal{F}/f^n\mathcal{F})$
は $f$-進位相に関する $H^{s - 1}(U, \mathcal{F})$ の通常の完備化である．
さらにこの加群は $f$ のある冪で零化されるので，その完備化は単に
$H^{s - 1}(U, \mathcal{F})$ 自身に等しい．これが望む結論である．
\end{proof}





\section{連結性への応用}
\label{algebraization-section-connected}

\noindent
本節では Grothendieck の連結性定理とその変種を論じる．原型は
\cite[Exposee XIII, 定理 2.1]{SGA2} にある．文献には Faltings の連結性定理と
呼ばれる版もあり，これは \cite[定理 6]{Faltings-some} を指すものと思われる．
\cite[定理 1.6]{Varbaro} に与えられた完備局所環に対する最適な形を述べ，証明しよう．

\begin{lemma}
\label{algebraization-lemma-punctured-still-connected}
\begin{reference}
\cite[定理 1.6]{Varbaro}
\end{reference}
$(A, \mathfrak m)$ を Noether 完備局所環とする．$I$ を $A$ の真のイデアルとする．
$X = \Spec(A)$ および $Y = V(I)$ と置く．次のように記す．
\begin{enumerate}
\item $d$ は $X$ の既約成分の次元の最小値である．
\item $c$ は，閉部分集合 $Z \subset X$ で $X \setminus Z$ が非連結となるものの次元の最小値である．
\end{enumerate}
このとき，閉部分集合 $Z \subset Y$ に対し，$Y \setminus Z$ は
$\dim(Z) < \min(c, d - 1) - \text{cd}(A, I)$ ならば連結である．特に，
$A/I$ の穴あきスペクトルは $\text{cd}(A, I) < \min(c, d - 1)$ ならば連結である．
\end{lemma}

\begin{proof}
まず最後の主張を証明する．第1の場合として，$A/I$ の穴あきスペクトルが空ならば，
Local Cohomology, 補題 \ref{local-cohomology-lemma-cd-bound-dim-local} により，
$X$ の各既約成分の次元は $\leq \text{cd}(A, I)$ である．したがって
$\min(c, d - 1) - \text{cd}(A, I) < 0$ を得て，この場合の補題が従う．ゆえに
$U \cap Y$ は空でないと仮定してよい．ここで $U = X \setminus \{\mathfrak m\}$ は
$A$ の穴あきスペクトルである．$A$ をその被約化で置き換えてよい．
$A$ は双対化複体をもち（Dualizing Complexes, 補題
\ref{dualizing-lemma-ubiquity-dualizing}），さらに $A$ は $I$ に関して完備である（Algebra, 補題
\ref{algebra-lemma-complete-by-sub}）ことに注意する．
$d - 1 > \text{cd}(A, I)$ と仮定すれば，補題 \ref{algebraization-lemma-application-theorem} により，写像
$$
\colim H^0(V, \mathcal{O}_V)
\longrightarrow
\lim H^0(U, \mathcal{O}_U/I^n\mathcal{O}_U)
$$
は同型である．ここで余極限は，$V \subset U$ で $U \cap Y$ を含む開集合全体にわたる．
$U \cap Y$ が非連結ならば，その $n$ 次無限小近傍も $U$ の中で，すべての $n$ に対して
非連結である．したがって右辺は非自明な冪等元をもつ（ここで $U \cap Y$ が空でないことを用いる）．
ゆえに非連結な $V$ が見つかる．$Z = X \setminus V$ と置く．Local Cohomology, 補題
\ref{local-cohomology-lemma-cd-bound-dim-local} により，$Z$ の各既約成分の次元は
$\leq \text{cd}(A, I)$ である．よって $c \leq \text{cd}(A, I)$ となり，最後の主張が証明された．

\medskip\noindent
今証明したことから，補題の主張を次のように導ける．閉部分集合 $Z \subset Y$ に対して
$Y \setminus Z$ が非連結であり，$\dim(Z) = e$ と仮定する．規約により連結空間は空でないことを想起する．
したがって，(a) $Y = Z$ であるか，または (b) $Y \setminus Z = W_1 \amalg W_2$ であり，
$W_i$ は $Y \setminus Z$ の空でない開かつ閉な部分集合である．場合 (b) では，
$w_i \in W_i$ を $U$ の閉点となるように選べる．Morphisms, 補題
\ref{morphisms-lemma-ubiquity-Jacobson-schemes} を参照されたい．このとき
$f_1, \ldots, f_e \in \mathfrak m$ で $V(f_1, \ldots, f_e) \cap Z = \{\mathfrak m\}$ を満たすものを
見つけられ，場合 (b) では $w_i \in V(f_1, \ldots, f_e)$ とも仮定できる．実際，素イデアル回避を
帰納的に用いて，$f_i$ を $\dim V(f_1, \ldots, f_i) \cap Z = e - i$ を満たし，かつ場合 (b) では
$w_1, w_2 \in V(f_i)$ となるように選べる．すると $A/I + (f_1, \ldots, f_e)$ の
穴あきスペクトルは非連結である（細部は省略する）．また
$\text{cd}(A, I + (f_1, \ldots, f_e)) \leq \text{cd}(A, I) + e$
が Local Cohomology, 補題 \ref{local-cohomology-lemma-cd-sum} および
\ref{local-cohomology-lemma-bound-cd} により成り立つから，第1部分より
$$
\text{cd}(A, I) + e \geq \min(c, d - 1)
$$
を得る．これは $e \geq \min(c, d - 1) - \text{cd}(A, I)$ を意味し，示すべきことである．
\end{proof}

\begin{lemma}
\label{algebraization-lemma-connected}
$I \subset \mathfrak a$ を Noether 環 $A$ のイデアルとする．次を仮定する．
\begin{enumerate}
\item $A$ は $I$-進完備であり，双対化複体をもつ．
\item $\mathfrak p \subset A$ が $V(I)$ に含まれない極小素イデアルで，
$\mathfrak q \in V(\mathfrak p) \cap V(\mathfrak a)$ ならば，
$\dim((A/\mathfrak p)_\mathfrak q) > \text{cd}(A, I) + 1$ である．
\item $V \subset \Spec(A)$ が空でない開集合で $V(I) \setminus V(\mathfrak a)$ を含むならば，
これは連結である\footnote{たとえば $A$ が整域の場合である．}．
\end{enumerate}
このとき $V(I) \setminus V(\mathfrak a)$ は空であるか，または連結である．
\end{lemma}

\begin{proof}
$A$ をその被約化で置き換えてよい．すると (2) の不等式は $A$ のすべての随伴素イデアルに対して成り立つ．
命題 \ref{algebraization-proposition-application-H0} により
$$
\colim H^0(V, \mathcal{O}_V) = \lim H^0(T_n, \mathcal{O}_{T_n})
$$
である．ここで余極限は (3) の開集合 $V$ 全体にわたり，$T_n$ は $n$ 次無限小近傍で，
$T = V(I) \setminus V(\mathfrak a)$ のものを
$U = \Spec(A) \setminus V(\mathfrak a)$ の中でとったものである．したがって $T$ は空であるか連結である．
実際，そうでなければ右辺は非自明な冪等元をもつが，左辺はもたないと仮定した．細部は省略する．
\end{proof}

\begin{lemma}
\label{algebraization-lemma-H0}
$A$ を，双対化複体をもち，非零元 $f \in A$ に関して完備な Noether 整域とする．
$f \in \mathfrak a \subset A$ をイデアルとする．$Z = V(\mathfrak a)$ の各既約成分が
余次元 $> 2$ を $X = \Spec(A)$ においてもつと仮定する．すなわち，$Z$ の各既約成分が
余次元 $> 1$ を $Y = V(f)$ においてもつと仮定する．このとき $Y \setminus Z$ は連結である．
\end{lemma}

\begin{proof}
これは補題 \ref{algebraization-lemma-connected} の特別な場合である（その証明は命題
\ref{algebraization-proposition-application-H0} に依存する）．以下では，より容易な命題
\ref{algebraization-proposition-alternative-colim-H0} を用いて証明する．

\medskip\noindent
$U = X \setminus Z$ と置く．命題 \ref{algebraization-proposition-alternative-colim-H0} により同型
$$
\colim \Gamma(V, \mathcal{O}_V) \to
\lim_n \Gamma(U, \mathcal{O}_U/f^n \mathcal{O}_U)
$$
を得る．ここで余極限は，$V \subset U$ で $U \cap Y$ を含む開集合全体にわたる．
$U \cap Y$ が非連結ならば，ある $V$ に対して $\Gamma(V, \mathcal{O}_V)$ に非自明な冪等元が存在する．
これは不可能である．実際，$V$ は整スキームであり，$X$ は整域のスペクトルだからである．
\end{proof}






\section{完備化関手}
\label{algebraization-section-completion}

\noindent
$X$ を Noether スキームとする．$Y \subset X$ を，準連接イデアル層
$\mathcal{I} \subset \mathcal{O}_X$ をもつ閉部分スキームとする．
本節では，連接 $\mathcal{O}_X$-加群の逆系 $(\mathcal{F}_n)$ であって，
$\mathcal{F}_n$ が $I^n$ により零化され，遷移写像が同型
$\mathcal{F}_{n + 1}/I^n\mathcal{F}_{n + 1} \to \mathcal{F}_n$ を誘導するものを考える．
これらの逆系の圏は
$$
\textit{Coh}(X, \mathcal{I})
$$
Cohomology of Schemes，第 \ref{coherent-section-coherent-formal} 節において
このように表された．この圏は，$X$ の $Y$ に沿う形式完備化上の連接加群の圏と
同値である．しかし，ここではまだ形式スキームもその上の連接加群も導入して
いないので，この用語を用いることはできない．特に関心があるのは完備化関手
$$
\textit{Coh}(\mathcal{O}_X)
\longrightarrow
\textit{Coh}(X, \mathcal{I}),\quad
\mathcal{F} \longmapsto \mathcal{F}^\wedge
$$
Cohomology of Schemes，式 (\ref{coherent-equation-completion-functor}) を参照されたい．

\begin{lemma}
\label{algebraization-lemma-completion-fully-faithful}
$X$ を Noether スキームとし，$Y \subset X$ を閉部分スキームとする．
$Y_n \subset X$ を第 $n$ 無限小近傍，すなわち $Y$ の $X$ における無限小近傍とする．
次の条件を考える．
\begin{enumerate}
\item $X$ は準アフィンであり，
$\Gamma(X, \mathcal{O}_X) \to \lim \Gamma(Y_n, \mathcal{O}_{Y_n})$
は同型である．
\item $X$ は豊富な可逆加群 $\mathcal{L}$ をもち，
$\Gamma(X, \mathcal{L}^{\otimes m}) \to
\lim \Gamma(Y_n, \mathcal{L}^{\otimes m}|_{Y_n})$
はすべての $m \gg 0$ に対して同型である．
\item 任意の有限局所自由 $\mathcal{O}_X$-加群 $\mathcal{E}$ に対して写像
$\Gamma(X, \mathcal{E}) \to \lim \Gamma(Y_n, \mathcal{E}|_{Y_n})$
は同型である．
\item 完備化関手
$\textit{Coh}(\mathcal{O}_X) \to \textit{Coh}(X, \mathcal{I})$
は，有限局所自由な対象からなる充満部分圏上で全忠実である．
\end{enumerate}
このとき (1) $\Rightarrow$ (2) $\Rightarrow$ (3) $\Rightarrow$ (4) かつ
(4) $\Rightarrow$ (3) である．
\end{lemma}

\begin{proof}
(3) $\Rightarrow$ (4) の証明．$\mathcal{F}$ と $\mathcal{G}$ が $X$ 上有限局所自由であるとする．
$\mathcal{H} = \SheafHom_{\mathcal{O}_X}(\mathcal{G}, \mathcal{F})$
を考え，Cohomology of Schemes，補題
\ref{coherent-lemma-completion-internal-hom} を用いると，(3) から (4) が従う．

\medskip\noindent
(2) $\rightarrow$ (3) の証明．すなわち，$\mathcal{L}$ を $X$ 上豊富とし，
$\mathcal{E}$ を有限局所自由 $\mathcal{O}_X$-加群とする．
ある普遍的完全列
$$
0 \to \mathcal{E} \to
(\mathcal{L}^{\otimes p})^{\oplus r} \to
(\mathcal{L}^{\otimes q})^{\oplus s}
$$
が，ある $r, s \geq 0$ と $0 \ll p \ll q$ に対して存在すると主張する．
これが成り立てば，完全列
$$
0 \to \lim \Gamma(\mathcal{E}|_{Y_n}) \to
\lim \Gamma((\mathcal{L}^{\otimes p})^{\oplus r}|_{Y_n}) \to
\lim \Gamma((\mathcal{L}^{\otimes q})^{\oplus s}|_{Y_n})
$$
と (2) の同型から，(3) の同型が得られる．この主張を証明するため，双対局所自由加群
$\SheafHom_{\mathcal{O}_X}(\mathcal{E}, \mathcal{O}_X)$
を考え，Properties，命題 \ref{properties-proposition-characterize-ample} を適用して，全射
$$
(\mathcal{L}^{\otimes -p})^{\oplus r}
\longrightarrow
\SheafHom_{\mathcal{O}_X}(\mathcal{E}, \mathcal{O}_X)
$$
を得る．双対を取ると完全列の最初の写像が得られる（全射であることは普遍的なので，
これは普遍的単射である）．余核に対して同じことを繰り返せば二番目の写像が得られる．
いくつかの細部は省略する．

\medskip\noindent
(1) $\Rightarrow$ (2) の証明．$X$ が準アフィンならば $\mathcal{O}_X$ は
豊富な可逆加群だからである．Properties，補題
\ref{properties-lemma-quasi-affine-O-ample} を参照されたい．

\medskip\noindent
(4) $\Rightarrow$ (3) の証明は省略する．
\end{proof}

\noindent
Noether スキームと準連接イデアル層 $\mathcal{I} \subset \mathcal{O}_X$ が与えられたとする．
対象 $(\mathcal{F}_n)$ が $\textit{Coh}(X, \mathcal{I})$ に属するとき，各 $\mathcal{F}_n$ が
有限局所自由 $\mathcal{O}_X/\mathcal{I}^n$-加群であるとき，{\it 有限局所自由}であるという．

\begin{lemma}
\label{algebraization-lemma-completion-fully-faithful-general}
$X$ を Noether スキームとし，$Y \subset X$ をイデアル層
$\mathcal{I} \subset \mathcal{O}_X$ をもつ閉部分スキームとする．
$Y_n \subset X$ を第 $n$ 無限小近傍，すなわち $Y$ の $X$ における無限小近傍とする．
$\mathcal{V}$ を，開部分スキーム $V \subset X$ で $Y$ を含むものの集合とし，
包含関係の逆で順序づける．
\begin{enumerate}
\item $X$ は準アフィンであり，
$$
\colim_\mathcal{V} \Gamma(V, \mathcal{O}_V)
\longrightarrow
\lim \Gamma(Y_n, \mathcal{O}_{Y_n})
$$
は同型である．
\item $X$ は豊富な可逆加群 $\mathcal{L}$ をもち，
$$
\colim_\mathcal{V} \Gamma(V, \mathcal{L}^{\otimes m})
\longrightarrow
\lim \Gamma(Y_n, \mathcal{L}^{\otimes m}|_{Y_n})
$$
はすべての $m \gg 0$ に対して同型である．
\item 任意の $V \in \mathcal{V}$ と任意の有限局所自由
$\mathcal{O}_V$-加群 $\mathcal{E}$ に対して写像
$$
\colim_{V' \geq V} \Gamma(V', \mathcal{E}|_{V'})
\longrightarrow
\lim \Gamma(Y_n, \mathcal{E}|_{Y_n})
$$
は同型である．
\item 完備化関手
$$
\colim_\mathcal{V} \textit{Coh}(\mathcal{O}_V)
\longrightarrow
\textit{Coh}(X, \mathcal{I}),
\quad
\mathcal{F} \longmapsto \mathcal{F}^\wedge
$$
は有限局所自由な対象からなる充満部分圏上で全忠実である（上の説明を参照）．
\end{enumerate}
このとき (1) $\Rightarrow$ (2) $\Rightarrow$ (3) $\Rightarrow$ (4) かつ
(4) $\Rightarrow$ (3) である．
\end{lemma}

\begin{proof}
$\mathcal{V}$ は有向集合なので，余極限は Categories，第
\ref{categories-section-directed-colimits} 節におけるものと同じであることに注意する．
残りの論証は補題 \ref{algebraization-lemma-completion-fully-faithful} の証明における論証と
ほとんどまったく同じであるため，読み飛ばすことを勧める．

\medskip\noindent
(3) $\Rightarrow$ (4) の証明．$\mathcal{F}$ と $\mathcal{G}$ が
$V \in \mathcal{V}$ 上有限局所自由であるとする．
$\mathcal{H} = \SheafHom_{\mathcal{O}_V}(\mathcal{G}, \mathcal{F})$
を考え，Cohomology of Schemes，補題
\ref{coherent-lemma-completion-internal-hom} を用いると，(3) から (4) が従う．

\medskip\noindent
(2) $\Rightarrow$ (3) の証明．$\mathcal{L}$ を $X$ 上豊富とし，
$\mathcal{E}$ を有限局所自由 $\mathcal{O}_V$-加群とする．ここで
$V \in \mathcal{V}$ である．ある普遍的完全列
$$
0 \to \mathcal{E} \to
(\mathcal{L}^{\otimes p})^{\oplus r}|_{V} \to
(\mathcal{L}^{\otimes q})^{\oplus s}|_{V}
$$
が，ある $r, s \geq 0$ と $0 \ll p \ll q$ に対して存在すると主張する．
これが成り立てば，(2) の同型から (3) の同型が従う．この主張を証明するため，
$\mathcal{L}|_V$ が豊富であることを思い出す．Properties，補題
\ref{properties-lemma-ample-on-locally-closed} を参照されたい．双対局所自由加群
$\SheafHom_{\mathcal{O}_V}(\mathcal{E}, \mathcal{O}_V)$
を考え，Properties，命題 \ref{properties-proposition-characterize-ample} を適用して，全射
$$
(\mathcal{L}^{\otimes -p})^{\oplus r}|_V \longrightarrow
\SheafHom_{\mathcal{O}_V}(\mathcal{E}, \mathcal{O}_V)
$$
を得る（全射であることは普遍的なので，これは普遍的単射である）．
双対を取ると完全列の最初の写像が得られる．余核に対して同じことを繰り返せば
二番目の写像が得られる．いくつかの細部は省略する．

\medskip\noindent
(1) $\Rightarrow$ (2) の証明．$X$ が準アフィンならば $\mathcal{O}_X$ は
豊富な可逆加群だからである．Properties，補題
\ref{properties-lemma-quasi-affine-O-ample} を参照されたい．

\medskip\noindent
(4) $\Rightarrow$ (3) の証明は省略する．
\end{proof}

\begin{lemma}
\label{algebraization-lemma-recognize-formal-coherent-modules}
$X$ を Noether スキームとする．$\mathcal{I} \subset \mathcal{O}_X$ を
準連接イデアル層とする．関手
$$
\textit{Coh}(X, \mathcal{I}) \longrightarrow \text{Pro-}\QCoh(\mathcal{O}_X)
$$
は全忠実である．Categories，注意 \ref{categories-remark-pro-category} を参照されたい．
\end{lemma}

\begin{proof}
$(\mathcal{F}_n)$ と $(\mathcal{G}_n)$ を $\textit{Coh}(X, \mathcal{I})$ の対象とする．
pro-対象の射 $\alpha$ で $(\mathcal{F}_n)$ から $(\mathcal{G}_n)$ へのものは，
写像の系
$\alpha_n : \mathcal{F}_{m(n)} \to \mathcal{G}_n$
によって与えられる．ここで，これらは遷移写像と両立し，
$\mathbf{N} \to \mathbf{N}$，$n \mapsto m(n)$ は増加関数である
（特に $m(n) \geq n$）．
$\mathcal{F}_n = \mathcal{F}_{m(n)}/\mathcal{I}^n\mathcal{F}_{m(n)}$
であり，$\mathcal{G}_n$ は $\mathcal{I}^n$ によって零化されるので，
$\alpha_n$ は写像 $\mathcal{F}_n \to \mathcal{G}_n$ を誘導する．
\end{proof}

\noindent
次に，本章で先に得た成果を用いて証明できる種類の全忠実性結果について，
いくつかの例を加える．

\begin{lemma}
\label{algebraization-lemma-fully-faithful}
$I \subset \mathfrak a$ を Noether 環 $A$ のイデアルとする．
$U = \Spec(A) \setminus V(\mathfrak a)$ とする．次を仮定する．
\begin{enumerate}
\item $A$ は $I$-進完備であり，双対化複体をもつ．
\item $\mathfrak p \subset A$ が随伴素イデアルで，
$\mathfrak p \not \in V(I)$ かつ
$V(\mathfrak p) \cap V(I) \not \subset V(\mathfrak a)$ であり，さらに
$\mathfrak q \in V(\mathfrak p) \cap V(\mathfrak a)$ ならば，
$\dim((A/\mathfrak p)_\mathfrak q) > \text{cd}(A, I) + 1$ である．
\item $\mathfrak p \subset A$ が $\mathfrak p \not \in V(I)$ および
$V(\mathfrak p) \cap V(I) \subset V(\mathfrak a)$ を満たすならば，
$\text{depth}(A_\mathfrak p) \geq 2$ である．
\end{enumerate}
このとき完備化関手
$$
\textit{Coh}(\mathcal{O}_U)
\longrightarrow
\textit{Coh}(U, I\mathcal{O}_U),
\quad
\mathcal{F} \longmapsto \mathcal{F}^\wedge
$$
は有限局所自由な対象からなる充満部分圏上で全忠実である．
\end{lemma}

\begin{proof}
補題 \ref{algebraization-lemma-completion-fully-faithful} により，
$$
\Gamma(U, \mathcal{O}_U) =
\lim \Gamma(U, \mathcal{O}_U/I^n\mathcal{O}_U)
$$
を示せば十分である．これは補題 \ref{algebraization-lemma-application-H0} から直ちに従う．
\end{proof}

\begin{lemma}
\label{algebraization-lemma-fully-faithful-alternative}
$A$ を Noether 環とする．$f \in \mathfrak a \subset A$ を $A$ のイデアルの元とする．
$U = \Spec(A) \setminus V(\mathfrak a)$ とする．次を仮定する．
\begin{enumerate}
\item $A$ は $f$-進完備である．
\item $H^1_\mathfrak a(A)$ と $H^2_\mathfrak a(A)$ は $f$ のある冪により零化される．
\end{enumerate}
このとき完備化関手
$$
\textit{Coh}(\mathcal{O}_U)
\longrightarrow
\textit{Coh}(U, I\mathcal{O}_U),
\quad
\mathcal{F} \longmapsto \mathcal{F}^\wedge
$$
は有限局所自由な対象からなる充満部分圏上で全忠実である．
\end{lemma}

\begin{proof}
補題 \ref{algebraization-lemma-completion-fully-faithful} により，
$$
\Gamma(U, \mathcal{O}_U) =
\lim \Gamma(U, \mathcal{O}_U/I^n\mathcal{O}_U)
$$
を示せば十分である．これは補題 \ref{algebraization-lemma-alternative-H0} から直ちに従う．
\end{proof}

\begin{lemma}
\label{algebraization-lemma-fully-faithful-simple-two}
$A$ を Noether 環とする．$f \in \mathfrak a$ を $A$ のイデアルの元とする．
$U = \Spec(A) \setminus V(\mathfrak a)$ とする．次を仮定する．
\begin{enumerate}
\item $A$ は双対化複体をもち，$f$ に関して完備である．
\item 任意の素イデアル $\mathfrak p \subset A$ に対し，$f \not \in \mathfrak p$ かつ
$\mathfrak q \in V(\mathfrak p) \cap V(\mathfrak a)$ ならば，
$\text{depth}(A_\mathfrak p) + \dim((A/\mathfrak p)_\mathfrak q) > 2$ である．
\end{enumerate}
このとき完備化関手
$$
\textit{Coh}(\mathcal{O}_U)
\longrightarrow
\textit{Coh}(U, I\mathcal{O}_U),
\quad
\mathcal{F} \longmapsto \mathcal{F}^\wedge
$$
は有限局所自由な対象からなる充満部分圏上で全忠実である．
\end{lemma}

\begin{proof}
結論は補題 \ref{algebraization-lemma-fully-faithful-alternative} および
Local Cohomology，命題 \ref{local-cohomology-proposition-annihilator} から従う．
\end{proof}

\begin{lemma}
\label{algebraization-lemma-fully-faithful-general}
$I \subset \mathfrak a \subset A$ を Noether 環 $A$ のイデアルとする．
$U = \Spec(A) \setminus V(\mathfrak a)$ とする．$\mathcal{V}$ を，
$U$ の開部分スキームで $U \cap V(I)$ を含むものの集合とし，包含関係の逆で順序づける．
次を仮定する．
\begin{enumerate}
\item $A$ は $I$-進完備であり，双対化複体をもつ．
\item 随伴素イデアル $\mathfrak p \subset A$ が
$I \not \subset \mathfrak p$ かつ
$V(\mathfrak p) \cap V(I) \not \subset V(\mathfrak a)$ を満たし，さらに
$\mathfrak q \in V(\mathfrak p) \cap V(\mathfrak a)$ ならば，
$\dim((A/\mathfrak p)_\mathfrak q) > \text{cd}(A, I) + 1$ である．
\end{enumerate}
このとき完備化関手
$$
\colim_\mathcal{V} \textit{Coh}(\mathcal{O}_V)
\longrightarrow
\textit{Coh}(U, I\mathcal{O}_U),
\quad
\mathcal{F} \longmapsto \mathcal{F}^\wedge
$$
は有限局所自由な対象からなる充満部分圏上で全忠実である．
\end{lemma}

\begin{proof}
補題 \ref{algebraization-lemma-completion-fully-faithful-general} により，
$$
\colim_\mathcal{V} \Gamma(V, \mathcal{O}_V) =
\lim \Gamma(U, \mathcal{O}_U/I^n\mathcal{O}_U)
$$
を示せば十分である．これは命題 \ref{algebraization-proposition-application-H0} から直ちに従う．
\end{proof}

\begin{lemma}
\label{algebraization-lemma-fully-faithful-general-alternative}
$A$ を Noether 環とする．$f \in \mathfrak a \subset A$ を $A$ のイデアルの元とする．
$U = \Spec(A) \setminus V(\mathfrak a)$ とする．$\mathcal{V}$ を，$U$ の開部分スキームで
$U \cap V(f)$ を含むものの集合とし，包含関係の逆で順序づける．次を仮定する．
\begin{enumerate}
\item $A$ は $f$-進完備である．
\item $f$ は非零因子である．
\item $H^1_\mathfrak a(A/fA)$ は有限 $A$-加群である．
\end{enumerate}
このとき完備化関手
$$
\colim_\mathcal{V} \textit{Coh}(\mathcal{O}_V)
\longrightarrow
\textit{Coh}(U, f\mathcal{O}_U),
\quad
\mathcal{F} \longmapsto \mathcal{F}^\wedge
$$
は有限局所自由な対象からなる充満部分圏上で全忠実である．
\end{lemma}

\begin{proof}
補題 \ref{algebraization-lemma-completion-fully-faithful-general} により，
$$
\colim_\mathcal{V} \Gamma(V, \mathcal{O}_V) =
\lim \Gamma(U, \mathcal{O}_U/I^n\mathcal{O}_U)
$$
を示せば十分である．これは補題 \ref{algebraization-lemma-alternative-colim-H0} から直ちに従う．
\end{proof}

\begin{lemma}
\label{algebraization-lemma-fully-faithful-very-general}
$I \subset \mathfrak a \subset A$ を Noether 環 $A$ のイデアルとする．
$U = \Spec(A) \setminus V(\mathfrak a)$ とする．$\mathcal{V}$ を，
$U$ の開部分スキームで $U \cap V(I)$ を含むものの集合とし，包含関係の逆で順序づける．
$\mathcal{F}$ と $\mathcal{G}$ を連接 $\mathcal{O}_V$-加群とする．ここで
ある $V \in \mathcal{V}$ を取る．写像
$$
\colim_{V' \geq V} \Hom_V(\mathcal{G}|_{V'}, \mathcal{F}|_{V'})
\longrightarrow
\Hom_{\textit{Coh}(U, I\mathcal{O}_U)}(\mathcal{G}^\wedge, \mathcal{F}^\wedge)
$$
は，次の仮定のもとで全単射である．
\begin{enumerate}
\item $A$ は $I$-進完備であり，双対化複体をもつ．
\item $x \in \text{Ass}(\mathcal{F})$，$x \not \in V(I)$，
$\overline{\{x\}} \cap V(I) \not \subset V(\mathfrak a)$，かつ
$z \in \overline{\{x\}} \cap V(\mathfrak a)$ ならば，
$\dim(\mathcal{O}_{\overline{\{x\}}, z}) > \text{cd}(A, I) + 1$ である．
\end{enumerate}
\end{lemma}

\begin{proof}
連接 $\mathcal{O}_U$-加群 $\mathcal{F}'$ と $\mathcal{G}'$ で，$V$ への制限が
$\mathcal{F}$ と $\mathcal{G}$ であるものを選べる．Properties，補題
\ref{properties-lemma-lift-finite-presentation} を参照されたい．さらに $\mathcal{F}'$ の選択を
変更して，$\text{Ass}(\mathcal{F}') \subset V$ とできる．例えば Local Cohomology，補題
\ref{local-cohomology-lemma-get-depth-1-along-Z} を参照されたい．したがって，$V$ を $U$ で，
$\mathcal{F}$ と $\mathcal{G}$ を $\mathcal{F}'$ と $\mathcal{G}'$ で置き換えてよく，実際そうする．
$\mathcal{H} = \SheafHom_{\mathcal{O}_U}(\mathcal{G}, \mathcal{F})$ とおく．
これは連接 $\mathcal{O}_U$-加群である．次が成り立つ．
$$
\Hom_V(\mathcal{G}|_V, \mathcal{F}|_V) =
H^0(V, \mathcal{H})
\quad\text{かつ}\quad
\lim H^0(U, \mathcal{H}/\mathcal{I}^n\mathcal{H}) =
\Mor_{\textit{Coh}(U, I\mathcal{O}_U)}
(\mathcal{G}^\wedge, \mathcal{F}^\wedge)
$$
Cohomology of Schemes，補題 \ref{coherent-lemma-completion-internal-hom} を参照されたい．
したがって，命題 \ref{algebraization-proposition-application-H0} の仮定が $\mathcal{H}$ に対して
成り立つことを示せば，証明は完了する．これは
$\text{Ass}(\mathcal{H}) \subset \text{Ass}(\mathcal{F})$ だから成り立つ．
Cohomology of Schemes，補題 \ref{coherent-lemma-hom-into-depth} を参照されたい．
\end{proof}








\section{連接形式加群の代数化 I}
\label{algebraization-section-algebraization-modules}

\noindent
完備化関手の本質的全射性（以下を参照）は，
\cite{SGA2}，\cite{MRaynaud-book}，および \cite{MRaynaud-paper} において
体系的に研究された．以下のアフィンな状況で考える．

\begin{situation}
\label{algebraization-situation-algebraize}
ここで $A$ は Noether 環であり，$I \subset \mathfrak a \subset A$ はイデアルである．
$X = \Spec(A)$，$Y = V(I) = \Spec(A/I)$，および
$Z = V(\mathfrak a) = \Spec(A/\mathfrak a)$ とおく．さらに $U = X \setminus Z$ とおく．
\end{situation}

\noindent
本節では，$\textit{Coh}(U, I\mathcal{O}_U)$ の対象が完備化関手
$\textit{Coh}(\mathcal{O}_U) \to \textit{Coh}(U, I\mathcal{O}_U)$ の像に入ることを
保証する条件を探す．Cohomology of Schemes，第
\ref{coherent-section-coherent-formal} 節および第 \ref{algebraization-section-completion} 節を参照されたい．

\begin{lemma}
\label{algebraization-lemma-system-of-modules}
状況 \ref{algebraization-situation-algebraize} において考える．逆系 $(M_n)$ で，
$A$-加群からなり次を満たすものを考える．
\begin{enumerate}
\item $M_n$ は有限 $A$-加群である．
\item $M_n$ は $I^n$ により零化される．
\item $M_{n + 1}/I^nM_{n + 1} \to M_n$ の核と余核は
$\mathfrak a$-冪捩れである．
\end{enumerate}
このとき $(\widetilde{M}_n|_U)$ は $\textit{Coh}(U, I\mathcal{O}_U)$ に属する．
逆に，$\textit{Coh}(U, I\mathcal{O}_U)$ の任意の対象はこのようにして得られる．
\end{lemma}

\begin{proof}
$(\widetilde{M}_n|_U)$ が $\textit{Coh}(U, I\mathcal{O}_U)$ に属することの確認は省略する．
$(\mathcal{F}_n)$ を $\textit{Coh}(U, I\mathcal{O}_U)$ の対象とする．
Local Cohomology，補題
\ref{local-cohomology-lemma-finiteness-pushforwards-and-H1-local} により，
$\mathcal{F}_n = \widetilde{M_n}$ となる，ある有限 $A/I^n$-加群 $M_n$ が存在する．
$M_n$ を $H^0_\mathfrak a(M_n)$ で割ることにより，
$M_n \subset H^0(U, \mathcal{F}_n)$ と仮定してよい．Dualizing Complexes，補題
\ref{dualizing-lemma-divide-by-torsion} および既に参照した補題を参照されたい．
$M_{n + 1}$ を，$M_n$ の，写像 $M_{n + 1} \to H^0(U, \mathcal{F}_{n + 1})
\to H^0(U, \mathcal{F}_n)$ による逆像で帰納的に置き換えることにより，
$M_{n + 1}$ が $M_n$ に写ると仮定してよい．こうして (1) と (2) を満たす逆系
$(M_n)$ で $\mathcal{F}_n = \widetilde{M_n}$ となるものが得られる．(3) を見るには，
$M_{n + 1}/I^nM_{n + 1} \to M_n$ が有限 $A$-加群の写像であり，
$\widetilde{\ }$ を施して $U$ に制限すると同型を誘導することを用いる
（ここでも最初に参照した補題を用いる）．
\end{proof}

\noindent
状況 \ref{algebraization-situation-algebraize} では，Cohomology of Schemes，式
(\ref{coherent-equation-completion-functor}) の完備化関手
\begin{equation}
\label{algebraization-equation-completion}
\textit{Coh}(\mathcal{O}_U)
\longrightarrow
\textit{Coh}(U, I\mathcal{O}_U),\quad
\mathcal{F} \longmapsto \mathcal{F}^\wedge
\end{equation}
を考察できる．$A$ が $I$-進完備ならば，形式切断の代数化に関する先の結果により，
この関手は適当な部分圏上で全忠実である．例については第
\ref{algebraization-section-completion} 節および補題 \ref{algebraization-lemma-fully-faithful-inequalities} を参照されたい．
次に，$(\mathcal{F}_n)$ を $\textit{Coh}(U, I\mathcal{O}_U)$ の対象とする．
引き続き $A$ が $I$-進完備であると仮定すると，次を問うことができる：
$(\mathcal{F}_n)$ が上に表示した完備化関手の本質像に入るのはいつか．

\begin{lemma}
\label{algebraization-lemma-essential-image-completion}
状況 \ref{algebraization-situation-algebraize} において，$(\mathcal{F}_n)$ を
$\textit{Coh}(U, I\mathcal{O}_U)$ の対象とする．次の条件を考える．
\begin{enumerate}
\item $(\mathcal{F}_n)$ は関手 (\ref{algebraization-equation-completion}) の本質像に属する．
\item $(\mathcal{F}_n)$ は連接 $\mathcal{O}_U$-加群の完備化である．
\item $(\mathcal{F}_n)$ は連接 $\mathcal{O}_V$-加群の完備化である．ここで
$U \cap Y \subset V \subset U$ は開である．
\item $(\mathcal{F}_n)$ は，$U$ への，連接 $\mathcal{O}_X$-加群の制限を完備化したものである．
\item $(\mathcal{F}_n)$ は，$U$ への，連接 $\mathcal{O}_X$-加群の完備化の制限である．
\item 対象 $(\mathcal{G}_n)$ が $\textit{Coh}(X, I\mathcal{O}_X)$ に存在し，$U$ への制限が
$(\mathcal{F}_n)$ であるものが存在する．
\end{enumerate}
このとき条件 (1)，(2)，(3)，(4)，(5) は同値であり，(6) を含意する．
$A$ が $I$-進完備ならば，条件 (6) は他の条件を含意する．
\end{lemma}

\begin{proof}
(1) と (2) は同値である．実際，連接 $\mathcal{O}_U$-加群 $\mathcal{F}$ の完備化は，
定義により $\mathcal{F}$ の関手 (\ref{algebraization-equation-completion}) による像である．
$V \subset U$ が $U \cap Y$ を含む開部分スキームならば，
$$
\textit{Coh}(V, I\mathcal{O}_V) =
\textit{Coh}(U, I\mathcal{O}_U)
$$
が成り立つ．なぜなら，連接 $\mathcal{O}_V$-加群で $V \cap Y$ に台をもつものの圏は，
連接 $\mathcal{O}_U$-加群で $U \cap Y$ に台をもつものの圏と同じだからである．したがって，
連接 $\mathcal{O}_V$-加群の完備化は $\textit{Coh}(U, I\mathcal{O}_U)$ の対象である．
以上を踏まえると，(2)，(3)，(4)，(5) の同値性は，関手
$\textit{Coh}(\mathcal{O}_X) \to \textit{Coh}(\mathcal{O}_U) \to
\textit{Coh}(\mathcal{O}_V)$ が本質的全射であることから従う．
Properties，補題 \ref{properties-lemma-lift-finite-presentation} を参照されたい．

\medskip\noindent
(5) は常に (6) を含意する．$A$ が $I$-進完備であると仮定する．このとき
$\textit{Coh}(X, I\mathcal{O}_X)$ の任意の対象は，Cohomology of Schemes，補題
\ref{coherent-lemma-inverse-systems-affine} により有限 $A$-加群に対応する．
したがって，この場合には (6) が (5) を含意する．
\end{proof}

\begin{example}
\label{algebraization-example-not-algebraizable}
$k$ を体とする．$A = k[x, y][[t]]$，$I = (t)$，
$\mathfrak a = (x, y, t)$ とする．状況 \ref{algebraization-situation-algebraize} の記法を用いる．
$U \cap Y = (D(x) \cap Y) \cup (D(y) \cap Y)$ はアフィン開被覆であることに注意する．
$n \geq 1$ に対し，可逆加群 $\mathcal{L}_n$ で
$\mathcal{O}_U/t^n\mathcal{O}_U$ 上のものを次のように考える．$A_x/t^nA_x$ と $A_y/t^nA_y$ を，
$A_{xy}/t^nA_{xy}$ の可逆元，すなわち次の形の任意の冪級数の像を用いて貼り合わせる：
$$
u = 1 + \frac{t}{xy} + \sum_{n \geq 2} a_n \frac{t^n}{(xy)^{\varphi(n)}}
$$
ここで $a_n \in k[x, y]$ かつ $\varphi(n) \in \mathbf{N}$ である．
このとき $(\mathcal{L}_n)$ は $\textit{Coh}(U, I\mathcal{O}_U)$ の可逆な対象であるが，
連接 $\mathcal{O}_U$-加群 $\mathcal{L}$ の完備化ではない．論証の概略だけを述べ，
細部の大半は省略する．$y \in U \cap Y$ とする．このとき茎 $\mathcal{L}_y$ の完備化は
可逆加群となるはずであるから，$\mathcal{L}_y$ は可逆である．したがって，開集合
$V \subset U$ で，$U \cap Y$ を含み，$\mathcal{L}|_V$ が可逆となるものが存在するはずである．
Divisors，補題 \ref{divisors-lemma-extend-invertible-module} により，可逆 $A$-加群 $M$ で
$\widetilde{M}|_V \cong \mathcal{L}|_V$ となるものが得られる．しかし環 $A$ は一意分解整域なので，
$M \cong A$ である．これは $\mathcal{L}_n \cong \mathcal{O}_U/I^n\mathcal{O}_U$ を含意する．
構成により $\mathcal{L}_2 \not \cong \mathcal{O}_U/I^2\mathcal{O}_U$ なので，望む矛盾を得る．

\medskip\noindent
$a_n = 0$ を $n \geq 2$ に対して取れば，
$\lim H^0(U, \mathcal{L}_n)$ は零でないことに注意する．この場合，関数 $x$（$D(x)$ 上）と
関数 $x + t/y$（$D(y)$ 上）は貼り合わさる．一方，$a_n = 1$ かつ
$\varphi(n) = 2^n$，またはさらに $\varphi(n) = n^2$ と取ると，
$\lim H^0(U, \mathcal{L}_n)$ が零であることを読者は示せる．これは，この場合に
$(\mathcal{L}_n)$ が代数化可能でないことの別の証明を与える．
\end{example}

\noindent
状況 \ref{algebraization-situation-algebraize} において環 $A$ が $I$-進完備でない場合，
補題 \ref{algebraization-lemma-essential-image-completion} は，$(\mathcal{F}_n)$ が制限関手
$$
\textit{Coh}(X, I\mathcal{O}_X)
\longrightarrow
\textit{Coh}(U, I\mathcal{O}_U)
$$
の本質像に入るかを問うべきことを示唆する．しかし，これは $(\mathcal{F}_n)$ が
代数化可能であることを意味するとは，もはやいえない．そこで次の用語を導入する．

\begin{definition}
\label{algebraization-definition-algebraizable}
状況 \ref{algebraization-situation-algebraize} において，$(\mathcal{F}_n)$ を
$\textit{Coh}(U, I\mathcal{O}_U)$ の対象とする．{\it $(\mathcal{F}_n)$ は $X$ へ延長される}とは，
対象 $(\mathcal{G}_n)$ が $\textit{Coh}(X, I\mathcal{O}_X)$ に存在し，$U$ への制限が
$(\mathcal{F}_n)$ と同型であることをいう．
\end{definition}

\noindent
この概念は，完備化上で代数化可能であることと同値である．

\begin{lemma}
\label{algebraization-lemma-algebraizable}
状況 \ref{algebraization-situation-algebraize} において，$(\mathcal{F}_n)$ を
$\textit{Coh}(U, I\mathcal{O}_U)$ の対象とする．$A', I', \mathfrak a'$ を
$I$-進完備化，すなわちそれぞれ $A, I, \mathfrak a$ の完備化とする．$X' = \Spec(A')$ および
$U' = X' \setminus V(\mathfrak a')$ とおく．次は同値である．
\begin{enumerate}
\item $(\mathcal{F}_n)$ は $X$ へ延長される．
\item $(\mathcal{F}_n)$ の $U'$ への引き戻しは，連接 $\mathcal{O}_{U'}$-加群の完備化である．
\end{enumerate}
\end{lemma}

\begin{proof}
$A \to A'$ は同型 $A/I \to A'/I'$ を誘導する平坦な環準同型であることを思い出す．
Algebra，補題 \ref{algebra-lemma-completion-flat} および
\ref{algebra-lemma-completion-complete} を参照されたい．したがって，$X' \to X$ は
同型 $Y' \to Y$ を誘導する平坦射である．ゆえに $U' \to U$ は，同型
$U' \cap Y' \to U \cap Y$ を誘導する平坦射である．これは，可換図式
$$
\xymatrix{
\textit{Coh}(X', I\mathcal{O}_{X'}) \ar[r] &
\textit{Coh}(U', I\mathcal{O}_{U'}) \\
\textit{Coh}(X, I\mathcal{O}_X) \ar[u] \ar[r] &
\textit{Coh}(U, I\mathcal{O}_U) \ar[u]
}
$$
における垂直関手が同値であることを含意する．Cohomology of Schemes，補題
\ref{coherent-lemma-inverse-systems-pullback-equivalence} を参照されたい．
この事実と補題 \ref{algebraization-lemma-essential-image-completion} の結果から，本補題は形式的に従う．
\end{proof}

\noindent
状況 \ref{algebraization-situation-algebraize} において，$(\mathcal{F}_n)$ を
$\textit{Coh}(U, I\mathcal{O}_U)$ の対象とする．$(\mathcal{F}_n)$ が $X$ へ延長されるかを
判定するためには，$A$-加群
\begin{equation}
\label{algebraization-equation-guess}
M = \lim H^0(U, \mathcal{F}_n)
\end{equation}
を見るのが自然である．$M$ は極限位相をもち，これは（先験的には）$I$-進位相より粗いことに
注意する．実際，$M \to H^0(U, \mathcal{F}_n)$ は $I^nM$ を零化する．
次の標準写像がある．
$$
\widetilde{M}|_U \to \widetilde{M/I^nM}|_U \to
\widetilde{H^0(U, \mathcal{F}_n)}|_U \to
\mathcal{F}_n
$$
$\widetilde{M}$ の $U$ への制限が連接加群となり，$(\mathcal{F}_n)$ がこの加群の完備化に
なることを期待したくなる．しかし，$A$ が完備でなければ，これが成り立つ見込みはほとんど
ないので，これは素朴すぎる．さらに $A$ が $I$-進完備であっても，この概念は扱うのが
非常に難しい．より素朴でない方法として，次の要件を考える．

\begin{definition}
\label{algebraization-definition-canonically-algebraizable}
状況 \ref{algebraization-situation-algebraize} において，$(\mathcal{F}_n)$ を
$\textit{Coh}(U, I\mathcal{O}_U)$ の対象とする．{\it $(\mathcal{F}_n)$ は $X$ へ
標準的に延長される}とは，逆系
$$
\{\widetilde{H^0(U, \mathcal{F}_n)}\}_{n \geq 1}
$$
が $\QCoh(\mathcal{O}_X)$ において対象 $(\mathcal{G}_n)$ と pro-同型であり，後者が
$\textit{Coh}(X, I\mathcal{O}_X)$ に属することをいう．
\end{definition}

\noindent
補題 \ref{algebraization-lemma-canonically-algebraizable} において，定義
\ref{algebraization-definition-canonically-algebraizable} の条件が定義
\ref{algebraization-definition-algebraizable} の条件より強いことを見る．

\begin{lemma}
\label{algebraization-lemma-canonically-algebraizable}
状況 \ref{algebraization-situation-algebraize} において，$(\mathcal{F}_n)$ を
$\textit{Coh}(U, I\mathcal{O}_U)$ の対象とする．$(\mathcal{F}_n)$ が $X$ へ
標準的に延長されるならば，次が成り立つ．
\begin{enumerate}
\item $(\widetilde{H^0(U, \mathcal{F}_n)})$ は対象 $(\mathcal{G}_n)$ と pro-同型であり，後者は
$\textit{Coh}(X, I \mathcal{O}_X)$ に属し，一意な同型を除いて一意である．
\item $(\mathcal{G}_n)$ の $U$ への制限は $(\mathcal{F}_n)$ と同型である．すなわち，
$(\mathcal{F}_n)$ は $X$ へ延長される．
\item 逆系 $\{H^0(U, \mathcal{F}_n)\}$ は Mittag--Leffler 条件を満たす．
\item (\ref{algebraization-equation-guess}) の加群 $M$ は $I$-進完備化した $A$ 上有限であり，
$M$ 上の極限位相は $I$-進位相である．
\end{enumerate}
\end{lemma}

\begin{proof}
(1) における $(\mathcal{G}_n)$ の存在は定義
\ref{algebraization-definition-canonically-algebraizable} から従う．(1) における $(\mathcal{G}_n)$ の一意性は補題
\ref{algebraization-lemma-recognize-formal-coherent-modules} から従う．
$\mathcal{G}_n = \widetilde{M_n}$ と書く．このとき $\{M_n\}$ は有限 $A$-加群の逆系であり，
$M_n = M_{n + 1}/I^n M_{n + 1}$ を満たす．定義
\ref{algebraization-definition-canonically-algebraizable} により，逆系
$\{H^0(U, \mathcal{F}_n)\}$ は $\{M_n\}$ と pro-同型である．したがって逆系
$\{H^0(U, \mathcal{F}_n)\}$ は Mittag--Leffler 条件を満たし，
$M = \lim M_n$ である（位相加群として）．ゆえに (4) における $M$ の性質は，
Algebra，補題 \ref{algebra-lemma-limit-complete}，
\ref{algebra-lemma-finite-over-complete-ring}，および
\ref{algebra-lemma-hathat-finitely-generated} から従う．$U$ は準アフィンなので，標準写像
$$
\widetilde{H^0(U, \mathcal{F}_n)}|_U \to \mathcal{F}_n
$$
は同型である（Properties，補題
\ref{properties-lemma-quasi-coherent-quasi-affine}）．したがって
$(\mathcal{G}_n|_U)$ と $(\mathcal{F}_n)$ は pro-同型であり，ゆえに補題
\ref{algebraization-lemma-recognize-formal-coherent-modules} により同型である．
\end{proof}

\begin{lemma}
\label{algebraization-lemma-canonically-extend-base-change}
状況 \ref{algebraization-situation-algebraize} において，$(\mathcal{F}_n)$ を
$\textit{Coh}(U, I\mathcal{O}_U)$ の対象とする．$A \to A'$ を平坦な環準同型とする．
$X' = \Spec(A')$ とおき，$U' \subset X'$ を $U$ の逆像とし，誘導される射を
$g : U' \to U$ と表す．$(\mathcal{F}'_n) = (g^*\mathcal{F}_n)$ とおく．
Cohomology of Schemes，補題 \ref{coherent-lemma-inverse-systems-pullback} を参照されたい．
$(\mathcal{F}_n)$ が $X$ へ標準的に延長されるならば，$(\mathcal{F}'_n)$ は $X'$ へ
標準的に延長される．さらに，補題 \ref{algebraization-lemma-canonically-algebraizable} で得られた
$(\mathcal{F}_n)$ の延長を引き戻すと，$(\mathcal{F}'_n)$ の延長になる．
\end{lemma}

\begin{proof}
$f : X' \to X$ を誘導される射とする．平坦基変換により
$H^0(U', \mathcal{F}'_n) = H^0(U, \mathcal{F}_n) \otimes_A A'$ である．
Cohomology of Schemes，補題 \ref{coherent-lemma-flat-base-change-cohomology} を参照されたい．
したがって，$(\mathcal{G}_n)$ が $\textit{Coh}(X, I\mathcal{O}_X)$ に属し，
$(\widetilde{H^0(U, \mathcal{F}_n)})$ と pro-同型ならば，
$(f^*\mathcal{G}_n)$ は
$$
(f^*\widetilde{H^0(U, \mathcal{F}_n)}) =
(\widetilde{H^0(U, \mathcal{F}_n) \otimes_A A'}) =
(\widetilde{H^0(U', \mathcal{F}'_n)})
$$
と pro-同型である．これで証明が完了する．
\end{proof}

\begin{lemma}
\label{algebraization-lemma-when-done}
状況 \ref{algebraization-situation-algebraize} において，$(\mathcal{F}_n)$ を
$\textit{Coh}(U, I\mathcal{O}_U)$ の対象とする．$M$ を (\ref{algebraization-equation-guess}) のものとする．
次を仮定する．
\begin{enumerate}
\item[(a)] 逆系 $H^0(U, \mathcal{F}_n)$ は Mittag--Leffler 条件を満たす．
\item[(b)] $M$ 上の極限位相は $I$-進位相と一致する．
\item[(c)] $M \to H^0(U, \mathcal{F}_n)$ の像は有限 $A$-加群であり，これはすべての $n$ に対して成り立つ．
\end{enumerate}
このとき $(\mathcal{F}_n)$ は $X$ へ標準的に延長される．特に，$A$ が $I$-進完備ならば，
$(\mathcal{F}_n)$ は連接 $\mathcal{O}_U$-加群の完備化である．
\end{lemma}

\begin{proof}
$H^0(U, \mathcal{F}_n)$ が Mittag--Leffler 条件を満たし，$M$ 上の極限位相が
$I$-進位相であるので，$\{M/I^nM\}$ と $\{H^0(U, \mathcal{F}_n)\}$ は
$A$-加群の pro-同型な逆系である．そこで
$$
\mathcal{G}_n = \widetilde{M/I^n M}
$$
とおく．定義 \ref{algebraization-definition-canonically-algebraizable} の条件を確認するには，$M$ が
$I$-進完備化した $A$ 上の有限加群であることを示せば十分である．これは，条件 (c) と
上の議論により $M/I^n M$ が有限であること，および Algebra，補題
\ref{algebra-lemma-finite-over-complete-ring} から従う．
\end{proof}

\noindent
次は，ある意味で上の補題 \ref{algebraization-lemma-when-done} の最も直接的な応用である．

\begin{lemma}
\label{algebraization-lemma-algebraization-principal-variant}
状況 \ref{algebraization-situation-algebraize} において，$(\mathcal{F}_n)$ を
$\textit{Coh}(U, I\mathcal{O}_U)$ の対象とする．次を仮定する．
\begin{enumerate}
\item $I = (f)$ は非零因子 $f \in \mathfrak a$ により生成される単項イデアルである．
\item $\mathcal{F}_n$ は有限局所自由 $\mathcal{O}_U/f^n\mathcal{O}_U$-加群である．
\item $H^1_\mathfrak a(A/fA)$ と $H^2_\mathfrak a(A/fA)$ は有限 $A$-加群である．
\end{enumerate}
このとき $(\mathcal{F}_n)$ は $X$ へ標準的に延長される．特に，$A$ が完備ならば，
$(\mathcal{F}_n)$ は連接 $\mathcal{O}_U$-加群の完備化である．
\end{lemma}

\begin{proof}
補題 \ref{algebraization-lemma-when-done} の仮定 (a)，(b)，(c) を確認することで証明する．

\medskip\noindent
$\mathcal{F}_n$ は $\mathcal{O}_U/f^n\mathcal{O}_U$ 上局所自由なので，短完全列
$0 \to \mathcal{F}_n \to \mathcal{F}_{n + 1} \to \mathcal{F}_1 \to 0$
がすべての $n$ に対して得られる．したがって，条件 (b) は Cohomology，補題
\ref{cohomology-lemma-topology-I-adic-f} により成り立つ．

\medskip\noindent
$f$ は非零因子なので，短完全列
$$
0 \to A/f^nA \xrightarrow{f} A/f^{n + 1}A \to A/fA \to 0
$$
を得て，これに対応する短完全列
$0 \to \mathcal{F}_n \to \mathcal{F}_{n + 1} \to \mathcal{F}_1 \to 0$ をもつ．
以下では Local Cohomology，補題
\ref{local-cohomology-lemma-finiteness-pushforwards-and-H1-local} を断りなく用いる．
仮定から $H^0(U, \mathcal{O}_U/f\mathcal{O}_U)$ と
$H^1(U, \mathcal{O}_U/f\mathcal{O}_U)$ は有限 $A$-加群である．したがって
$\mathcal{F}_1$ についても同じことが成り立つ．Local Cohomology，補題
\ref{local-cohomology-lemma-finiteness-for-finite-locally-free} を参照されたい．
帰納法と短完全列を用いると，$H^0(U, \mathcal{F}_n)$ は有限 $A$-加群であり，
これはすべての $n$ に対して成り立つ．こうして仮定 (c) が満たされる．

\medskip\noindent
最後に，$H^1(U, \mathcal{F}_1)$ は有限 $A$-加群なので，Cohomology，補題
\ref{cohomology-lemma-ML} を適用して仮定 (a) が成り立つことが分かる．
\end{proof}

\begin{remark}
\label{algebraization-remark-interesting-case-variant}
補題 \ref{algebraization-lemma-algebraization-principal-variant} において，$A$ が普遍鎖状で
Cohen--Macaulay な形式的ファイバーをもつならば（例えば $A$ が双対化複体をもつ場合），
条件
$H^1_\mathfrak a(A/fA)$ と $H^2_\mathfrak a(A/fA)$
が有限 $A$-加群であることは，
$$
\text{depth}((A/f)_\mathfrak p) + \dim((A/\mathfrak p)_\mathfrak q) > 2
$$
がすべての $\mathfrak p \in V(f) \setminus V(\mathfrak a)$ および
$\mathfrak q \in V(\mathfrak p) \cap V(\mathfrak a)$ に対して成り立つことと同値である．
これは Local Cohomology，定理 \ref{local-cohomology-theorem-finiteness} による．

\medskip\noindent
例えば，$A/fA$ が $(S_2)$ であり，$Z = V(\mathfrak a)$ の各既約成分が
余次元 $\geq 3$ をもち，ここで周囲の空間が $Y = \Spec(A/fA)$ ならば，
$H^1_\mathfrak a(A/fA)$ と $H^2_\mathfrak a(A/fA)$ の有限性が得られる．
これは補題 \ref{algebraization-lemma-algebraization-principal} で得られた，わずかに弱い条件と
対比されるべきである（注意 \ref{algebraization-remark-interesting-case} も参照）．
\end{remark}






\section{連接形式加群の代数化 II}
\label{algebraization-section-algebraization-modules-general}

\noindent
第 \ref{algebraization-section-algebraization-modules} 節で始めた議論を続ける．
初読時には本節を飛ばしてよい．

\begin{lemma}
\label{algebraization-lemma-map-kernel-cokernel-on-closed}
状況 \ref{algebraization-situation-algebraize} において考える．
$(\mathcal{F}_n) \to (\mathcal{F}'_n)$ を $\textit{Coh}(U, I\mathcal{O}_U)$ の射で，
その核と余核が $I$ のある冪によって零化されるものとする．このとき
\begin{enumerate}
\item $(\mathcal{F}_n)$ が $X$ へ延長されることと，$(\mathcal{F}'_n)$ が $X$ へ延長されることは同値である．
\item $(\mathcal{F}_n)$ が連接 $\mathcal{O}_U$-加群の完備化であることと，
$(\mathcal{F}'_n)$ がそうであることは同値である．
\end{enumerate}
\end{lemma}

\begin{proof}
(2) は Cohomology of Schemes，補題 \ref{coherent-lemma-existence-easy} から直ちに従う．
(1) を見るには，まず補題 \ref{algebraization-lemma-algebraizable} を用いて，$A$ が $I$-進完備である
場合に帰着する．しかしこの場合，補題 \ref{algebraization-lemma-essential-image-completion} により
(1) は (2) に帰着する．
\end{proof}

\noindent
次の二つの補題は，もともと補題 \ref{algebraization-lemma-when-done} の証明で用いられていた．
どのような中間結果が得られるかに関心のある読者のため，ここに残しておく．

\begin{lemma}
\label{algebraization-lemma-when-ML}
状況 \ref{algebraization-situation-algebraize} において，$(\mathcal{F}_n)$ を
$\textit{Coh}(U, I\mathcal{O}_U)$ の対象とする．逆系
$H^0(U, \mathcal{F}_n)$ が Mittag--Leffler 条件を満たすならば，標準写像
$$
\widetilde{M/I^nM}|_U \to \mathcal{F}_n
$$
はすべての $n$ に対して全射である．ここで $M$ は (\ref{algebraization-equation-guess}) のものとする．
\end{lemma}

\begin{proof}
全射性は，ある点 $y \in Y \setminus Z$ における茎上で確認できる．
$y$ が素イデアル $\mathfrak q \subset A$ に対応するならば，
$f \in \mathfrak a$ かつ $f \not \in \mathfrak q$ となるものを選べる．このとき
$$
M_f \longrightarrow H^0(U, \mathcal{F}_n)_f = H^0(D(f), \mathcal{F}_n)
$$
が全射であることを示せば十分である．実際，$D(f)$ はアフィンである（等号は Properties，
補題 \ref{properties-lemma-invert-f-sections} による）．Mittag--Leffler 条件により，
$$
\Im(M \to H^0(U, \mathcal{F}_n)) =
\Im(H^0(U, \mathcal{F}_m) \to H^0(U, \mathcal{F}_n))
$$
がある $m \geq n$ に対して成り立つ．コホモロジーの長完全列を用いると，
$$
\Coker(H^0(U, \mathcal{F}_m) \to H^0(U, \mathcal{F}_n))
\subset
H^1(U, \Ker(\mathcal{F}_m \to \mathcal{F}_n))
$$
$U = X \setminus V(\mathfrak a)$ なので，この $H^1$ は $\mathfrak a$-冪捩れである．
したがって $f$ を可逆にすると余核は零になる．
\end{proof}

\begin{lemma}
\label{algebraization-lemma-when-topology}
状況 \ref{algebraization-situation-algebraize} において，$(\mathcal{F}_n)$ を
$\textit{Coh}(U, I\mathcal{O}_U)$ の対象とする．$M$ を (\ref{algebraization-equation-guess}) のものとする．
$$
\mathcal{G}_n = \widetilde{M/I^nM}.
$$
とおく．$M$ 上の極限位相が $I$-進位相と一致するならば，
$\mathcal{G}_n|_U$ は連接 $\mathcal{O}_U$-加群であり，逆系の写像
$$
(\mathcal{G}_n|_U) \longrightarrow (\mathcal{F}_n)
$$
はアーベル圏 $\textit{Coh}(U, I\mathcal{O}_U)$ において単射である．
\end{lemma}

\begin{proof}
まず，$\mathcal{G}_n$ は準連接 $\mathcal{O}_X$-加群で，$I^n$ により零化され，
$\mathcal{G}_{n + 1}/I^n\mathcal{G}_{n + 1} = \mathcal{G}_n$.
であることに注意する．次を考える．
$$
M_n = \Im(M \longrightarrow H^0(U, \mathcal{F}_n))
$$
仮定は，逆系 $(M_n)$ と $(M/I^nM)$ が $\text{Mod}_A$ の pro-対象として同型であることをいう．
$f \in \mathfrak a$ を，$D(f) \subset U$ がアフィン開集合となるように選ぶ．このとき
$$
(M_n)_f \subset H^0(U, \mathcal{F}_n)_f = H^0(D(f), \mathcal{F}_n)
$$
等号は Properties，補題 \ref{properties-lemma-invert-f-sections} により成り立つ．
したがって $\widetilde{M_n}|_U \to \mathcal{F}_n$ は単射である．ゆえに
$\widetilde{M_n}|_U$ は連接加群である（Cohomology of Schemes，補題
\ref{coherent-lemma-coherent-Noetherian-quasi-coherent-sub-quotient}）．
$M \to M/I^nM$ は全射であり，$M_{n'} \to M/I^nM$ を経由する．ここで
ある $n' \geq n$ を取る．したがって，$\mathcal{G}_n|_U$ は連接加群の商として連接である．
これを証明冒頭の注意と合わせると，$(\mathcal{G}_n|_U)$ は確かに
$\textit{Coh}(U, I\mathcal{O}_U)$ の対象をなす．最後に，写像の単射性を示すには，
$$
\lim (M/I^nM)_f = \lim H^0(D(f), \mathcal{G}_n) \to
\lim H^0(D(f), \mathcal{F}_n)
$$
が単射であることを示せば十分である．Cohomology of Schemes，補題
\ref{coherent-lemma-inverse-systems-abelian} および
\ref{coherent-lemma-inverse-systems-affine} を参照されたい．
$\lim (M_n)_f \to \lim H^0(D(f), \mathcal{F}_n)$ の単射性は明らかであり（上を参照），
pro-系についての注意により $\lim (M_n)_f = \lim (M/I^nM)_f$ である．
これで証明が完了する．
\end{proof}





\section{距離関数}
\label{algebraization-section-distance}

\noindent
$Y$ を Noether スキームとし，$Z \subset Y$ を閉部分集合とする．
次の関数を定義する．
\begin{equation}
\label{algebraization-equation-delta-Z}
\delta^Y_Z = \delta_Z : Y \longrightarrow \mathbf{Z}_{\geq 0} \cup \{\infty\}
\end{equation}
これは $Y$ の点の $Z$ からの「距離」を測る．
非形式的な議論については，注意 \ref{algebraization-remark-discussion} を参照されたい．
$y \in Y$ とする．$\delta_Z(y) = \infty$ とおくのは，$y$ が $Y$ の，
$Z$ と交わらない連結成分に含まれる場合である．
$y$ が $Y$ の，$Z$ と交わる連結成分に含まれるならば，
$k \geq 0$ と，次の系を取ることができる．
$$
V_0 \subset W_0 \supset V_1 \subset W_1 \supset \ldots \supset
V_k \subset W_k
$$
これは $Y$ の整な閉部分スキームからなる系で，$V_0 \subset Z$ であり，
$y \in W_k$ は生成点である．$c_i = \text{codim}(V_i, W_i)$ を
$i = 0, \ldots, k$ に対して定め，$b_i = \text{codim}(V_{i + 1}, W_i)$ を
$i = 0, \ldots, k - 1$ に対して定める．このような系に対して
$$
\delta(V_0, W_0, V_1, \ldots, W_k) =
k +
\max_{i = 0, 1, \ldots, k}
(c_i + c_{i + 1} + \ldots + c_k - b_i - b_{i + 1} - \ldots - b_{k - 1})
$$
これは $\geq k$ である．実際，$i = k$ と取ることができ，$c_k \geq 0$ である．
最後に，
$$
\delta_Z(y) = \min \delta(V_0, W_0, V_1, \ldots, W_k)
$$
とおく．ここで最小値は，上のような $Y$ の整な閉部分スキームのすべての系を走る．

\begin{lemma}
\label{algebraization-lemma-discussion}
$Y$ を Noether スキームとし，$Z \subset Y$ を閉部分集合とする．
\begin{enumerate}
\item $y \in Y$ に対して，$\delta_Z(y) = 0 \Leftrightarrow y \in Z$ である．
\item 部分集合 $\{y \in Y \mid \delta_Z(y) \leq k\}$ は特殊化で安定である．
\item $y \in Y$ および $z \in \overline{\{y\}} \cap Z$ に対して，
$\dim(\mathcal{O}_{\overline{\{y\}}, z}) \geq \delta_Z(y)$ である．
\item $\delta$ が $Y$ 上の次元関数ならば，
$\delta_Z(y) \geq \delta(y) - \delta_{max}$ である．ここで $\delta_{max}$ は
$\delta$ が $Z$ 上でとる最大値である．
\item $Y = \Spec(A)$ が鎖状 Noether 局所環のスペクトルで，その極大イデアルが
$\mathfrak m$ かつ $Z = \{\mathfrak m\}$ ならば，
$\delta_Z(y) = \dim(\overline{\{y\}})$ である．
\item $Y' \subset Y$ が開部分スキームならば，
$\delta^{Y'}_{Y' \cap Z}(y') \geq \delta^Y_Z(y')$ が $y' \in Y'$ に対して成り立つ．
\end{enumerate}
$Y$ が鎖状であると仮定する．このとき，
\begin{enumerate}
\setcounter{enumi}{6}
\item $y' \leadsto y$ を $Y$ の点の直後の特殊化とする．
$Y$ が鎖状ならば，$\delta_Z(y') \leq \delta_Z(y) + 1$ である．
\item 次の特殊化のパターンが与えられたとする．
$$
\xymatrix{
& y'_0 \ar@{~>}[ld] \ar@{~>}[rd] &
& y'_1 \ar@{~>}[ld] & \ldots
& y'_{k - 1} \ar@{~>}[rd] &
\\
y_0 & &
y_1 & &
\ldots & &
y_k = y
}
$$
これは $Y$ の点の間のもので，$y_0 \in Z$ かつ $y_i' \leadsto y_i$ は
直後の特殊化である．このとき $\delta_Z(y_k) \leq k$ である．
\end{enumerate}
\end{lemma}

\begin{proof}
(1) の証明．$y \in Z$ ならば，$k = 0$ および
$V_0 = W_0 = \overline{\{y\}}$ と取ることができ，$\delta(V_0, W_0) = 0$ を得るので，
$\delta_Z(y) = 0$ である．$y \not \in Z$ ならば，任意の系
$V_0 \subset W_0 \supset V_1 \subset W_1 \supset \ldots \subset W_k$
で $y$ に対するものについて，$k = 0$ かつ $V_0 \not = W_0$ であるか，または $k > 0$ である．
いずれの場合も $\delta(V_0, W_0, \ldots, W_k) > 0$ である．したがって $\delta_Z(y) > 0$ である．

\medskip\noindent
(2) の証明．$y \leadsto y'$ を非自明な特殊化とし，
$V_0 \subset W_0 \supset V_1 \subset W_1 \supset \ldots \subset W_k$
$y$ に対する系とする．ここで二つの場合がある．
場合 I：$V_k = W_k$，すなわち $c_k = 0$ とする．この場合，
$V'_k = W'_k = \overline{\{y'\}}$ とおくことができる．
簡単な計算により，
$\delta(V_0, W_0, \ldots, V'_k, W'_k) \leq
\delta(V_0, W_0, \ldots, V_k, W_k)$
である．これは $b_{k - 1}$ だけがより大きな整数に変わるためである．
場合 II：$V_k \not = W_k$，すなわち $c_k > 0$ とする．
この場合，$V_{k + 1} = W_{k + 1} = \overline{\{y'\}}$ とおくと，
$V_0 \subset W_0 \supset \ldots \subset W_k \supset V_{k + 1}
\subset W_{k + 1}$ は $y'$ に対する系となる．このとき $c_{k + 1} = 0$ かつ
$b_k > 0$ なので，次を得る．
\begin{align*}
& \delta(V_0, \ldots, W_{k + 1}) \\
& =
k + 1 +
\max_{i = 0, 1, \ldots, k + 1}
(c_i + c_{i + 1} + \ldots + c_k + c_{k + 1}
- b_i - b_{i + 1} - \ldots - b_{k - 1} - b_k) \\
& =
k + 1 +
\max_{i = 0, 1, \ldots, k + 1}
(c_i + c_{i + 1} + \ldots + c_k
- b_i - b_{i + 1} - \ldots - b_{k - 1} - b_k) \\
& \leq
k + \max_{i = 0, 1, \ldots, k}
(c_i + c_{i + 1} + \ldots + c_k - b_i - b_{i + 1} - \ldots - b_{k - 1}) \\
& = \delta(V_0, \ldots, W_k)
\end{align*}
この不等式が成り立つのは，$c_k > 0$ かつ $b_k > 0$ であり，とりわけ
$\delta(V_0, \ldots, W_k) \geq k + c_k \geq k + 1$ が従うためである．

\medskip\noindent
(3) の証明．$y \in Y$ および $z \in \overline{\{y\}} \cap Z$ が与えられると，
次の系を得る．
$$
V_0 = \overline{\{z\}} \subset W_0 = \overline{\{y\}}
$$
また，$c_0 = \text{codim}(V_0, W_0) = \dim(\mathcal{O}_{\overline{\{y\}}, z})$
である（Properties，補題 \ref{properties-lemma-codimension-local-ring}）．
したがって $\delta(V_0, W_0) = 0 + c_0 = c_0$ となり，主張が従う．

\medskip\noindent
(4) の証明．$\delta$ を $Y$ 上の次元関数とする．
$V_0 \subset W_0 \supset V_1 \subset W_1 \supset \ldots \subset W_k$
を $y$ に対する系とする．$y'_i \in W_i$ および $y_i \in V_i$ をそれぞれ
生成点とする．したがって $y_0 \in Z$ かつ $y_k = y$ である．このとき，
$$
\delta(y_i) - \delta(y_{i - 1}) =
\delta(y'_{i - 1}) - \delta(y_{i - 1}) - \delta(y'_{i - 1}) + \delta(y_i) =
c_{i - 1} - b_{i - 1}
$$
最後に，$\delta(y'_k) - \delta(y_{k - 1}) = c_k$ である．
したがって，
$$
\delta(y) - \delta(y_0) =
c_0 + \ldots + c_k - b_0 - \ldots - b_{k - 1}
$$
ゆえに
$\delta(V_0, W_0, \ldots, W_k) \geq k + \delta(y) - \delta(y_0)$
であり，主張が従う．

\medskip\noindent
(5) の証明．関数 $\delta(y) = \dim(\overline{\{y\}})$ は次元関数である．
したがって (4) により $\delta(y) \leq \delta_Z(y)$ である．
(3) により $\delta_Z(y) \leq \delta(y)$ でもあるから，証明は完了する．

\medskip\noindent
(6) の証明．$\delta^{Y'}_{Y' \cap Z}$ の計算で考えるべき系が少ないので，これは明らかである．

\medskip\noindent
(7) の証明．
$V_0 \subset W_0 \supset V_1 \subset W_1 \supset \ldots \subset W_k$
を $y$ に対する系とする．$W'_k = \overline{\{y'\}}$ とおく．
$Y$ は鎖状なので，
$\text{codim}(V_k, W'_k) = \text{codim}(V_k, W_k) + 1$.
である．これから容易に $\delta(V_0, \ldots, V_k, W'_k) \leq
\delta(V_0, \ldots, V_k, W_k) + 1$ が従い，主張が証明される．

\medskip\noindent
(8) の証明．(7) と (2) を組み合わせればよい．
\end{proof}

\begin{lemma}
\label{algebraization-lemma-change-distance-function}
$Y$ を普遍鎖状 Noether スキームとする．$Z \subset Y$ を閉部分スキームとする．
$f : Y' \to Y$ を，すべてのファイバーの次元が $\leq e$ である有限型射とする．
$Z' = f^{-1}(Z)$ とおく．このとき，
$$
\delta_Z(y) \leq \delta_{Z'}(y') + e - \text{trdeg}_{\kappa(y)}(\kappa(y'))
$$
が成り立つ．ここで $y' \in Y'$ の像は $y \in Y$ である．
\end{lemma}

\begin{proof}
$\delta_{Z'}(y') = \infty$ ならば，証明すべきことはない．
$\delta_{Z'}(y') < \infty$ ならば，整な閉部分スキームの系
$$
V'_0 \subset W'_0 \supset V'_1 \subset W'_1 \supset \ldots \subset W'_k
$$
を $Y'$ の中で選び，$V'_0 \subset Z'$，$y'$ は $W'_k$ の生成点，かつ
$\delta_{Z'}(y') = \delta(V'_0, W'_0, \ldots, W'_k)$ となるようにする．
次を
$$
V_0 \subset W_0 \supset V_1 \subset W_1 \supset \ldots \subset W_k
$$
上のスキームの $Y$ におけるスキーム論的像とする．$y$ は $W_k$ の生成点であり，
$V_0 \subset Z$ であることに注意する．各 $i$ について次の図式を考える．
$$
\xymatrix{
V'_i \ar[r] \ar[d] & W'_i \ar[d] & V'_{i + 1} \ar[l] \ar[d] \\
V_i \ar[r] & W_i & V_{i + 1} \ar[l]
}
$$
相対次元 $n_i$ を $V'_i/V_i$ のものとし，相対次元
$m_i$ を $W'_i/W_i$ のものとする．より正確には，これらは対応する
関数体の拡大の超越次数である．
$c_i = \text{codim}(V_i, W_i)$,
$c'_i = \text{codim}(V'_i, W'_i)$,
$b_i = \text{codim}(V_{i + 1}, W_i)$，および
$b'_i = \text{codim}(V'_{i + 1}, W'_i)$ とおく．
次元公式により，
$$
c_i = c'_i + n_i - m_i
\quad\text{かつ}\quad
b_i = b'_i + n_{i + 1} - m_i
$$
Morphisms，補題 \ref{morphisms-lemma-dimension-formula} を参照されたい．
したがって $c_i - b_i = c'_i - b'_i + n_i - n_{i + 1}$ である．ゆえに，
\begin{align*}
& c_i + c_{i + 1} + \ldots + c_k - b_i - b_{i + 1} - \ldots - b_{k - 1} \\
& =
c'_i + c'_{i + 1} + \ldots + c'_k - b'_i - b'_{i + 1} - \ldots - b'_{k - 1}
+ n_i - n_k + c_k - c'_k \\
& =
c'_i + c'_{i + 1} + \ldots + c'_k - b'_i - b'_{i + 1} - \ldots - b'_{k - 1}
+ n_i - m_k
\end{align*}
したがって，
\begin{align*}
\max_{i = 0, \ldots, k}
& (c_i + c_{i + 1} + \ldots + c_k - b_i - b_{i + 1} - \ldots - b_{k - 1}) \\
& =
\max_{i = 0, \ldots, k}
(c'_i + c'_{i + 1} + \ldots + c'_k - b'_i - b'_{i + 1} - \ldots - b'_{k - 1}
+ n_i - m_k) \\
& =
\max_{i = 0, \ldots, k}
(c'_i + c'_{i + 1} + \ldots + c'_k - b'_i - b'_{i + 1} - \ldots - b'_{k - 1}
+ n_i) - m_k \\
& \leq
\max_{i = 0, \ldots, k}
(c'_i + c'_{i + 1} + \ldots + c'_k - b'_i - b'_{i + 1} - \ldots - b'_{k - 1})
+ e - m_k
\end{align*}
$m_k = \text{trdeg}_{\kappa(y)}(\kappa(y'))$ であるから，
$$
\delta(V_0, W_0, \ldots, W_k) \leq
\delta(V'_0, W'_0, \ldots, W'_k) + e - \text{trdeg}_{\kappa(y)}(\kappa(y'))
$$
を得る．これは求める不等式である．
\end{proof}

\begin{remark}
\label{algebraization-remark-discussion}
$Y$ を鎖状 Noether スキームとし，$Z \subset Y$ を閉部分集合とする．
補題 \ref{algebraization-lemma-discussion} により，
$$
\delta_Z(y) \leq \min
\left\{ k \middle|
\begin{matrix}
\text{次の特殊化を }Y\text{ において取れる} \\
y_0 \leftarrow y'_0 \rightarrow y_1 \leftarrow y'_1 \rightarrow \ldots
\leftarrow y'_{k - 1} \rightarrow y_k = y \\
\text{ただし }y_0 \in Z\text{ かつ }y_i' \leadsto y_i
\text{ は直後の特殊化}
\end{matrix}
\right\}
$$
である．$Y$ が体上有限型ならば，等号が成り立つと主張する．
この結果が必要になれば，正確な主張を定式化してここで証明することにする．
しかし一般には，この不等式の右辺によって $\delta_Z$ を定義したとき，
補題 \ref{algebraization-lemma-change-distance-function} がなお成り立つかどうかは分からない．
\end{remark}

\begin{example}
\label{algebraization-example-distance}
$k$ を体とし，$Y = \mathbf{A}^n_k$ とする．
$\delta : Y \to \mathbf{Z}_{\geq 0}$ を通常の次元関数とする．
\begin{enumerate}
\item $Z = \{z\}$ で，$z$ がある閉点ならば，
\begin{enumerate}
\item $\delta_Z(y) = \delta(y)$ であるのは $y \leadsto z$ の場合であり，
\item $\delta_Z(y) = \delta(y) + 1$ であるのは $y \not \leadsto z$ の場合である．
\end{enumerate}
\item $Z$ が閉部分多様体で，$W = \overline{\{y\}}$ ならば，
\begin{enumerate}
\item $\delta_Z(y) = 0$ であるのは $W \subset Z$ の場合，
\item $\delta_Z(y) = \dim(W) - \dim(Z)$ であるのは $Z$ が $W$ に含まれる場合，
\item $\delta_Z(y) = 1$ であるのは $\dim(W) \leq \dim(Z)$ かつ $W \not \subset Z$ の場合，
\item $\delta_Z(y) = \dim(W) - \dim(Z) + 1$ であるのは
$\dim(W) > \dim(Z)$ かつ $Z \not \subset W$ の場合である．
\end{enumerate}
\end{enumerate}
場合 (1) の一般化として，$Y$ が体上有限型で，$Z = \{z\}$ が閉点である場合がある．
このとき $\delta_Z(y) = \delta(y) + t$ であり，$t$ は $z$ と
$\overline{\{y\}}$ の閉点とを結ぶ曲線の鎖の長さの最小値である．
\end{example}







\section{連接形式加群の代数化 III}
\label{algebraization-section-uniqueness}

\noindent
第 \ref{algebraization-section-algebraization-modules} 節および
第 \ref{algebraization-section-algebraization-modules-general} 節で始めた議論を続ける．
第 \ref{algebraization-section-distance} 節の距離関数を用いて，
状況 \ref{algebraization-situation-algebraize} における連接形式加群に対する
いくつかの自然な条件を定式化する．

\medskip\noindent
状況 \ref{algebraization-situation-algebraize} において点 $y \in U \cap Y$ が与えられると，
$I$-進完備化
$$
\mathcal{O}_{X, y}^\wedge = \lim \mathcal{O}_{X, y}/I^n\mathcal{O}_{X, y}
$$
を考えることができる．これは $I\mathcal{O}_{X, y}^\wedge$ に関して完備な
Noether 局所環であり，その極大イデアルは $\mathfrak m_y^\wedge$ である．
Algebra，第 \ref{algebra-section-completion-noetherian} 節を参照されたい．
$(\mathcal{F}_n)$ を $\textit{Coh}(U, I\mathcal{O}_U)$ の対象とする．
$(\mathcal{F}_n)$ の $y$ における「茎」を次の式で定義する．
$$
\mathcal{F}_y^\wedge = \lim \mathcal{F}_{n, y}
$$
これは $\mathcal{O}_{X, y}^\wedge$ 上の有限加群である．Algebra，補題
\ref{algebra-lemma-limit-complete} および
\ref{algebra-lemma-finite-over-complete-ring} を参照されたい．

\begin{definition}
\label{algebraization-definition-s-d-inequalities}
状況 \ref{algebraization-situation-algebraize} において，$(\mathcal{F}_n)$ を
$\textit{Coh}(U, I\mathcal{O}_U)$ の対象とし，$a, b$ を整数とする．
$\delta^Y_Z$ は (\ref{algebraization-equation-delta-Z}) のものとする．
$(\mathcal{F}_n)$ が {\it $(a, b)$-不等式を満たす}とは，
$y \in U \cap Y$ および素イデアル $\mathfrak p \subset \mathcal{O}_{X, y}^\wedge$ で
$\mathfrak p \not \in V(I\mathcal{O}_{X, y}^\wedge)$ を満たすものに対して，
\begin{enumerate}
\item $V(\mathfrak p) \cap V(I\mathcal{O}_{X, y}^\wedge) \not =
\{\mathfrak m_y^\wedge\}$ ならば，
$$
\text{depth}((\mathcal{F}^\wedge_y)_\mathfrak p) + \delta^Y_Z(y) \geq a
\quad\text{または}\quad
\text{depth}((\mathcal{F}^\wedge_y)_\mathfrak p) +
\dim(\mathcal{O}_{X, y}^\wedge/\mathfrak p) + \delta^Y_Z(y) > b
$$
\item $V(\mathfrak p) \cap V(I\mathcal{O}_{X, y}^\wedge) =
\{\mathfrak m_y^\wedge\}$ ならば，
$$
\text{depth}((\mathcal{F}^\wedge_y)_\mathfrak p) + \delta^Y_Z(y) > a
$$
\end{enumerate}
$(\mathcal{F}_n)$ が {\it 狭義の $(a, b)$-不等式を満たす}とは，
$y \in U \cap Y$ および素イデアル
$\mathfrak p \subset \mathcal{O}_{X, y}^\wedge$ で
$\mathfrak p \not \in V(I\mathcal{O}_{X, y}^\wedge)$ を満たすものに対して，
次が成り立つことをいう．
$$
\text{depth}((\mathcal{F}^\wedge_y)_\mathfrak p) + \delta^Y_Z(y) > a
\quad\text{または}\quad
\text{depth}((\mathcal{F}^\wedge_y)_\mathfrak p) +
\dim(\mathcal{O}_{X, y}^\wedge/\mathfrak p) + \delta^Y_Z(y) > b
$$
\end{definition}

\noindent
いくつかの初等的な観察を述べる．

\begin{lemma}
\label{algebraization-lemma-elementary}
状況 \ref{algebraization-situation-algebraize} において，$(\mathcal{F}_n)$ を
$\textit{Coh}(U, I\mathcal{O}_U)$ の対象とし，$a, b$ を整数とする．
\begin{enumerate}
\item $(\mathcal{F}_n)$ が $I$ のある冪で零化されるならば，
$(\mathcal{F}_n)$ は $(a, b)$-不等式を任意の $a, b$ に対して満たす．
\item $(\mathcal{F}_n)$ が $(a + 1, b)$-不等式を満たすならば，
$(\mathcal{F}_n)$ は狭義の $(a, b)$-不等式を満たす．
\end{enumerate}
$\text{cd}(A, I) \leq d$ で，$A$ が双対化複体をもつならば，
\begin{enumerate}
\item[(3)] $(\mathcal{F}_n)$ が $(s, s + d)$-不等式を満たすための必要十分条件は，
すべての $y \in U \cap Y$ に対して組
$\mathcal{O}_{X, y}^\wedge, I\mathcal{O}_{X, y}^\wedge,
\{\mathfrak m_y^\wedge\}, \mathcal{F}_y^\wedge, s - \delta^Y_Z(y), d$
が状況 \ref{algebraization-situation-bootstrap} のとおりであることである．
\item[(4)]
$(\mathcal{F}_n)$ が狭義の $(s, s + d)$-不等式を満たすならば，
$(\mathcal{F}_n)$ は $(s, s + d)$-不等式を満たす．
\end{enumerate}
\end{lemma}

\begin{proof}
(4) を除いて直ちに従う．(4) は補題 \ref{algebraization-lemma-bootstrap-bis-bis} と
(3) の対応から従う．
\end{proof}

\begin{lemma}
\label{algebraization-lemma-explain-2-3-cd-1}
状況 \ref{algebraization-situation-algebraize} において，$(\mathcal{F}_n)$ を
$\textit{Coh}(U, I\mathcal{O}_U)$ の対象とする．$\text{cd}(A, I) = 1$ ならば，
$\mathcal{F}$ が $(2, 3)$-不等式を満たすための必要十分条件は，
$$
\text{depth}((\mathcal{F}^\wedge_y)_\mathfrak p) +
\dim(\mathcal{O}_{X, y}^\wedge/\mathfrak p) + \delta^Y_Z(y) > 3
$$
すべての $y \in U \cap Y$ および $\mathfrak p \subset \mathcal{O}_{X, y}^\wedge$ で
$\mathfrak p \not \in V(I\mathcal{O}_{X, y}^\wedge)$ を満たすものに対して成り立つことである．
\end{lemma}

\begin{proof}
素イデアル $\mathfrak p \subset \mathcal{O}_{X, y}^\wedge$ で
$\mathfrak p \not \in V(I\mathcal{O}_{X, y}^\wedge)$ を満たすものに対して，
$V(\mathfrak p) \cap V(I\mathcal{O}_{X, y}^\wedge) =
\{\mathfrak m_y^\wedge\}
\Leftrightarrow \dim(\mathcal{O}_{X, y}^\wedge/\mathfrak p) = 1$
である．これは $\text{cd}(A, I) = 1$ による．Local Cohomology，補題
\ref{local-cohomology-lemma-cd-change-rings} および
\ref{local-cohomology-lemma-cd-bound-dim-local} を参照されたい．
さて，三つの数
$\alpha = \text{depth}((\mathcal{F}^\wedge_y)_\mathfrak p) \geq 0$,
$\beta = \dim(\mathcal{O}_{X, y}^\wedge/\mathfrak p) \geq 1$，および
$\gamma = \delta^Y_Z(y) \geq 1$ を考える．
このとき定義 \ref{algebraization-definition-s-d-inequalities} が要求するのは，
\begin{enumerate}
\item $\beta > 1$ ならば，
$\alpha + \gamma \geq 2$ または $\alpha + \beta + \gamma > 3$ であり，
\item $\beta = 1$ ならば，$\alpha + \gamma > 2$ であること
\end{enumerate}
である．これが $\alpha + \beta + \gamma > 3$ と同値であることは直ちに分かる．
\end{proof}

\noindent
この節の残りは，初読時には飛ばすことを勧める．ここでは，$A$ が $I$-進完備なとき，
$(\mathcal{F}_n)$ の圏で，$\textit{Coh}(U, I\mathcal{O}_U)$ の対象として $X$ へ延長され，
狭義の $(1, 1 + \text{cd}(A, I))$-不等式を満たすものの圏が，
連接 $\mathcal{O}_U$-加群の圏のある充満部分圏と同値であることを示す．

\begin{lemma}
\label{algebraization-lemma-sanity}
状況 \ref{algebraization-situation-algebraize} において，$\mathcal{F}$ を連接
$\mathcal{O}_U$-加群とし，$d \geq 1$ とする．次を仮定する．
\begin{enumerate}
\item $A$ は $I$-進完備で，双対化複体をもち，
$\text{cd}(A, I) \leq d$ である．
\item 完備化 $\mathcal{F}^\wedge$ は $\mathcal{F}$ のもので，
狭義の $(1, 1 + d)$-不等式を満たす．
\end{enumerate}
$x \in X$ を点とし，$W = \overline{\{x\}}$ とおく．
$W \cap Y$ が $Z$ に含まれる既約成分と，含まれない既約成分とをもつならば，
$\text{depth}(\mathcal{F}_x) \geq 1$ である．
\end{lemma}

\begin{proof}
$W \cap Y = W_1 \cup \ldots \cup W_n$ を既約成分への分解とする．
仮定により，番号を付け替えれば，$0 < m < n$ で
$W_1, \ldots, W_m  \subset Z$ かつ
$W_{m + 1}, \ldots, W_n \not \subset Z$ となるものを取れる．したがって，
$$
W \cap Y \setminus
\left((W_1 \cup \ldots \cup W_m) \cap (W_{m + 1} \cup \ldots \cup W_n)\right)
$$
は非連結である．補題 \ref{algebraization-lemma-connected} により，
$1 \leq i \leq m < j \leq n$ と $z \in W_i \cap W_j$ で
$\dim(\mathcal{O}_{W, z}) \leq d + 1$ を満たすものが存在する．
直後の特殊化 $y \leadsto z$ で，$y \in W_j$ かつ $y \not \in Z$ となるものを選ぶ．
$y$ の存在は Properties，補題
\ref{properties-lemma-complement-closed-point-Jacobson} から従う．
$\delta^Y_Z(y) = 1$ かつ $\dim(\mathcal{O}_{W, y}) \leq d$ であることに注意する．
$\mathfrak p \subset \mathcal{O}_{X, y}$ を $x$ に対応する素イデアルとする．
$\mathfrak p' \subset \mathcal{O}_{X, y}^\wedge$ を
$\mathfrak p\mathcal{O}_{X, y}^\wedge$ の上にある極小素イデアルとする．このとき，
$$
\text{depth}(\mathcal{F}_x) =
\text{depth}((\mathcal{F}^\wedge_y)_{\mathfrak p'})
\quad\text{かつ}\quad
\dim(\mathcal{O}_{W, y}) = \dim(\mathcal{O}_{X, y}^\wedge/\mathfrak p')
$$
Algebra，補題 \ref{algebra-lemma-apply-grothendieck-module} および
Local Cohomology，補題 \ref{local-cohomology-lemma-change-completion} を参照されたい．
これで，与えられた不等式から結論を読み取ればよい．
\end{proof}

\begin{lemma}
\label{algebraization-lemma-recover}
状況 \ref{algebraization-situation-algebraize} において，$\mathcal{F}$ を連接
$\mathcal{O}_U$-加群とし，$d \geq 1$ とする．次を仮定する．
\begin{enumerate}
\item $A$ は $I$-進完備で，双対化複体をもち，
$\text{cd}(A, I) \leq d$ である．
\item 完備化 $\mathcal{F}^\wedge$ は $\mathcal{F}$ のもので，
狭義の $(1, 1+ d)$-不等式を満たす．
\item $x \in U$ で $\overline{\{x\}} \cap Y \subset Z$ を満たすものに対して，
$\text{depth}(\mathcal{F}_x) \geq 2$ である．
\end{enumerate}
このとき $H^0(U, \mathcal{F}) \to \lim H^0(U, \mathcal{F}/I^n\mathcal{F})$
は同型である．
\end{lemma}

\begin{proof}
補題 \ref{algebraization-lemma-application-H0} が適用できることを示す．
そこで，$x \in \text{Ass}(\mathcal{F})$ かつ $x \not \in Y$ とする．
$W = \overline{\{x\}}$ とおく．条件 (3) により $W \cap Y \not \subset Z$ である．
補題 \ref{algebraization-lemma-sanity} により，$W \cap Y$ の既約成分はどれも $Z$ に含まれない．
したがって $z \in W \cap Z$ ならば，直後の特殊化 $y \leadsto z$ で
$y \in W \cap Y$ かつ $y \not \in Z$ を満たすものが存在する．
$y$ の存在については，Properties，補題
\ref{properties-lemma-complement-closed-point-Jacobson} を用いる．
このとき $\delta^Y_Z(y) = 1$ であり，仮定から
$\dim(\mathcal{O}_{W, y}) > d$ が従う．したがって
$\dim(\mathcal{O}_{W, z}) > 1 + d$ となり，証明は完了する．
\end{proof}

\begin{lemma}
\label{algebraization-lemma-fully-faithful-inequalities}
状況 \ref{algebraization-situation-algebraize} において，$\mathcal{F}$ を連接
$\mathcal{O}_U$-加群とし，$d \geq 1$ とする．次を仮定する．
\begin{enumerate}
\item $A$ は $I$-進完備で，双対化複体をもち，
$\text{cd}(A, I) \leq d$ である．
\item 完備化 $\mathcal{F}^\wedge$ は $\mathcal{F}$ のもので，
狭義の $(1, 1 + d)$-不等式を満たす．
\item $x \in U$ で $\overline{\{x\}} \cap Y \subset Z$ を満たすものに対して，
$\text{depth}(\mathcal{F}_x) \geq 2$ である．
\end{enumerate}
このとき写像
$$
\Hom_U(\mathcal{G}, \mathcal{F})
\longrightarrow
\Hom_{\textit{Coh}(U, I\mathcal{O}_U)}(\mathcal{G}^\wedge, \mathcal{F}^\wedge)
$$
は，任意の連接 $\mathcal{O}_U$-加群 $\mathcal{G}$ に対して全単射である．
\end{lemma}

\begin{proof}
$\mathcal{H} = \SheafHom_{\mathcal{O}_U}(\mathcal{G}, \mathcal{F})$ とおく．
Cohomology of Schemes，補題
\ref{coherent-lemma-hom-into-depth} または
More on Algebra，補題 \ref{more-algebra-lemma-hom-into-depth} を用いると，
$\mathcal{H}$ の完備化は狭義の $(1, 1 + d)$-不等式を満たし，
$x \in U$ で $\overline{\{x\}} \cap Y \subset Z$ を満たすものに対して，
$\text{depth}(\mathcal{H}_x) \geq 2$ であることが分かる．詳細は省略する．
したがって補題 \ref{algebraization-lemma-recover} により，
$$
\Hom_U(\mathcal{G}, \mathcal{F}) =
H^0(U, \mathcal{H}) =
\lim H^0(U, \mathcal{H}/\mathcal{I}^n\mathcal{H}) =
\Mor_{\textit{Coh}(U, I\mathcal{O}_U)}
(\mathcal{G}^\wedge, \mathcal{F}^\wedge)
$$
最後の等号については，Cohomology of Schemes，補題
\ref{coherent-lemma-completion-internal-hom} を参照されたい．
\end{proof}

\begin{lemma}
\label{algebraization-lemma-construct-unique}
状況 \ref{algebraization-situation-algebraize} において，$(\mathcal{F}_n)$ を
$\textit{Coh}(U, I\mathcal{O}_U)$ の対象とし，$d \geq 1$ とする．次を仮定する．
\begin{enumerate}
\item $A$ は $I$-進完備で，双対化複体をもち，
$\text{cd}(A, I) \leq d$ である．
\item $(\mathcal{F}_n)$ はある連接 $\mathcal{O}_U$-加群の完備化である．
\item $(\mathcal{F}_n)$ は狭義の $(1, 1 + d)$-不等式を満たす．
\end{enumerate}
このとき一意な連接 $\mathcal{O}_U$-加群 $\mathcal{F}$ が存在し，その完備化は
$(\mathcal{F}_n)$ であり，$x \in U$ で $\overline{\{x\}} \cap Y \subset Z$
を満たすものに対して $\text{depth}(\mathcal{F}_x) \geq 2$ である．
\end{lemma}

\begin{proof}
$\mathcal{O}_U$-加群 $\mathcal{F}$ で，その完備化が $(\mathcal{F}_n)$ である連接なものを選ぶ．
次をおく．
$T = \{x \in U \mid \overline{\{x\}} \cap Y \subset Z\}$.
$\mathcal{F}$ を構成するため，Local Cohomology，補題
\ref{local-cohomology-lemma-make-S2-along-T-simple} を $\mathcal{F}$ と $T$ に適用する．
一意性は補題 \ref{algebraization-lemma-fully-faithful-inequalities} の写像性質から従う．

\medskip\noindent
$T$ は $U$ において特殊化で安定なので，確認すべきことは次だけである．
$x' \leadsto x$ が $U$ の点の直後の特殊化で，$x \in T$ かつ
$x' \not \in T$ ならば，$\text{depth}(\mathcal{F}_{x'}) \geq 1$ である．
$W = \overline{\{x\}}$ および $W' = \overline{\{x'\}}$ とおく．
$x' \not \in T$ なので，$W' \cap Y$ は $Z$ に含まれない．
$W' \cap Y$ が $Z$ に含まれる既約成分をもつならば，
補題 \ref{algebraization-lemma-sanity} により証明は完了する．そうでなければ，
$W_1$ を $W \cap Y$ の既約成分，$W'_1$ を $W' \cap Y$ の既約成分として，
$W_1 \subset W'_1$ を満たすものを選ぶ．$z \in W_1$ を生成点とする．
$y \leadsto z$，$y \in W'_1$ が直後の特殊化で，$y \not \in Z$ となるものを取る．
$y$ の存在は $W'_1 \not \subset Z$（上を参照）と Properties，補題
\ref{properties-lemma-complement-closed-point-Jacobson} から従う．
このとき $z \in Z$，$x \leadsto z$，
$x' \leadsto y \leadsto z$，$y \in Y \setminus Z$，かつ $\delta^Y_Z(y) = 1$ である．
Local Cohomology，補題 \ref{local-cohomology-lemma-cd-bound-dim-local} と，$z$ が
$W \cap Y$ の生成点であることにより，$\dim(\mathcal{O}_{W, z}) \leq d$ である．
$x' \leadsto x$ は直後の特殊化なので，$\dim(\mathcal{O}_{W', z}) \leq d + 1$ である．
$y \not = z$ なので，$\dim(\mathcal{O}_{W', y}) \leq d$ を得る．
$\text{depth}(\mathcal{F}_{x'}) = 0$ ならば仮定 (3) と矛盾する．
$\mathcal{O}_{X, y}$ からその完備化への移行の詳細は省略する．これで証明が完了する．
\end{proof}








\section{連接形式加群の代数化 IV}
\label{algebraization-section-algebraization-modules-local}

\noindent
本節では，局所的な場合に補題
\ref{algebraization-lemma-algebraization-principal-variant} のより強い二つの版，すなわち補題
\ref{algebraization-lemma-algebraization-principal} および
\ref{algebraization-lemma-algebraization-principal-bis} を証明する．
これらの補題は，より一般的な命題 \ref{algebraization-proposition-cd-1} によって
不要になるが，その証明はかなり容易である．

\begin{lemma}
\label{algebraization-lemma-algebraization-principal}
状況 \ref{algebraization-situation-algebraize} において，$(\mathcal{F}_n)$ を
$\textit{Coh}(U, I\mathcal{O}_U)$ の対象とする．次を仮定する．
\begin{enumerate}
\item $A$ は局所環で，$\mathfrak a = \mathfrak m$ はその極大イデアルである；
\item $A$ は双対化複体をもつ；
\item $I = (f)$ は非零因子 $f \in \mathfrak m$ によって生成される単項イデアルである；
\item $\mathcal{F}_n$ は有限局所自由
$\mathcal{O}_U/f^n\mathcal{O}_U$-加群である；
\item $\mathfrak p \in V(f) \setminus \{\mathfrak m\}$ ならば
$\text{depth}((A/f)_\mathfrak p) + \dim(A/\mathfrak p) > 1$ である；
\item $\mathfrak p \not \in V(f)$ かつ
$V(\mathfrak p) \cap V(f) \not = \{\mathfrak m\}$ ならば
$\text{depth}(A_\mathfrak p) + \dim(A/\mathfrak p) > 3$.
\end{enumerate}
このとき $(\mathcal{F}_n)$ は標準的に $X$ 上へ延長される．特に，$A$ が
完備ならば，$(\mathcal{F}_n)$ はある連接
$\mathcal{O}_U$-加群の完備化である．
\end{lemma}

\begin{proof}
補題 \ref{algebraization-lemma-when-done} の仮定 (a)，(b)，(c) を
検証することによって証明する．

\medskip\noindent
各 $\mathcal{F}_n$ は $\mathcal{O}_U/f^n\mathcal{O}_U$ 上局所自由なので，
短完全列
$0 \to \mathcal{F}_n \to \mathcal{F}_{n + 1} \to \mathcal{F}_1 \to 0$
がすべての $n$ に対してある．したがって条件 (b) は Cohomology，補題
\ref{cohomology-lemma-topology-I-adic-f} により成り立つ．

\medskip\noindent
$n$ に関する帰納法と短完全列
$0 \to A/f^n \to A/f^{n + 1} \to A/f \to 0$ により，
$A/f^nA$ の随伴素イデアルは $A/fA$ の随伴素イデアルと一致する．
仮定 (3) により，$\mathcal{F}_n$ の随伴点は $A/f^nA$ の随伴素イデアルであって
$\mathfrak m$ と異なるものに対応するので，(5) と Local Cohomology，命題
\ref{local-cohomology-proposition-kollar} から
$M_n = H^0(U, \mathcal{F}_n)$ は有限 $A$-加群である．
したがって仮定 (c) が成り立つ．

\medskip\noindent
証明を終えるには，ある $n > 1$ に対して次の写像の像が
$$
H^1(U, \mathcal{F}_n) \longrightarrow H^1(U, \mathcal{F}_1)
$$
有限長の $A$-加群であることを示せば十分である．実際，Cohomology，補題
\ref{cohomology-lemma-ML-better} により，これは仮定 (a) を導く．補題
\ref{algebraization-lemma-ML-local} により，十分大きい $n$ に対してこの像は $n$ に依存しない．
$\omega_A^\bullet$ を $A$ の正規化された双対化複体とする．局所双対定理と
Matlis 双対性（Dualizing Complexes，補題
\ref{dualizing-lemma-special-case-local-duality} および命題
\ref{dualizing-proposition-matlis}）により，主張は次の写像の像が
$$
\text{Ext}^{-2}_A(M_1, \omega_A^\bullet) \to
\text{Ext}^{-2}_A(M_n, \omega_A^\bullet)
$$
有限長となることと同値である（$n \gg 1$）．ここに現れる加群は有限
$A$-加群であり，$V(f)$ 上に台をもつ．したがって，素イデアル
$\mathfrak q$ であって $f$ を含み $\mathfrak m$ と異なるものにおいて局所化した後にこの写像が
零になることを示せば十分である．$\omega_{A_\mathfrak q}^\bullet$ を
$A_\mathfrak q$ 上の正規化された双対化複体とする．
$\omega_{A_\mathfrak q}^\bullet =
(\omega_A^\bullet)_\mathfrak q[\dim(A/\mathfrak q)]$ であることを思い出そう
（Dualizing Complexes，補題 \ref{dualizing-lemma-dimension-function}）．
(4) で与えられた $\mathcal{F}_n$ の局所構造を用いると，次の写像が
$$
\text{Ext}^{-2 + \dim(A/\mathfrak q)}_{A_\mathfrak q}(
A_\mathfrak q/f, \omega_{A_\mathfrak q}^\bullet)
\to
\text{Ext}^{-2 + \dim(A/\mathfrak q)}_{A_\mathfrak q}(
A_\mathfrak q/f^n, \omega_{A_\mathfrak q}^\bullet)
$$
十分大きい $n$ に対して零になることを示せば十分だと分かる．
$\dim(A/\mathfrak q) > 3$ ならば，これは Local Cohomology，補題
\ref{local-cohomology-lemma-sitting-in-degrees} から直ちに従う．
それ以外の場合には長完全列
$$
\ldots
\xrightarrow{f^n}
H^{-1}(\omega_{A_\mathfrak q}^\bullet)
\to
\text{Ext}^{-1}_{A_\mathfrak q}(
A_\mathfrak q/f^n, \omega_{A_\mathfrak q}^\bullet) \to
H^0(\omega_{A_\mathfrak q}^\bullet)
\xrightarrow{f^n}
H^0(\omega_{A_\mathfrak q}^\bullet)
\to
\text{Ext}^0_{A_\mathfrak q}(
A_\mathfrak q/f^n, \omega_{A_\mathfrak q}^\bullet) \to 0
$$
を用いる．$\dim(A/\mathfrak q) = 2$ ならば，
$H^0(\omega_{A_\mathfrak q}^\bullet) = 0$
である．実際，$\text{depth}(A_\mathfrak q) \geq 1$ である．これは
$f$ が非零因子だからである．
したがって長完全列から，必要な条件は
$$
f^{n - 1} :
H^{-1}(\omega_{A_\mathfrak q}^\bullet)/f \to
H^{-1}(\omega_{A_\mathfrak q}^\bullet)/f^n
$$
が零であることだと分かる．ここで $H^{-1}(\omega^\bullet_\mathfrak q)$ は，
次を満たす素イデアル $\mathfrak p \subset A_\mathfrak q$ 上に台をもつ有限加群である：
$\text{depth}(A_\mathfrak p) + \dim((A/\mathfrak p)_\mathfrak q) \leq 1$.
$\dim((A/\mathfrak p)_\mathfrak q) = \dim(A/\mathfrak p) - 2$ なので，
条件 (6) からこれらの素イデアルは $V(f)$ に含まれる．
したがって，十分大きい $n$ に対する所望の消滅．
最後に $\dim(A/\mathfrak q) = 1$ とする．条件 (5) と $f$ が非零因子であることから，
$A_\mathfrak q$ の深さは少なくとも $2$ である．したがって
$H^0(\omega_{A_\mathfrak q}^\bullet) =
H^{-1}(\omega_{A_\mathfrak q}^\bullet) = 0$
であり，長完全列から主張は次の写像の消滅と同値だと分かる：
$$
f^{n - 1} :
H^{-2}(\omega_{A_\mathfrak q}^\bullet)/f \to
H^{-2}(\omega_{A_\mathfrak q}^\bullet)/f^n
$$
ここで $H^{-2}(\omega^\bullet_\mathfrak q)$ は，素イデアル
$\mathfrak p \subset A_\mathfrak q$ であって
$\text{depth}(A_\mathfrak p) + \dim((A/\mathfrak p)_\mathfrak q)
\leq 2$ を満たすものの上に台をもつ有限加群である．条件 (6) により，
これらの素イデアルはすべて $V(f)$ に含まれる．
したがって，十分大きい $n$ に対する所望の消滅．
\end{proof}

\begin{remark}
\label{algebraization-remark-interesting-case}
$(A, \mathfrak m)$ を次元 $\geq 4$ の完備 Noether 正規局所整域とし，
$f \in \mathfrak m$ は零でないとする．このとき補題
\ref{algebraization-lemma-algebraization-principal} の仮定 (1)，(2)，(3)，(5)，(6) が
満たされる．したがって，$U$ の $U \cap V(f)$ に沿う形式的完備化上の
ベクトル束は代数化できる．補題 \ref{algebraization-lemma-algebraization-principal-bis} では，
これをより一般の連接形式加群へ一般化する．注意
\ref{algebraization-remark-interesting-case-bis} とも比較されたい．
\end{remark}

\begin{lemma}
\label{algebraization-lemma-helper-algebraize}
状況 \ref{algebraization-situation-algebraize} において，$(M_n)$ を補題
\ref{algebraization-lemma-system-of-modules} のような $A$-加群の逆系とし，
$(\mathcal{F}_n)$ を対応する
$\textit{Coh}(U, I\mathcal{O}_U)$ の対象とする．
$d \geq \text{cd}(A, I)$ および $s \geq 0$ を整数とする．
上の記法のもとで次を仮定する．
\begin{enumerate}
\item $A$ は局所環で，極大イデアルは $\mathfrak m = \mathfrak a$ である；
\item $A$ は双対化複体をもつ；
\item $(\mathcal{F}_n)$ は $(s, s + d)$-不等式を満たす
（定義 \ref{algebraization-definition-s-d-inequalities}）．
\end{enumerate}
$E$ を $A$ の剰余体の移入包とする．このとき $i \leq s$ に対して，有限
$A$-加群 $N$ で $I$ のある冪により消されるものと，$n \gg 0$ に対する両立する写像
$$
H^i_\mathfrak m(M_n) \to \Hom_A(N, E)
$$
が存在し，その余核は有限長の $A$-加群であり，その核 $K_n$ は，
$\Im(K_{n''} \to K_{n'})$ が $n'' \gg n' \gg 0$ に対して有限長となるような
逆系をなす．
\end{lemma}

\begin{proof}
$\omega_A^\bullet$ を正規化された双対化複体とする．このとき
$\delta^Y_Z = \delta$ はこの双対化複体に付随する次元関数である．
$\Ext^{-i}_A(M_n, \omega_A^\bullet)$ は有限 $A$-加群であり，$I^n$ で消されることに
注意する．$0 \leq i \leq s$ を固定する．以下では $n_1 > n_0 > 0$ を，
$$
N = \Im(\Ext^{-i}_A(M_{n_0}, \omega_A^\bullet) \to
\Ext^{-i}_A(M_{n_1}, \omega_A^\bullet))
$$
とおいたとき，写像
$$
N \to \Ext^{-i}_A(M_n, \omega_A^\bullet),\quad n \geq n_1
$$
の核が有限長の $A$-加群となり，余核 $Q_n$ が，
$\Im(Q_{n'} \to Q_{n''})$ が $n'' \gg n' \gg n_1$ に対して有限長となるような
系をなすように選ぶ．
これは，系 $\{\Ext^{-i}_A(M_n, \omega_A^\bullet)\}_{n \geq 1}$ が，
有限 $A$-加群の圏を有限長 $A$-加群の Serre 部分圏で割った商圏において
本質的に定値であるという主張と同値である．局所双対定理
（Dualizing Complexes，補題
\ref{dualizing-lemma-special-case-local-duality}）と Matlis 双対性
（Dualizing Complexes，命題 \ref{dualizing-proposition-matlis}）から，
補題の主張にある写像
$$
H^i_\mathfrak m(M_n) \to \Hom_A(N, E),\quad n \geq n_1
$$
が得られる．

\medskip\noindent
$f \in \mathfrak m$ を取る．$B = A_f^\wedge$ を，$I$-進完備化によって
局所化 $A_f$ から得られる環とする．
$\omega_{A_f}^\bullet = \omega_A^\bullet \otimes_A A_f$
および $\omega_B^\bullet = \omega_A^\bullet \otimes_A B$ は双対化複体である
（Dualizing Complexes，補題 \ref{dualizing-lemma-dualizing-localize} および
\ref{dualizing-lemma-completion-henselization-dualizing}）．$M$ を有限 $B$-加群
$\lim M_{n, f}$ とする（Cohomology of Schemes，補題
\ref{coherent-lemma-inverse-systems-affine} の議論と比較されたい）．このとき
$$
\Ext^{-i}_A(M_n, \omega_A^\bullet)_f =
\Ext^{-i}_{A_f}(M_{n, f}, \omega_{A_f}^\bullet) =
\Ext^{-i}_B(M/I^n M, \omega_B^\bullet)
$$
$\mathfrak m$ は有限個の $f \in \mathfrak m$ によって生成できるので，各 $f$ に対して系
$$
\{\Ext^{-i}_B(M/I^n M, \omega_B^\bullet)\}_{n \geq 1}
$$
が本質的に定値であることを示せば十分である．詳細の一部を省略する．

\medskip\noindent
$\mathfrak q \subset IB$ を素イデアルとする．すると $\mathfrak q$ は点
$y \in U \cap Y$ に対応する．$\delta(\mathfrak q) = \dim(\overline{\{y\}})$ は
$\omega_B^\bullet$ に付随する次元関数の値でもあることに注意する
（詳細は省略する．$\omega_B^\bullet$ が $\omega_A^\bullet$ と $B$ との
テンソル積によって得られることを用いよ）．仮定 (3) と補題
\ref{algebraization-lemma-elementary} により，補題
\ref{algebraization-lemma-algebraize-local-cohomology-bis} を
$B_\mathfrak q, IB_\mathfrak q, \mathfrak qB_\mathfrak q, M_\mathfrak q$
に適用できる．ただし $s$ は $s - \delta(y)$ で置き換える．これにより
$$
H^{i - \delta(\mathfrak q)}_{\mathfrak qB_\mathfrak q}(M_\mathfrak q) =
\lim H^{i - \delta(\mathfrak q)}_{\mathfrak qB_\mathfrak q}(
(M/I^nM)_\mathfrak q)
$$
を得て，この加群は $I$ のある冪で消される．補題
\ref{algebraization-lemma-terrific} により，逆系
$H^{i - \delta(\mathfrak q)}_{\mathfrak qB_\mathfrak q}((M/I^nM)_\mathfrak q)$
は本質的に定値であり，その値は
$H^{i - \delta(\mathfrak q)}_{\mathfrak qB_\mathfrak q}(M_\mathfrak q)$ である．
$(\omega_B^\bullet)_\mathfrak q[-\delta(\mathfrak q)]$ は $B_\mathfrak q$ 上の
正規化された双対化複体なので，局所双対定理から系
$$
\Ext^{-i}_B(M/I^n M, \omega_B^\bullet)_\mathfrak q
$$
は本質的に定値であり，その値は
$\Ext^{-i}_B(M, \omega_B^\bullet)_\mathfrak q$ であると分かる．

\medskip\noindent
証明を終えるため，補題 \ref{algebraization-lemma-bootstrap} の証明と同様に大域化する．
ここでは系の値が既に分かっているので，議論はより容易である．すなわち，写像
$$
\alpha_n :
\Ext^{-i}_B(M/I^n M, \omega_B^\bullet)
\longrightarrow
\Ext^{-i}_B(M, \omega_B^\bullet)
$$
を，$n$ を動かして考える．上の結果により，各 $\mathfrak q$ に対して，ある $n$ を，
$\alpha_n$ が $\mathfrak q$ で局所化した後に全射となるように見つけられる．
$B$ は Noether 環であり，$\Ext^{-i}_B(M, \omega_B^\bullet)$ は有限加群なので，
ある $n$ を，$\alpha_n$ が全射となるように見つけられる．任意の $n$ で
$\alpha_n$ が全射となるものと，素イデアル $\mathfrak q \in V(IB)$ を与えると，ある $n' > n$ を，
$\Ker(\alpha_n)$ が $\Ext^{-i}(M/I^{n'}M, \omega_B^\bullet)$ へ零写像で移るように
見つけられる．ただし少なくとも $\mathfrak q$ で局所化した後についてである．
$\Ker(\alpha_n)$ は有限 $A$-加群であり，切断の台は準コンパクトなので，ある $n'$ を，
$\Ker(\alpha_n)$ が $\Ext^{-i}(M/I^{n'}M, \omega_B^\bullet)$ へ零写像で移るように
見つけられる．
このようにして，$\Ext^{-i}(M/I^n M, \omega_B^\bullet)$ は本質的に定値で，
その値は $\Ext^{-i}(M, \omega_B^\bullet)$ であると分かる．これで証明が終わる．
\end{proof}

\noindent
補題 \ref{algebraization-lemma-algebraization-principal} のより一般的な版を与える．

\begin{lemma}
\label{algebraization-lemma-algebraization-principal-bis}
状況 \ref{algebraization-situation-algebraize} において，$(\mathcal{F}_n)$ を
$\textit{Coh}(U, I\mathcal{O}_U)$ の対象とする．次を仮定する．
\begin{enumerate}
\item $A$ は局所環で，$\mathfrak a = \mathfrak m$ はその極大イデアルである；
\item $A$ は双対化複体をもつ；
\item $I = (f)$ は単項イデアルである；
\item $(\mathcal{F}_n)$ は $(2, 3)$-不等式を満たす．
\end{enumerate}
このとき $(\mathcal{F}_n)$ は $X$ 上へ延長される．特に，$A$ が
$I$-進完備ならば，$(\mathcal{F}_n)$ はある連接
$\mathcal{O}_U$-加群の完備化である．
\end{lemma}

\begin{proof}
$\textit{Coh}(U, I\mathcal{O}_U)$ はアーベル圏であることを思い出そう．
Cohomology of Schemes，補題
\ref{coherent-lemma-inverse-systems-abelian} を参照されたい．$U$ のアフィン開集合上で，
対象 $(\mathcal{F}_n)$ は Noether 環上の有限加群に対応する
（Cohomology of Schemes，補題
\ref{coherent-lemma-inverse-systems-affine}）．したがって写像
$f^N : (\mathcal{F}_n) \to (\mathcal{F}_n)$ の核は，十分大きい $N$ に対して
安定する．補題 \ref{algebraization-lemma-map-kernel-cokernel-on-closed} および
\ref{algebraization-lemma-essential-image-completion} により，補題を証明するにあたって
$(\mathcal{F}_n)$ をそのような写像の像で置き換えてよい．したがって，
$f$ は $(\mathcal{F}_n)$ 上単射であると仮定してよい．この置換の後，逆系
$(\mathcal{F}_n)$ は $U$ 上で補題 \ref{algebraization-lemma-equivalent-f-good} の同値な条件を満たす．
以下の証明では，これを断りなく用いる．

\medskip\noindent
補題 \ref{algebraization-lemma-when-done} の仮定 (a)，(b)，(c) を検証する．
仮定 (b) は Cohomology，補題
\ref{cohomology-lemma-topology-I-adic-f} により成り立つ．

\medskip\noindent
補題 \ref{algebraization-lemma-system-of-modules} のような加群の逆系 $\{M_n\}$ を取る．
必要ならば $H^0_\mathfrak m(M_n) = 0$ と仮定してよい．これは $M_n$ を
$M_n/H^0_\mathfrak m(M_n)$ で置き換えることにより達成できる．すると短完全列
$$
0 \to M_n \to H^0(U, \mathcal{F}_n) \to H^1_\mathfrak m(M_n) \to 0
$$
をすべての $n$ に対して得る．$E$ を $A$ の剰余体の移入包とする．補題
\ref{algebraization-lemma-helper-algebraize} と現在の仮定 (4) により，整数 $m \geq 0$，有限
$A$-加群 $N_1$，$N_2$ であって，$f^c$ で消されるもの（ある $c \geq 0$ に対して），および
両立する写像系
$$
H^i_\mathfrak m(M_n) \to \Hom_A(N_i, E), \quad i = 1, 2
$$
を $n \geq m$ に対して，補題で述べた性質をもつように選べる．

\medskip\noindent
$M = \lim H^0(U, \mathcal{F}_n)$ は $A$-加群であり，その極限位相は $f$-進位相であることが
分かっている．したがって，$n$ を与えると，加群 $M/f^nM$ は
$H^0(U, \mathcal{F}_N)$ の部分商である．ここである $N \gg n$ を選んでいる．
上で得た情報を見ると，$f^cM/f^nM$ は有限 $A$-加群である．$f$ は $M$ 上の
非零因子なので，$M/f^{n - c}M$ は有限 $A$-加群であると結論できる．
このようにして，補題 \ref{algebraization-lemma-when-done} の仮定 (c) が成り立つと分かる．

\medskip\noindent
次に，加群
$$
Ob = \lim H^1(U, \mathcal{F}_n) = \lim H^2_\mathfrak m(M_n)
$$
を調べる．$n \geq m$ に対し，$K_n$ を写像
$H^2_\mathfrak m(M_n) \to \Hom_A(N_2, E)$ の核とする．
$K = \lim K_n$ とおくと，完全列
$$
0 \to K \to Ob \to \Hom_A(N_2, E)
$$
を得る．上の議論により，$Ob = \lim H^2_\mathfrak m(M_n)$ 上の極限位相は
$f$-進位相である．$N_2$ は $f^c$ で消されるので，$K = \lim K_n$ 上の
極限位相についても同じことが成り立つ．したがって $K/fK$ は $K_n$ の部分商である
（$n \gg 1$）．一方，補題 \ref{algebraization-lemma-helper-algebraize} の結論により
$\{K_n\}$ は有限長 $A$-加群の逆系と pro-同型なので，$K/fK$ は有限長
$A$-加群の部分商である．よって $K$ は有限 $A$-加群である．Algebra，補題
\ref{algebra-lemma-finite-over-complete-ring} を参照されたい．
（実際，$\dim(\text{Supp}(K)) = 1$ であることさえ分かるが，これは用いない．）

\medskip\noindent
$n \geq 1$ を与え，境界写像
$$
\delta_n :
H^0(U, \mathcal{F}_n)
\longrightarrow
\lim_N H^1(U, f^n\mathcal{F}_N) \xrightarrow{f^{-n}} Ob
$$
を考える（第2の写像は同型である）．これは短完全列
$$
0 \to f^n\mathcal{F}_N \to \mathcal{F}_N \to \mathcal{F}_n \to 0
$$
から得られる．各 $n$ に対して
$$
P_n = \Im(H^0(U, \mathcal{F}_{n + m}) \to H^0(U, \mathcal{F}_n))
$$
とおく．ここで $m$ は上で選んだものである．$\{P_n\}$ は逆系をなし，大域切断上の写像
$f : \mathcal{F}_n \to \mathcal{F}_{n + 1}$ は $P_n$ を $P_{n + 1}$ へ写すことに
注意する．$p \in P_n$ ならば，$\delta_n(p) \in K \subset Ob$ である．実際，
$\delta_n(p)$ は
$H^1(U, f^n\mathcal{F}_{n + m}) = H^2_\mathfrak m(M_m)$
において零に写り，また $\delta_n$ と $Ob \to \Hom_A(N_2, E)$ の合成は
$H^2_\mathfrak m(M_m)$ を経由するからである．これは $m$ の選び方による．
したがって
$$
\bigoplus\nolimits_{n \geq 0} \Im(P_n \to Ob)
$$
は有限生成次数付き $A[T]$-加群であり，$T$ は $f$ 倍として作用する．実際，これは
$K[T]$ の次数付き部分加群であり，$K$ は $A$ 上有限である．Cohomology，補題
\ref{cohomology-lemma-ML-general} の証明と同様に議論すると\footnote{表示した加群の
$\delta_{n_j}(p_j)$ という形の斉次生成元を選ぶ．$k = \max(n_j)$ とおけば，
$n \geq k$ と任意の $p \in P_n$ に対して，ある $a_j \in A$ を，
$p - \sum a_j f^{n - n_j} p_j$ が $\delta_n$ の核に入り，したがって
$P_{n'}$ の像に入るように見つけられる．これはすべての $n' \geq n$ に対して成り立つ．
ゆえに $\Im(P_n \to P_{n - k}) = \Im(P_{n'} \to P_{n - k})$ が
すべての $n' \geq n$ に対して成り立つ．}，
逆系 $\{P_n\}$ は ML を満たすと分かる．$\{P_n\}$ は
$\{H^0(U, \mathcal{F}_n)\}$ と pro-同型なので，
$\{H^0(U, \mathcal{F}_n)\}$ は ML をもつと結論できる．したがって補題
\ref{algebraization-lemma-when-done} の仮定 (a) が成り立ち，証明が完了する．
\end{proof}

\noindent
補題 \ref{algebraization-lemma-algebraization-principal-bis} の条件は，次のように具体化できる．

\begin{lemma}
\label{algebraization-lemma-unwinding-conditions}
状況 \ref{algebraization-situation-algebraize} において，$(\mathcal{F}_n)$ を
$\textit{Coh}(U, I\mathcal{O}_U)$ の対象とする．次を仮定する．
\begin{enumerate}
\item $A$ は極大イデアル $\mathfrak a = \mathfrak m$ をもつ局所環である；
\item $\text{cd}(A, I) = 1$.
\end{enumerate}
このとき，$(\mathcal{F}_n)$ が $(2, 3)$-不等式を満たすための必要十分条件は，
すべての $y \in U \cap Y$ で $\dim(\{y\}) = 1$ を満たすものと，すべての素イデアル
$\mathfrak p \subset \mathcal{O}_{X, y}^\wedge$ で
$\mathfrak p \not \in V(I\mathcal{O}_{X, y}^\wedge)$ を満たすものに対して
$$
\text{depth}((\mathcal{F}_y^\wedge)_\mathfrak p) +
\dim(\mathcal{O}_{X, y}^\wedge/\mathfrak p) > 2
$$
が成り立つことである．
\end{lemma}

\begin{proof}
補題 \ref{algebraization-lemma-explain-2-3-cd-1} を以下では断りなく用いる．特に，
この条件が必要であることが分かる．逆に，条件が成り立つと仮定する．補題
\ref{algebraization-lemma-discussion} により
$\delta^Y_Z(y) = \dim(\overline{\{y\}})$ であることに注意し，この関数を
$\delta$ と書く．$y \in U \cap Y$ とする．$\delta(y) > 2$ ならば，補題
\ref{algebraization-lemma-explain-2-3-cd-1} の不等式が成り立つ．最後に
$\delta(y) = 2$ と仮定する．示すべきことは
$$
\text{depth}((\mathcal{F}_y^\wedge)_\mathfrak p) +
\dim(\mathcal{O}_{X, y}^\wedge/\mathfrak p) > 1
$$
である．特殊化 $y \leadsto y'$ で $\delta(y') = 1$ を満たすものを選ぶ．すると環準同型
$\mathcal{O}_{X, y'}^\wedge \to \mathcal{O}_{X, y}^\wedge$ があり，標的は
$\mathcal{O}_{X, y'}^\wedge$ をある素イデアル $\mathfrak q$ で局所化して完備化したものと
同一視される．ここで $\dim(\mathcal{O}_{X, y'}^\wedge/\mathfrak q) = 1$ である．
さらに
$$
\mathcal{F}_y^\wedge =
\mathcal{F}_{y'}^\wedge
\otimes_{\mathcal{O}_{X, y'}^\wedge}
\mathcal{O}_{X, y}^\wedge
$$
を得る．$\mathfrak p' \subset \mathcal{O}_{X, y'}^\wedge$ を
$\mathfrak p$ の像とする．Local Cohomology，補題
\ref{local-cohomology-lemma-change-completion} により
\begin{align*}
\text{depth}((\mathcal{F}_y^\wedge)_\mathfrak p) +
\dim(\mathcal{O}_{X, y}^\wedge/\mathfrak p)
& =
\text{depth}((\mathcal{F}_{y'}^\wedge)_{\mathfrak p'}) +
\dim((\mathcal{O}_{X, y}^\wedge/\mathfrak p)_{\mathfrak p'}) \\
& =
\text{depth}((\mathcal{F}_{y'}^\wedge)_{\mathfrak p'}) +
\dim(\mathcal{O}_{X, y}^\wedge/\mathfrak p') - 1
\end{align*}
が成り立つ．最後の等式は，特殊化が直後のものだからである．
したがって，仮定した $y', \mathfrak p'$ に対する不等式により補題の証明．
\end{proof}

\begin{lemma}
\label{algebraization-lemma-unwinding-conditions-bis}
状況 \ref{algebraization-situation-algebraize} において，$(\mathcal{F}_n)$ を
$\textit{Coh}(U, I\mathcal{O}_U)$ の対象とする．次を仮定する．
\begin{enumerate}
\item $A$ は極大イデアル $\mathfrak a = \mathfrak m$ をもつ局所環である；
\item $A$ は双対化複体をもつ；
\item $\text{cd}(A, I) = 1$,
\item $y \in U \cap Y$ に対して，加群 $\mathcal{F}_y^\wedge$ は
$V(I\mathcal{O}_{X, y}^\wedge)$ の外で有限局所自由である．たとえば
$\mathcal{F}_n$ が有限局所自由 $\mathcal{O}_U/I^n\mathcal{O}_U$-加群ならばよい；
\item 次のいずれか一方が成り立つ：
\begin{enumerate}
\item $A_f$ は $(S_2)$ であり，$X$ のすべての既約成分で $Y$ に含まれないものは
次元 $\geq 4$ である；または
\item $\mathfrak p \not \in V(f)$ かつ
$V(\mathfrak p) \cap V(f) \not = \{\mathfrak m\}$ ならば
$\text{depth}(A_\mathfrak p) + \dim(A/\mathfrak p) > 3$.
\end{enumerate}
\end{enumerate}
このとき $(\mathcal{F}_n)$ は $(2, 3)$-不等式を満たす．
\end{lemma}

\begin{proof}
補題 \ref{algebraization-lemma-unwinding-conditions} の判定法を用いる．
$y \in U \cap Y$ で $\dim(\overline{\{y\}} = 1$ を満たすものを取り，素イデアル
$\mathfrak p$ で $\mathfrak p \subset \mathcal{O}_{X, y}^\wedge$ かつ
$\mathfrak p \not \in V(I\mathcal{O}_{X, y}^\wedge)$ を満たすものを取る．条件 (4) から
$\text{depth}((\mathcal{F}_y^\wedge)_\mathfrak p) =
\text{depth}((\mathcal{O}_{X, y}^\wedge)_\mathfrak p)$.
したがって，示すべきことは
$$
\text{depth}((\mathcal{O}_{X, y}^\wedge)_\mathfrak p) +
\dim(\mathcal{O}_{X, y}^\wedge/\mathfrak p) > 2
$$
である．$\mathfrak p_0 \subset A$ を $\mathfrak p$ の像とし，
$\mathfrak q \subset A$ を $y$ に対応する素イデアルとする．
Local Cohomology，補題 \ref{local-cohomology-lemma-change-completion} により
\begin{align*}
\text{depth}((\mathcal{O}_{X, y}^\wedge)_\mathfrak p) +
\dim(\mathcal{O}_{X, y}^\wedge/\mathfrak p)
& =
\text{depth}(A_{\mathfrak p_0}) + \dim((A/\mathfrak p_0)_\mathfrak q) \\
& =
\text{depth}(A_{\mathfrak p_0}) + \dim(A/\mathfrak p_0) - 1
\end{align*}
(5)(a) が成り立つならば，これは
$$
\geq \min(2, \dim(A_{\mathfrak p_0})) + \dim(A/\mathfrak p_0) - 1
$$
以上である．いずれにせよ $\dim(A/\mathfrak p_0) \geq 2$ であることに注意する．
したがって最小値が $2$ ならば証明は終わる．そうでなければ
$$
\dim(A_{\mathfrak p_0}) + \dim(A/\mathfrak p_0) - 1 \geq 4 - 1
$$
を得る．実際，$\Spec(A)$ の各成分で $\mathfrak p_0$ を通るものは次元
$\geq 4$ である．(5)(b) が成り立つならば，結論は直ちに従う．
\end{proof}

\begin{remark}
\label{algebraization-remark-interesting-case-bis}
$(A, \mathfrak m)$ を，双対化複体をもち $f \in \mathfrak m$ に関して完備な
Noether 局所環とする．$(\mathcal{F}_n)$ を $\textit{Coh}(U, f\mathcal{O}_U)$ の対象とする．
ここで $U$ は $A$ の穴あきスペクトルである．$Y = V(f) \subset X = \Spec(A)$ とおく．
$y \in U \cap V(f)$ が $U$ において閉じている，すなわち
$\dim(\overline{\{y\}}) = 1$ を満たすとき，$\mathcal{O}_{X, y}^\wedge$-加群
$\mathcal{F}_y^\wedge$ が次の二条件を満たすと仮定する．
\begin{enumerate}
\item $\mathcal{F}_y^\wedge[1/f]$ は $(S_2)$ である
$\mathcal{O}_{X, y}^\wedge[1/f]$-加群である；
\item $\mathfrak p \in \text{Ass}(\mathcal{F}_y^\wedge[1/f])$ に対して
$\dim(\mathcal{O}_{X, y}^\wedge/\mathfrak p) \geq 3$ である．
\end{enumerate}
このとき $(\mathcal{F}_n)$ は $U$ 上の連接加群の完備化である．これは補題
\ref{algebraization-lemma-algebraization-principal-bis} および
\ref{algebraization-lemma-unwinding-conditions} から従う．
\end{remark}
































\section{連接形式加群の改良}
\label{algebraization-section-improving-formal-modules}

\noindent
Noether スキーム $X$ を取る．$Y \subset X$ を，準連接イデアル層
$\mathcal{I} \subset \mathcal{O}_X$ をもつ閉部分スキームとする．
$(\mathcal{F}_n)$ を $\textit{Coh}(X, \mathcal{I})$ の対象とする．本節では，
$(\mathcal{F}_n) \to (\mathcal{F}'_n)$
という写像で，連接加群に対して Local Cohomology，第
\ref{local-cohomology-section-improve} 節で構成した写像と同様のものを構成する．
点 $y \in Y$ に対して
$$
\mathcal{O}_{X, y}^\wedge = \lim \mathcal{O}_{X, y}/\mathcal{I}^n_y, \quad
\mathcal{I}_y^\wedge = \lim \mathcal{I}_y/\mathcal{I}^n_y
\quad\text{かつ}\quad
\mathfrak m_y^\wedge = \lim \mathfrak m_y/\mathcal{I}_y^n
$$
とおく．すると $\mathcal{O}_{X, y}^\wedge$ は，極大イデアル
$\mathfrak m_y^\wedge$ をもつ Noether 局所環で，
$\mathcal{I}_y^\wedge = \mathcal{I}_y\mathcal{O}_{X, y}^\wedge$ に関して完備である．
また
$$
\mathcal{F}_y^\wedge = \lim \mathcal{F}_{n, y}
$$
とおく．すると $\mathcal{F}_y^\wedge$ は $\mathcal{O}_{X, y}^\wedge$ 上の有限加群であり，
$\mathcal{F}_y^\wedge/(\mathcal{I}_y^\wedge)^n\mathcal{F}_y^\wedge =
\mathcal{F}_{n, y}$ がすべての $n$ に対して成り立つ．Algebra，補題
\ref{algebra-lemma-limit-complete} および
\ref{algebra-lemma-finite-over-complete-ring} を参照されたい．

\begin{lemma}
\label{algebraization-lemma-divide-torsion-formal-coherent-module}
上の状況で，$X$ は局所的に双対化複体をもつと仮定する．
$T \subset Y$ を特殊化で安定な部分集合とする．$y \in T$ と，非極大素イデアル
$\mathfrak p \subset \mathcal{O}_{X, y}^\wedge$ で
$V(\mathfrak p) \cap V(\mathcal{I}^\wedge_y) = \{\mathfrak m_y^\wedge\}$ を満たすものに対して
$$
\text{depth}_{(\mathcal{O}_{X, y})_\mathfrak p}
((\mathcal{F}^\wedge_y)_\mathfrak p) > 0
$$
が成り立つと仮定する．このとき標準的な写像
$(\mathcal{F}_n) \to (\mathcal{F}_n')$
で，連接 $\mathcal{O}_X$-加群の逆系の間のものが存在し，次の性質をもつ．
\begin{enumerate}
\item $y \in T$ に対して $\text{depth}(\mathcal{F}'_{n, y}) \geq 1$ である；
\item $(\mathcal{F}'_n)$ は pro-系としてある対象 $(\mathcal{G}_n)$ と同型であり，
後者は $\textit{Coh}(X, \mathcal{I})$ の対象である；
\item 誘導される射 $(\mathcal{F}_n) \to (\mathcal{G}_n)$ は
$\textit{Coh}(X, \mathcal{I})$ の射として全射であり，その核は
$\mathcal{I}$ のある冪で消される．
\end{enumerate}
\end{lemma}

\begin{proof}
各 $n$ に対して，$\mathcal{F}_n \to \mathcal{F}'_n$ を Local Cohomology，補題
\ref{local-cohomology-lemma-get-depth-1-along-Z} で構成した全射とする．これは
$\mathcal{F}_n$ を，$T$ 上に台をもつ切断の部分層で割った商なので，標準的な写像
$\mathcal{F}'_{n + 1} \to \mathcal{F}'_n$ が得られる．これにより，写像
$(\mathcal{F}_n) \to (\mathcal{F}_n')$ を連接 $\mathcal{O}_X$-加群の逆系の間に得る．
性質 (1) は構成から成り立つ．

\medskip\noindent
性質 (2)，(3) を証明するため，$X = \Spec(A_0)$ はアフィンで，$A_0$ は双対化複体を
もつと仮定してよい．$I_0 \subset A_0$ を $Y$ に対応するイデアルとする．
$A, I$ を，$I_0$-進完備化によって $A_0, I_0$ から得られるものとする．
後で用いるため，$A$ は双対化複体をもつことに注意しておく
（Dualizing Complexes，補題 \ref{dualizing-lemma-ubiquity-dualizing}）．
$M$ を有限 $A$-加群で $(\mathcal{F}_n)$ に対応するものとする．Cohomology of Schemes，補題
\ref{coherent-lemma-inverse-systems-affine} を参照されたい．すると
$\mathcal{F}_n$ は $M_n = M/I^nM$ に対応する．また，$\mathcal{F}'_n$ は商
$M'_n = M_n / H^0_T(M_n)$ に対応することを思い出そう．Local Cohomology，補題
\ref{local-cohomology-lemma-get-depth-1-along-Z} とその証明を参照されたい．

\medskip\noindent
$s = 0$ および $d = \text{cd}(A, I)$ とおく．$A, I, T, M, s, d$ は状況
\ref{algebraization-situation-bootstrap} の仮定 (1)，(3)，(4)，(6) を満たすと主張する．
実際，(1)，(3) は上から直ちに従い，$s = 0$ なので (4) は空条件であり，
(6) は補題の主張で置いた仮定である．

\medskip\noindent
定理 \ref{algebraization-theorem-final-bootstrap} により，$\{H^0_T(M_n)\}$ は本質的に定値で，
その値は $H^0_T(M)$ である．したがって短完全列
$0 \to H^0_T(M_n) \to M_n \to M'_n \to 0$ の極限は短完全列
$$
0 \to H^0_T(M) \to M \to \lim M'_n \to 0
$$
である．$M' = \lim M'_n = M/H^0_T(M)$ とおく．これは有限 $A$-加群である．
$(\mathcal{G}_n)$ を $\textit{Coh}(X, \mathcal{I})$ の対象で，$M'$ に対応するものとする．
証明を終えるには，標準的な写像 $\{M'/I^nM'\} \to \{M'_n\}$ が pro-同型であることを
示さなければならない．これは，
$\{H^0_T(M) + I^nM\} \to \{\Ker(M \to M'_n)\}$ が pro-同型であることと同値である．
$I^nM$ で割れば，$\{H^0_T(M)/H^0_T(M) \cap I^nM\} \to \{H^0_T(M_n)\}$ が
pro-同型であることを示せば十分である．$H^0_T(M)$ は $I$ のある冪で消されるので，
Artin--Rees により $H^0_T(M) \cap I^nM = 0$ がすべての $n \gg 0$ に対して成り立つ．
したがって所望の pro-同型を得る．実際，上で $\{H^0_T(M_n)\}$ は本質的に定値で，
その値が $H^0_T(M)$ であることを見た．
\end{proof}

\begin{lemma}
\label{algebraization-lemma-improvement-formal-coherent-module-better}
上の状況で，$X$ は局所的に双対化複体をもつと仮定する．
$T' \subset T \subset Y$ を特殊化で安定な部分集合とする．
$d \geq 0$ を整数とする．次を仮定する．
\begin{enumerate}
\item[(a)] アフィン局所的に $X = \Spec(A_0)$，$Y = V(I_0)$ であり，
$\text{cd}(A_0, I_0) \leq d$ である；
\item[(b)] $y \in T$ と，非極大素イデアル
$\mathfrak p \subset \mathcal{O}_{X, y}^\wedge$ で
$V(\mathfrak p) \cap V(\mathcal{I}_y^\wedge) = \{\mathfrak m_y^\wedge\}$ を満たすものに対して
$$
\text{depth}_{(\mathcal{O}_{X, y})_\mathfrak p}
((\mathcal{F}^\wedge_y)_\mathfrak p) > 0
$$
である；
\item[(c)] $y \in T'$ と，素イデアル
$\mathfrak p \subset \mathcal{O}_{X, y}^\wedge$ で
$\mathfrak p \not \in V(\mathcal{I}_y^\wedge)$ かつ
$V(\mathfrak p) \cap V(\mathcal{I}_y^\wedge) \not =
\{\mathfrak m_y^\wedge\}$ を満たすものに対して
$$
\text{depth}_{(\mathcal{O}_{X, y})_\mathfrak p}
((\mathcal{F}^\wedge_y)_\mathfrak p) \geq 1
\quad\text{または}\quad
\text{depth}_{(\mathcal{O}_{X, y})_\mathfrak p}
((\mathcal{F}^\wedge_y)_\mathfrak p) +
\dim(\mathcal{O}_{X, y}^\wedge/\mathfrak p) > 1 + d
$$
である；
\item[(d)] $y \in T'$ と，非極大素イデアル
$\mathfrak p \subset \mathcal{O}_{X, y}^\wedge$ で
$V(\mathfrak p) \cap V(\mathcal{I}_y^\wedge) = \{\mathfrak m_y^\wedge\}$ を満たすものに対して
$$
\text{depth}_{(\mathcal{O}_{X, y})_\mathfrak p}
((\mathcal{F}^\wedge_y)_\mathfrak p) > 1
$$
である；
\item[(e)] $y \leadsto y'$ が直後の特殊化であり，$y' \in T'$ ならば
$y \in T$ である．
\end{enumerate}
このとき標準的な写像 $(\mathcal{F}_n) \to (\mathcal{F}_n'')$ で，連接
$\mathcal{O}_X$-加群の逆系の間のものが存在し，次の性質をもつ．
\begin{enumerate}
\item $y \in T$ に対して $\text{depth}(\mathcal{F}''_{n, y}) \geq 1$ である；
\item $y' \in T'$ に対して $\text{depth}(\mathcal{F}''_{n, y'}) \geq 2$ である；
\item $(\mathcal{F}''_n)$ は pro-系としてある対象 $(\mathcal{H}_n)$ と同型であり，
後者は $\textit{Coh}(X, \mathcal{I})$ の対象である；
\item 誘導される射 $(\mathcal{F}_n) \to (\mathcal{H}_n)$ は
$\textit{Coh}(X, \mathcal{I})$ の射として，$\mathcal{I}$ のある冪で消される核と余核をもつ．
\end{enumerate}
\end{lemma}

\begin{proof}
補題 \ref{algebraization-lemma-divide-torsion-formal-coherent-module} とその証明と同様に，各 $n$ に対して
$\mathcal{F}_n \to \mathcal{F}'_n$ を，$T$ を用いて Local Cohomology，補題
\ref{local-cohomology-lemma-get-depth-1-along-Z} で構成した全射とする．次に，
$\mathcal{F}'_n \to \mathcal{F}''_n$ を Local Cohomology，補題
\ref{local-cohomology-lemma-make-S2-along-T} とその証明で構成した単射とする．
これらの構成から標準的な写像 $\mathcal{F}''_{n + 1} \to \mathcal{F}''_n$ が得られ，したがって写像
$$
(\mathcal{F}_n) \longrightarrow (\mathcal{F}_n') \longrightarrow
(\mathcal{F}''_n)
$$
を連接 $\mathcal{O}_X$-加群の逆系の間に得る．性質 (1)，(2) は構成から成り立つ．

\medskip\noindent
性質 (3)，(4) を証明するため，$X = \Spec(A_0)$ はアフィンで，$A_0$ は双対化複体を
もつと仮定してよい．$I_0 \subset A_0$ を $Y$ に対応するイデアルとする．
$A, I$ を，$I_0$-進完備化によって $A_0, I_0$ から得られるものとする．
後で用いるため，$A$ は双対化複体をもつことに注意しておく
（Dualizing Complexes，補題 \ref{dualizing-lemma-ubiquity-dualizing}）．
$M$ を有限 $A$-加群で $(\mathcal{F}_n)$ に対応するものとする．Cohomology of Schemes，補題
\ref{coherent-lemma-inverse-systems-affine} を参照されたい．すると
$\mathcal{F}_n$ は $M_n = M/I^nM$ に対応する．また，$\mathcal{F}'_n$ は商
$M'_n = M_n / H^0_T(M_n)$ に対応することを思い出そう．さらに，
$M' = \lim M'_n$ は $M$ を $H^0_T(M)$ で割った商であり，$H^0_T(M)$ は
$I$-冪捩れであり，$\{M'_n\}$ と $\{M'/I^nM'\}$ は pro-系として同型である．
最後に，$\mathcal{F}''_n$ は拡大
$$
0 \to M'_n \to M''_n \to H^1_{T'}(M'_n) \to 0
$$
に対応する．Local Cohomology，補題
\ref{local-cohomology-lemma-make-S2-along-T} の証明を参照されたい．

\medskip\noindent
$s = 1$ とおく．$A, I, T', M', s, d$ は状況 \ref{algebraization-situation-bootstrap} の仮定
(1)，(3)，(4)，(6) を満たすと主張する．実際，(1)，(3) は直ちに従い，
(c) から (4) が従い，(d) から (6) が従う．(c) $\Rightarrow$ (4) および
(d) $\Rightarrow$ (6) の証明は省略する．

\medskip\noindent
定理 \ref{algebraization-theorem-final-bootstrap} により，$\{H^1_{T'}(M'/I^nM')\}$ は本質的に定値で，
その値は $H^1_{T'}(M')$ である．したがって $\{H^1_{T'}(M'_n)\}$ も本質的に定値で，
その値は $H^1_{T'}(M')$ である．よって，上に表示した短完全列の極限は短完全列
$$
0 \to M' \to \lim M''_n \to H^1_{T'}(M') \to 0
$$
である．$M'' = \lim M''_n$ とおく．Local Cohomology，命題
\ref{local-cohomology-proposition-finiteness} により，$H^1_{T'}(M')$，したがって
$M''$ は有限 $A$-加群である（条件を検証する際には，$M'$ が少なくとも $1$ の深さを
$T - T'$ の素イデアルでもつことを用いる）．$(\mathcal{H}_n)$ を
$\textit{Coh}(X, \mathcal{I})$ の対象で，有限 $A$-加群 $M''$ に対応するものとする．
証明を終えるには，標準的な写像 $\{M''/I^nM''\} \to \{M''_n\}$ が pro-同型であることを
示さなければならない．$\{M'/I^nM'\}$ が $\{M'_n\}$ と pro-同型であることは既に
分かっているので，これは
$\{H^1_{T'}(M')/I^nH^1_{T'}(M')\} \to \{H^1_{T'}(M'_n)\}$ が pro-同型であることと
同値である（検証は読者に委ね，省略する）．
$H^1_{T'}(M')/I^nH^1_{T'}(M') = H^1_{T'}(M')$ が十分大きい $n$ に対して成り立つので，
これは上で見たように $\{H^1_{T'}(M'_n)\}$ が本質的に定値で，その値が
$H^1_{T'}(M')$ であることから従う．
\end{proof}

\begin{lemma}
\label{algebraization-lemma-improvement-application}
状況 \ref{algebraization-situation-algebraize} において，$A$ は双対化複体をもつと仮定する．
$d \geq \text{cd}(A, I)$ とする．$(\mathcal{F}_n)$ を
$\textit{Coh}(U, I\mathcal{O}_U)$ の対象とする．$(\mathcal{F}_n)$ は
$(2, 2 + d)$-不等式を満たすと仮定する．定義
\ref{algebraization-definition-s-d-inequalities} を参照されたい．このとき標準的な写像
$(\mathcal{F}_n) \to (\mathcal{F}_n'')$ で，連接 $\mathcal{O}_U$-加群の逆系の間のものが
存在し，次の性質をもつ．
\begin{enumerate}
\item $\text{depth}(\mathcal{F}''_{n, y}) + \delta^Y_Z(y) \geq 3$
がすべての $y \in U \cap Y$ に対して成り立つ；
\item $(\mathcal{F}''_n)$ は pro-系としてある対象 $(\mathcal{H}_n)$ と同型であり，
後者は $\textit{Coh}(U, I\mathcal{O}_U)$ の対象である；
\item 誘導される射 $(\mathcal{F}_n) \to (\mathcal{H}_n)$ は
$\textit{Coh}(U, I\mathcal{O}_U)$ の射として，$I$ のある冪で消される核と余核をもつ；
\item 加群 $H^0(U, \mathcal{F}''_n)$ および $H^1(U, \mathcal{F}''_n)$ は有限
$A$-加群であり，これはすべての $n$ に対して成り立つ．
\end{enumerate}
\end{lemma}

\begin{proof}
存在と性質 (1)，(2)，(3) は，補題
\ref{algebraization-lemma-improvement-formal-coherent-module-better} を
$U$，$U \cap Y$，$T = \{y \in U \cap Y : \delta^Y_Z(y) \leq 2\}$，
$T' = \{y \in U \cap Y : \delta^Y_Z(y) \leq 1\}$，$(\mathcal{F}_n)$ に適用して従う．
加群 $H^0(U, \mathcal{F}''_n)$ および $H^1(U, \mathcal{F}''_n)$ の有限性は，
Local Cohomology，補題 \ref{local-cohomology-lemma-finiteness-Rjstar} と，補題
\ref{algebraization-lemma-discussion} で証明した関数 $\delta^Y_Z(-)$ の基本的性質から従う．
\end{proof}











\section{連接形式加群の代数化 V}
\label{algebraization-section-algebraization-modules-conclusion}

\noindent
本節では，連接形式加群の代数化に関する最も一般的な結果を証明する．まず，
イデアルのコホモロジー次元が $1$ である場合にこれを証明する．次に，この結果を
爆発に適用して，より一般的な結果を証明する．

\begin{lemma}
\label{algebraization-lemma-cd-1-canonical}
状況 \ref{algebraization-situation-algebraize} において，$(\mathcal{F}_n)$ を
$\textit{Coh}(U, I\mathcal{O}_U)$ の対象とする．次を仮定する．
\begin{enumerate}
\item $A$ は双対化複体をもち，$\text{cd}(A, I) = 1$ である；
\item $(\mathcal{F}_n)$ は，逆系 $(\mathcal{F}_n'')$ であって，連接
$\mathcal{O}_U$-加群からなり，
$\text{depth}(\mathcal{F}''_{n, y}) + \delta^Y_Z(y) \geq 3$
がすべての $y \in U \cap Y$ に対して成り立つものと pro-同型である．
\end{enumerate}
このとき $(\mathcal{F}_n)$ は $X$ へ標準的に延長される．定義
\ref{algebraization-definition-canonically-algebraizable} を参照されたい．
\end{lemma}

\begin{proof}
補題 \ref{algebraization-lemma-when-done} の仮定 (a)，(b)，(c) を検証する．
始める前に，加群
$H^0(U, \mathcal{F}''_n)$ および $H^1(U, \mathcal{F}''_n)$
が有限 $A$-加群であり，これはすべての $n$ に対して成り立つことを指摘しておく．
Local Cohomology，補題 \ref{local-cohomology-lemma-finiteness-Rjstar} による．

\medskip\noindent
各 $p \geq 0$ に対し，$\lim H^p(U, \mathcal{F}_n)$ 上の極限位相は，
補題 \ref{algebraization-lemma-topology-I-adic} により $I$-進位相であることに注意する．
特に，仮定 (b) が成り立つ．

\medskip\noindent
$M = \lim H^0(U, \mathcal{F}_n)$ は $A$-加群で，その極限位相は
$I$-進位相であることが分かっている．したがって，$n$ を与えると，加群
$M/I^nM$ は $H^0(U, \mathcal{F}_N)$ の部分商であり，これはある $N \gg n$ に対して成り立つ．
逆系 $\{H^0(U, \mathcal{F}_N)\}$ は有限 $A$-加群の逆系，すなわち
$\{H^0(U, \mathcal{F}''_N)\}$ と pro-同型なので，$M/I^nM$ は有限である．
したがって $M$ は有限である．Algebra，補題
\ref{algebra-lemma-finite-over-complete-ring} を参照されたい．
特に，仮定 (c) が成り立つ．

\medskip\noindent
各 $n \geq 0$ に対して $Ob_n = \lim_N H^1(U, I^n\mathcal{F}_N)$ と書く．
特別な場合が $Ob = Ob_0 = \lim_N H^1(U, \mathcal{F}_N)$ である．
前段落とまったく同じ議論により，$Ob$ は有限 $A$-加群であることが分かる．
（実際，$Ob/I Ob$ が $\mathfrak a$ のある冪で消されることも分かっているが，
これを利用するのはいささか難しいようである．）

\medskip\noindent
$\mathcal{F} = \lim \mathcal{F}_n$ とおき，生成元
$f_1, \ldots, f_r$ を $I$ に対して選び，$c \geq 1$ を選び，
$\Phi_\mathcal{F}$ を補題 \ref{algebraization-lemma-cd-is-one-for-system} にあるものとして選ぶ．
以後，補題 \ref{algebraization-lemma-properties-system} の結果を断りなく用いる．
特に，各 $n \geq 1$ に対して写像
$$
\delta_n :
H^0(U, \mathcal{F}_n)
\longrightarrow
H^1(U, I^n\mathcal{F})
\longrightarrow
Ob_n
$$
がある．第1の写像は短完全列
$0 \to I^n\mathcal{F} \to \mathcal{F} \to \mathcal{F}_n \to 0$
から，第2の写像は $I^n\mathcal{F} = \lim I^n\mathcal{F}_N$ から得られる．
後で次を用いる．$\delta_n(s) = 0$ が $s \in H^0(U, \mathcal{F}_n)$ に対して
成り立つならば，各 $n' \geq n$ に対して
$s' \in H^0(U, \mathcal{F}_{n'})$ であって $s$ へ写るものを見つけることができる．
可換図式
$$
\xymatrix{
H^0(U, \mathcal{F}_{nc}) \ar[r] \ar[dd] &
H^1(U, I^{nc}\mathcal{F}) \ar[dd] \ar[rd]^{\Phi_\mathcal{F}} \\
& &
\bigoplus_{e_1 + \ldots + e_r = n}
H^1(U, \mathcal{F}) \cdot T_1^{e_1} \ldots T_r^{e_r} \ar[ld] \\
H^0(U, \mathcal{F}_n) \ar[r] &
H^1(U, I^n\mathcal{F})
}
$$
があることに注意する．障害写像
$H^0(U, \mathcal{F}_n) \to Ob_n$
は，写像
$H^0(U, \mathcal{F}_{nc}) \to H^0(U, \mathcal{F}_n)$
の像を部分加群
$$
Ob'_n =
\Im\left(
\bigoplus\nolimits_{e_1 + \ldots + e_r = n}
Ob \cdot T_1^{e_1} \ldots T_r^{e_r} \to Ob_n
\right)
$$
へ送ると結論できる．ここで直和因子
$Ob \cdot T_1^{e_1} \ldots T_r^{e_r}$ 上では，$I$ の冪を法とする簡約によって，乗法写像
$f_1^{e_1} \ldots f_r^{e_r} : \mathcal{F} \to I^n\mathcal{F}$ から得られる
コホモロジー上の写像を用いる．構成により
$$
\bigoplus\nolimits_{n \geq 0} Ob'_n
$$
は Rees 代数 $\bigoplus_{n \geq 0} I^n$ 上の有限次数付き加群である．
各 $n$ に対して
$$
M_n = \{s \in H^0(U, \mathcal{F}_n) \mid \delta_n(s) \in Ob'_n\}
$$
とおく．$\{M_n\}$ は逆系であり，大域切断上の
$f_j : \mathcal{F}_n \to \mathcal{F}_{n + 1}$ は $M_n$ を $M_{n + 1}$ へ写すことに
注意する．Cohomology，補題 \ref{cohomology-lemma-ML-general} の証明とまったく
同じ議論により，$\{M_n\}$ は ML であることが分かる．実際，Rees 代数は
Noether なので，形 $\delta_{n_j}(z_j)$ の斉次生成元を，$z_j \in M_{n_j}$ として，
次数付き部分加群 $\bigoplus_{n \geq 0} \Im(M_n \to Ob'_n)$ に対して有限個選ぶことができる．
そこで $k = \max(n_j)$ とおくと，$n \geq k$ および任意の $z \in M_n$ に対し，
$a_j \in I^{n - n_j}$ であって $z - \sum a_j z_j$ が $\delta_n$ の核に入るものを
見つけることができ，したがってこの元は $M_{n'}$ の像に，すべての
$n' \geq n$ に対して入る（実際，$\delta_n$ が消えることから，
$z - \sum a_j z_j$ を元 $z' \in H^0(U, \mathcal{F}_{n'c})$ へ，すべての
$n' \ge n$ に対して持ち上げることができ，上で証明したことにより，$z'$ の
$H^0(U, \mathcal{F}_{n'})$ における像は $M_{n'}$ に属する）．
したがって $\Im(M_n \to M_{n - k}) = \Im(M_{n'} \to M_{n - k})$ であり，
これはすべての $n' \geq n$ に対して成り立つ．

\medskip\noindent
$n$ を選ぶ．いま証明した $\{M_n\}$ の Mittag--Leffler 性により，
$n' \geq n$ であって，$M_{n'} \to M_n$ の像が $M' \to M_n$ の像と
一致するものを見つけることができる．上の議論から，$M' \to M_n$ の像は
$H^0(U, \mathcal{F}_{n'c}) \to H^0(U, \mathcal{F}_n)$ の像を含むことが分かる．
したがって $\{M_n\}$ と $\{H^0(U, \mathcal{F}_n)\}$ は pro-同型である．
それゆえ $\{H^0(U, \mathcal{F}_n)\}$ は ML であり，最後に仮定 (a) が成り立つと
結論できる．これで証明が完了する．
\end{proof}

\begin{proposition}[コホモロジー次元 1 における代数化]
\label{algebraization-proposition-cd-1}
\begin{reference}
この結果の局所的な場合は \cite[IV Corollaire 2.9]{MRaynaud-book} である．
\end{reference}
状況 \ref{algebraization-situation-algebraize} において，$(\mathcal{F}_n)$ を
$\textit{Coh}(U, I\mathcal{O}_U)$ の対象とする．次を仮定する．
\begin{enumerate}
\item $A$ は双対化複体をもち，$\text{cd}(A, I) = 1$ である；
\item $(\mathcal{F}_n)$ は $(2, 3)$-不等式を満たす．定義
\ref{algebraization-definition-s-d-inequalities} を参照されたい．
\end{enumerate}
このとき $(\mathcal{F}_n)$ は $X$ へ延長される．特に，$A$ が
$I$-進完備ならば，$(\mathcal{F}_n)$ はある連接 $\mathcal{O}_U$-加群の完備化である．
\end{proposition}

\begin{proof}
補題 \ref{algebraization-lemma-map-kernel-cokernel-on-closed} により，$(\mathcal{F}_n)$ を，
補題 \ref{algebraization-lemma-improvement-application} で得られた対象 $(\mathcal{H}_n)$ であって
$\textit{Coh}(U, I\mathcal{O}_U)$ に属するものに置き換えてよい．
したがって，$(\mathcal{F}_n)$ は，補題 \ref{algebraization-lemma-improvement-application} に述べた
性質をもつ逆系 $(\mathcal{F}_n'')$ と pro-同型であると仮定してよい．
補題 \ref{algebraization-lemma-cd-1-canonical} で，$(\mathcal{F}_n)$ が $X$ へ標準的に延長されることを
証明した．最後の主張は補題 \ref{algebraization-lemma-canonically-algebraizable} から従う．
\end{proof}

\begin{lemma}
\label{algebraization-lemma-blowup}
状況 \ref{algebraization-situation-algebraize} において，$(\mathcal{F}_n)$ を
$\textit{Coh}(U, I\mathcal{O}_U)$ の対象とする．次を仮定する．
\begin{enumerate}
\item $A$ は双対化複体をもつ；
\item 爆発 $b : X' \to X$ であって $I$ のもののすべてのファイバーは次元 $\leq d - 1$ をもつ；
\item 次のいずれかが成り立つ：
\begin{enumerate}
\item $(\mathcal{F}_n)$ は $(d + 1, d + 2)$-不等式
（定義 \ref{algebraization-definition-s-d-inequalities}）を満たす；または
\item $y \in U \cap Y$ および素イデアル
$\mathfrak p \subset \mathcal{O}_{X, y}^\wedge$ が
$\mathfrak p \not \in V(I\mathcal{O}_{X, y}^\wedge)$ を満たすならば，
$$
\text{depth}((\mathcal{F}^\wedge_y)_\mathfrak p) +
\dim(\mathcal{O}_{X, y}^\wedge/\mathfrak p) + \delta^Y_Z(y) > d + 2
$$
\end{enumerate}
\end{enumerate}
が成り立つ．このとき $(\mathcal{F}_n)$ は $X$ へ延長される．
\end{lemma}

\begin{proof}
$Y' \subset X'$ を例外因子とする．$Z' \subset Y'$ を $Z \subset Y$ の逆像とする．
すると $U' = X' \setminus Z'$ は $U$ の逆像である．
(\ref{algebraization-equation-delta-Z}) にある $\delta^{Y'}_{Z'}$ を用いて
$$
T' = \{y' \in Y' \mid \delta^{Y'}_{Z'}(y') = 1\}
\subset
T = \{y' \in Y' \mid \delta^{Y'}_{Z'}(y') = 1\text{ または }2\}
$$
とおく．補題 \ref{algebraization-lemma-discussion} の (1)，(2) により，これらは
$U' \cap Y' = Y' \setminus Z'$ の特殊化で安定な部分集合である．
対象 $(b|_{U'}^*\mathcal{F}_n)$ であって
$\textit{Coh}(U', I\mathcal{O}_{U'})$ に属するものを考える．
Cohomology of Schemes，補題 \ref{coherent-lemma-inverse-systems-pullback} を
参照されたい．$y' \in U' \cap Y'$ に対して
$$
\mathcal{F}_{y'}^\wedge = \lim (b|_{U'}^*\mathcal{F}_n)_{y'}
$$
を，この引き戻しの $y'$ における「茎」と表す．補題
\ref{algebraization-lemma-improvement-formal-coherent-module-better} の条件
(a)，(b)，(c)，(d)，(e) が，対象 $(b|_{U'}^*\mathcal{F}_n)$ に対し，$U'$ 上で，
$d$ を $1$ で置き換え，部分集合を $T' \subset T \subset U' \cap Y'$ とすれば
成り立つと主張する．$Y'$ は有効 Cartier 因子であり，したがって局所的に $1$ 本の
方程式で切り出されるので，条件 (a) は成り立つ．条件 (e) は補題
\ref{algebraization-lemma-discussion} の (7) により成り立つ．(b)，(c)，(d) を証明するには
少し準備が必要である．

\medskip\noindent
$y' \in U' \cap Y'$ とし，$\mathfrak p' \subset \mathcal{O}_{X', y'}^\wedge$ を
$V(I\mathcal{O}_{X', y'}^\wedge)$ に含まれない素イデアルとする．
$y = b(y') \in U \cap Y$ と表す．$f \in I$ を，$y'$ がアフィン爆発代数
$A[\frac{I}{f}]$ のスペクトルに含まれるように選ぶ．Divisors，補題
\ref{divisors-lemma-blowing-up-affine} を参照されたい．任意の $A$-代数 $B$ に対し，
$B' = B[\frac{IB}{f}]$ を対応するアフィン爆発代数と表す．$I$-進完備化を
${\ }^\wedge$ で表す．$f$ の選び方により，環準同型
$(\mathcal{O}_{X, y}^\wedge)' \to \mathcal{O}_{X', y'}^\wedge$.
を得る．$\mathfrak q' \subset (\mathcal{O}_{X, y}^\wedge)'$ を
$\mathfrak m_{y'}^\wedge$ の逆像とすると，
$((\mathcal{O}_{X, y}^\wedge)'_{\mathfrak q'})^\wedge =
\mathcal{O}_{X', y'}^\wedge$.
であることが分かる．$\mathfrak p \subset \mathcal{O}_{X, y}^\wedge$ を対応する
素イデアルとする．この時点で可換図式
$$
\xymatrix{
\mathcal{O}_{X, y}^\wedge \ar[d] \ar[r] &
(\mathcal{O}_{X, y}^\wedge)' \ar[d]_\alpha \ar[r] &
(\mathcal{O}_{X, y}^\wedge)'_{\mathfrak q'} \ar[d] \ar[r]_\beta &
\mathcal{O}_{X', y'}^\wedge \ar[d] \\
\mathcal{O}_{X, y}^\wedge/\mathfrak p \ar[r] &
(\mathcal{O}_{X, y}^\wedge/\mathfrak p)' \ar[r] &
(\mathcal{O}_{X, y}^\wedge/\mathfrak p)'_{\mathfrak q'} \ar[r]^\gamma &
((\mathcal{O}_{X, y}^\wedge/\mathfrak p)'_{\mathfrak q'})^\wedge \ar[d] \\
 & & &
\mathcal{O}_{X', y'}^\wedge/\mathfrak p'
}
$$
を得る．その縦の矢印は全射である．More on Algebra，補題
\ref{more-algebra-lemma-completion-dimension} と次元公式
（Algebra，補題 \ref{algebra-lemma-dimension-formula}）により，
$$
\dim(((\mathcal{O}_{X, y}^\wedge/\mathfrak p)'_{\mathfrak q'})^\wedge) =
\dim((\mathcal{O}_{X, y}^\wedge/\mathfrak p)'_{\mathfrak q'}) =
\dim(\mathcal{O}_{X, y}^\wedge/\mathfrak p)
- \text{trdeg}(\kappa(y')/\kappa(y))
$$
を得る．引き戻し，茎，局所化，完備化の定義をたどると，
$$
(\mathcal{F}_y^\wedge)_{\mathfrak p}
\otimes_{(\mathcal{O}_{X, y}^\wedge)_\mathfrak p}
(\mathcal{O}_{X', y'}^\wedge)_{\mathfrak p'}
=
(\mathcal{F}_{y'}^\wedge)_{\mathfrak p'}
$$
を得る．詳細は省略する．図式中の環準同型 $\beta$ および $\gamma$ は，
Gorenstein（したがって Cohen--Macaulay）ファイバーをもつ平坦写像である．実際，
これらは双対化複体をもつ環の完備化だからである．Dualizing Complexes，補題
\ref{dualizing-lemma-formal-fibres-gorenstein} および
\ref{dualizing-lemma-dualizing-gorenstein-formal-fibres}
および More on Algebra，節 \ref{more-algebra-section-properties-formal-fibres} の議論を
参照されたい．
$(\mathcal{O}_{X, y}^\wedge)_\mathfrak p =
(\mathcal{O}_{X, y}^\wedge)'_{\tilde{\mathfrak p}}$
であることに注意する．ここで $\tilde{\mathfrak p}$ は図式中の $\alpha$ の核である．一方，
$(\mathcal{O}_{X, y}^\wedge)'_{\tilde{\mathfrak p}}
\to (\mathcal{O}_{X', y'}^\wedge)_{\mathfrak p'}$
は，上の議論により Cohen--Macaulay ファイバーをもつ平坦写像である．したがって
$(\mathcal{O}_{X, y}^\wedge)_\mathfrak p  \to
(\mathcal{O}_{X', y'}^\wedge)_{\mathfrak p'}$ は Cohen--Macaulay ファイバーをもつ平坦写像である．
Algebra，補題 \ref{algebra-lemma-apply-grothendieck-module} を用いると，
$$
\text{depth}((\mathcal{F}_{y'}^\wedge)_{\mathfrak p'}) =
\text{depth}((\mathcal{F}_y^\wedge)_{\mathfrak p}) +
\dim(F_\mathfrak r)
$$
を得る．ここで $F$ は
$(\mathcal{O}_{X, y}^\wedge/\mathfrak p)'_{\mathfrak q'}$
の一般形式ファイバーであり，$\mathfrak r$ は $\mathfrak p'$ に対応する素イデアルである．
$(\mathcal{O}_{X, y}^\wedge/\mathfrak p)'_{\mathfrak q'}$ は普遍鎖状局所整域なので，
その $I$-進完備化は，Ratliff の定理（More on Algebra，命題
\ref{more-algebra-proposition-ratliff}）により等次元かつ（普遍）鎖状である．したがって
$$
\dim(((\mathcal{O}_{X, y}^\wedge/\mathfrak p)'_{\mathfrak q'})^\wedge) =
\dim(F_\mathfrak r) + \dim(\mathcal{O}_{X', y'}^\wedge/\mathfrak p')
$$
を得る．これと補題 \ref{algebraization-lemma-change-distance-function} を組み合わせると，
\begin{equation}
\label{algebraization-equation-one}
\begin{aligned}
&
\text{depth}((\mathcal{F}_{y'}^\wedge)_{\mathfrak p'}) +
\delta^{Y'}_{Z'}(y') \\
& =
\text{depth}((\mathcal{F}_y^\wedge)_{\mathfrak p}) +
\dim(F_\mathfrak r) + \delta^{Y'}_{Z'}(y') \\
& \geq
\text{depth}((\mathcal{F}_y^\wedge)_{\mathfrak p}) + \delta^Y_Z(y) +
\dim(F_\mathfrak r) + \text{trdeg}(\kappa(y')/\kappa(y)) - (d - 1) \\
& =
\text{depth}((\mathcal{F}_y^\wedge)_{\mathfrak p}) + \delta^Y_Z(y) - (d - 1)
+ \dim(\mathcal{O}_{X, y}^\wedge/\mathfrak p) -
\dim(\mathcal{O}_{X', y'}^\wedge/\mathfrak p')
\end{aligned}
\end{equation}
ここで
$\dim(\mathcal{O}_{X, y}^\wedge/\mathfrak p) \geq
\dim(\mathcal{O}_{X', y'}^\wedge/\mathfrak p')$ であることに注意されたい．
これを書き換えると，
\begin{equation}
\label{algebraization-equation-two}
\begin{aligned}
&
\text{depth}((\mathcal{F}_{y'}^\wedge)_{\mathfrak p'}) +
\dim(\mathcal{O}_{X', y'}^\wedge/\mathfrak p') +
\delta^{Y'}_{Z'}(y') \\
& \geq
\text{depth}((\mathcal{F}_y^\wedge)_{\mathfrak p}) +
\dim(\mathcal{O}_{X, y}^\wedge/\mathfrak p) +
\delta^Y_Z(y) - (d - 1)
\end{aligned}
\end{equation}
を得る．この不等式により，残りの条件を検証できる．

\medskip\noindent
補題 \ref{algebraization-lemma-improvement-formal-coherent-module-better} の条件 (b) および (d) を考える．
次を仮定する．
$V(\mathfrak p') \cap V(I\mathcal{O}_{X', y'}^\wedge) =
\{\mathfrak m_{y'}^\wedge\}$.
これは，$\dim(\mathcal{O}_{X', y'}^\wedge/\mathfrak p') = 1$ であることを意味する．
実際，$Z'$ は有効 Cartier 因子だからである．
(b) と (d) を合わせた条件は
$$
\text{depth}((\mathcal{F}_{y'}^\wedge)_{\mathfrak p'}) + \delta^{Y'}_{Z'}(y')
> 2
$$
と同値である．$(\mathcal{F}_n)$ が (3)(b) の不等式を満たすならば，
(\ref{algebraization-equation-two}) を適用して，これが成り立つと直ちに結論できる．
$(\mathcal{F}_n)$ が (3)(a)，すなわち $(d + 1, d + 2)$-不等式を満たすならば，
いずれにせよ
$$
\text{depth}((\mathcal{F}_y^\wedge)_{\mathfrak p}) + \delta^Y_Z(y)
\geq d + 1
\quad\text{または}\quad
\text{depth}((\mathcal{F}_y^\wedge)_{\mathfrak p}) +
\dim(\mathcal{O}_{X, y}^\wedge/\mathfrak p) +
\delta^Y_Z(y) > d + 2
$$
であることが分かる．上の (\ref{algebraization-equation-one}) および (\ref{algebraization-equation-two}) を見ると，
これは，場合によっては
$\dim(\mathcal{O}_{X, y}^\wedge/\mathfrak p) = 1$.
のときを除いて，求める主張を与える．しかし，
$\dim(\mathcal{O}_{X, y}^\wedge/\mathfrak p) = 1$ ならば，
$V(\mathfrak p) \cap V(I\mathcal{O}_{X, y}^\wedge) = \{\mathfrak m_y^\wedge\}$
であり，実際
$$
\text{depth}((\mathcal{F}_y^\wedge)_{\mathfrak p}) + \delta^Y_Z(y) > d + 1
$$
であることが分かる．実際，$(\mathcal{F}_n)$ は $(d + 1, d + 2)$-不等式を
満たすからである．よって再び結論を得る．

\medskip\noindent
補題 \ref{algebraization-lemma-improvement-formal-coherent-module-better} の条件 (c) を考える．
次を仮定する．
$V(\mathfrak p') \cap V(I\mathcal{O}_{X', y'}^\wedge) \not =
\{\mathfrak m_{y'}^\wedge\}$. このとき条件 (c) は
$$
\text{depth}((\mathcal{F}_{y'}^\wedge)_{\mathfrak p'}) + \delta^{Y'}_{Z'}(y')
\geq 2
\quad\text{または}\quad
\text{depth}((\mathcal{F}_{y'}^\wedge)_{\mathfrak p'}) +
\dim(\mathcal{O}_{X', y'}^\wedge/\mathfrak p') +
\delta^{Y'}_{Z'}(y') > 3
$$
と同値である．$(\mathcal{F}_n)$ が (3)(b) の不等式を満たすならば，
(\ref{algebraization-equation-two}) を適用して，表示された二つの不等式のうち第2のものが
成り立つことが分かる．$(\mathcal{F}_n)$ が (3)(a)，すなわち
$(d + 1, d + 2)$-不等式を満たすならば，これは
(\ref{algebraization-equation-one}) および (\ref{algebraization-equation-two}) から直ちに従う．
これで主張の証明が完了する．

\medskip\noindent
$(b|_{U'}^*\mathcal{F}_n) \to (\mathcal{F}_n'')$ と
$(\mathcal{H}_n)$ を $\textit{Coh}(U', I\mathcal{O}_{U'})$ において，
補題 \ref{algebraization-lemma-improvement-formal-coherent-module-better} にあるように選ぶ．
任意のアフィン開集合 $W \subset X'$ に対し，
$\delta^{W \cap Y'}_{W \cap Z'}(y') \geq \delta^{Y'}_{Z'}(y')$ であることに注意する．これは
補題 \ref{algebraization-lemma-discussion} の (6) であることに注意する．したがって
$(\mathcal{H}_n|_W)$ は
補題 \ref{algebraization-lemma-cd-1-canonical} の仮定を満たす．それゆえ
$(\mathcal{H}_n|_W)$ は $W$ へ標準的に延長される．
$(\mathcal{G}_{W, n})$ を，$\textit{Coh}(W, I\mathcal{O}_W)$ に属する標準的延長で，
補題 \ref{algebraization-lemma-canonically-algebraizable} にあるものとする．
補題 \ref{algebraization-lemma-canonically-extend-base-change} により，$W' \subset W$ に対して
一意な同型
$$
(\mathcal{G}_{W, n}|_{W'}) \longrightarrow
(\mathcal{G}_{W', n})
$$
があり，これは与えられた同型
$(\mathcal{G}_{W, n}|_{W \cap U}) \cong (\mathcal{H}_n|_{W \cap U})$.
と両立する．したがって，対象 $(\mathcal{G}_n)$ であって
$\textit{Coh}(X', I\mathcal{O}_{X'})$ に属し，その $U$ への制限が $(\mathcal{H}_n)$ と
同型であるものが存在すると結論する．

\medskip\noindent
$A$ が $I$-進完備ならば，次のように証明を終えることができる．
Grothendieck の存在定理（Cohomology of Schemes，補題
\ref{coherent-lemma-existence-projective}）により，$(\mathcal{G}_n)$ は
ある連接 $\mathcal{O}_{X'}$-加群の完備化である．次に，Cohomology of Schemes，補題
\ref{coherent-lemma-existence-easy} により，$(b|_{U'}^*\mathcal{F}_n)$ は
ある連接 $\mathcal{O}_{U'}$-加群 $\mathcal{F}'$ の完備化であることが分かる．
Cohomology of Schemes，補題
\ref{coherent-lemma-inverse-systems-push-pull} により，写像
$$
(\mathcal{F}_n) \longrightarrow ((b|_{U'})_*\mathcal{F}')^\wedge
$$
があり，その核と余核は $I$ のある冪で消される．最後に，補題
\ref{algebraization-lemma-map-kernel-cokernel-on-closed} を適用すればよい．

\medskip\noindent
$A$ が完備でないならば，証明を始める前に $A$ をその完備化で置き換えてよい．
補題 \ref{algebraization-lemma-algebraizable} を参照されたい．完備化後にも仮定は成り立つ．
条件 (3) については明らかであり，条件 (1) については Dualizing Complexes，補題
\ref{dualizing-lemma-ubiquity-dualizing} から，条件 (2) については Divisors，補題
\ref{divisors-lemma-flat-base-change-blowing-up} から従う．
したがって，完備な場合から一般の場合が従う．
\end{proof}

\begin{proposition}[生成元の少ないイデアルに対する代数化]
\label{algebraization-proposition-d-generators}
状況 \ref{algebraization-situation-algebraize} において，$(\mathcal{F}_n)$ を
$\textit{Coh}(U, I\mathcal{O}_U)$ の対象とする．次を仮定する．
\begin{enumerate}
\item $A$ は双対化複体をもつ；
\item $V(I) = V(f_1, \ldots, f_d)$ が，ある $d \geq 1$ および
$f_1, \ldots, f_d \in A$ に対して成り立つ；
\item 次のいずれかが成り立つ：
\begin{enumerate}
\item $(\mathcal{F}_n)$ は $(d + 1, d + 2)$-不等式
（定義 \ref{algebraization-definition-s-d-inequalities}）を満たす；または
\item $y \in U \cap Y$ および素イデアル
$\mathfrak p \subset \mathcal{O}_{X, y}^\wedge$ が
$\mathfrak p \not \in V(I\mathcal{O}_{X, y}^\wedge)$ を満たすならば，
$$
\text{depth}((\mathcal{F}^\wedge_y)_\mathfrak p) +
\dim(\mathcal{O}_{X, y}^\wedge/\mathfrak p) + \delta^Y_Z(y) > d + 2
$$
\end{enumerate}
\end{enumerate}
が成り立つ．このとき $(\mathcal{F}_n)$ は $X$ へ延長される．特に，$A$ が
$I$-進完備ならば，$(\mathcal{F}_n)$ はある連接 $\mathcal{O}_U$-加群の完備化である．
\end{proposition}

\begin{proof}
$I = (f_1, \ldots, f_d)$ と仮定してよい．Cohomology of Schemes，補題
\ref{coherent-lemma-inverse-systems-ideals-equivalence} を参照されたい．すると，
$X$ の $I$ における爆発のすべてのファイバーは次元が高々 $d - 1$ であることが分かる．
したがって，補題 \ref{algebraization-lemma-blowup} から延長を得る．最後の主張は補題
\ref{algebraization-lemma-essential-image-completion} から従う．
\end{proof}

\noindent
次の補題を，注意 \ref{algebraization-remark-interesting-case-variant}，
\ref{algebraization-remark-interesting-case},
\ref{algebraization-remark-interesting-case-bis}，および
\ref{algebraization-remark-interesting-case-ter} と比較されたい．

\begin{lemma}
\label{algebraization-lemma-interesting-case-final}
状況 \ref{algebraization-situation-algebraize} において，$(\mathcal{F}_n)$ を
$\textit{Coh}(U, I\mathcal{O}_U)$ の対象とする．次を仮定する．
\begin{enumerate}
\item $A$ は双対化複体をもつ局所環である；
\item $X$ のすべての既約成分は同じ次元をもつ；
\item スキーム $X \setminus Y$ は Cohen--Macaulay である；
\item $I$ は $d$ 個の元で生成される；
\item $\dim(X) - \dim(Z) > d + 2$ である；
\item $y \in U \cap Y$ に対し，加群 $\mathcal{F}_y^\wedge$ は
$V(I\mathcal{O}_{X, y}^\wedge)$ の外で有限局所自由である．例えば，
$\mathcal{F}_n$ が有限局所自由な $\mathcal{O}_U/I^n\mathcal{O}_U$-加群ならばそうである．
\end{enumerate}
このとき $(\mathcal{F}_n)$ は $X$ へ延長される．特に $A$ が $I$-進完備ならば，
$(\mathcal{F}_n)$ はある連接 $\mathcal{O}_U$-加群の完備化である．
\end{lemma}

\begin{proof}
命題 \ref{algebraization-proposition-d-generators} の仮定 (1)，(2)，(3)(b) が満たされることを示す．
(1) と (2) については明らかである．

\medskip\noindent
$y \in U \cap Y$ とし，$\mathfrak p$ を素イデアル
$\mathfrak p \subset \mathcal{O}_{X, y}^\wedge$ であって，
$\mathfrak p \not \in V(I\mathcal{O}_{X, y}^\wedge)$ を満たすものとする．
最後の条件から
$\text{depth}((\mathcal{F}_y^\wedge)_\mathfrak p) =
\text{depth}((\mathcal{O}_{X, y}^\wedge)_\mathfrak p)$
である．
である．$X \setminus Y$ は Cohen--Macaulay なので，
$(\mathcal{O}_{X, y}^\wedge)_\mathfrak p$ は Cohen--Macaulay である．したがって
\begin{align*}
& \text{depth}((\mathcal{F}^\wedge_y)_\mathfrak p) +
\dim(\mathcal{O}_{X, y}^\wedge/\mathfrak p) + \delta^Y_Z(y) \\
& =
\dim((\mathcal{O}_{X, y}^\wedge)_\mathfrak p) +
\dim(\mathcal{O}_{X, y}^\wedge/\mathfrak p) + \delta^Y_Z(y) \\
& =
\dim(\mathcal{O}_{X, y}^\wedge) + \delta^Y_Z(y)
\end{align*}
を得る．最後の等式は，第2の条件により $\mathcal{O}_{X, y}$ が等次元だからである．
$\delta(y) = \dim(\overline{\{y\}})$ とおく．$A$ は鎖状局所環なので，これは次元関数である．
補題 \ref{algebraization-lemma-discussion} により
$\delta^Y_Z(y) \geq \delta(y) - \dim(Z)$ である．$X$ は等次元なので，
$$
\dim(\mathcal{O}_{X, y}^\wedge) + \delta^Y_Z(y)
\geq \dim(\mathcal{O}_{X, y}^\wedge) + \delta(y) - \dim(Z)
= \dim(X) - \dim(Z)
$$
を得る．したがって，求める不等式が得られ，証明が完了する．
\end{proof}

\begin{remark}
\label{algebraization-remark-question}
命題 \ref{algebraization-proposition-d-generators} において，$I$ が $d$ 個の生成元をもつという仮定を
$\text{cd}(A, I) \leq d$ という仮定で置き換えた類似の主張を，われわれは証明も
反証もできていない．証明または反例をご存じならば，
\href{mailto:stacks.project@gmail.com}{stacks.project@gmail.com}.
宛てにメールをいただきたい．もう一つの明らかな問題は，命題
\ref{algebraization-proposition-d-generators} の条件がどの程度必要であるか，というものである．
\end{remark}




\section{連接形式加群の代数化 VI}
\label{algebraization-section-algebraization-modules-yet-more}

\noindent
本節では，さらにいくつかの証明しやすい場合を加える．

\begin{proposition}
\label{algebraization-proposition-algebraization-regular-sequence}
状況 \ref{algebraization-situation-algebraize} において，$(\mathcal{F}_n)$ を
$\textit{Coh}(U, I\mathcal{O}_U)$ の対象とする．次を仮定する．
\begin{enumerate}
\item $f_1, \ldots, f_d \in I$ が存在し，$y \in U \cap Y$ に対して
イデアル $I\mathcal{O}_{X, y}$ は $f_1, \ldots, f_d$ で生成され，
$f_1, \ldots, f_d$ は $\mathcal{F}_y^\wedge$-正則列をなす；
\item $H^0(U, \mathcal{F}_1)$ および $H^1(U, \mathcal{F}_1)$ は有限
$A$-加群である．
\end{enumerate}
このとき $(\mathcal{F}_n)$ は $X$ へ標準的に延長される．特に，$A$ が
完備ならば，$(\mathcal{F}_n)$ はある連接 $\mathcal{O}_U$-加群の完備化である．
\end{proposition}

\begin{proof}
補題 \ref{algebraization-lemma-when-done} の仮定 (a)，(b)，(c) を検証して証明する．
すべての $n$ に対して短完全列
$$
0 \to I^n\mathcal{F}_{n + 1} \to \mathcal{F}_{n + 1} \to \mathcal{F}_n \to 0
$$
がある．$f_1, \ldots, f_d$ は正則列（したがって準正則列；Algebra，補題
\ref{algebra-lemma-regular-quasi-regular} を参照）を，各「茎」
$\mathcal{F}_y^\wedge$ 上でなし，また
$I\mathcal{F}_n = (f_1, \ldots, f_d)\mathcal{F}_n$ がすべての $n$ に対して
成り立つので，
$$
I^n\mathcal{F}_{n + 1} =
\bigoplus\nolimits_{e_1 + \ldots + e_d = n} \mathcal{F}_1 \cdot
f_1^{e_1} \ldots f_d^{e_d}
$$
を得る．これは茎上で検証できる．$H^0(U, \mathcal{F}_1)$ の有限性の仮定と
帰納法を用いると，まず $M_n = H^0(U, \mathcal{F}_n)$ は有限 $A$-加群であり，
これはすべての $n$ に対して成り立つと結論できる．このようにして，補題
\ref{algebraization-lemma-when-done} の条件 (c) が成り立つことが分かる．さらに
$$
\bigoplus\nolimits_{n \geq 0} H^1(U, I^n\mathcal{F}_{n + 1})
$$
は有限次数付き $R = \bigoplus I^n/I^{n +1}$-加群であることも分かる．
Cohomology，補題 \ref{cohomology-lemma-ML-general} により，補題
\ref{algebraization-lemma-when-done} の条件 (a) が満たされると結論する．最後に，補題
\ref{algebraization-lemma-when-done} の条件 (b) も満たされる．実際，
$\bigoplus H^0(U, I^n\mathcal{F}_{n + 1})$ は有限次数付き $R$-加群であり，
Cohomology，補題 \ref{cohomology-lemma-topology-I-adic-general} を適用できるからである．
\end{proof}

\begin{remark}
\label{algebraization-remark-interesting-case-ter}
命題 \ref{algebraization-proposition-algebraization-regular-sequence} の状況において，
$A$ が双対化複体をもつとさらに仮定すれば，
$H^0(U, \mathcal{F}_1)$ および $H^1(U, \mathcal{F}_1)$ が有限であるという条件は，
$$
\text{depth}(\mathcal{F}_{1, y}) +
\dim(\mathcal{O}_{\overline{\{y\}}, z}) > 2
$$
がすべての $y \in U \cap Y$ および $z \in Z \cap \overline{\{y\}}$ に対して
成り立つことと同値である．Local Cohomology，補題
\ref{local-cohomology-lemma-finiteness-Rjstar} を参照されたい．例えば，
$\mathcal{F}_1$ が有限局所自由な $\mathcal{O}_{U \cap Y}$-加群であり，
$Y$ が $(S_2)$ であり，$\text{codim}(Z', Y') \geq 3$ が，既約成分
$Y'$（$Y$ のもの）と $Z'$（$Z$ のもの）の各対で $Z' \subset Y'$ を
満たすものに対して成り立つならば，これは成り立つ．
\end{remark}

\begin{proposition}
\label{algebraization-proposition-algebraization-flat}
状況 \ref{algebraization-situation-algebraize} において，$(\mathcal{F}_n)$ を
$\textit{Coh}(U, I\mathcal{O}_U)$ の対象とする．Noether 局所環
$(R, \mathfrak m)$ と環準同型 $R \to A$ で，次を満たすものが存在すると仮定する．
\begin{enumerate}
\item $I = \mathfrak m A$ である；
\item $y \in U \cap Y$ に対して茎 $\mathcal{F}_y^\wedge$ は $R$-平坦である；
\item $H^0(U, \mathcal{F}_1)$ および $H^1(U, \mathcal{F}_1)$ は有限
$A$-加群である．
\end{enumerate}
このとき $(\mathcal{F}_n)$ は $X$ へ標準的に延長される．特に，$A$ が
完備ならば，$(\mathcal{F}_n)$ はある連接 $\mathcal{O}_U$-加群の完備化である．
\end{proposition}

\begin{proof}
証明は命題 \ref{algebraization-proposition-algebraization-regular-sequence} の証明とまったく同じである．
実際，$\kappa = R/\mathfrak m$ とすれば，$n \geq 0$ に対して同型
$$
I^n \mathcal{F}_{n + 1} \cong
\mathcal{F}_1 \otimes_\kappa \mathfrak m^n/\mathfrak m^{n + 1}
$$
があり，右辺は $\mathcal{F}_1$ のコピーの有限直和である．これは茎を見れば
検証できる．残りはすべてまったく同じである．
\end{proof}

\begin{remark}
\label{algebraization-remark-interesting-case-quater}
命題 \ref{algebraization-proposition-algebraization-flat} は
\cite[定理 2.10 (i)]{Baranovsky} の局所版である．局所的な結果から大域的な結果を
導くのは直接的である；その議論を概説しよう．$(R, \mathfrak m)$ を完備 Noether 局所環，
$X \to \Spec(R)$ を固有射とする．$n \geq 1$ に対して
$X_n = X \times_{\Spec(R)} \Spec(R/\mathfrak m^n)$ とおく．
$Z \subset X_1$ を特殊ファイバーの閉部分集合とする．$U = X \setminus Z$ とおき，
包含射を $j : U \to X$ と表す．対象
$$
(\mathcal{F}_n) \text{ は } \textit{Coh}(U, \mathfrak m\mathcal{O}_U)
$$
が与えられ，これが $R$ 上平坦である，すなわち $\mathcal{F}_n$ が
$R/\mathfrak m^n$ 上平坦であり，これはすべての $n$ に対して成り立つと仮定する．
$j_*\mathcal{F}_1$ および $R^1j_*\mathcal{F}_1$ は連接加群であると仮定する．
すると，$X$ 上アフィン局所的に，$(\mathcal{F}_n)$ の標準的延長を，命題
\ref{algebraization-proposition-algebraization-flat} により得る．
さらに，この延長の形成は局所化と可換である（補題
\ref{algebraization-lemma-algebraization-principal-variant} による）．したがって，標準的な大域対象
$(\mathcal{G}_n)$ であって $\textit{Coh}(X, \mathfrak m\mathcal{O}_X)$ に属し，
その $U$ への制限が $(\mathcal{F}_n)$ であるものを得る．Grothendieck の存在定理
（Cohomology of Schemes，命題 \ref{coherent-proposition-existence-proper}）により，
連接 $\mathcal{O}_X$-加群 $\mathcal{G}$ であって，その完備化が
$(\mathcal{G}_n)$ であるものが存在する．このようにして $(\mathcal{F}_n)$ は
代数化可能，すなわちある連接 $\mathcal{O}_U$-加群の完備化であることが分かる．

\medskip\noindent
さらに，$j_*\mathcal{F}_1$ および $R^1j_*\mathcal{F}_1$ の連接性は
特殊ファイバー上の条件である．実際，$j_1 : U_1 \to X_1$ を
$j : U \to X$ の特殊ファイバーと表せば，$\mathcal{F}_1$ を $U_1$ 上の連接層とみなすことができ，
$j_*\mathcal{F}_1 = j_{1, *}\mathcal{F}_1$ および
$R^1j_*\mathcal{F}_1 = R^1j_{1, *}\mathcal{F}_1$.
を得る．したがって，例えば $X_1$ が $(S_2)$ かつ既約であり，
$\dim(X_1) - \dim(Z) \geq 3$ であり，$\mathcal{F}_1$ が局所自由な
$\mathcal{O}_{U_1}$-加群ならば，$j_{1, *}\mathcal{F}_1$ および
$R^1j_{1, *}\mathcal{F}_1$ は連接加群である．
\end{remark}








\section{完備化関手への応用}
\label{algebraization-section-completion-application}

\noindent
本節では，参照しやすいように，既に得られたいくつかの結果を組み合わせるだけである．
ここにはほかにも多くの（より強い）結果を述べることができる．

\begin{lemma}
\label{algebraization-lemma-equivalence-better}
状況 \ref{algebraization-situation-algebraize} において，次を仮定する．
\begin{enumerate}
\item $A$ は双対化複体をもち，$I$-進完備である；
\item $I = (f)$ は一つの元で生成される；
\item $A$ は局所環で，極大イデアルは $\mathfrak a = \mathfrak m$ である；
\item 次のいずれかが成り立つ：
\begin{enumerate}
\item $A_f$ は $(S_2)$ であり，$\mathfrak p \subset A$，
$f \not \in \mathfrak p$ を満たす極小なものに対して
$\dim(A/\mathfrak p) \geq 4$ である；または
\item $\mathfrak p \not \in V(f)$ かつ
$V(\mathfrak p) \cap V(f) \not = \{\mathfrak m\}$ ならば，
$\text{depth}(A_\mathfrak p) + \dim(A/\mathfrak p) > 3$ である．
\end{enumerate}
\end{enumerate}
このとき，$U_0 = U \cap V(f)$ とすれば，完備化関手
$$
\colim_{U_0 \subset U' \subset U\text{ 開}}
\textit{Coh}(\mathcal{O}_{U'})
\longrightarrow
\textit{Coh}(U, f\mathcal{O}_U)
$$
は有限局所自由対象からなる充満部分圏上の同値である．
\end{lemma}

\begin{proof}
補題 \ref{algebraization-lemma-fully-faithful-general} から，この関手は充満忠実であることが従う
（詳細は省略する）．本質的全射性を証明しよう．$(\mathcal{F}_n)$ を有限局所自由な
$\textit{Coh}(U, f\mathcal{O}_U)$ の対象とする．補題
\ref{algebraization-lemma-algebraization-principal-bis} または命題 \ref{algebraization-proposition-cd-1} により，
連接 $\mathcal{O}_U$-加群 $\mathcal{F}$ であって，$(\mathcal{F}_n)$ が
$\mathcal{F}$ の完備化であるものが存在する．実際，いずれの結果を適用する場合も，
検証すべきことは $(\mathcal{F}_n)$ が $(2, 3)$-不等式を満たすことだけである．
これは補題 \ref{algebraization-lemma-unwinding-conditions-bis} で行われている．$y \in U_0$ ならば，
$f$-進完備化された茎 $\mathcal{F}_y$ は，$f$-進完備化された
$\mathcal{O}_{U, y}$ 上の有限自由加群と同型である．したがって
$\mathcal{F}$ は開近傍 $U'$ 上で有限局所自由であり，この近傍は $U_0$ を含む．
これで証明が完了する．
\end{proof}

\begin{lemma}
\label{algebraization-lemma-equivalence}
状況 \ref{algebraization-situation-algebraize} において，次を仮定する．
\begin{enumerate}
\item $I = (f)$ は単項である；
\item $A$ は $f$-進完備である；
\item $f$ は非零因子である；
\item $H^1_\mathfrak a(A/fA)$ および $H^2_\mathfrak a(A/fA)$ は有限
$A$-加群である．
\end{enumerate}
このとき，$U_0 = U \cap V(f)$ とすれば，完備化関手
$$
\colim_{U_0 \subset U' \subset U\text{ 開}}
\textit{Coh}(\mathcal{O}_{U'})
\longrightarrow
\textit{Coh}(U, f\mathcal{O}_U)
$$
は有限局所自由対象からなる充満部分圏上の同値である．
\end{lemma}

\begin{proof}
この関手は補題 \ref{algebraization-lemma-fully-faithful-general-alternative} により充満忠実である．
本質的全射性は補題 \ref{algebraization-lemma-algebraization-principal-variant} から従う．
\end{proof}








\section{連接三つ組}
\label{algebraization-section-coherent-triples}

\noindent
$(A, \mathfrak m)$ を Noether 局所環とする．
$f \in \mathfrak m$ を非零因子とする．
$X = \Spec(A)$，$X_0 = \Spec(A/fA)$，$U = X \setminus V(\mathfrak m)$，
$U_0 = U \cap X_0$ とおく．
$(\mathcal{F}, \mathcal{F}_0, \alpha)$ が{\it 連接三つ組}であるとは，
次が成り立つことをいう．
\begin{enumerate}
\item $\mathcal{F}$ は連接 $\mathcal{O}_U$-加群であり，
$f : \mathcal{F} \to \mathcal{F}$ は単射である；
\item $\mathcal{F}_0$ は連接 $\mathcal{O}_{X_0}$-加群である；
\item $\alpha : \mathcal{F}/f\mathcal{F} \to \mathcal{F}_0|_{U_0}$
は同型である．
\end{enumerate}
{\it 連接三つ組の射}の概念は明らかに定まり，これによってすべての連接三つ組の
集まりは圏となる．

\medskip\noindent
連接三つ組の圏は加法的であるがアーベル圏ではない．
しかし，連接三つ組の短完全列が何を意味するかは明らかである．

\medskip\noindent
二つの連接三つ組 $(\mathcal{F}, \mathcal{F}_0, \alpha)$ および
$(\mathcal{G}, \mathcal{G}_0, \beta)$ が与えられても，
$(\mathcal{F} \otimes_{\mathcal{O}_U} \mathcal{G},
\mathcal{F}_0  \otimes_{\mathcal{O}_{X_0}} \mathcal{G}_0,
\alpha \otimes \beta)$ が連接三つ組になるとは限らない\footnote{実際，
$f$ は $\mathcal{F} \otimes_{\mathcal{O}_U} \mathcal{G}$ 上で
単射であるとは限らない．}．
しかし，茎 $\mathcal{G}_x$ がすべての $x \in U_0$ に対して
自由ならば，これは成り立つ．

\medskip\noindent
連接三つ組 $(\mathcal{G}, \mathcal{G}_0, \beta)$ が{\it 局所自由}，
それぞれ{\it 可逆}であるとは，$\mathcal{G}$ と $\mathcal{G}_0$ が
局所自由加群，それぞれ可逆加群であることをいう．この場合，
$(\mathcal{G}, \mathcal{G}_0, \beta)$ とのテンソル積をとることには
意味があり（上を参照），連接三つ組の短完全列を連接三つ組の短完全列へ移す．

\begin{lemma}
\label{algebraization-lemma-prepare-chi-triple}
任意の連接三つ組 $(\mathcal{F}, \mathcal{F}_0, \alpha)$ に対して，
連接 $\mathcal{O}_X$-加群 $\mathcal{F}'$ であって，
$f : \mathcal{F}' \to \mathcal{F}'$ が単射となるもの，同型
$\alpha' : \mathcal{F}'|_U \to \mathcal{F}$，および写像
$\alpha'_0 : \mathcal{F}'/f\mathcal{F}' \to \mathcal{F}_0$
であって $\alpha \circ (\alpha' \bmod f) = \alpha'_0|_{U_0}$ を
満たすものが存在する．
\end{lemma}

\begin{proof}
有限 $A$-加群 $M$ を，$\mathcal{F}$ が $U$ への，ある連接
$\mathcal{O}_X$-加群の制限であり，その加群が $M$ に対応するように選ぶ．
局所コホモロジー，補題
\ref{local-cohomology-lemma-finiteness-pushforwards-and-H1-local}．
$\mathcal{F}$ は $f$-捩れをもたないので，$M$ をその $f$ の冪で
消える捩れ部分による商で置き換えてよい．
一方，$M_0 = \Gamma(X_0, \mathcal{F}_0)$ とおくと，
$\mathcal{F}_0$ は連接 $\mathcal{O}_{X_0}$-加群であり，有限
$A/fA$-加群 $M_0$ に対応する．
スキームのコホモロジー，補題 \ref{coherent-lemma-homs-over-open} により，
ある $n$ が存在し，同型 $\alpha_0$ は $A/fA$-加群準同型
$\mathfrak m^n M/fM \to M_0$ に対応する
（その核と余核は $\mathfrak m$ のある冪で消えるが，これは必要ない）．
したがって $M' = \mathfrak m^n M$ とし，$\mathcal{F}'$ を連接
$\mathcal{O}_X$-加群で $M'$ に対応するものとすれば，補題は明らかである．
\end{proof}

\noindent
$(\mathcal{F}, \mathcal{F}_0, \alpha)$ を連接三つ組とする．
補題 \ref{algebraization-lemma-prepare-chi-triple} のように
$\mathcal{F}', \alpha', \alpha'_0$ を選ぶ．
\begin{equation}
\label{algebraization-equation-chi-triple}
\chi(\mathcal{F}, \mathcal{F}_0, \alpha) =
\text{length}_A(\Coker(\alpha'_0)) - 
\text{length}_A(\Ker(\alpha'_0))
\end{equation}
とおく．右辺の式には意味がある．実際，$\alpha'_0$ は $U_0$ 上の
同型であり，したがってその核および連接は $\{\mathfrak m\}$ 上に台をもつ
連接加群であるから有限長をもつ
（代数，補題 \ref{algebra-lemma-support-point}）．

\begin{lemma}
\label{algebraization-lemma-well-defined-chi-triple}
(\ref{algebraization-equation-chi-triple}) の量
$\chi(\mathcal{F}, \mathcal{F}_0, \alpha)$ は，補題
\ref{algebraization-lemma-prepare-chi-triple} における
$\mathcal{F}', \alpha', \alpha'_0$ の選択に依存しない．
\end{lemma}

\begin{proof}
$\mathcal{F}', \alpha', \alpha'_0$ および
$\mathcal{F}'', \alpha'', \alpha''_0$ をそのような二つの選択とする．
$n > 0$ に対し $\mathcal{F}'_n = \mathfrak m^n \mathcal{F}'$ とおく．
スキームのコホモロジー，補題 \ref{coherent-lemma-homs-over-open} により，
ある $n$ に対して $\mathcal{O}_X$-加群写像
$\mathcal{F}'_n \to \mathcal{F}''$ で，同一視
$\mathcal{F}''|_U = \mathcal{F}'|_U$ と一致するものが存在する．
この同一視は $\alpha'$ と $\alpha''$ により定まる．
すると図式
$$
\xymatrix{
\mathcal{F}'_n/f\mathcal{F}'_n \ar[r] \ar[d] &
\mathcal{F}'/f\mathcal{F}' \ar[d]^{\alpha_0'} \\
\mathcal{F}''/f\mathcal{F}'' \ar[r]^{\alpha_0''} &
\mathcal{F}_0
}
$$
は $U_0$ に制限すると可換である．したがってスキームのコホモロジー，補題
\ref{coherent-lemma-homs-over-open} により，
$\mathfrak m^l(\mathcal{F}'_n/f\mathcal{F}'_n)$ に制限すれば可換となる
ような $l > 0$ が存在する．
$\mathcal{F}'_{n + l}/f\mathcal{F}'_{n + l} \to \mathcal{F}'_n/f\mathcal{F}'_n$
は $\mathfrak m^l(\mathcal{F}'_n/f\mathcal{F}'_n)$ を経由するので，
$n$ を $n + l$ で置き換えればこの図式が可換になることがわかる．
言い換えると，第三の選択
$\mathcal{F}''', \alpha''', \alpha'''_0$
であって，写像 $\mathcal{F}''' \to \mathcal{F}''$ および
$\mathcal{F}''' \to \mathcal{F}'$ が $X$ 上に存在し，$U$ および $X_0$ 上の
写像と両立するものを得た．これにより次の段落で扱う場合へ帰着される．

\medskip\noindent
$\mathcal{F}'' \to \mathcal{F}'$ を $X$ 上の写像で，$\alpha', \alpha''$ と
$U$ 上で，$\alpha'_0, \alpha''_0$ と $X_0$ 上で
両立するものがあると仮定する．$\mathcal{F}'' \to \mathcal{F}'$ は
$U$ 上の同型であり，かつ $f : \mathcal{F}'' \to \mathcal{F}''$ が
単射であるから，この写像も単射であることに注意する．明らかに
$\mathcal{F}'/\mathcal{F}''$ は $\{\mathfrak m\}$ 上に台をもち，
したがって有限長である．連接 $\mathcal{O}_{X_0}$-加群の写像
$$
\mathcal{F}''/f\mathcal{F}'' \to
\mathcal{F}'/f\mathcal{F}' \xrightarrow{\alpha'_0}
\mathcal{F}_0
$$
があり，その合成は $\alpha''_0$ で，これらは $U_0$ 上の同型である．
初等的なホモロジー代数により，$6$ 項完全列
$$
\begin{matrix}
0 \to
\Ker(\mathcal{F}''/f\mathcal{F}'' \to \mathcal{F}'/f\mathcal{F}') \to
\Ker(\alpha''_0) \to
\Ker(\alpha'_0) \to \\
\Coker(\mathcal{F}''/f\mathcal{F}'' \to \mathcal{F}'/f\mathcal{F}') \to
\Coker(\alpha''_0) \to
\Coker(\alpha'_0) \to 0
\end{matrix}
$$
を得る．長さの加法性（代数，補題 \ref{algebra-lemma-length-additive}）により，
次を示せば十分であることがわかる．
$$
\text{length}_A(
\Coker(\mathcal{F}''/f\mathcal{F}'' \to \mathcal{F}'/f\mathcal{F}')) -
\text{length}_A(
\Ker(\mathcal{F}''/f\mathcal{F}'' \to \mathcal{F}'/f\mathcal{F}')) = 0
$$
これは図式
$$
\xymatrix{
0 \ar[r] &
\mathcal{F}'' \ar[r]_f \ar[d] &
\mathcal{F}'' \ar[r] \ar[d] &
\mathcal{F}''/f\mathcal{F}'' \ar[r] \ar[d] &
0 \\
0 \ar[r] &
\mathcal{F}' \ar[r]^f &
\mathcal{F}' \ar[r] &
\mathcal{F}'/f\mathcal{F}' \ar[r] &
0
}
$$
に蛇の補題を適用し，$\mathcal{F}'/\mathcal{F}''$ が有限長であることを
用いれば従う．
\end{proof}

\begin{lemma}
\label{algebraization-lemma-ses-chi-triple}
等式
$\chi(\mathcal{G}, \mathcal{G}_0, \beta) =
\chi(\mathcal{F}, \mathcal{F}_0, \alpha) +
\chi(\mathcal{H}, \mathcal{H}_0, \gamma)$ は，列
$$
0 \to
(\mathcal{F}, \mathcal{F}_0, \alpha) \to
(\mathcal{G}, \mathcal{G}_0, \beta) \to
(\mathcal{H}, \mathcal{H}_0, \gamma)
\to 0
$$
が連接三つ組の短完全列ならば成り立つ．
\end{lemma}

\begin{proof}
$\mathcal{G}', \beta', \beta'_0$ を補題
\ref{algebraization-lemma-prepare-chi-triple} のように，三つ組
$(\mathcal{G}, \mathcal{G}_0, \beta)$ に対して選ぶ．
包含射を $j : U \to X$ と書く．$\mathcal{F}' \subset \mathcal{G}'$ を
合成
$$
\mathcal{G}' \xrightarrow{\beta'} j_*\mathcal{G} \to j_*\mathcal{H}
$$
の核とする．$\mathcal{H}' = \mathcal{G}'/\mathcal{F}'$ は
$j_*\mathcal{H}$ の連接部分層であり，したがって
$f : \mathcal{H}' \to \mathcal{H}'$ は単射であることに注意する．
よって蛇の補題により短完全列
$$
0 \to \mathcal{F}'/f\mathcal{F}' \to
\mathcal{G}'/f\mathcal{G}' \to
\mathcal{H}'/f\mathcal{H}' \to 0
$$
を得る．構成により同型
$\alpha' : \mathcal{F}'|_U \to \mathcal{F}$,
$\beta' : \mathcal{G}'|_U \to \mathcal{G}$，および
$\gamma' : \mathcal{H}'|_U \to \mathcal{H}$ がある．
証明を完了するには，補題 \ref{algebraization-lemma-prepare-chi-triple} のような写像
$\alpha'_0 : \mathcal{F}'/f\mathcal{F}' \to \mathcal{F}_0$ および
$\gamma'_0 : \mathcal{H}'/f\mathcal{H}' \to \mathcal{H}_0$ を構成し，
これらが可換図式
$$
\xymatrix{
0 \ar[r] &
\mathcal{F}'/f\mathcal{F}' \ar[r] \ar@{..>}[d]^{\alpha'_0} &
\mathcal{G}'/f\mathcal{G}' \ar[r] \ar[d]^{\beta'_0} &
\mathcal{H}'/f\mathcal{H}' \ar[r] \ar@{..>}[d]^{\gamma'_0} &
0 \\
0 \ar[r] &
\mathcal{F}_0 \ar[r] &
\mathcal{G}_0 \ar[r] &
\mathcal{H}_0 \ar[r] &
0
}
$$
に入るようにする必要がある．しかし，$\mathcal{G}'$ の最初の選択では
これは不可能かもしれない．表示した図式から，障害はちょうど合成
$$
\delta :
\mathcal{F}'/f\mathcal{F}' \to
\mathcal{G}'/f\mathcal{G}' \xrightarrow{\beta'_0}
\mathcal{G}_0 \to
\mathcal{H}_0
$$
であることがわかる．$\delta$ の $U_0$ への制限は，
$\mathcal{F}'$ と $\mathcal{H}'$ の選択により零であることに注意する．したがって
スキームのコホモロジー，補題 \ref{coherent-lemma-homs-over-open} により，
ある $k > 0$ が存在し，$\delta$ は
$\mathfrak m^k \cdot (\mathcal{F}'/f\mathcal{F}')$ 上で消える．
$n > k$ に対して
$\mathcal{G}'_n = \mathfrak m^n \mathcal{G}'$，
$\mathcal{F}'_n = \mathcal{G}'_n \cap \mathcal{F}'$，および
$\mathcal{H}'_n = \mathcal{G}'_n/\mathcal{F}'_n$ とおく．
$\beta'_0$ と
$\mathcal{G}'_n/f\mathcal{G}'_n \to \mathcal{G}'/f\mathcal{G}'$
を合成して写像
$\beta'_{n, 0} : \mathcal{G}'_n/f\mathcal{G}'_n \to \mathcal{G}_0$
を得ることができ，これは補題 \ref{algebraization-lemma-prepare-chi-triple} のような写像である．
Artin--Rees（代数，補題 \ref{algebra-lemma-Artin-Rees}）により，
$n$ を，$\mathcal{F}'_n \subset \mathfrak m^k \mathcal{F}'$ と
なるように選べる．上と同様，写像
$f : \mathcal{F}'_n \to \mathcal{F}'_n$,
$f : \mathcal{G}'_n \to \mathcal{G}'_n$，および
$f : \mathcal{H}'_n \to \mathcal{H}'_n$ は単射であり，やはり蛇の補題により
短完全列
$$
0 \to \mathcal{F}'_n/f\mathcal{F}'_n \to
\mathcal{G}'_n/f\mathcal{G}'_n \to
\mathcal{H}'_n/f\mathcal{H}'_n \to 0
$$
を得る．上と同様に同型
$\alpha'_n : \mathcal{F}'_n|_U \to \mathcal{F}$,
$\beta'_n : \mathcal{G}'_n|_U \to \mathcal{G}$，および
$\gamma'_n : \mathcal{H}'_n|_U \to \mathcal{H}$ がある．
先ほどと同じ障害
$$
\delta_n :
\mathcal{F}'_n/f\mathcal{F}'_n \to
\mathcal{G}'_n/f\mathcal{G}'_n
\xrightarrow{\beta'_{n, 0}}
\mathcal{G}_0 \to
\mathcal{H}_0
$$
を考える．しかし，可換図式
$$
\xymatrix{
\mathcal{F}'_n/f\mathcal{F}'_n \ar[r] \ar[d] &
\mathcal{G}'_n/f\mathcal{G}'_n \ar[r]_{\beta'_{n, 0}} \ar[d] &
\mathcal{G}_0 \ar[r] \ar[d] &
\mathcal{H}_0 \ar[d] \\
\mathcal{F}'/f\mathcal{F}' \ar[r] &
\mathcal{G}'/f\mathcal{G}' \ar[r]^{\beta'_0} &
\mathcal{G}_0 \ar[r] &
\mathcal{H}_0
}
$$
と，$n$ の選択および $\delta$ に関する観察から
$\delta_n = 0$ がわかる．これにより求める写像
$\alpha'_{n, 0} : \mathcal{F}'_n/f\mathcal{F}'_n \to \mathcal{F}_0$，および
$\gamma'_{n, 0} : \mathcal{H}'_n/f\mathcal{H}'_n \to \mathcal{H}_0$
を得る．したがって，
$\mathcal{F}'_n, \alpha'_n, \alpha'_{n, 0}$,
$\mathcal{G}'_n, \beta'_n, \beta'_{n, 0}$，および
$\mathcal{H}'_n, \gamma'_n, \gamma'_{n, 0}$
を用いて
$\chi(\mathcal{F}, \mathcal{F}_0, \alpha)$,
$\chi(\mathcal{G}, \mathcal{G}_0, \beta)$，および
$\chi(\mathcal{H}, \mathcal{H}_0, \gamma)$ を計算してよい．最後に，図式
$$
\xymatrix{
0 \ar[r] &
\mathcal{F}'_n/f\mathcal{F}'_n \ar[r] \ar[d] &
\mathcal{G}'_n/f\mathcal{G}'_n \ar[r] \ar[d] &
\mathcal{H}'_n/f\mathcal{H}'_n \ar[r] \ar[d] &
0 \\
0 \ar[r] &
\mathcal{F}_0 \ar[r] &
\mathcal{G}_0 \ar[r] &
\mathcal{H}_0 \ar[r] &
0
}
$$
に蛇の補題を適用し，長さの加法性
（代数，補題 \ref{algebra-lemma-length-additive}）を用いれば補題が従う．
\end{proof}

\begin{proposition}
\label{algebraization-proposition-hilbert-triple}
$(\mathcal{F}, \mathcal{F}_0, \alpha)$ を連接三つ組とし，
$(\mathcal{L}, \mathcal{L}_0, \lambda)$ を可逆な連接三つ組とする．
このとき，関数
$$
\mathbf{Z} \longrightarrow \mathbf{Z},\quad
n \longmapsto
\chi((\mathcal{F}, \mathcal{F}_0, \alpha) \otimes
(\mathcal{L}, \mathcal{L}_0, \lambda)^{\otimes n})
$$
は次数 $\leq \dim(\text{Supp}(\mathcal{F}))$ の多項式である．
\end{proposition}

\noindent
より正確には，$\mathcal{F} = 0$ ならばこの関数は定数である．
$\mathcal{F}$ が $U$ 内に有限な台をもつならば，この関数は定数である．
$\mathcal{F}$ の $U$ 内での台が次元 $1$，すなわち
$\mathcal{F}$ の台の $X$ における閉包が次元 $2$ ならば，
この関数は一次式である，等々．

\begin{proof}
これは $\mathcal{F}$ の台の次元に関する帰納法で証明する．

\medskip\noindent
基底の場合は $\mathcal{F} = 0$ のときである．このとき
$\mathcal{F}_0$ は零であるか，その台は $\{\mathfrak m\}$ である．
この場合，
$$
(\mathcal{F}, \mathcal{F}_0, \alpha) \otimes
(\mathcal{L}, \mathcal{L}_0, \lambda)^{\otimes n} =
(0, \mathcal{F}_0 \otimes \mathcal{L}_0^{\otimes n}, 0) \cong
(0, \mathcal{F}_0, 0)
$$
が成り立つ．したがって補題の関数は定数であり，その値は
$\mathcal{F}_0$ の長さに等しい．

\medskip\noindent
帰納段階．$\mathcal{F}$ の台が空でないと仮定する．
$\mathcal{G}_0 \subset \mathcal{F}_0$ を，$\{\mathfrak m\}$ 上に
台をもつ切断からなる部分加群とする．すると短完全列
$$
0 \to (0, \mathcal{G}_0, 0) \to
(\mathcal{F}, \mathcal{F}_0, \alpha) \to
(\mathcal{F}, \mathcal{F}_0/\mathcal{G}_0, \alpha) \to 0
$$
を得る．この列は，可逆な連接三つ組
$(\mathcal{L}, \mathcal{L}_0, \lambda)$ をテンソルしても完全なままである
（上の議論を参照）．したがって $\chi$ の加法性
（補題 \ref{algebraization-lemma-ses-chi-triple}）と上で説明した基底の場合により，
帰納段階を
$(\mathcal{F}, \mathcal{F}_0/\mathcal{G}_0, \alpha)$
に対して証明すれば十分である．このようにして，$\mathfrak m$ は
$\mathcal{F}_0$ の随伴点でないと仮定してよいことがわかる．

\medskip\noindent
$T = \text{Ass}(\mathcal{F}) \cup \text{Ass}(\mathcal{F}/f\mathcal{F})$
とおく．$U$ は準アフィンなので，$s \in \Gamma(U, \mathcal{L})$ で，
どの $u \in T$ においても消えないものを見つけることができる．
性質，補題
\ref{properties-lemma-quasi-affine-invertible-nonvanishing-section}．
$s$ に $\mathfrak m$ の適当な元を掛けることにより，
$\lambda(s \bmod f) = s_0|_{U_0}$ がある
$s_0 \in \Gamma(X_0, \mathcal{L}_0)$ に対して成り立つと仮定してよい．
詳細は省略する．
連接三つ組の圏において射
$$
(s, s_0) :
(\mathcal{O}_U, \mathcal{O}_{X_0}, 1)
\longrightarrow
(\mathcal{L}, \mathcal{L}_0, \lambda)
$$
を得る．
$\mathcal{G} = \Coker(s : \mathcal{F} \to \mathcal{F} \otimes \mathcal{L})$
および
$\mathcal{G}_0 = \Coker(s_0 : \mathcal{F}_0 \to
\mathcal{F}_0 \otimes \mathcal{L}_0)$ とおく．写像
$s_0 : \mathcal{F}_0 \to \mathcal{F}_0 \otimes \mathcal{L}_0$ は，
$U_0$ 上で $s$ の選択により単射であり，$\mathfrak m$ は
$\mathcal{F}_0$ の随伴点ではないので，単射であることに注意する．
したがって，同型
$\beta : \mathcal{G}/f\mathcal{G} \to \mathcal{G}_0|_{U_0}$
であって，短完全列
$$
0 \to
(\mathcal{F}, \mathcal{F}_0, \alpha) \to
(\mathcal{F}, \mathcal{F}_0, \alpha) \otimes
(\mathcal{L}, \mathcal{L}_0, \lambda) \to
(\mathcal{G}, \mathcal{G}_0, \beta) \to 0
$$
を得るようなものが存在する．台の次元に関する帰納法により，命題は
連接三つ組 $(\mathcal{G}, \mathcal{G}_0, \beta)$ に対して成り立つ．
補題 \ref{algebraization-lemma-ses-chi-triple} の加法性を用いると，
$$
n \longmapsto 
\chi((\mathcal{F}, \mathcal{F}_0, \alpha) \otimes
(\mathcal{L}, \mathcal{L}_0, \lambda)^{\otimes n + 1})
-
\chi((\mathcal{F}, \mathcal{F}_0, \alpha) \otimes
(\mathcal{L}, \mathcal{L}_0, \lambda)^{\otimes n})
$$
が多項式であることがわかる．すべての整数上で定義された関数に対する
代数，補題 \ref{algebra-lemma-numerical-polynomial} の変形により結論する
（詳細は省略する）．
\end{proof}

\begin{lemma}
\label{algebraization-lemma-nonnegative-chi-triple}
$\text{depth}(A) \geq 3$，または同値なこととして
$\text{depth}(A/fA) \geq 2$ を仮定する．
$(\mathcal{L}, \mathcal{L}_0, \lambda)$ を可逆な連接三つ組とする．このとき
$$
\chi(\mathcal{L}, \mathcal{L}_0, \lambda) =
\text{length}_A \Coker(\Gamma(U, \mathcal{L}) \to \Gamma(U_0, \mathcal{L}_0))
$$
であり，特にこれは $\geq 0$ である．さらに，
$\chi(\mathcal{L}, \mathcal{L}_0, \lambda) = 0$ であることと
$\mathcal{L} \cong \mathcal{O}_U$ であることは同値である．
\end{lemma}

\begin{proof}
深さに関する二条件の同値性は，代数，補題
\ref{algebra-lemma-depth-drops-by-one} から従う．
深さの条件により
$\Gamma(U, \mathcal{O}_U) = A$ および
$\Gamma(U_0, \mathcal{O}_{U_0}) = A/fA$ である．双対化複体，補題
\ref{dualizing-lemma-depth} および局所コホモロジー，補題
\ref{local-cohomology-lemma-finiteness-pushforwards-and-H1-local}．
局所コホモロジー，補題
\ref{local-cohomology-lemma-finiteness-for-finite-locally-free}
を用いると，$M = \Gamma(U, \mathcal{L})$ は有限 $A$-加群であることが
わかる．すると，例えば局所コホモロジー，補題
\ref{local-cohomology-lemma-finiteness-pushforwards-and-H1-local} の (4)，
または因子，補題 \ref{divisors-lemma-depth-pushforward} により，
$\text{depth}(M) \geq 2$ が従う．また，
$\mathcal{L}_0 \cong \mathcal{O}_{X_0}$ である．実際，$X_0$ は
局所スキームである．したがって
$M_0 = \Gamma(X_0, \mathcal{L}_0) = \Gamma(U_0, \mathcal{L}_0|_{U_0})$
であり，この加群は $A/fA$ と同型であることもわかる．

\medskip\noindent
上により，$\mathcal{F}' = \widetilde{M}$ は連接
$\mathcal{O}_X$-加群であり，その $U$ への制限は $\mathcal{L}$ と同型である．
同型 $\lambda : \mathcal{L}/f\mathcal{L} \to \mathcal{L}_0|_{U_0}$ は，
大域切断上の写像 $M/fM \to M_0$ を定め，これは $U_0$ 上の同型である．
$\text{depth}(M) \geq 2$ なので $H^0_\mathfrak m(M/fM) = 0$ であり，
したがって $M/fM \to M_0$ は単射である．よって定義により
$$
\chi(\mathcal{L}, \mathcal{L}_0, \lambda) =
\text{length}_A \Coker(M/fM \to M_0)
$$
となり，補題の最初の主張を得る．

\medskip\noindent
最後に，この長さが $0$ ならば，$M \to M_0$ は全射である．
したがって，$s \in M = \Gamma(U, \mathcal{L})$ で
$\mathcal{L}_0$ の自明化切断へ写るものを見つけることができる．
有限 $A$-加群 $K$，$Q$ を完全列
$$
0 \to K \to A \xrightarrow{s} M \to Q \to 0
$$
により定める．$K$ と $Q$ の台は $U_0$ と交わらない．実際，$s$ は
$U_0$ の各点で零でない．代数，補題
\ref{algebra-lemma-depth-in-ses} を用いると，
$\text{depth}(K) \geq 2$ であることがわかる
（$As \subset M$ は $\text{depth} \geq 1$ をもち，$M$ の部分加群で
あることに注意する）．したがって代数，補題
\ref{algebra-lemma-bound-depth} により，$K$ の台は空でなければ
次元 $\geq 2$ をもつ．これは
$\text{Supp}(M) \cap V(f) \subset \{\mathfrak m\}$ に矛盾するので，
$K = 0$ でなければならない．
$K = 0$ のとき $\text{depth}(Q) \geq 2$ がわかり，
前と同様に $Q = 0$ と結論する．したがって $A \cong M$ であり，
$\mathcal{L}$ は自明である．
\end{proof}







\section{穿孔スペクトル上の可逆加群 I}
\label{algebraization-section-local-lefschetz-for-pic}

\noindent
本節ではピカール群に対するいくつかの局所レフシェッツ定理を証明する．
着想の一部は
\cite{Kollar-pic}，\cite{Bhatt-local}，および \cite{Kollar-map-pic}
による．

\begin{lemma}
\label{algebraization-lemma-injective-torsion-in-pic}
$(A, \mathfrak m)$ を Noether 局所環とする．$f \in \mathfrak m$ を
非零因子とし，$\text{depth}(A/fA) \geq 2$，または同値なこととして
$\text{depth}(A) \geq 3$ を仮定する．$U$，それぞれ $U_0$ を
$A$，それぞれ $A/fA$ の穿孔スペクトルとする．写像
$$
\Pic(U) \to \Pic(U_0)
$$
は捩れ部分上で単射である．
\end{lemma}

\begin{proof}
$\mathcal{L}$ を可逆 $\mathcal{O}_U$-加群とする．
$\mathcal{L}$ が $0$ へ写ること（$\Pic(U_0)$ において）と，
$\mathcal{L}$ を節 \ref{algebraization-section-coherent-triples} のような可逆な
連接三つ組 $(\mathcal{L}, \mathcal{L}_0, \lambda)$ へ拡張できることは
同値である．命題 \ref{algebraization-proposition-hilbert-triple} により，関数
$$
n \longmapsto \chi((\mathcal{L}, \mathcal{L}_0, \lambda)^{\otimes n})
$$
は多項式である．補題 \ref{algebraization-lemma-nonnegative-chi-triple} により，
この多項式の値が零であることと
$\mathcal{L}^{\otimes n}$ が自明であることは同値である．
したがって $\mathcal{L}$ が捩れ元ならば，この多項式は無限個の零点をもち，
ゆえに恒等的に零であり，したがって $\mathcal{L}$ は自明である．
\end{proof}

\begin{proposition}[Koll\'ar]
\label{algebraization-proposition-injective-pic}
\begin{reference}
\cite[定理 1.9]{Kollar-map-pic}
\end{reference}
$(A, \mathfrak m)$ を Noether 局所環とする．$f \in \mathfrak m$ とする．
次を仮定する．
\begin{enumerate}
\item $A$ は双対化複体をもつ；
\item $f$ は非零因子である；
\item $\text{depth}(A/fA) \geq 2$，または同値なこととして
$\text{depth}(A) \geq 3$ である；
\item $f \in \mathfrak p \subset A$ が
$\dim(A/\mathfrak p) = 2$ を満たす素イデアルならば，
$\text{depth}(A_\mathfrak p) \geq 2$ である．
\end{enumerate}
$U$，それぞれ $U_0$ を $A$，それぞれ $A/fA$ の穿孔スペクトルとする．写像
$$
\Pic(U) \to \Pic(U_0)
$$
は単射である．最後に，(1)，(2)，(3) が成り立ち，$A$ が $(S_2)$ で，
$\dim(A) \geq 4$ ならば，(4) が成り立つ．
\end{proposition}

\begin{proof}
$\mathcal{L}$ を可逆 $\mathcal{O}_U$-加群とする．
$\mathcal{L}$ が $0$ へ写ること（$\Pic(U_0)$ において）と，
$\mathcal{L}$ を節 \ref{algebraization-section-coherent-triples} のような可逆な
連接三つ組 $(\mathcal{L}, \mathcal{L}_0, \lambda)$ へ拡張できることは
同値である．命題 \ref{algebraization-proposition-hilbert-triple} により，関数
$$
n \longmapsto \chi((\mathcal{L}, \mathcal{L}_0, \lambda)^{\otimes n})
$$
は多項式 $P$ である．補題 \ref{algebraization-lemma-nonnegative-chi-triple} により，
$P(n) \geq 0$ がすべての $n \in \mathbf{Z}$ に対して成り立ち，等号が
成り立つことと $\mathcal{L}^{\otimes n}$ が自明であることは同値である．
特に $P(0) = 0$ であり，$P$ は恒等的に零であって結論を得るか，
または $P$ は偶数次数 $\geq 2$ をもつ．

\medskip\noindent
$M = \Gamma(U, \mathcal{L})$ および
$M_0 = \Gamma(X_0, \mathcal{L}_0) = \Gamma(U_0, \mathcal{L}_0)$ とおく．
すると $M$ は有限 $A$-加群で，深さ $\geq 2$ であり，
$M_0 \cong A/fA$ である．補題 \ref{algebraization-lemma-nonnegative-chi-triple} の
証明を参照せよ．$H^2_\mathfrak m(M)$ は有限 $A$-加群である．これは
局所コホモロジー，補題
\ref{local-cohomology-lemma-local-finiteness-for-finite-locally-free}
と，$H^i_\mathfrak m(A) = 0$ が $i = 0, 1, 2$ に対して成り立つことから
従う．ここで $\text{depth}(A) \geq 3$ を用いた．短完全列
$$
0 \to M/fM \to M_0 \to Q \to 0
$$
を考える．補題 \ref{algebraization-lemma-nonnegative-chi-triple} により，$Q$ は
$\chi(\mathcal{L}, \mathcal{L}_0, \lambda)$ に等しい有限長をもつ．
$Q = H^1_\mathfrak m(M/fM)$ および
$H^i_\mathfrak m(M/fM) = H^i_\mathfrak m(M_0) \cong H^i_\mathfrak m(A/fA)$
を，表示した短完全列に付随する局所コホモロジーの長完全列から
$i > 1$ に対して得る．列
$0 \to M \to M \to M/fM \to 0$ に付随する局所コホモロジーの
長完全列を考える．これは
$$
0 \to Q \to H^2_\mathfrak m(M) \to H^2_\mathfrak m(M) \to
H^2_\mathfrak m(A/fA)
$$
で始まる．長さの加法性を用いると，
$\chi(\mathcal{L}, \mathcal{L}_0, \lambda)$
は $H^2_\mathfrak m(M) \to H^2_\mathfrak m(A/fA)$ の像の長さに
等しいことがわかる．

\medskip\noindent
証明の残りを明らかにするため，まず特殊な場合に補題を証明しよう．
しばらく $H^2_\mathfrak m(A/fA)$ が有限長加群であると仮定する．すると
$P(1) \leq \text{length}_A H^2_\mathfrak m(A/fA)$ となるはずである．
まったく同じ議論を $\mathcal{L}^{\otimes n}$ に適用すると，
$P(n) \leq \text{length}_A H^2_\mathfrak m(A/fA)$ がすべての $n$ に対して
わかる．したがって $P$ は正の次数をもつことができず，結論を得る．
証明の残りでは，この議論を修正して $P(n)$ の一次の上界を与える．
これで十分である．

\medskip\noindent
写像
$H^2_\mathfrak m(M) \to H^2_\mathfrak m(M_0) \cong H^2_\mathfrak m(A/fA)$
を調べよう．正規化された双対化複体 $\omega_A^\bullet$ を $A$ に対して
選ぶ．局所双対性
（双対化複体，補題 \ref{dualizing-lemma-special-case-local-duality}）により，
この写像は写像
$$
\text{Ext}^{-2}_A(M, \omega_A^\bullet)
\longleftarrow
\text{Ext}^{-2}_A(M_0, \omega_A^\bullet)
$$
の Matlis 双対であり，したがってその像は同じ有限長をもつ．
有限 $A$-加群
$\text{Ext}^{-2}_A(M_0, \omega_A^\bullet)$ の台は，空でなければ
$\mathfrak m$ と，有限個の素イデアル
$\mathfrak p_1, \ldots, \mathfrak p_r$ からなる．これらは $f$ を含み，
$\dim(A/\mathfrak p_i) = 1$ を満たす．実際，局所コホモロジー，補題
\ref{local-cohomology-lemma-sitting-in-degrees} により，この台は
次を満たす素イデアル $\mathfrak p \subset A$ の集合に含まれる：
$\text{depth}_{A_\mathfrak p}(M_{0, \mathfrak p}) + \dim(A/\mathfrak p) \leq 2$．
したがって，素イデアル $\mathfrak p$ で，$f$ を含み，
$\dim(A/\mathfrak p) = 2$ および
$\text{depth}_{A_\mathfrak p}(M_{0, \mathfrak p}) = 0$ を満たすものが
存在しないことを示せば十分である．しかし，
$M_{0, \mathfrak p} \cong (A/fA)_\mathfrak p$ なので，このような素イデアルは
$\text{depth}(A_\mathfrak p) = 1$ を与え，仮定 (4) に矛盾する．
$t \in \Gamma(U, \mathcal{L}^{\otimes -1})$ を，各点
$\mathfrak p_1, \ldots, \mathfrak p_r$ で消えないように選ぶ．性質，補題
\ref{properties-lemma-quasi-affine-invertible-nonvanishing-section}．
大域切断上での $t$ による乗法は写像 $t : M \to A$ を定め，これは
同型 $M_{\mathfrak p_i} \to A_{\mathfrak p_i}$ を
$i = 1, \ldots, r$ に対して定める．
$t_0 = t|_{U_0}$ を，対応する
$\Gamma(U_0, \mathcal{L}_0^{\otimes -1})$ の切断と書く．これは同様に
写像 $t_0 : M_0 \to A/fA$ を定め，$t$ と両立する．
したがって可換図式
$$
\xymatrix{
\text{Ext}^{-2}_A(M, \omega_A^\bullet) &
\text{Ext}^{-2}_A(M_0, \omega_A^\bullet) \ar[l] \\
\text{Ext}^{-2}_A(A, \omega_A^\bullet) \ar[u]^t &
\text{Ext}^{-2}_A(A/fA, \omega_A^\bullet) \ar[l] \ar[u]_{t_0}
}
$$
がある．上の水平写像の像の長さは
$\text{Ext}^{-2}_A(A/fA, \omega_A^\bullet)$ の長さと
$t_0$ の余核の長さとの和以下であることが従う．

\medskip\noindent
ここで，$\mathcal{L}$ を $\mathcal{L}^n$ で $n > 1$ に対して
置き換えるならば，
$$
t^n :
M_n = \Gamma(U, \mathcal{L}^{\otimes n})
\longrightarrow
\Gamma(U, \mathcal{O}_U) = A
$$
を $t$ の代わりに用いることができる．これは $t_0$ をその $n$ 乗で
置き換える．したがって写像
$\text{Ext}^{-2}_A(M_n, \omega_A^\bullet) \leftarrow
\text{Ext}^{-2}_A(M_{n, 0}, \omega_A^\bullet)$
の像の長さは，$\text{Ext}^{-2}_A(A/fA, \omega_A^\bullet)$ の長さと，
写像
$$
t_0^n : 
\text{Ext}^{-2}_A(A/fA, \omega_A^\bullet)
\longrightarrow
\text{Ext}^{-2}_A(M_{n, 0}, \omega_A^\bullet)
$$
の余核の長さとの和以下である．同型 $M_0 \cong A/fA$ を介して，
写像 $t_0$ は $g : A/fA \to A/fA$ となる．ここで
$g \in A/fA$ である．対応する同型
$M_{n, 0} \cong A/fA$ を介して，写像 $t_0^n$ は
$g^n : A/fA \to A/fA$ となる．したがって上の余核の長さは
$\text{Ext}^{-2}_A(A/fA, \omega_A^\bullet)$ の $g^n$ による商の長さである．
$\text{Ext}^{-2}_A(A/fA, \omega_A^\bullet)$ は有限 $A$-加群であり，
その台 $T$ は次元 $1$ をもち，また $V(g) \cap T$ は $t$ の選択により
閉点のみからなる．したがって代数，補題
\ref{algebra-lemma-support-dimension-d} により，この長さは $n$ に
関して一次的に増大する．

\medskip\noindent
証明を完了するため最後の主張を証明する．
$f \in \mathfrak m \subset A$ が (1)，(2)，(3) を満たし，$A$ は $(S_2)$ で，
$\dim(A) \geq 4$ であると仮定する．条件 (1) から $A$ は鎖状である．
双対化複体，補題 \ref{dualizing-lemma-universally-catenary} を参照せよ．
すると局所コホモロジー，補題
\ref{local-cohomology-lemma-catenary-S2-equidimensional}
により $\Spec(A)$ は等次元である．したがって
$\dim(A_\mathfrak p) + \dim(A/\mathfrak p) \geq 4$ が任意の素イデアル
$\mathfrak p$（$A$ の素イデアル）に対して成り立つ．すると
$\text{depth}(A_\mathfrak p) \geq \min(2, \dim(A_\mathfrak p))
\geq \min(2, 4 - \dim(A/\mathfrak p))$ であり，したがって (4) が成り立つ．
\end{proof}

\begin{remark}
\label{algebraization-remark-compare-SGA2}
SGA2 には次の結果がある．$(A, \mathfrak m)$ を Noether 局所環とし，
$f \in \mathfrak m$ とする．$A$ は正則環の商であり，元
$f$ は非零因子であり，さらに次を仮定する．
\begin{enumerate}
\item[(a)] $\mathfrak p \subset A$ が
$\dim(A/\mathfrak p) = 1$ を満たす素イデアルならば，
$\text{depth}(A_\mathfrak p) \geq 2$ である；
\item[(b)] $\text{depth}(A/fA) \geq 3$，または同値なこととして
$\text{depth}(A) \geq 4$ である．
\end{enumerate}
$U$，それぞれ $U_0$ を $A$，それぞれ $A/fA$ の穿孔スペクトルとする．
このとき写像
$$
\Pic(U) \to \Pic(U_0)
$$
は単射である．これは \cite[Exposee XI, 補題 3.16]{SGA2} である
\footnote{条件 (a) は条件 (b) から従う．代数，補題
\ref{algebra-lemma-depth-localization} を参照せよ．}．
SGA2 のこの結果は命題 \ref{algebraization-proposition-injective-pic} から従う．実際，
\begin{enumerate}
\item 正則環の商は双対化複体をもつ
（双対化複体，補題 \ref{dualizing-lemma-regular-gorenstein} および
命題 \ref{dualizing-proposition-dualizing-essentially-finite-type} を参照）；
\item $\text{depth}(A) \geq 4$ ならば，
$\text{depth}(A_\mathfrak p) \geq 2$ がすべての素イデアル
$\mathfrak p$ で $\dim(A/\mathfrak p) = 2$ を満たすものに対して成り立つ．
代数，補題 \ref{algebra-lemma-depth-localization} を参照せよ．
\end{enumerate}
\end{remark}






\section{穿孔スペクトル上の可逆加群 II}
\label{algebraization-section-local-lefschetz-for-pic-surjective}

\noindent
次に，ピカール群に対する局所レフシェッツの全射性を扱う．
まず，$U_0$ 上の可逆加群をある開近傍へ延長するために，
次の簡単な判定条件がある．

\begin{lemma}
\label{algebraization-lemma-surjective-Pic-first}
$(A, \mathfrak m)$ を Noether 局所環とし，$f \in \mathfrak m$ とする．
次を仮定する．
\begin{enumerate}
\item $A$ は $f$ 進完備である．
\item $f$ は非零因子である．
\item $H^1_\mathfrak m(A/fA)$ および $H^2_\mathfrak m(A/fA)$
は有限 $A$-加群である．
\item $H^3_\mathfrak m(A/fA) = 0$\footnote{(3) と (4) は，
$\text{depth}(A/fA) \geq 4$ のとき，または同値なことだが $\text{depth}(A) \geq 5$ のときに成り立つことに注意せよ．}．
\end{enumerate}
$U$，$U_0$ をそれぞれ $A$，$A/fA$ の穿孔スペクトルとする．
このとき
$$
\colim_{U_0 \subset U' \subset U\text{ は開}} \Pic(U')
\longrightarrow
\Pic(U_0)
$$
は全射である．
\end{lemma}

\begin{proof}
$U_0 \subset U_n \subset U$ を，第 $n$ 次の $U_0$ の無限小近傍とする．
$U_n$ の $U_{n + 1}$ におけるイデアル層は $\mathcal{O}_{U_0}$ と同型であることに注意する．
実際，$U_0 \subset U$ は非零因子 $f$ で切り出される主閉部分スキームである．
したがって，可換群の完全列
$$
\Pic(U_{n + 1}) \to \Pic(U_n) \to
H^2(U_0, \mathcal{O}_{U_0}) = H^3_\mathfrak m(A/fA) = 0
$$
を得る．More on Morphisms, 補題
\ref{more-morphisms-lemma-picard-group-first-order-thickening} を参照せよ．
したがって，任意の可逆 $\mathcal{O}_{U_0}$-加群は，ある可逆な連接形式加群，
すなわち $\textit{Coh}(U, f\mathcal{O}_U)$ の可逆対象の制限である．
補題 \ref{algebraization-lemma-equivalence} を適用すれば結論を得る．
\end{proof}

\begin{remark}
\label{algebraization-remark-surjective-Pic-second}
$(A, \mathfrak m)$ を Noether 局所環とし，$f \in \mathfrak m$ とする．
次を仮定すれば，補題 \ref{algebraization-lemma-surjective-Pic-first} の結論が成り立つ．
\begin{enumerate}
\item $A$ は双対化複体をもつ．
\item $A$ は $f$ 進完備である．
\item $f$ は非零因子である．
\item 次のいずれかが成り立つ．
\begin{enumerate}
\item $A_f$ は $(S_2)$ であり，$\mathfrak p \subset A$ が
$f \not \in \mathfrak p$ を満たす極小なものならば $\dim(A/\mathfrak p) \geq 4$ である．または，
\item $\mathfrak p \not \in V(f)$ かつ
$V(\mathfrak p) \cap V(f) \not = \{\mathfrak m\}$ ならば，
$\text{depth}(A_\mathfrak p) + \dim(A/\mathfrak p) > 3$ である．
\end{enumerate}
\item $H^3_{\mathfrak m}(A/fA) = 0$.
\end{enumerate}
証明は補題 \ref{algebraization-lemma-surjective-Pic-first} の証明とまったく同じであり，
補題 \ref{algebraization-lemma-equivalence-better} を補題 \ref{algebraization-lemma-equivalence} の代わりに用いればよい．
ここでは二点述べておく必要がある．(a)
$H^3_{\mathfrak m}(A/fA) = 0$ と分かる例を，
$\text{depth}(A/fA) \geq 4$ を仮定せずに見つけるのは難しいように思われる．
(b) 補題 \ref{algebraization-lemma-equivalence-better} の証明は，
補題 \ref{algebraization-lemma-equivalence} の証明よりもかなり難しい．
\end{remark}

\begin{lemma}
\label{algebraization-lemma-surjective-Pic-first-better}
$(A, \mathfrak m)$ を Noether 局所環とし，$f \in \mathfrak m$ とする．
次を仮定する．
\begin{enumerate}
\item 補題 \ref{algebraization-lemma-surjective-Pic-first} の条件が成り立つ．
\item 任意の極大イデアル $\mathfrak p \subset A_f$ に対して，
$(A_f)_\mathfrak p$ の穿孔スペクトルのピカール群は自明である．
\end{enumerate}
$U$，$U_0$ をそれぞれ $A$，$A/fA$ の穿孔スペクトルとする．
このとき
$$
\Pic(U) \longrightarrow \Pic(U_0)
$$
は全射である．
\end{lemma}

\begin{proof}
$\mathcal{L}_0 \in \Pic(U_0)$ とする．補題
\ref{algebraization-lemma-surjective-Pic-first} により，開集合 $U_0 \subset U' \subset U$ と
$\mathcal{L}' \in \Pic(U')$ が存在し，その $U_0$ への制限は $\mathcal{L}_0$ である．
$U' \supset U_0$ なので，$U \setminus U'$ は (2) に現れる素イデアル
$\mathfrak p_1, \ldots, \mathfrak p_n$ に対応する点からなる．
仮定により，可逆加群 $\mathcal{L}'_i$ を $\Spec(A_{\mathfrak p_i})$ 上に見つけられ，
$\mathcal{L}'$ と穿孔スペクトル $U' \times_U \Spec(A_{\mathfrak p_i})$ 上で一致させられる．
実際，自明な可逆加群は常に延長できる．
Limits, 補題 \ref{limits-lemma-glueing-near-closed-point-modules} を
$n$ 回適用すると，$\mathcal{L}'$ は $U$ 上の可逆加群へ延長されることが分かる．
\end{proof}

\begin{lemma}
\label{algebraization-lemma-local-pic-to-completion}
$(A, \mathfrak m)$ を深さ $\geq 2$ の Noether 局所環とする．
$A^\wedge$ をその完備化とする．$U$，$U^\wedge$ をそれぞれ
$A$，$A^\wedge$ の穿孔スペクトルとする．このとき
$\Pic(U) \to \Pic(U^\wedge)$ は単射である．
\end{lemma}

\begin{proof}
$\mathcal{L}$ を可逆 $\mathcal{O}_U$-加群とし，その引き戻し
$\mathcal{L}^\wedge$ を $U^\wedge$ 上で考える．
$H^0(U, \mathcal{O}_U) = A$ である．これは深さに関する仮定と
Dualizing Complexes, 補題 \ref{dualizing-lemma-depth} および
Local Cohomology, 補題
\ref{local-cohomology-lemma-finiteness-pushforwards-and-H1-local} による．
したがって $\mathcal{L}$ が自明であることと，
$M = H^0(U, \mathcal{L})$ が $A$-加群として $A$ と同型であることは同値である．
（詳細は省略する．）$A \to A^\wedge$ は平坦なので，平坦基底変換により
$M \otimes_A A^\wedge = \Gamma(U^\wedge, \mathcal{L}^\wedge)$ である．
Cohomology of Schemes, 補題 \ref{coherent-lemma-flat-base-change-cohomology} を参照せよ．
最後に，$M \cong A$ であることと
$M \otimes_A A^\wedge \cong A^\wedge$ であることが同値なのは容易に分かる．
\end{proof}

\begin{lemma}
\label{algebraization-lemma-trivial-local-pic-regular}
$(A, \mathfrak m)$ を正則局所環とする．このとき $A$ の
穿孔スペクトルのピカール群は自明である．
\end{lemma}

\begin{proof}
Divisors, 補題 \ref{divisors-lemma-extend-invertible-module} と
More on Algebra, 補題 \ref{more-algebra-lemma-regular-local-UFD} を組み合わせればよい．
\end{proof}

\noindent
これで先の結果をブートストラップし，次元 $\geq 4$ の完全交叉の
穿孔スペクトルではピカール群が自明であることを証明できる．
Noether 局所環は，その完備化が正則局所環を正則列の生成する
イデアルで割った商であるとき，完全交叉と呼ばれることを思い出そう．
Divided Power Algebra, 節 \ref{dpa-section-lci} の議論を参照せよ．

\begin{proposition}[Grothendieck]
\label{algebraization-proposition-trivial-local-pic-complete-intersection}
$(A, \mathfrak m)$ を Noether 局所環とする．$A$ が次元 $\geq 4$ の
完全交叉ならば，$A$ の穿孔スペクトルのピカール群は自明である．
\end{proposition}

\begin{proof}
補題 \ref{algebraization-lemma-local-pic-to-completion} により，$A$ は完備局所環であると仮定してよい．
仮定により，$A = B/(f_1, \ldots, f_r)$ と書ける．ここで $B$ は完備正則局所環であり，
$f_1, \ldots, f_r$ は正則列である．$r$ に関する帰納法で証明を終える．
基底の場合 $r = 0$ は補題 \ref{algebraization-lemma-trivial-local-pic-regular} から従う．

\medskip\noindent
$A = B/(f_1, \ldots, f_r)$ と仮定し，命題が $r - 1$ に対して成り立つとする．
$A' = B/(f_1, \ldots, f_{r - 1})$ と置き，補題
\ref{algebraization-lemma-surjective-Pic-first-better} を $f_r \in A'$ に適用する．
これは次の理由で可能である．
\begin{enumerate}
\item 局所環が完備なので，補題 \ref{algebraization-lemma-surjective-Pic-first} の条件 (1) が成り立つ．
\item $f_1, \ldots, f_r$ は正則列なので，補題
\ref{algebraization-lemma-surjective-Pic-first} の条件 (2) が成り立つ．
\item $A = A'/f_r A'$ は次元 $\dim(A) \geq 4$ の Cohen--Macaulay 環なので，
補題 \ref{algebraization-lemma-surjective-Pic-first} の条件 (3) と (4) が成り立つ．
\item 帰納法の仮定により，補題 \ref{algebraization-lemma-surjective-Pic-first-better} の条件 (2) が成り立つ．
実際，$\dim((A'_{f_r})_\mathfrak p) \geq 4$ である．ここで $\mathfrak p$ は
$A'_{f_r}$ の極大素イデアルである．また，
$(A'_{f_r})_\mathfrak p = B_\mathfrak q/(f_1, \ldots, f_{r - 1})$ となる
素イデアル $\mathfrak q \subset B$ があり，$B_\mathfrak q$ は正則である．
\end{enumerate}
これで証明が完了する．
\end{proof}

\begin{example}
\label{algebraization-example-grothendieck-sharp}
命題 \ref{algebraization-proposition-trivial-local-pic-complete-intersection} の
次元の下界は鋭い．例えば，$A = k[[x, y, z, w]]/(xy - zw)$ の
穿孔スペクトルのピカール群は非自明である．実際，イデアル $I = (x, z)$ が
切り出す有効 Cartier 因子を $D$ とすると，これは穿孔スペクトル $U$，すなわち
$A$ の穿孔スペクトル上の因子である．$I_x, I_y, I_z, I_w$ が
$A_x, A_y, A_z, A_w$ の可逆イデアルであることから，これは容易に分かる．一方，
$A/I$ の深さは $\geq 1$（実際には $2$）なので，$I$ の深さは $\geq 2$
（実際には $3$）である．したがって
$I = \Gamma(U, \mathcal{O}_U(-D))$ である．ゆえに $\mathcal{O}_U(-D)$ が自明ならば，
$I \cong \Gamma(U, \mathcal{O}_U) = A$ となるはずである．しかし，$I$ は $1$ 元で
生成されないので，これは正しくない．
\end{example}

\begin{example}
\label{algebraization-example-grothendieck-sharp-bis}
命題 \ref{algebraization-proposition-trivial-local-pic-complete-intersection} は，商
$$
A = B/(f_1, \ldots, f_r)
$$
で $B$ が正則かつ $\dim(B) - r \geq 4$ であるものへは拡張できない．
言い換えると，$f_1, \ldots, f_r$ が正則列であるという条件は，
$A$ の穿孔スペクトルのピカール群が消えるために一般には必要である．
実際，$k$ を体とし，次のように置く．
$$
A = k[[a, b, x, y, z, u, v, w]]/(a^3, b^3, xa^2 + yab + zb^2, w^2)
$$
ここで $A = A_0[w]/(w^2)$ であり，
$A_0 = k[[a, b, x, y, z, u, v]]/(a^3, b^3, xa^2 + yab + zb^2)$ である．
以下で $A_0$ の深さが $2$ であることを示す．
$U$ を $A$ の穿孔スペクトル，$U_0$ を $A_0$ の穿孔スペクトルと記す．
短完全列 $0 \to A_0 \to A \to A_0 \to 0$ があり，その最初の射は
$w$ による乗法で与えられることに注意する．
More on Morphisms, 補題
\ref{more-morphisms-lemma-picard-group-first-order-thickening}
により，完全列
$$
H^0(U, \mathcal{O}_U^*) \to
H^0(U_0, \mathcal{O}_{U_0}^*) \to
H^1(U_0, \mathcal{O}_{U_0}) \to
\Pic(U)
$$
を得る．$A_0$ の深さ，したがって $A$ の深さは $2$ なので，
$H^0(U_0, \mathcal{O}_{U_0}) = A_0$ および $H^0(U, \mathcal{O}_U) = A$ が成り立ち，
$H^1(U_0, \mathcal{O}_{U_0})$ は非零である．Dualizing Complexes, 補題
\ref{dualizing-lemma-depth} と Local Cohomology, 補題
\ref{local-cohomology-lemma-local-cohomology} を参照せよ．
したがって上に表示した最後の射は非零であり，$\Pic(U)$ は非零であると結論する．

\medskip\noindent
$A_0$ の深さが $2$ であることを示すには，
$A_1 = k[[a, b, x, y, z]]/(a^3, b^3, xa^2 + yab + zb^2)$ の深さが
$0$ であることを示せば十分である．実際，$a^2b^2$ は $A_1$ の非零元へ写り，
各変数 $a, b, x, y, z$ によって零化される．例えば
$ya^2b^2 = (yab)(ab) = - (xa^2 + zb^2)(ab) = -xa^3b - yab^3 = 0$ が $A_1$ で成り立つ．
他の場合も同様である．
\end{example}










\section{レフシェッツ定理への応用}
\label{algebraization-section-lefschetz}

\noindent
本節では，射影スキーム $P$ 上の連接層と，豊富因子 $Q$ に沿う
形式完備化上の連接加群との関係を論じる．

\medskip\noindent
$k$ を体とする．$P$ を $k$ 上の固有スキームとする．
$\mathcal{L}$ を豊富な可逆 $\mathcal{O}_P$-加群とする．
$s \in \Gamma(P, \mathcal{L})$ を切断とし\footnote{$s$ が正則切断であるとは仮定しない．
したがって，$Q$ は $P$ の局所主閉部分スキームにすぎず，
有効 Cartier 因子であるとは限らない．}，
$Q = Z(s)$ をその零点スキームとする．Divisors, 定義
\ref{divisors-definition-zero-scheme-s} を参照せよ．
任意の $n \geq 1$ に対し，$Q_n = Z(s^n)$ を第 $n$ 次の，$Q$ に沿う無限小近傍と記す．
$\mathcal{F}$ が連接 $\mathcal{O}_P$-加群ならば，その制限を
$\mathcal{F}_n = \mathcal{F}|_{Q_n}$ と記す．すなわち，これは $\mathcal{F}$ の，閉埋め込み
$Q_n \to P$ による引き戻しである．

\begin{proposition}
\label{algebraization-proposition-lefschetz}
上の状況で，整数 $\sigma$ が存在し，任意の点 $p \in P \setminus Q$ に対して
$$
\text{depth}(\mathcal{F}_p) + \dim(\overline{\{p\}}) > \sigma
$$
が成り立つと仮定する．このとき写像
$$
H^i(P, \mathcal{F}) \longrightarrow \lim H^i(Q_n, \mathcal{F}_n)
$$
は $0 \leq i < \sigma$ に対して同型である．
\end{proposition}

\begin{proof}
More on Morphisms, 補題 \ref{more-morphisms-lemma-apply-proj-spec} を用い，
さらに More on Morphisms, 節 \ref{more-morphisms-section-proj-spec} の
記法と結果を断りなく用いる．少なくともこの補題の主張を理解しなければ，
以下の証明は意味をなさない．この場合，
$A = \bigoplus_{m \geq 0} \Gamma(P, \mathcal{L}^{\otimes m})$
は有限型 $k$-代数であり，その各次数部分は有限次元 $k$-ベクトル空間であることに注意する．
Cohomology of Schemes, 補題 \ref{coherent-lemma-coherent-proper-ample} を参照せよ．

\medskip\noindent
ここでは $s$ を元 $f \in A_1 \subset A$，すなわち次数 $1$ の $A$ の斉次元とみなす．
$Y = V(f) \subset X$ を $f$ が定める閉部分スキームと記す．するとスキーム論的に
$U \cap Y = (\pi|_U)^{-1}(Q)$ である．記法
$\mathcal{F}_U = \pi^*\mathcal{F}|_U = (\pi|_U)^*\mathcal{F}$
を思い出そう．これは連接 $\mathcal{O}_U$-加群である．
次を満たす有限 $A$-加群 $M$ を選ぶ．
$\mathcal{F}_U = \widetilde{M}|_U$
（存在については Local Cohomology, 補題
\ref{local-cohomology-lemma-finiteness-pushforwards-and-H1-local} を参照せよ）．
$H^i_Z(M)$ は $f$ のある冪で零化されると主張する．これは $i \leq \sigma + 1$ に対してである．

\medskip\noindent
この主張を証明するために Local Cohomology, 命題
\ref{local-cohomology-proposition-annihilator} を適用する．幾何学的に言い換えると，
$u \in U$，$u \not \in Y$ および $z \in \overline{\{u\}} \cap Z$ に対して
$$
\text{depth}(\mathcal{F}_{U, u}) +
\dim(\mathcal{O}_{\overline{\{u\}}, z}) > \sigma + 1
$$
を証明すれば十分である．これには少し考えるだけでよい．

\medskip\noindent
$Z = \Spec(A_0)$ は $X$ の閉点からなる有限集合であることに注意する．
実際，$A_0$ は有限次元 $k$-代数である．
（$Z$ を一点集合にしたい読者は，有限 $k$-代数 $A_0$ を $k$ で置き換えればよい．
これは証明の他の部分には何も影響しない．）

\medskip\noindent
射 $\pi : L \to P$ とその制限 $\pi|_U : U \to P$ は，相対次元 $1$ の滑らかな射である．
$u \in U$，$u \not \in Y$ および $z \in \overline{\{u\}} \cap Z$ とする．
$p = \pi(u) \in P \setminus Q$ をその像とする．このとき $u$ は，
$\pi$ の $p$ 上のファイバーの生成点であるか，そのファイバーの閉点である．
$u$ がファイバーの生成点ならば，
$\text{depth}(\mathcal{F}_{U, u}) = \text{depth}(\mathcal{F}_p)$
かつ
$\dim(\overline{\{u\}}) = \dim(\overline{\{p\}}) + 1$.
$u$ がファイバーの閉点ならば，
$\text{depth}(\mathcal{F}_{U, u}) = \text{depth}(\mathcal{F}_p) + 1$
かつ
$\dim(\overline{\{u\}}) = \dim(\overline{\{p\}})$.
どちらの場合にも
$\dim(\overline{\{u\}}) = \dim(\mathcal{O}_{\overline{\{u\}}, z})$
である．実際，$Z$ のすべての点は閉点である．したがって，求める不等式は
命題の仮定から従う．

\medskip\noindent
$A'$ を $f$ 進の，$A$ の完備化とする．Algebra, 補題
\ref{algebra-lemma-completion-flat} により $A \to A'$ は平坦である．
$U' \subset X' = \Spec(A')$ を $U$ の逆像とし，$Y'$ と $Z'$ も同様に定める．
$\mathcal{F}'$ を $U'$ 上の $\mathcal{F}_U$ の引き戻しとし，
$M' = M \otimes_A A'$ と置く．局所コホモロジーの平坦基底変換
（Local Cohomology, 補題 \ref{local-cohomology-lemma-torsion-change-rings}）により
$$
H^i_{Z'}(M') = H^i_Z(M) \otimes_A A'
$$
であり，$i \leq \sigma + 1$ ならば，これらは $f$ のある冪で零化される．次の図式を考える．
$$
\xymatrix{
& H^i(U, \mathcal{F}_U) \ar[ld] \ar[d] \ar[r] &
\lim H^i(U, \mathcal{F}_U/f^n\mathcal{F}_U) \ar@{=}[d] \\
H^i(U, \mathcal{F}_U) \otimes_A A' \ar@{=}[r] &
H^i(U', \mathcal{F}') \ar[r] &
\lim H^i(U', \mathcal{F}'/f^n\mathcal{F}')
}
$$
補題 \ref{algebraization-lemma-alternative-higher} と今証明した捩れ性により，
下の水平射は $i < \sigma$ に対して同型である．水平な等号は平坦基底変換
（Cohomology of Schemes, 補題 \ref{coherent-lemma-flat-base-change-cohomology}）であり，
垂直な等号は，$U \cap Y$ と $U' \cap Y'$，およびそれらの第 $n$ 無限小近傍が
互いに同型に写されるためである（$f$ に関して完備化しているからである）．

\medskip\noindent
More on Morphisms, 式 (\ref{more-morphisms-equation-cohomology-torsor}) を適用すると，
両立する直和分解
$$
\lim H^i(U, \mathcal{F}_U/f^n\mathcal{F}_U) =
\lim
\left(
\bigoplus\nolimits_{m \in \mathbf{Z}}
H^i(Q_n, \mathcal{F}_n \otimes \mathcal{L}^{\otimes m})
\right)
$$
および
$$
H^i(U, \mathcal{F}_U) =
\bigoplus\nolimits_{m \in \mathbf{Z}}
H^i(P, \mathcal{F} \otimes \mathcal{L}^{\otimes m})
$$
を得る．したがって Algebra, 補題 \ref{algebra-lemma-daniel-litt} により結論を得る．
\end{proof}

\begin{lemma}
\label{algebraization-lemma-lefschetz-addendum}
$k$ を体とする．$X$ を $k$ 上の固有スキームとする．
$\mathcal{L}$ を豊富な可逆 $\mathcal{O}_X$-加群とする．
$s \in \Gamma(X, \mathcal{L})$ とする．$Y = Z(s)$ を $s$ の零点スキームとし，
その第 $n$ 無限小近傍を $Y_n = Z(s^n)$ とする．
$\mathcal{F}$ を連接 $\mathcal{O}_X$-加群とする．
任意の $x \in X \setminus Y$ に対して
$$
\text{depth}(\mathcal{F}_x) + \dim(\overline{\{x\}}) > 1
$$
が成り立つと仮定する．このとき $\Gamma(V, \mathcal{F}) \to \lim \Gamma(Y_n, \mathcal{F}|_{Y_n})$ は，
任意の開部分スキーム $V \subset X$ で $Y$ を含むものに対して同型である．
\end{lemma}

\begin{proof}
命題 \ref{algebraization-proposition-lefschetz} により，これは $V = X$ に対して成り立つ．
したがって，写像
$\Gamma(V, \mathcal{F}) \to \lim \Gamma(Y_n, \mathcal{F}|_{Y_n})$
が単射であることを示せば十分である．$\sigma \in \Gamma(V, \mathcal{F})$ が
零に写るならば，その台は $Y$ と交わらない（詳細は省略する．ヒント：
Krull の共通部分定理を用いる）．したがって，閉包 $T \subset X$，すなわち
$\text{Supp}(\sigma)$ の閉包は $Y$ と交わらない．ゆえに $T$ は $k$ 上固有であり
（$X$ の閉部分だから），かつアフィンである（アフィンスキーム $X \setminus Y$ の
閉部分だからである．Morphisms, 補題 \ref{morphisms-lemma-proper-ample-delete-affine} を参照せよ）．
したがって $k$ 上有限である
（Morphisms, 補題 \ref{morphisms-lemma-finite-proper}）．
ゆえに $T$ は $X$ の閉点からなる有限集合である．仮定により，
$\text{depth}(\mathcal{F}_x) \geq 2$ であり，特に少なくとも $1$ である．これは $x \in T$ に対して成り立つ．
したがって
$\Gamma(V, \mathcal{F}) \to \Gamma(V \setminus T, \mathcal{F})$
は単射であり，望みどおり $\sigma = 0$ である．
\end{proof}


\begin{example}
\label{algebraization-example-lefschetz}
$k$ を体とし，$X$ を $k$ 上の固有多様体とする．
$Y \subset X$ を $\mathcal{O}_X(Y)$ が豊富となる有効 Cartier 因子とし，
$Y_n$ をその第 $n$ 無限小近傍と記す．
$\mathcal{E}$ を有限局所自由 $\mathcal{O}_X$-加群とする．
命題 \ref{algebraization-proposition-lefschetz} のいくつかの特別な場合を挙げる．
\begin{enumerate}
\item $X$ が曲線ならば，何も得られない．
\item $X$ が Cohen--Macaulay（例えば正規）曲面ならば，
$$
H^0(X, \mathcal{E}) \to \lim H^0(Y_n, \mathcal{E}|_{Y_n})
$$
は同型である．
\item $X$ が Cohen--Macaulay 三次元多様体ならば，
$$
H^0(X, \mathcal{E}) \to \lim H^0(Y_n, \mathcal{E}|_{Y_n})
\quad\text{および}\quad
H^1(X, \mathcal{E}) \to \lim H^1(Y_n, \mathcal{E}|_{Y_n})
$$
はいずれも同型である．
\end{enumerate}
このパターンは明らかであろう．$X$ が正規三次元多様体ならば，
$H^0$ については結論できるが，$H^1$ については結論できない．
\end{example}

\noindent
次の主要結果を証明する前に，一つ補題が必要である．

\begin{lemma}
\label{algebraization-lemma-Gm-equivariant-extend-canonically}
状況 \ref{algebraization-situation-algebraize} において，$(\mathcal{F}_n)$ を
$\textit{Coh}(U, I\mathcal{O}_U)$ の対象とする．次を仮定する．
\begin{enumerate}
\item $A$ は次数付き環であり，$\mathfrak a = A_+$ であり，
$I$ は斉次イデアルである．
\item $(\mathcal{F}_n) = (\widetilde{M_n}|_U)$ である．ここで $(M_n)$ は
次数付き $A$-加群の逆系である．
\item $(\mathcal{F}_n)$ は $X$ へ標準的に延長される．
\end{enumerate}
このとき，次を満たす有限次数付き $A$-加群 $N$ が存在する．
\begin{enumerate}
\item[(a)] 逆系 $(N/I^nN)$ と $(M_n)$ は，次数付き $A$-加群を
$A_+$-冪捩れ加群で割った圏においてプロ同型である．
\item[(b)] $(\mathcal{F}_n)$ は，$N$ に付随する連接加群の完備化である．
\end{enumerate}
\end{lemma}

\begin{proof}
$(\mathcal{G}_n)$ を補題 \ref{algebraization-lemma-canonically-algebraizable} における
標準的延長とする．$A$ と $M_n$ の次数付けは作用
$$
a : \mathbf{G}_m \times X \longrightarrow X
$$
を定める．これは群スキーム $\mathbf{G}_m$ の $X$ 上の作用であって，
$(\widetilde{M_n})$ は $\mathbf{G}_m$-同変準連接 $\mathcal{O}_X$-加群の逆系となる．
Groupoids, 例 \ref{groupoids-example-Gm-on-affine} を参照せよ．
$\mathfrak a$ と $I$ は斉次イデアルなので，閉部分スキーム $Z$，$Y$ と
開部分スキーム $U$ は，$\mathbf{G}_m$-不変な閉部分スキームおよび開部分スキームである．
制限 $(\mathcal{F}_n)$ は $(\widetilde{M_n})$ の制限であり，$\mathbf{G}_m$-同変な
連接 $\mathcal{O}_U$-加群の逆系である．言い換えると，$(\mathcal{F}_n)$ は
$\mathbf{G}_m$-同変連接形式加群である．すなわち，同型
$$
\alpha : (a^*\mathcal{F}_n) \longrightarrow (p^*\mathcal{F}_n)
$$
が $\mathbf{G}_m \times U$ 上にあり，適切な余サイクル条件を満たす．
$a$ と $p$ はアフィンスキームの平坦射なので，補題
\ref{algebraization-lemma-canonically-extend-base-change} により，一意な同型
$$
\beta : (a^*\mathcal{G}_n) \longrightarrow (p^*\mathcal{G}_n)
$$
が $\mathbf{G}_m \times X$ 上に存在し，$\alpha$ に $\mathbf{G}_m \times U$ 上で
制限される．一意性から，$\beta$ は対応する余サイクル条件を満たす．
このようにして，各 $\mathcal{G}_n$ は遷移射と両立する形で
$\mathbf{G}_m$-同変連接 $\mathcal{O}_X$-加群となる．

\medskip\noindent
Groupoids, 補題 \ref{groupoids-lemma-Gm-equivariant-module} により，
$\mathcal{G}_n$ とその $\mathbf{G}_m$-同変構造は，次数付き $A$-加群 $N_n$ に対応する．
遷移射 $N_{n + 1} \to N_n$ は次数付き加群の射である．$N_n$ は有限 $A$-加群であり，
$N_n = N_{n + 1}/I^n N_{n + 1}$ であることに注意する．これは
$(\mathcal{G}_n)$ が $\textit{Coh}(X, I\mathcal{O}_X)$ の対象だからである．
$N$ を Algebra, 補題 \ref{algebra-lemma-finiteness-graded} で得られる有限次数付き
$A$-加群とする．このとき $N_n = N/I^nN$ であるから，$(\mathcal{G}_n)$ は
$N$ に付随する連接加群の完備化であり，したがって特に (b) が成り立つ．

\medskip\noindent
(a) を見るには，上で述べた状況をもう少し展開する必要がある．まず，
$M_n \to H^0(U, \mathcal{F}_n)$ の核と余核は $A_+$-冪捩れであることに注意する
（Local Cohomology, 補題
\ref{local-cohomology-lemma-finiteness-pushforwards-and-H1-local}）．
$H^0(U, \mathcal{F}_n)$ には自然な次数付けがあり，これらの射と系の遷移射は
次数付き $A$-加群の射となることにも注意する．例えば，
$(U \to X)_*\mathcal{F}_n$ が $\mathbf{G}_m$-同変加群であり，$X$ 上にあることと，
Groupoids, 補題 \ref{groupoids-lemma-Gm-equivariant-module} を用いればよい．
次に，定義 \ref{algebraization-definition-canonically-algebraizable} と補題
\ref{algebraization-lemma-canonically-algebraizable} により，$(N_n)$ と $(H^0(U, \mathcal{F}_n))$ は
プロ同型であることを思い出そう．このプロ同型を定める射が次数付き加群の射であることの
確認は省略する．したがって，$(N_n)$ と $(M_n)$ は，次数付き $A$-加群を
$A_+$-冪捩れ加群で割った圏においてプロ同型である．
\end{proof}

\noindent
$k$ を体とする．$P$ を $k$ 上の固有スキームとする．
$\mathcal{L}$ を豊富な可逆 $\mathcal{O}_P$-加群とする．
$s \in \Gamma(P, \mathcal{L})$ を切断とし，$Q = Z(s)$ をその零点スキームとする．
Divisors, 定義 \ref{divisors-definition-zero-scheme-s} を参照せよ．
$\mathcal{I} \subset \mathcal{O}_P$ を $Q$ のイデアル層とする．
$\textit{Coh}(P, \mathcal{I})$ により，Cohomology of Schemes, 節
\ref{coherent-section-coherent-formal} で導入した連接形式加群の圏を表す．

\begin{proposition}
\label{algebraization-proposition-lefschetz-existence}
上の状況で，$(\mathcal{F}_n)$ を $\textit{Coh}(P, \mathcal{I})$ の対象とする．
任意の $q \in Q$ と，任意の素イデアル $\mathfrak p \in \mathcal{O}_{P, q}^\wedge$ で
$\mathfrak p \not \in V(\mathcal{I}_q^\wedge)$ を満たすものに対して
$$
\text{depth}((\mathcal{F}_q^\wedge)_\mathfrak p) +
\dim(\mathcal{O}_{P, q}^\wedge/\mathfrak p) +
\dim(\overline{\{q\}}) > 2
$$
が成り立つと仮定する．このとき $(\mathcal{F}_n)$ は連接
$\mathcal{O}_P$-加群の完備化である．
\end{proposition}

\begin{proof}
命題を証明するため，Cohomology of Schemes, 補題
\ref{coherent-lemma-existence-easy} により，$(\mathcal{F}_n)$ を，それと
$\mathcal{I}$-捩れだけ異なる対象で置き換えてよい（より正確には以下を参照）．
$T' = \{q \in Q \mid \dim(\overline{\{q\}}) = 0\}$ および
$T = \{q \in Q \mid \dim(\overline{\{q\}}) \leq 1\}$ と置く．
命題の仮定は，$Q \subset P$，$(\mathcal{F}_n)$ および $T' \subset T \subset Q$ が，
補題 \ref{algebraization-lemma-improvement-formal-coherent-module-better} の条件を $d = 1$ として
満たすことにほかならない．不等式の自明な変形に加えて，
$V(\mathfrak p) \cap V(\mathcal{I}^\wedge_y) = \{\mathfrak m^\wedge_y\}
\Leftrightarrow \dim(\mathcal{O}_{P, q}^\wedge/\mathfrak p) = 1$
であることを用いる．実際，$\mathcal{I}_y^\wedge$ は $1$ 元で生成される．
この二つの注意を合わせると，$(\mathcal{F}_n)$ を補題
\ref{algebraization-lemma-improvement-formal-coherent-module-better} で得られる対象
$(\mathcal{H}_n)$，すなわち $\textit{Coh}(P, \mathcal{I})$ の対象で置き換えてよい．
したがって，$(\mathcal{F}_n)$ は逆系 $(\mathcal{F}_n'')$，すなわち連接
$\mathcal{O}_P$-加群の逆系とプロ同型であり，
$\text{depth}(\mathcal{F}''_{n, q}) + \dim(\overline{\{q\}}) \geq 2$ を
任意の $q \in Q$ に対して満たすと仮定してよい．

\medskip\noindent
More on Morphisms, 補題 \ref{more-morphisms-lemma-apply-proj-spec} を用い，
さらに More on Morphisms, 節 \ref{more-morphisms-section-proj-spec} の
記法と結果を断りなく用いる．少なくともこの補題の主張を理解しなければ，
以下の証明は意味をなさない．この場合，
$A = \bigoplus_{m \geq 0} \Gamma(P, \mathcal{L}^{\otimes m})$
は有限型 $k$-代数であり，その各次数部分は有限次元 $k$-ベクトル空間であることに注意する．
Cohomology of Schemes, 補題 \ref{coherent-lemma-coherent-proper-ample} を参照せよ．

\medskip\noindent
Cohomology of Schemes, 補題 \ref{coherent-lemma-inverse-systems-pullback} により，
$\pi|_U : U \to P$ による引き戻し $(\pi|_U^*\mathcal{F}_n)$ は
$\textit{Coh}(U, f\mathcal{O}_U)$ の対象であり，逆系
$(\pi|_U^*\mathcal{F}_n'')$，すなわち連接 $\mathcal{O}_U$-加群の逆系とプロ同型である．
次を主張する．
$$
\text{depth}(\pi|_U^*\mathcal{F}''_{n, y}) + \delta_Z^Y(y) \geq 3
$$
これは任意の $y \in U \cap Y$ に対して成り立つ．$Z$ のすべての点は閉点なので，
$\delta_Z^Y(y) \geq \dim(\overline{\{y\}})$ であり，これは任意の
$y \in U \cap Y$ に対して成り立つ．補題 \ref{algebraization-lemma-discussion} を参照せよ．
$q \in Q$ を $y$ の像とする．射 $\pi : U \to P$ は相対次元 $1$ の滑らかな射なので，
$y$ は $\pi$ のあるファイバーの閉点または生成点である．したがって
$$
\text{depth}(\pi^*\mathcal{F}''_{n, y}) + \delta_Z^Y(y)
\geq
\text{depth}(\pi^*\mathcal{F}''_{n, y}) + \dim(\overline{\{y\}}) =
\text{depth}(\mathcal{F}''_{n, q}) + \dim(\overline{\{q\}}) + 1
$$
である．これは深さまたは次元のいずれかが $1$ 増えるためである．
これで主張が証明された．

\medskip\noindent
補題 \ref{algebraization-lemma-cd-1-canonical} により，$(\pi|_U^*\mathcal{F}_n)$ は
$X$ へ標準的に延長されると結論する．次に
$$
M_n = \Gamma(U, \pi|_U^*\mathcal{F}_n) =
\bigoplus\nolimits_{m \in \mathbf{Z}}
\Gamma(P, \mathcal{F}_n \otimes_{\mathcal{O}_P} \mathcal{L}^{\otimes m})
$$
は標準的に次数付き $A$-加群である．More on Morphisms, 式
(\ref{more-morphisms-equation-cohomology-torsor}) を参照せよ．
Properties, 補題 \ref{properties-lemma-quasi-coherent-quasi-affine} により
$\pi|_U^*\mathcal{F}_n = \widetilde{M_n}|_U$ である．
したがって補題 \ref{algebraization-lemma-Gm-equivariant-extend-canonically} を適用して，
有限次数付き $A$-加群 $N$ で，$(M_n)$ と $(N/I^nN)$ が次数付き
$A$-加群を $A_+$-捩れ加群で割った圏においてプロ同型となるものを得る．
$\mathcal{F}$ を連接 $\mathcal{O}_P$-加群で，$N$ に付随するものとする．
Cohomology of Schemes, 命題
\ref{coherent-proposition-coherent-modules-on-proj-general} を参照せよ．
同じ命題により，$(\mathcal{F}/\mathcal{I}^n\mathcal{F})$ は
$(\mathcal{F}_n)$ とプロ同型である．両者は $\textit{Coh}(P, \mathcal{I})$ の対象なので，
補題 \ref{algebraization-lemma-recognize-formal-coherent-modules} により証明が完了する．
\end{proof}

\begin{example}
\label{algebraization-example-lefschetz-existence}
$k$ を体とし，$X$ を $k$ 上の固有多様体とする．
$Y \subset X$ を $\mathcal{O}_X(Y)$ が豊富となる有効 Cartier 因子とし，
$\mathcal{I} \subset \mathcal{O}_X$ を対応するイデアル層と記す．
$(\mathcal{E}_n)$ を $\textit{Coh}(X, \mathcal{I})$ の対象とし，各 $\mathcal{E}_n$ は
有限局所自由であるとする．命題 \ref{algebraization-proposition-lefschetz-existence} の
いくつかの特別な場合を挙げる．
\begin{enumerate}
\item $X$ が曲線または曲面ならば，何も得られない．
\item $X$ が Cohen--Macaulay 三次元多様体ならば，$(\mathcal{E}_n)$ は
連接 $\mathcal{O}_X$-加群 $\mathcal{E}$ の完備化である．
\item より一般に，$\dim(X) \geq 3$ かつ $X$ が $(S_3)$ ならば，
$(\mathcal{E}_n)$ は連接 $\mathcal{O}_X$-加群 $\mathcal{E}$ の完備化である．
\end{enumerate}
もちろん，$\mathcal{E}$ が存在すれば，$\mathcal{E}$ は $Y$ のある開近傍上で有限局所自由である．
\end{example}

\begin{proposition}
\label{algebraization-proposition-lefschetz-equivalence}
$k$ を体とする．$X$ を $k$ 上の固有スキームとする．
$\mathcal{L}$ を豊富な可逆 $\mathcal{O}_X$-加群とし，
$s \in \Gamma(X, \mathcal{L})$ とする．$Y = Z(s)$ を $s$ の零点スキームとし，
$\mathcal{I} \subset \mathcal{O}_X$ を対応するイデアル層と記す．
$\mathcal{V}$ を $X$ の開部分スキーム全体の集合で，$Y$ を含むものとし，
包含関係の逆で順序を入れる．任意の $x \in X \setminus Y$ に対して
$$
\text{depth}(\mathcal{O}_{X, x}) + \dim(\overline{\{x\}}) > 2
$$
が成り立つと仮定する．このとき完備化関手
$$
\colim_\mathcal{V}
\textit{Coh}(\mathcal{O}_V)
\longrightarrow
\textit{Coh}(X, \mathcal{I})
$$
は有限局所自由対象からなる充満部分圏の間の同値である．
\end{proposition}

\begin{proof}
充満忠実性を証明するには，
$$
\colim_\mathcal{V} \Gamma(V, \mathcal{L}^{\otimes m})
\longrightarrow
\lim \Gamma(Y_n, \mathcal{L}^{\otimes m}|_{Y_n})
$$
がすべての $m$ に対して同型であることを示せば十分である．
補題 \ref{algebraization-lemma-completion-fully-faithful-general} を参照せよ．
これは補題 \ref{algebraization-lemma-lefschetz-addendum} から従う．

\medskip\noindent
本質的全射性を示す．$(\mathcal{F}_n)$ を $\textit{Coh}(X, \mathcal{I})$ の
有限局所自由対象とする．このとき $y \in Y$ に対して
$\mathcal{F}_y^\wedge = \lim \mathcal{F}_{n, y}$ は有限自由
$\mathcal{O}_{X, y}^\wedge$-加群である．
$\mathfrak p \subset \mathcal{O}_{X, y}^\wedge$ を，
$\mathfrak p \not \in V(\mathcal{I}_y^\wedge)$ を満たす素イデアルとする．
すると $\mathfrak p$ は素イデアル $\mathfrak p_0 \subset \mathcal{O}_{X, y}$ の上にあり，
これは特殊化 $x \leadsto y$ に対応し，$x \not \in Y$ を満たす．
Local Cohomology, 補題 \ref{local-cohomology-lemma-change-completion} と
若干の次元論（Varieties, 節 \ref{varieties-section-algebraic-schemes} を参照）により
$$
\text{depth}((\mathcal{O}_{X, y}^\wedge)_\mathfrak p) +
\dim(\mathcal{O}_{X, y}^\wedge/\mathfrak p) =
\text{depth}(\mathcal{O}_{X, x}) +
\dim(\overline{\{x\}}) - \dim(\overline{\{y\}})
$$
である．したがって仮定から，命題 \ref{algebraization-proposition-lefschetz-existence} の
仮定が満たされ，$(\mathcal{F}_n)$ は連接 $\mathcal{O}_X$-加群
$\mathcal{F}$ の完備化であることが分かる．さらに，$\mathcal{F}_y$ は有限自由であり，
これはすべての $y \in Y$ に対して成り立つから，$\mathcal{F}$ はある開近傍 $V$，すなわち $Y$ の開近傍上で
有限局所自由である．これで証明が完了する．
\end{proof}
